% Generated mechanically from frozen input p07/ko/algebraization.tex.
% Do not edit this generated file; edit the bound locale source instead.

% OK, start here.
%

\setcounter{chapter}{51}
\chapter{대수적 기하와 형식기하}



\phantomsection
\label{algebraization-section-phantom}


\section{서론}
\label{algebraization-section-introduction}

\noindent
이 장에서는 형식 대수기하에 관한 연구를 계속하며, 특히 형식적 대상이
대수적 대상의 완비화인지라는 문제를 다룬다. 기본 참고문헌은 \cite{SGA2}이다.
Stacks 프로젝트에서 이미 논의한 결과를 나열하면 다음과 같다.
\begin{enumerate}
\item 형식함수 정리. 스킴의 코호몰로지의 절
\ref{coherent-section-theorem-formal-functions}을 보라.
\item 연접 형식 가군. 스킴의 코호몰로지의 절
\ref{coherent-section-coherent-formal}을 보라.
\item Grothendieck 존재 정리. 스킴의 코호몰로지의 절
\ref{coherent-section-existence}, \ref{coherent-section-existence-proper},
\ref{coherent-section-existence-proper-support}을 보라.
\item Grothendieck 대수화 정리. 스킴의 코호몰로지의 절
\ref{coherent-section-algebraization}을 보라.
\item 더 일반적인 Grothendieck 존재 정리. 평탄성 심화의 절
\ref{flat-section-existence}과 \ref{flat-section-existence-derived}를 보라.
\end{enumerate}
이 장의 내용을 개관하자.

\medskip\noindent
스킴 $X$와 유한형 준연접 아이디얼 층
$\mathcal{I} \subset \mathcal{O}_X$를 잡자. 이 장에서 다루는 많은 문제는
역계 $(\mathcal{F}_n)$과 관련된다. 이 역계의 항들은 준연접
$\mathcal{O}_X$-가군이고 다음을 만족한다.
$\mathcal{F}_n = \mathcal{F}_{n + 1}/\mathcal{I}^n\mathcal{F}_{n + 1}$.
중요한 특수한 경우는 $X$가 뇌터 환 $A$ 위의 스킴이고
$\mathcal{I} = I \mathcal{O}_X$인 경우이다. 여기서 $I \subset A$는 아이디얼이다.
층 코호몰로지의 절 \ref{cohomology-section-inverse-systems},
\ref{cohomology-section-inverse-systems-bis},
\ref{cohomology-section-inverse-systems-tri}에는 몇 가지 일반적인 결과가 있다.
이 장의 절 \ref{algebraization-section-ML-degree-zero}과
\ref{algebraization-section-formal-functions-principal}에는 스킴과 준연접 가군에 특유한
결과가 들어 있다. 절 \ref{algebraization-section-formal-sections-cd-one}에서는
$\lim H^p(X, \mathcal{F}_n)$ 위의 극한 위상이 $I$-진 위상임을
$\text{cd}(A, I) = 1$인 경우에 증명한다. 이 장의 한 주제는 주 아이디얼의 경우
$I = (f)$에 증명된 결과가 $\text{cd}(A, I) = 1$만을 가정해도 성립함을
보이는 것이다.

\medskip\noindent
절 \ref{algebraization-section-derived-completion}에서는 환 달린 사이트
$(\mathcal{C}, \mathcal{O})$ 위의 가군을 유한형 아이디얼 층
$\mathcal{I}$에 관하여 유도 완비하는 것을 논의한다. 이 절은 대수학 심화의
절 \ref{more-algebra-section-derived-completion}에서 설명한 가환대수의
유도 완비 이론을 자연스럽게 잇는다. 첫째 주요 결과는 유도 완비가 존재한다는
것이다. 둘째 주요 결과는 환 달린 사이트의 사상 $f$에 대하여 유도 완비가
유도 직접상과 가환한다는 것이다.
$$
(Rf_*K)^\wedge = Rf_*(K^\wedge)
$$
여기서 위쪽의 아이디얼 층이 아래쪽 아이디얼에서 오는 단면들로 국소적으로
생성된다고 가정한다. 보조정리
\ref{algebraization-lemma-pushforward-commutes-with-derived-completion}을 보라. 문제의
아이디얼들이 대역적으로 유한 생성이면 두 주요 결과 모두 매우 초등적임을
강조한다. 이 장에서 이 이론을 적용하는 모든 경우가 여기에 해당한다. 위에
표시한 등식은 형식함수 정리의 ``올바른'' 형태이다. 절
\ref{algebraization-section-formal-functions}의 논의를 보라.

\medskip\noindent
$A$를 뇌터 환이라 하고 $I, J$를 $A$의 두 아이디얼이라 하자. $M$을 유한
$A$-가군이라 하자. 이 장의 다음 주제는 사상
$$
R\Gamma_J(M) \longrightarrow R\Gamma_J(M)^\wedge
$$
이다. 이는 $M$의 국소 코호몰로지에서 그 유도 $I$-진 완비로 가는 사상이다.
적절한 깊이 조건을 부과하면 이 사상은 일정한 차수 범위의 코호몰로지에서
동형사상이 된다. 절 \ref{algebraization-section-algebraization-sections-general}에서는
본질적으로 방금 말한 일반성에서 작업한다. 절
\ref{algebraization-section-algebraization-punctured}에서는 $A$가 국소환이고
$J = \mathfrak m$이 극대 아이디얼이라고 가정한다. 독자에게 이 절을 이 부분의
다른 두 절보다 먼저 읽기를 권한다. 마지막으로 절 \ref{algebraization-section-bootstrap}에서는
국소적인 경우를 발판으로 삼아 일반적인 경우에 대한 더 강한 결과를 얻는다.

\medskip\noindent
이 장의 다음 부분에서는 국소 코호몰로지의 완비화에 관한 결과를 사용하여
연접 가군의 완비화의 코호몰로지에 관한 몇 가지 결과를 얻는다. 이 목록이
망라적이지는 않다. 더 정확히 말해 $A$를 뇌터 환, $I \subset A$를 아이디얼,
$U \subset \Spec(A)$를 열린 부분스킴이라 하자. $\mathcal{F}$가 연접
$\mathcal{O}_U$-가군이면 다음 사상들을 생각할 수 있다.
$$
H^i(U, \mathcal{F}) \longrightarrow \lim H^i(U, \mathcal{F}/I^n\mathcal{F})
$$
그리고 일정한 차수 범위에서 동형사상을 얻는지 물을 수 있다. 절
\ref{algebraization-section-algebraization-sections}에서는 $U$가 국소환의 천공 스펙트럼인
몇 가지 예를 계산한다. 절 \ref{algebraization-section-algebraization-sections-coherent}에서는
일반적인 경우를 논의한다. 절 \ref{algebraization-section-connected}에서는 얻은 결과 가운데
일부를 대수기하의 연결성 문제에 적용한다.

\medskip\noindent
이 장의 나머지 절들은 연접 형식 가군의 대수화를 논의하는 데 할애한다.
다시 말해 연접 가군의 역계 $(\mathcal{F}_n)$을 위와 같은 $U$ 위에 잡고
$\mathcal{F}_n = \mathcal{F}_{n + 1}/I^n\mathcal{F}_{n + 1}$
를 만족한다고 하자. 연접 $\mathcal{O}_U$-가군 $\mathcal{F}$로서
$\mathcal{F}_n = \mathcal{F}/I^n\mathcal{F}$
가 모든 $n$에 대하여 성립하는 것이 존재하는지 묻는다. 절
\ref{algebraization-section-algebraization-modules}을 읽기를 권한다. 거기에는 이 문제의
정확한 진술, 유용한 일반 결과(보조정리 \ref{algebraization-lemma-when-done}), 그리고
자명하지 않은 응용(보조정리 \ref{algebraization-lemma-algebraization-principal-variant})이
있다. 이 경우를 본질적으로
넘어서는 결과를 증명하려면 훨씬 더 많은 이론을 전개해야 한다. 이 장에서 얻은
이 유형의 가장 강한 결과는 절 \ref{algebraization-section-algebraization-modules-conclusion}을
보라.



\section{형식 단면, I}
\label{algebraization-section-ML-degree-zero}

\noindent
먼저 층 코호몰로지의 절 \ref{cohomology-section-inverse-systems}을
살펴보기를 권한다.

\begin{lemma}
\label{algebraization-lemma-properties-system}
$X$를 스킴이라 하고 $\mathcal{I} \subset \mathcal{O}_X$를 준연접
아이디얼 층이라 하자. 다음을
$$
\ldots \to \mathcal{F}_3 \to \mathcal{F}_2 \to \mathcal{F}_1
$$
준연접 $\mathcal{O}_X$-가군의 역계라 하고, 다음을 만족한다고 하자.
$\mathcal{F}_n = \mathcal{F}_{n + 1}/\mathcal{I}^n\mathcal{F}_{n + 1}$.
$\mathcal{F} = \lim \mathcal{F}_n$로 놓자. 그러면
\begin{enumerate}
\item $\mathcal{F} = R\lim \mathcal{F}_n$,
\item 임의의 아핀 열린집합 $U \subset X$에 대하여
$H^p(U, \mathcal{F}) = 0$이 $p > 0$일 때 성립하고,
\item 각 $p$에 대하여 다음 짧은 완전열이 있다.
$0 \to R^1\lim H^{p - 1}(X, \mathcal{F}_n) \to
H^p(X, \mathcal{F}) \to \lim H^p(X, \mathcal{F}_n) \to 0$.
\end{enumerate}
더욱이 $\mathcal{I}$가 유한형이면
\begin{enumerate}
\item[(4)]
$\mathcal{F}_n = \mathcal{F}/\mathcal{I}^n\mathcal{F}$이고,
\item[(5)]
$\mathcal{I}^n \mathcal{F} = \lim_{m \geq n} \mathcal{I}^n\mathcal{F}_m$.
\end{enumerate}
\end{lemma}

\begin{proof}
항목 (1), (2), (3)은 전이 사상이 전사인 준연접 가군의 역계에 관한
일반적인 사실이다. 스킴의 유도 범주의 보조정리
\ref{perfect-lemma-Rlim-quasi-coherent}와 층 코호몰로지의 보조정리
\ref{cohomology-lemma-RGamma-commutes-with-Rlim}를 보라. 이제
$\mathcal{I}$가 유한형이라고 가정하자. $U \subset X$를 아핀 열린집합이라
하고 $U = \Spec(A)$라 쓰자. 또한 $\mathcal{I}|_U$가 $I \subset A$에
대응한다고 하자. $I$는 유한 생성 아이디얼임을 주목하라. 준연접
$\mathcal{O}_U$-가군과 $A$-가군 사이의 범주 동치(스킴의 보조정리
\ref{schemes-lemma-equivalence-quasi-coherent})에 의해
$M_n = \mathcal{F}_n(U)$은 $A$-가군의 역계이고
$M_n = M_{n + 1}/I^nM_{n + 1}$을 만족한다. 따라서
$$
M = \mathcal{F}(U) = \lim \mathcal{F}_n(U) = \lim M_n
$$
는 대수학의 보조정리 \ref{algebra-lemma-limit-complete}에 의해 $I$-진 완비
가군이며 $M/I^nM = M_n$을 만족한다. 이로써 (4)가 증명된다.
항목 (5)는 다음 명제로 옮겨진다.
$\lim_{m \geq n} I^nM/I^mM = I^nM$.
$I^mM = I^{m - n} \cdot I^nM$이므로, 이는 $I^mM$이 $I$-진 완비라는
명제일 뿐이다. 이는 대수학의 보조정리
\ref{algebra-lemma-hathat-finitely-generated}와 $M$이 완비라는 사실에서
따른다.
\end{proof}





\section{형식 단면, II}
\label{algebraization-section-formal-functions-principal}

\noindent
먼저 층 코호몰로지의 절 \ref{cohomology-section-inverse-systems-bis}와
\ref{cohomology-section-inverse-systems-tri}를 살펴보기를 권한다.

\begin{lemma}
\label{algebraization-lemma-equivalent-f-good}
$X$를 스킴이라 하고 $f \in \Gamma(X, \mathcal{O}_X)$라 하자. 다음을
$$
\ldots \to \mathcal{F}_3 \to \mathcal{F}_2 \to \mathcal{F}_1
$$
준연접 $\mathcal{O}_X$-가군의 역계라 하자. 다음 조건들은 서로 동치이다.
\begin{enumerate}
\item 모든 $n \geq 1$에 대하여 사상
$f : \mathcal{F}_{n + 1} \to \mathcal{F}_{n + 1}$은
$\mathcal{F}_{n + 1} \to \mathcal{F}_n$을 통하여 인수분해되어 다음 짧은
완전열을 준다.
$0 \to \mathcal{F}_n \to \mathcal{F}_{n + 1} \to \mathcal{F}_1 \to 0$,
\item 모든 $n \geq 1$에 대하여 사상
$f^n : \mathcal{F}_{n + 1} \to \mathcal{F}_{n + 1}$
은 $\mathcal{F}_{n + 1} \to \mathcal{F}_1$을 통하여 인수분해되어 다음 짧은
완전열을 준다.
$0 \to \mathcal{F}_1 \to \mathcal{F}_{n + 1} \to \mathcal{F}_n \to 0$
\item $\mathcal{O}_X$-가군 $\mathcal{G}$가 존재하며, 이는
$f$-가분이고 $\mathcal{F}_n = \mathcal{G}[f^n]$을 만족한다.
\item $\mathcal{O}_X$-가군 $\mathcal{F}$가 존재하며, 이는
$f$-꼬임이 없고 $\mathcal{F}_n = \mathcal{F}/f^n\mathcal{F}$을 만족한다.
\end{enumerate}
\end{lemma}

\begin{proof}
(1), (2), (3)의 동치와 함의 (4) $\Rightarrow$ (1)은 층 코호몰로지의
보조정리 \ref{cohomology-lemma-equivalent-f-good}에서 증명했다. (1)이
성립한다고 가정하자. $\mathcal{F} = \lim \mathcal{F}_n$으로 놓자.
보조정리 \ref{algebraization-lemma-properties-system}의 항목 (4)에 의해
$\mathcal{F}_n = \mathcal{F}/f^n\mathcal{F}$이다. $U \subset X$를
열린집합이라 하고 $s = (s_n) \in \mathcal{F}(U) = \lim \mathcal{F}_n(U)$라
하자. $n \geq 1$을 택하자. $fs = 0$이면 조건 (1)에 의해 $s_{n + 1}$은
$\mathcal{F}_{n + 1} \to \mathcal{F}_n$의 핵에 속한다. 따라서 $s_n = 0$이다.
$n$이 임의였으므로 $s = 0$이다. 그러므로 $\mathcal{F}$에는 $f$-꼬임이 없다.
\end{proof}

\begin{lemma}
\label{algebraization-lemma-formal-functions-principal}
\begin{reference}
\cite[보조정리 1.6]{Bhatt-local}을 약간 개선한 형태이다.
\end{reference}
$A$를 환이라 하고 $f \in A$라 하자. $X$를 $A$ 위의 스킴이라 하고
$\mathcal{F}$를 준연접 $\mathcal{O}_X$-가군이라 하자.
$\mathcal{F}[f^n] = \Ker(f^n : \mathcal{F} \to \mathcal{F})$가
안정화된다고 가정하자. 그러면
$$
R\Gamma(X, \lim \mathcal{F}/f^n\mathcal{F}) =
R\Gamma(X, \mathcal{F})^\wedge
$$
이다. 여기서 오른쪽은 아이디얼 $(f) \subset A$에 관한 유도 완비를 나타낸다.
따라서 $p \in \mathbf{Z}$에 대하여 다음 가환도식을 얻는다.
$$
\xymatrix{
& 0 & 0 \\
0 \ar[r] &
\widehat{H^p(X, \mathcal{F})} \ar[r] \ar[u] &
\lim H^p(X, \mathcal{F}/f^n\mathcal{F}) \ar[r] \ar[u] &
T_f(H^{p + 1}(X, \mathcal{F})) \ar[r] &
0 \\
0 \ar[r] &
H^0(H^p(X, \mathcal{F})^\wedge) \ar[r] \ar[u] &
H^p(X, \lim \mathcal{F}/f^n\mathcal{F}) \ar[r] \ar[u] &
T_f(H^{p + 1}(X, \mathcal{F})) \ar[r] \ar@{=}[u] &
0 \\
&
R^1\lim H^p(X, \mathcal{F})[f^n] \ar[u] \ar[r]^\cong &
R^1\lim H^{p - 1}(X, \mathcal{F}/f^n\mathcal{F}) \ar[u] \\
& 0 \ar[u] & 0 \ar[u]
}
$$
이 도식의 행과 열은 완전하며,
$\widehat{H^p(X, \mathcal{F})} =
\lim H^p(X, \mathcal{F})/f^n H^p(X, \mathcal{F})$
는 통상적인 $f$-진 완비이고 $T_f(-)$는 대수학 심화의 예
\ref{more-algebra-example-spectral-sequence-principal}에서처럼 $f$-진
Tate 가군을 나타낸다.
\end{lemma}

\begin{proof}
보조정리 \ref{algebraization-lemma-properties-system}에 의해
$\lim \mathcal{F}/f^n\mathcal{F} = R\lim \mathcal{F}/f^n \mathcal{F}$.
나머지는 모두 층 코호몰로지의 예
\ref{cohomology-example-formal-functions-principal}에서 따른다.
\end{proof}









\section{형식 단면, III}
\label{algebraization-section-formal-sections-cd-one}

\noindent
이 절에서는 보조정리 \ref{algebraization-lemma-topology-I-adic}을 증명한다. 이는 (뇌터
스킴과 연접 가군의 설정에서) 아이디얼 $I$가 주 아이디얼이라고 가정하지 않고
$\text{cd}(A, I) = 1$이라는 성질을 가정하는 경우의 층 코호몰로지의 보조정리
\ref{cohomology-lemma-topology-I-adic-f}에 해당한다.

\begin{lemma}
\label{algebraization-lemma-cd-one}
$I = (f_1, \ldots, f_r)$를 뇌터 환 $A$의 아이디얼이라 하자.
$\text{cd}(A, I) = 1$이면, $c \geq 1$과 사상
$\varphi_j : I^c \to A$가 존재하여 $\sum f_j \varphi_j : I^c \to I$가
포함사상이 된다.
\end{lemma}

\begin{proof}
$\text{cd}(A, I) = 1$이므로 여집합 $U = \Spec(A) \setminus V(I)$는
아핀이다(국소 코호몰로지의 보조정리
\ref{local-cohomology-lemma-cd-is-one}). $U = \Spec(B)$라 쓰자. 그러면
$IB = B$이고 $1 = \sum_{j = 1, \ldots, r} f_j b_j$로 쓸 수 있는데, 여기서
어떤 $b_j \in B$들이 존재한다. 스킴의 코호몰로지의 보조정리
\ref{coherent-lemma-homs-over-open}에 의해 $b_j$를 사상
$\varphi_j : I^c \to A$로 나타낼 수 있으며, 이는 어떤 $c \geq 0$에 대하여
성립한다. 그러면 $\sum f_j \varphi_j : I^c \to I \subset A$는 같은
보조정리에 따라 필요하면 $c$를 더 큰 정수로 바꾼 뒤 표준 매장이 된다.
\end{proof}

\begin{lemma}
\label{algebraization-lemma-cd-one-extend}
$I = (f_1, \ldots, f_r)$를 뇌터 환 $A$의 아이디얼이라 하고
$\text{cd}(A, I) = 1$이라 하자. $c \geq 1$과 $\varphi_j : I^c \to A$,
$j = 1, \ldots, r$를 보조정리 \ref{algebraization-lemma-cd-one}의 것이라 하자. 그러면
유일한 등급 $A$-대수 사상
$$
\Phi : \bigoplus\nolimits_{n \geq 0} I^{nc} \to A[T_1, \ldots, T_r]
$$
이 존재하여 $\Phi(g) = \sum \varphi_j(g) T_j$를 $g \in I^c$에 대하여
만족한다. 더욱이 $\Phi$와 사상
$A[T_1, \ldots, T_r] \to \bigoplus_{n \geq 0} I^n$,
$T_j \mapsto f_j$의 합성은 포함사상
$\bigoplus_{n \geq 0} I^{nc} \to \bigoplus_{n \geq 0} I^n$.
이다.
\end{lemma}

\begin{proof}
각 $j$와 $m \geq c$에 대하여 $\varphi_j$를 $I^m$에 제한하면 사상
$\varphi_j : I^m \to I^{m - c}$를 얻는다.
$j_1, \ldots, j_n \in \{1, \ldots, r\}$가 주어졌을 때 합성
$$
\varphi_{j_1} \ldots \varphi_{j_n} :
I^{nc} \to I^{(n - 1)c} \to \ldots \to I^c \to A
$$
은 지표 $j_1, \ldots, j_n$의 순서와 무관하다고 주장한다. 실제로
$g = g_1 \ldots g_n$이고 $g_i \in I^c$이면,
$$
(\varphi_{j_1} \ldots \varphi_{j_n})(g) =
\varphi_{j_1}(g_1) \ldots \varphi_{j_n}(g_n)
$$
이며, $A$에서의 곱셈이 가환이므로 이는 순서와 무관하다. 따라서
$\Phi$를, $g \in I^{nc}$를 다음 원소로 보내도록 정의할 수 있다.
$$
\Phi(g) = \sum\nolimits_{e_1 + \ldots + e_r = n}
(\varphi_1^{e_1} \circ \ldots \circ \varphi_r^{e_r})(g)
T_1^{e_1} \ldots T_r^{e_r}
$$
이것이 원하는 성질을 지닌 등급 $A$-대수 준동형임은 직접 확인된다.
유일성과 마지막 성질도 즉시 따른다. 이로써 보조정리가 증명되었다.
\end{proof}

\begin{lemma}
\label{algebraization-lemma-cd-one-extend-to-module}
$I = (f_1, \ldots, f_r)$를 뇌터 환 $A$의 아이디얼이라 하고
$\text{cd}(A, I) = 1$이라 하자. $c \geq 1$과 $\varphi_j : I^c \to A$,
$j = 1, \ldots, r$를 보조정리 \ref{algebraization-lemma-cd-one}의 것이라 하자.
$A \to B$를 환 준동형이라 하되 $B$는 뇌터 환이라 하고, $N$을 유한
$B$-가군이라 하자. 그러면 필요하면 $c$를 키우고 그에 따라
$\varphi_j$를 조정한 뒤 유일한 등급 $B$-가군 사상
$$
\Phi_N : \bigoplus\nolimits_{n \geq 0} I^{nc}N \to N[T_1, \ldots, T_r]
$$
이 존재하여 $\Phi_N(g x) = \Phi(g) x$를 $g \in I^{nc}$와 $x \in N$에
대하여 만족한다. 여기서 $\Phi$는 보조정리
\ref{algebraization-lemma-cd-one-extend}의 것이다. $\Phi_N$과 사상
$N[T_1, \ldots, T_r] \to \bigoplus_{n \geq 0} I^nN$,
$T_j \mapsto f_j$의 합성은 포함사상
$\bigoplus_{n \geq 0} I^{nc}N \to \bigoplus_{n \geq 0} I^nN$.
이다.
\end{lemma}

\begin{proof}
공식과 보조정리 \ref{algebraization-lemma-cd-one-extend}에서의 $\Phi$의 유일성으로부터
유일성은 명백하다. 뇌터 $A$-대수 $B' = B \oplus N$을 생각하자. 여기서
$N$은 제곱 영 아이디얼이다. $\Phi_N$의 존재를 보이려면 (보조정리
\ref{algebraization-lemma-cd-one}에 의해) $\varphi_j$가 사상
$\varphi'_j : I^cB' \to B'$로 확장됨을 보이면 충분하다. 이때 필요하면
$c$를 어떤 $c'$까지 키운다.
이때 $\varphi_j$는 포함사상 $I^{c'} \to I^c$와 $\varphi_j$의 합성으로
바꾼다. $\varphi_j$가 단면
$$
h_j \in \Gamma(\Spec(A) \setminus V(I), \mathcal{O}_{\Spec(A)})
$$
에 대응함을 상기하자. 스킴의 코호몰로지의 보조정리
\ref{coherent-lemma-homs-over-open}을 보라. (실제로 보조정리
\ref{algebraization-lemma-cd-one}의 증명에서 $\varphi_j$를 바로 이렇게 택했다.) 같은
보조정리를 사용하여 다음 당김을
$$
h'_j \in \Gamma(\Spec(B') \setminus V(IB'), \mathcal{O}_{\Spec(B')})
$$
즉 $h_j$의 당김을 $B'$-선형 사상
$\varphi'_j : I^{c'}B' \to B'$로 나타내자. 이는 어떤 $c' \geq c$에
대하여 가능하다. 이 보조정리를 한 번 더 적용하면 $\varphi_j$와의 일치는
$c'$가 충분히 클 때 성립한다. 실제로
$I^{c'}$의 유한한 생성원 목록에서 일치를 검사할 수 있다. 사소한
세부사항은 생략한다.
\end{proof}

\begin{lemma}
\label{algebraization-lemma-cd-is-one-for-system}
$I = (f_1, \ldots, f_r)$를 뇌터 환 $A$의 아이디얼이라 하고
$\text{cd}(A, I) = 1$이라 하자. $c \geq 1$과 $\varphi_j : I^c \to A$,
$j = 1, \ldots, r$를 보조정리 \ref{algebraization-lemma-cd-one}의 것이라 하자.
$X$를 $\Spec(A)$ 위의 뇌터 스킴이라 하자. 다음을
$$
\ldots \to \mathcal{F}_3 \to \mathcal{F}_2 \to \mathcal{F}_1
$$
연접 $\mathcal{O}_X$-가군의 역계라 하고
$\mathcal{F}_n = \mathcal{F}_{n + 1}/I^n\mathcal{F}_{n + 1}$을 만족한다고
하자. $\mathcal{F} = \lim \mathcal{F}_n$으로 놓자. 그러면 필요하면 $c$를
키우고 그에 따라 $\varphi_j$를 조정한 뒤 유일한 등급
$\mathcal{O}_X$-가군 사상
$$
\Phi_\mathcal{F} :
\bigoplus\nolimits_{n \geq 0} I^{nc}\mathcal{F}
\longrightarrow
\mathcal{F}[T_1, \ldots, T_r]
$$
이 존재하여 $\Phi_\mathcal{F}(g s) = \Phi(g) s$를 만족한다. 여기서
$g \in I^{nc}$이고 $s$는 $\mathcal{F}$의 국소 단면이며, $\Phi$는 보조정리
\ref{algebraization-lemma-cd-one-extend}의 것이다. $\Phi_\mathcal{F}$와 사상
$\mathcal{F}[T_1, \ldots, T_r] \to \bigoplus_{n \geq 0} I^n\mathcal{F}$,
$T_j \mapsto f_j$
의 합성은 표준 포함사상
$\bigoplus_{n \geq 0} I^{nc}\mathcal{F} \to
\bigoplus_{n \geq 0} I^n\mathcal{F}$.
이다.
\end{lemma}

\begin{proof}
$\mathcal{O}_X$-선형성과 $\Phi_\mathcal{F}(g s) = \Phi(g) s$라는 조건으로부터
유일성이 즉시 따른다. 여기서 $g \in I^{nc}$이고 $s$는 $\mathcal{F}$의
국소 단면이다. 따라서 $X = \Spec(B)$가 아핀이라고 가정할 수 있다.
$(\mathcal{F}_n)$은 스킴의 코호몰로지의 절
\ref{coherent-section-coherent-formal}에서 도입한 범주
$\textit{Coh}(X, I\mathcal{O}_X)$의 대상임을 주목하라. $B' = B^\wedge$를
$I$에 관한 $B$의 진 완비라 하자. 스킴의 코호몰로지의 보조정리
\ref{coherent-lemma-inverse-systems-affine}에 의해 대상 $(\mathcal{F}_n)$은
유한 $B'$-가군 $N$에 대응한다. 이는 $\mathcal{F}_n$이 유한 $B$-가군
$N/I^n N$에 결부된 연접 가군이라는 뜻이다. 보조정리
\ref{algebraization-lemma-cd-one-extend-to-module}을 $I \subset A \to B'$와 $N$에
적용하면, 필요하면 $c$를 키우고 그에 따라 $\varphi_j$를 조정한 뒤 유일한
사상
$$
\Phi_N : \bigoplus\nolimits_{n \geq 0} I^{nc}N \to N[T_1, \ldots, T_r]
$$
을 얻으며, 이는 대응하는 성질들을 만족한다. 차수 $n$에서 사상들의 역계
$N/I^mN \to \bigoplus_{e_1 + \ldots + e_r = n}
N/I^{m - nc}N \cdot T_1^{e_1} \ldots T_r^{e_r}$을 $m \geq nc$에 대하여
얻는다는 점을 주목하라. 이를 다시 연접 층으로 옮기면 $\Phi_N$은 사상계
$$
\Phi^n_m :
I^{nc}\mathcal{F}_m
\longrightarrow
\bigoplus\nolimits_{e_1 + \ldots + e_r = n}
\mathcal{F}_{m - nc} \cdot T_1^{e_1} \ldots T_r^{e_r}
$$
에 대응함을 알 수 있다. 여기서 $m \geq nc$와 $n \geq 1$은 변한다.
$m$에 관하여 이 사상들의 역극한을 취하면
$\Phi_\mathcal{F} = \bigoplus_n \lim_m \Phi^n_m$을 얻는다. 예를 들어 아핀
열린집합에서 계산하면 $\lim_m I^t\mathcal{F}_m = I^t \mathcal{F}$임을 알 수
있다. 그러나 실제로는 사상
$I^t\mathcal{F} \to \lim_m I^t\mathcal{F}_m$이 있음이 명백하므로 이 사실은
필요하지 않다.
\end{proof}

\begin{lemma}
\label{algebraization-lemma-topology-I-adic}
$I$를 뇌터 환 $A$의 아이디얼이라 하자. $X$를 $\Spec(A)$ 위의 뇌터
스킴이라 하자. 다음을
$$
\ldots \to \mathcal{F}_3 \to \mathcal{F}_2 \to \mathcal{F}_1
$$
연접 $\mathcal{O}_X$-가군의 역계라 하고
$\mathcal{F}_n = \mathcal{F}_{n + 1}/I^n\mathcal{F}_{n + 1}$을 만족한다고
하자. $\text{cd}(A, I) = 1$이면, 모든 $p \in \mathbf{Z}$에 대하여
$\lim H^p(X, \mathcal{F}_n)$ 위의 극한 위상은 $I$-진 위상이다.
\end{lemma}

\begin{proof}
먼저 $I^t \lim H^p(X, \mathcal{F}_n)$이 $H^p(X, \mathcal{F}_t)$에서 영으로
간다는 것은 명백하다. 따라서 $I$-진 위상은 극한 위상보다 세밀하다.
역방향을 보이기 위해 $\mathcal{F} = \lim \mathcal{F}_n$으로 놓고,
$f_1, \ldots, f_r$를 택하되 이들이 $I$를 생성하게 하고, $c \geq 1$을
택한 다음,
보조정리 \ref{algebraization-lemma-cd-is-one-for-system}에서와 같이
$\Phi_\mathcal{F}$를 택한다. 이제부터 보조정리
\ref{algebraization-lemma-properties-system}의 결과를 별도로 언급하지 않고 사용한다. 특히
다음 짧은 완전열이 있다.
$$
0 \to R^1\lim H^{p - 1}(X, \mathcal{F}_n) \to H^p(X, \mathcal{F})
\to \lim H^p(X, \mathcal{F}_n) \to 0
$$
따라서 $\xi$를 $\lim H^p(X, \mathcal{F}_n)$의 임의의 원소라 하면 이를 원소
$\xi' \in H^p(X, \mathcal{F})$로 올릴 수 있다. $\xi$가
$H^p(X, \mathcal{F}_{nc})$에서 영으로 가는 어떤 $n$이 있다고 가정하자.
다시 말해 $\xi$가
극한 위상에서 ``작다''고 가정하자. 다음 짧은 완전열이 있다.
$$
0 \to I^{nc}\mathcal{F} \to \mathcal{F} \to \mathcal{F}_{nc} \to 0
$$
따라서 이 가정은 $\xi'$를 원소
$\xi'' \in H^p(X, I^{nc}\mathcal{F})$로 올릴 수 있음을 뜻한다.
$\Phi_\mathcal{F}$를 적용하면
$$
\Phi_\mathcal{F}(\xi'') = \sum\nolimits_{e_1 + \ldots + e_r = n}
\xi'_{e_1, \ldots, e_r} \cdot T_1^{e_1} \ldots T_r^{e_r}
$$
을 얻는다. 여기서 어떤
$\xi'_{e_1, \ldots, e_r} \in H^p(X, \mathcal{F})$들이 존재한다. 그 상들을
$\xi_{e_1, \ldots, e_r} \in \lim H^p(X, \mathcal{F}_n)$라 하고 보조정리
\ref{algebraization-lemma-cd-is-one-for-system}의 마지막 주장을 사용하면
$$
\xi = \sum f_1^{e_1} \ldots f_r^{e_r} \xi_{e_1, \ldots, e_r}
$$
가 원하는 대로 $I^n \lim H^p(X, \mathcal{F}_n)$에 속함을 얻는다.
\end{proof}

\begin{example}
\label{algebraization-example-not-I-adic}
$k$를 체라 하자. $A = k[x, y][[s, t]]/(xs - yt)$라 하자.
$I = (s, t)$ 및 $\mathfrak a = (x, y, s, t)$라 하자.
$X = \Spec(A) - V(\mathfrak a)$ 및
$\mathcal{F}_n = \mathcal{O}_X/I^n\mathcal{O}_X$라 하자. 유리함수
$$
g = \frac{t}{x} = \frac{s}{y}
$$
는 어떤 열린집합 $V \subset X$에서 정칙이며, 이 열린집합은
$V(I\mathcal{O}_X)$의 근방이다. 따라서 각
거듭제곱 $g^e$는 단면 $g^e \in M = \lim H^0(X, \mathcal{F}_n)$을 정한다.
$g^e \to 0$이 $e \to \infty$일 때 $M$ 위의 극한 위상에서 성립함을
주목하라. 이는 $g^e$가 $\mathcal{F}_e$에서 영으로 가기 때문이다. 한편
$g^e \not \in IM$는 모든 $e$에 대하여 성립한다. 독자는
$H^0(U, \mathcal{F}_n)$을 계산하여 이를 확인할 수 있다. 계산은 생략한다.
$\text{cd}(A, I) = 2$임을
주목하라. 따라서 보조정리 \ref{algebraization-lemma-topology-I-adic}의 결과는 최선이다.
\end{example}








\section{Mittag--Leffler 조건}
\label{algebraization-section-ML}

\noindent
뇌터 국소환의 극대 아이디얼에 관한 국소 코호몰로지를 취할 때에는 흔히
Mittag--Leffler 조건을 자동으로 얻는다. 이로부터 천공 스펙트럼 위에서 전이
사상이 전사인 연접 층 역계의 고차 코호몰로지 군에도 같은 사실이 성립한다.

\begin{lemma}
\label{algebraization-lemma-descending-chain}
$(A, \mathfrak m)$을 뇌터 국소환이라 하자.
\begin{enumerate}
\item $M$을 유한 $A$-가군이라 하자. 그러면 $A$-가군
$H^i_\mathfrak m(M)$은 모든 $i$에 대하여 내림사슬 조건을 만족한다.
\item $U = \Spec(A) \setminus \{\mathfrak m\}$를 $A$의 천공 스펙트럼이라
하자. $\mathcal{F}$를 연접 $\mathcal{O}_U$-가군이라 하자. 그러면 $A$-가군
$H^i(U, \mathcal{F})$은 $i > 0$에 대하여 내림사슬 조건을 만족한다.
\end{enumerate}
\end{lemma}

\begin{proof}
항목 (1)을 $M$의 지지집합의 차원에 대한 귀납법으로 증명한다.
$M = 0$이면 명제가 성립하므로, $M$은 영이 아니라고 가정한다.

\medskip\noindent
귀납법의 기초 단계.
$\dim(\text{Supp}(M)) = 0$이면 $M$의 지지집합은 $\{\mathfrak m\}$이다.
국소 코호몰로지의 구성에서 명백하듯
$H^0_\mathfrak m(M) = M$이고 $H^i_\mathfrak m(M) = 0$이며, 후자는
$i > 0$에 대하여 성립한다. 쌍대화 복합체의 절
\ref{dualizing-section-local-cohomology}을 보라. $M$은 유한 길이를 가지므로
(대수학의 보조정리 \ref{algebra-lemma-length-finite}) 내림사슬 조건을
만족한다.

\medskip\noindent
귀납 단계. $\dim(\text{Supp}(M)) > 0$이라 가정하자. 기초 단계에 의해 유한
가군 $H^0_\mathfrak m(M) \subset M$은 내림사슬 조건을 만족한다. 쌍대화
복합체의 보조정리 \ref{dualizing-lemma-divide-by-torsion}에 의해 $M$을
$M/H^0_\mathfrak m(M)$으로 바꿀 수 있다. 그러면
$H^0_\mathfrak m(M) = 0$이다. 즉, $M$의 깊이는 $\geq 1$이다. 쌍대화
복합체의 보조정리 \ref{dualizing-lemma-depth}를 보라.
$x \in \mathfrak m$을 택하여 $x : M \to M$이 단사이게 하자. 대수학의
보조정리 \ref{algebra-lemma-one-equation-module}에 의해
$\dim(\text{Supp}(M/xM)) = \dim(\text{Supp}(M)) - 1$이고 귀납가설을 적용할
수 있다. 지표 $i$를 택하고 완전열
$$
H^{i - 1}_\mathfrak m(M/xM) \to H^i_\mathfrak m(M) \xrightarrow{x}
H^i_\mathfrak m(M)
$$
을 생각하자. 이는 짧은 완전열
$0 \to M \xrightarrow{x} M \to M/xM \to 0$에서 나온다. 따라서
$x$-꼬임 $H^i_\mathfrak m(M)[x]$는 내림사슬 조건을 만족하는 가군의
몫가군이므로 그 자체도 내림사슬 조건을 만족한다. 그러므로
$\mathfrak m$-꼬임 부분가군 $H^i_\mathfrak m(M)[\mathfrak m]$은 내림사슬
조건을 만족한다(따라서 $A/\mathfrak m$ 위에서 유한 차원이다). 이에 따라
$\mathfrak m$-멱 꼬임 가군 $H^i_\mathfrak m(M)$은 쌍대화 복합체의
보조정리 \ref{dualizing-lemma-describe-categories}에 의해 내림사슬 조건을
만족한다.

\medskip\noindent
항목 (2)는 항목 (1)과 국소 코호몰로지의 보조정리
\ref{local-cohomology-lemma-finiteness-pushforwards-and-H1-local}에서 따른다.
\end{proof}

\begin{lemma}
\label{algebraization-lemma-ML-local}
$(A, \mathfrak m)$을 뇌터 국소환이라 하자.
\begin{enumerate}
\item $(M_n)$을 유한 $A$-가군의 역계라 하자. 그러면 역계
$H^i_\mathfrak m(M_n)$은 모든 $i$에 대하여 Mittag--Leffler 조건을
만족한다.
\item $U = \Spec(A) \setminus \{\mathfrak m\}$를 $A$의 천공 스펙트럼이라
하자. $\mathcal{F}_n$을 연접 $\mathcal{O}_U$-가군의 역계라 하자. 그러면
역계 $H^i(U, \mathcal{F}_n)$은 $i > 0$에 대하여 Mittag--Leffler 조건을
만족한다.
\end{enumerate}
\end{lemma}

\begin{proof}
보조정리 \ref{algebraization-lemma-descending-chain}에서 즉시 따른다.
\end{proof}

\begin{lemma}
\label{algebraization-lemma-terrific}
$(A, \mathfrak m)$을 뇌터 국소환이라 하자. $(M_n)$을 유한 $A$-가군의
역계라 하자. $M \to \lim M_n$을 사상이라 하되, 여기서 $M$은 유한
$A$-가군이고 어떤 $i$에 대하여 사상
$H^i_\mathfrak m(M) \to \lim H^i_\mathfrak m(M_n)$
이 동형사상이라고 하자. 그러면 역계 $H^i_\mathfrak m(M_n)$은 본질적으로
상수이며 그 값은 $H^i_\mathfrak m(M)$이다.
\end{lemma}

\begin{proof}
보조정리 \ref{algebraization-lemma-ML-local}에 의해 역계 $H^i_\mathfrak m(M_n)$은
Mittag--Leffler 조건을 만족한다. $E_n \subset H^i_\mathfrak m(M_n)$을
$H^i_\mathfrak m(M_{n'})$의 상이라 하되 $n' \gg n$이라 하자. 그러면 $(E_n)$은
전이 사상이 전사인 역계이고 $H^i_\mathfrak m(M) = \lim E_n$이다.
$H^i_\mathfrak m(M)$은 보조정리 \ref{algebraization-lemma-descending-chain}에 의해
내림사슬 조건을 만족하므로, 전사
$H^i_\mathfrak m(M) \to E_n$ 가운데 자명하지 않은 핵을 갖는 것은 유한
개뿐이다. 따라서 원하는 대로 $E_n \to E_{n - 1}$은 모든 $n \gg 0$에
대하여 동형사상이다.
\end{proof}

\begin{lemma}
\label{algebraization-lemma-local-cohomology-derived-completion}
$(A, \mathfrak m)$을 뇌터 국소환이라 하자. $I \subset A$를 아이디얼이라
하자. $M$을 유한 $A$-가군이라 하자. 그러면
$$
H^i(R\Gamma_\mathfrak m(M)^\wedge) = \lim H^i_\mathfrak m(M/I^nM)
$$
가 모든 $i$에 대하여 성립한다. 여기서 $R\Gamma_\mathfrak m(M)^\wedge$는
유도 $I$-진 완비를 나타낸다.
\end{lemma}

\begin{proof}
쌍대화 복합체의 보조정리 \ref{dualizing-lemma-completion-local}과 보조정리
\ref{algebraization-lemma-ML-local}을 적용하면 $R^1\lim$ 항들이 소멸함을 알 수 있다.
\end{proof}








\section{환 달린 사이트 위의 유도 완비}
\label{algebraization-section-derived-completion}

\noindent
처음 읽을 때에는 이 절을 건너뛰기를 강하게 권한다.

\medskip\noindent
이 내용의 대수적 판은 대수학 심화의 절
\ref{more-algebra-section-derived-completion}에서 찾을 수 있다.
$\mathcal{O}$를 사이트 $\mathcal{C}$ 위의 환층이라 하자. $f$를
$\mathcal{O}$의 대역 단면이라 하자. $\mathcal{O}_f$로 나타낼 것은 국소화의
준층 $U \mapsto \mathcal{O}(U)_f$에 결부된 층이다.

\begin{lemma}
\label{algebraization-lemma-map-twice-localize}
$(\mathcal{C}, \mathcal{O})$를 환 달린 사이트라 하자. $f$를
$\mathcal{O}$의 대역 단면이라 하자.
\begin{enumerate}
\item $L, N \in D(\mathcal{O}_f)$에 대하여
$R\SheafHom_\mathcal{O}(L, N) = R\SheafHom_{\mathcal{O}_f}(L, N)$.
이다. 특히 두 $\mathcal{O}_f$-구조는
$R\SheafHom_\mathcal{O}(L, N)$ 위에서 일치한다.
\item $K \in D(\mathcal{O})$와 $L \in D(\mathcal{O}_f)$에 대하여
$$
R\SheafHom_\mathcal{O}(L, K) =
R\SheafHom_{\mathcal{O}_f}(L, R\SheafHom_\mathcal{O}(\mathcal{O}_f, K))
$$
이다. 특히
$R\SheafHom_\mathcal{O}(\mathcal{O}_f,
R\SheafHom_\mathcal{O}(\mathcal{O}_f, K)) =
R\SheafHom_\mathcal{O}(\mathcal{O}_f, K)$.
\item $g$가 $\mathcal{O}$의 또 다른 대역 단면이면
$$
R\SheafHom_\mathcal{O}(\mathcal{O}_f, R\SheafHom_\mathcal{O}(\mathcal{O}_g, K))
= R\SheafHom_\mathcal{O}(\mathcal{O}_{gf}, K).
$$
\end{enumerate}
\end{lemma}

\begin{proof}
(1)의 증명. $\mathcal{J}^\bullet$을 $\mathcal{O}_f$-가군의 K-단사
복합체로 잡되 $N$을 나타내게 하자. 사이트 위의 코호몰로지의
보조정리 \ref{sites-cohomology-lemma-K-injective-flat}에 의해
$\mathcal{J}^\bullet$은 $\mathcal{O}$-가군의 K-단사 복합체이기도 하다.
$\mathcal{F}^\bullet$을 $\mathcal{O}_f$-가군의 복합체로 잡되 $L$을
나타내게 하자. 그러면
$$
R\SheafHom_\mathcal{O}(L, N) =
R\SheafHom_\mathcal{O}(\mathcal{F}^\bullet, \mathcal{J}^\bullet) =
R\SheafHom_{\mathcal{O}_f}(\mathcal{F}^\bullet, \mathcal{J}^\bullet)
$$
이다. 이는 사이트 위의 가군의 보조정리
\ref{sites-modules-lemma-epimorphism-modules}에서 따르는데,
$\mathcal{J}^\bullet$이 $\mathcal{O}$-가군 및 $\mathcal{O}_f$-가군의 K-단사
복합체이기 때문이다.

\medskip\noindent
(2)의 증명. $\mathcal{I}^\bullet$을 $\mathcal{O}$-가군의 K-단사
복합체로 잡되 $K$를 나타내게 하자. 그러면
$R\SheafHom_\mathcal{O}(\mathcal{O}_f, K)$는
$\SheafHom_\mathcal{O}(\mathcal{O}_f, \mathcal{I}^\bullet)$로 나타내어진다.
이는 사이트 위의 코호몰로지의 보조정리
\ref{sites-cohomology-lemma-hom-K-injective}와
\ref{sites-cohomology-lemma-K-injective-flat}에 의해 $\mathcal{O}_f$-가군 및
$\mathcal{O}$-가군의 K-단사 복합체이다. $\mathcal{F}^\bullet$을
$\mathcal{O}_f$-가군의 복합체로 잡되 $L$을 나타내게 하자. 그러면
$$
R\SheafHom_\mathcal{O}(L, K) =
R\SheafHom_\mathcal{O}(\mathcal{F}^\bullet, \mathcal{I}^\bullet) =
R\SheafHom_{\mathcal{O}_f}(\mathcal{F}^\bullet,
\SheafHom_\mathcal{O}(\mathcal{O}_f, \mathcal{I}^\bullet))
$$
이다. 이는 사이트 위의 가군의 보조정리
\ref{sites-modules-lemma-adjoint-hom-restrict}과
$\SheafHom_\mathcal{O}(\mathcal{O}_f, \mathcal{I}^\bullet)$이
$\mathcal{O}_f$-가군의 K-단사 복합체라는 사실에서 따른다.

\medskip\noindent
(3)의 증명. 이는
$R\SheafHom_\mathcal{O}(\mathcal{O}_g, \mathcal{I}^\bullet)$
이 $\mathcal{O}$-가군의 복합체로서 K-단사라는 사실과 다음 등식
$\SheafHom_\mathcal{O}(\mathcal{O}_f,
\SheafHom_\mathcal{O}(\mathcal{O}_g, \mathcal{H})) = 
\SheafHom_\mathcal{O}(\mathcal{O}_{gf}, \mathcal{H})$
이 모든 $\mathcal{O}$-가군층 $\mathcal{H}$에 대하여 성립한다는 사실에서
따른다.
\end{proof}

\noindent
$K \in D(\mathcal{O})$라 하자. $T(K, f)$로 나타낼 것은 역계
$$
\ldots \to K \xrightarrow{f} K \xrightarrow{f} K
$$
의 $D(\mathcal{O})$ 안에서의 유도 극한이다(유도 범주의 정의
\ref{derived-definition-derived-limit}).

\begin{lemma}
\label{algebraization-lemma-hom-from-Af}
$(\mathcal{C}, \mathcal{O})$를 환 달린 사이트라 하자. $f$를
$\mathcal{O}$의 대역 단면이라 하자. $K \in D(\mathcal{O})$라 하자.
다음 조건들은 서로 동치이다.
\begin{enumerate}
\item $R\SheafHom_\mathcal{O}(\mathcal{O}_f, K) = 0$,
\item $R\SheafHom_\mathcal{O}(L, K) = 0$이 모든 $L$에 대하여 성립한다.
여기서 이 대상은 $D(\mathcal{O}_f)$에 속한다,
\item $T(K, f) = 0$.
\end{enumerate}
\end{lemma}

\begin{proof}
(2)가 (1)을 함의함은 명백하다. 함의 (1) $\Rightarrow$ (2)는 보조정리
\ref{algebraization-lemma-map-twice-localize}에서 따른다. $\mathcal{O}$-가군
$\mathcal{O}_f$의 자유 분해는 다음과 같다.
$$
0 \to \bigoplus\nolimits_{n \in \mathbf{N}} \mathcal{O} \to
\bigoplus\nolimits_{n \in \mathbf{N}} \mathcal{O}
\to \mathcal{O}_f \to 0
$$
여기서 첫째 사상은 국소 단면 $(x_0, x_1, \ldots)$을
$(x_0, x_1 - fx_0, x_2 - fx_1, \ldots)$로 보내고, 둘째 사상은
$(x_0, x_1, \ldots)$을 $x_0 + x_1/f + x_2/f^2 + \ldots$로 보낸다.
$\SheafHom_\mathcal{O}(-, \mathcal{I}^\bullet)$을 적용하자. 여기서
$\mathcal{I}^\bullet$은 $\mathcal{O}$-가군의 K-단사 복합체이며 $K$를
나타낸다. 그러면 복합체의 짧은 완전열
$$
0 \to \SheafHom_\mathcal{O}(\mathcal{O}_f, \mathcal{I}^\bullet) \to
\prod \mathcal{I}^\bullet \to \prod \mathcal{I}^\bullet \to 0
$$
을 얻는다. 이는 $\mathcal{I}^n$이 단사 $\mathcal{O}$-가군이기 때문이다.
이 곱들은 $D(\mathcal{O})$ 안에서의 곱이다. 단사 대상의 보조정리
\ref{injectives-lemma-derived-products}를 보라. 이는 대상 $T(K, f)$가
$R\SheafHom_\mathcal{O}(\mathcal{O}_f, K)$의 한 대표이며 그 대표가
$D(\mathcal{O})$ 안에 있음을 뜻한다. 따라서 (1)과 (3)은 동치이다.
\end{proof}

\begin{lemma}
\label{algebraization-lemma-ideal-of-elements-complete-wrt}
$(\mathcal{C}, \mathcal{O})$를 환 달린 사이트라 하자. $K \in D(\mathcal{O})$라
하자. 각 $U$에 집합 $\mathcal{I}(U)$을 대응시키자. 이 집합은 단면
$f \in \mathcal{O}(U)$ 가운데 $T(K|_U, f) = 0$을 만족하는 것들의 집합이다.
그러면 이 규칙은 $\mathcal{O}$ 안의 아이디얼 층을 이룬다.
\end{lemma}

\begin{proof}
보조정리 \ref{algebraization-lemma-hom-from-Af}의 결과를 더 언급하지 않고 사용한다.
$f \in \mathcal{I}(U)$이고 $g \in \mathcal{O}(U)$이면
$\mathcal{O}_{U, gf}$는 $\mathcal{O}_{U, f}$-가군이므로
$R\SheafHom_\mathcal{O}(\mathcal{O}_{U, gf}, K|_U) = 0$이고, 따라서
$gf \in \mathcal{I}(U)$이다. $f, g \in \mathcal{O}(U)$라 가정하자. 그러면
짧은 완전열
$$
0 \to \mathcal{O}_{U, f + g} \to
\mathcal{O}_{U, f(f + g)} \oplus \mathcal{O}_{U, g(f + g)} \to
\mathcal{O}_{U, gf(f + g)} \to 0
$$
이 있다. 이는 $f, g$가 $\mathcal{O}(U)_{f + g}$ 안의 단위 아이디얼을
생성하기 때문이다. 이는 대수학의 보조정리
\ref{algebra-lemma-standard-covering}와 마지막 화살표가 전사라는 쉬운
사실에서 따른다. $R\SheafHom_\mathcal{O}( - , K|_U)$는 삼각 범주의
완전 함자이므로 다음 세 대상의 소멸은
$R\SheafHom_{\mathcal{O}_U}(\mathcal{O}_{U, f(f + g)}, K|_U)$,
$R\SheafHom_{\mathcal{O}_U}(\mathcal{O}_{U, g(f + g)}, K|_U)$, 그리고
$R\SheafHom_{\mathcal{O}_U}(\mathcal{O}_{U, gf(f + g)}, K|_U)$,
다음 대상의 소멸을 함의한다.
$R\SheafHom_{\mathcal{O}_U}(\mathcal{O}_{U, f + g}, K|_U)$.
층 조건의 검증은 생략한다.
\end{proof}

\noindent
임의의 환 달린 사이트에 대하여 다음 정의를 할 수 있다.

\begin{definition}
\label{algebraization-definition-derived-complete}
$(\mathcal{C}, \mathcal{O})$를 환 달린 사이트라 하자.
$\mathcal{I} \subset \mathcal{O}$를 아이디얼 층이라 하자.
$K \in D(\mathcal{O})$라 하자. $K$가 {\it $\mathcal{I}$에 관해 유도
완비}라는 것은, $U$를 $\mathcal{C}$의 임의의 대상으로 하고
$f \in \mathcal{I}(U)$에 대하여 $T(K|_U, f)$가
$D(\mathcal{O}_U)$의 영대상이라는 뜻이다.
\end{definition}

\noindent
유도 완비 대상들로 이루어진 충만한 부분범주
$D_{comp}(\mathcal{O}) = D_{comp}(\mathcal{O}, \mathcal{I}) \subset
D(\mathcal{O})$는 포화 삼각 부분범주임이 명백하다. 유도 범주의 정의
\ref{derived-definition-triangulated-subcategory}와
\ref{derived-definition-saturated}를 보라. 이 부분범주는 $D(\mathcal{O})$
안의 곱과 호모토피 극한 아래 보존된다. 그러나 일반적으로 가산 직합 아래서는
보존되지 않는다.

\begin{lemma}
\label{algebraization-lemma-derived-complete-internal-hom}
$(\mathcal{C}, \mathcal{O})$를 환 달린 사이트라 하자.
$\mathcal{I} \subset \mathcal{O}$를 아이디얼 층이라 하자.
$K \in D(\mathcal{O})$이고 $L \in D_{comp}(\mathcal{O})$이면
$R\SheafHom_\mathcal{O}(K, L) \in D_{comp}(\mathcal{O})$.
\end{lemma}

\begin{proof}
$U$를 $\mathcal{C}$의 대상이라 하고 $f \in \mathcal{I}(U)$라 하자. 다음을
상기하라.
$$
\Hom_{D(\mathcal{O}_U)}(\mathcal{O}_{U, f}, R\SheafHom_\mathcal{O}(K, L)|_U)
=
\Hom_{D(\mathcal{O}_U)}(
K|_U \otimes_{\mathcal{O}_U}^\mathbf{L} \mathcal{O}_{U, f}, L|_U)
$$
이는 사이트 위의 코호몰로지의 보조정리
\ref{sites-cohomology-lemma-internal-hom}에서 따른다. 우변은 보조정리
\ref{algebraization-lemma-hom-from-Af}와 내부 Hom 및 실제 Hom 사이의 관계에 의해 영이다.
사이트 위의 코호몰로지의 보조정리
\ref{sites-cohomology-lemma-section-RHom-over-U}를 보라. 같은 소멸은 모든
$U'/U$에 대하여 성립한다. 따라서 대상
$R\SheafHom_{\mathcal{O}_U}(\mathcal{O}_{U, f},
R\SheafHom_\mathcal{O}(K, L)|_U)$은 $D(\mathcal{O}_U)$에 속하고, 그 $0$차
코호몰로지 층은 소멸한다(상기
인용). 다른 코호몰로지 층들에도 마찬가지이므로,
$R\SheafHom_{\mathcal{O}_U}(\mathcal{O}_{U, f},
R\SheafHom_\mathcal{O}(K, L)|_U)$는 $D(\mathcal{O}_U)$에서 영이다.
보조정리 \ref{algebraization-lemma-hom-from-Af}에 의해 결론을 얻는다.
\end{proof}

\begin{lemma}
\label{algebraization-lemma-restriction-derived-complete}
$\mathcal{C}$를 사이트라 하자. $\mathcal{O} \to \mathcal{O}'$을 환층의
준동형이라 하자. $\mathcal{I} \subset \mathcal{O}$를 아이디얼 층이라 하자.
$D_{comp}(\mathcal{O}, \mathcal{I})$의 스칼라 제한 함자
$D(\mathcal{O}') \to D(\mathcal{O})$ 아래 역상은
$D_{comp}(\mathcal{O}', \mathcal{I}\mathcal{O}')$이다.
\end{lemma}

\begin{proof}
보조정리 \ref{algebraization-lemma-ideal-of-elements-complete-wrt}을 사용하면
$K' \in D(\mathcal{O}')$이
$D_{comp}(\mathcal{O}', \mathcal{I}\mathcal{O}')$에 속할 필요충분조건은
$T(K'|_U, f)$가 모든 국소 단면 $f \in \mathcal{I}(U)$에 대하여 영인 것이다.
$T(K'|_U, f)$의 코호몰로지 층들은 아벨 층의 범주에서 계산되므로, $f$를
$\mathcal{O}$의 단면으로 보든 $f$의 상을 $\mathcal{O}'$의 단면으로 보든
상관없다. 이 사실과 유도 완비 대상의 정의에서 보조정리가 즉시 따른다.
\end{proof}

\begin{lemma}
\label{algebraization-lemma-pushforward-derived-complete}
$f : (\Sh(\mathcal{D}), \mathcal{O}') \to (\Sh(\mathcal{C}), \mathcal{O})$를
환 달린 토포스의 사상이라 하자. $\mathcal{I} \subset \mathcal{O}$와
$\mathcal{I}' \subset \mathcal{O}'$을 아이디얼 층이라 하고, $f^\sharp$가
$f^{-1}\mathcal{I}$를 $\mathcal{I}'$ 안으로 보낸다고 하자. 그러면 $Rf_*$는
$D_{comp}(\mathcal{O}', \mathcal{I}')$을
$D_{comp}(\mathcal{O}, \mathcal{I})$ 안으로 보낸다.
\end{lemma}

\begin{proof}
$f$가 연속 함자 $\mathcal{C} \to \mathcal{D}$에 대응하는 환 달린 사이트의
사상으로 주어진다고 가정할 수 있다(사이트 위의 가군의 보조정리
\ref{sites-modules-lemma-morphism-ringed-topoi-comes-from-morphism-ringed-sites}
).
$U$를 $\mathcal{C}$의 대상이라 하고 $g$를 $\mathcal{I}$의 단면이라 하되
$U$ 위에서 잡자. 다음을 보이면 된다.
$\Hom_{D(\mathcal{O}_U)}(\mathcal{O}_{U, g}, Rf_*K|_U) = 0$
여기서는 $K$가 $\mathcal{I}'$에 관해 유도 완비라고 가정한다. 실제로 사이트
위의 코호몰로지의 보조정리
\ref{sites-cohomology-lemma-section-RHom-over-U}
에 의해 이를 $U$ 위의 모든 대상과 $K$의 모든 옮김에 적용하면
$R\SheafHom_{\mathcal{O}_U}(\mathcal{O}_{U, g}, Rf_*K|_U)$가 영임을 얻는다.
이는 $T(Rf_*K|_U, g)$가 영임을 함의하며(보조정리
\ref{algebraization-lemma-hom-from-Af}), 이것이 보일 바이다(정의
\ref{algebraization-definition-derived-complete}). $V$를 $\mathcal{D}$ 안에서 $U$의 상이라
하자. 그러면
$$
\Hom_{D(\mathcal{O}_U)}(\mathcal{O}_{U, g}, Rf_*K|_U) =
\Hom_{D(\mathcal{O}'_V)}(\mathcal{O}'_{V, g'}, K|_V) = 0
$$
이다. 여기서 $g' = f^\sharp(g) \in \mathcal{I}'(V)$이다. 둘째 등식은 $K$가
유도 완비이므로 성립하고, 첫째 등식은 $\mathcal{O}_{U, g}$의 유도 당김이
$\mathcal{O}'_{V, g'}$라는 사실과 사이트 위의 코호몰로지의 보조정리
\ref{sites-cohomology-lemma-adjoint}에서 성립한다.
\end{proof}

\noindent
다음 보조정리는 유도 완비가 존재하는 가장 간단한 경우이다.

\begin{lemma}
\label{algebraization-lemma-derived-completion}
$(\mathcal{C}, \mathcal{O})$를 환 달린 사이트라 하자. $f_1, \ldots, f_r$를
$\mathcal{O}$의 대역 단면이라 하자. $\mathcal{I} \subset \mathcal{O}$를
$f_1, \ldots, f_r$가 생성하는 아이디얼 층이라 하자. 그러면 포함 함자
$D_{comp}(\mathcal{O}) \to D(\mathcal{O})$는 왼쪽 수반을 갖는다. 즉,
임의의 대상 $K$가 $D(\mathcal{O})$ 안에서 주어지면 사상 $K \to K^\wedge$가
존재하며 $K^\wedge$는 $D_{comp}(\mathcal{O})$에 속하고, 사상
$$
\Hom_{D(\mathcal{O})}(K^\wedge, E) \longrightarrow \Hom_{D(\mathcal{O})}(K, E)
$$
은 $E$가 $D_{comp}(\mathcal{O})$에 속할 때마다 전단사이다. 실제로
$$
K^\wedge =
R\SheafHom_\mathcal{O}
(\mathcal{O} \to \prod\nolimits_{i_0} \mathcal{O}_{f_{i_0}} \to
\prod\nolimits_{i_0 < i_1} \mathcal{O}_{f_{i_0}f_{i_1}} \to
\ldots \to \mathcal{O}_{f_1\ldots f_r}, K)
$$
가 $K$에 대하여 함자적으로 성립한다.
\end{lemma}

\begin{proof}
$K^\wedge$를 보조정리의 마지막 표시식으로 정의하자. 복합체의 사상
$$
(\mathcal{O} \to \prod\nolimits_{i_0} \mathcal{O}_{f_{i_0}} \to
\prod\nolimits_{i_0 < i_1} \mathcal{O}_{f_{i_0}f_{i_1}} \to
\ldots \to \mathcal{O}_{f_1\ldots f_r}) \longrightarrow \mathcal{O}
$$
이 있고, 이는 사상 $K \to K^\wedge$를 유도한다. $K^\wedge$가 유도 완비이며,
$K \to K^\wedge$가 $K$가 유도 완비일 때 동형사상임을 보이면 충분하다.

\medskip\noindent
 $f$를 $\mathcal{O}$의 대역 단면이라 하자. 보조정리
\ref{algebraization-lemma-map-twice-localize}에 의해 대상
$R\SheafHom_\mathcal{O}(\mathcal{O}_f, K^\wedge)$
은 다음과 같다.
$$
R\SheafHom_\mathcal{O}(
(\mathcal{O}_f \to \prod\nolimits_{i_0} \mathcal{O}_{ff_{i_0}} \to
\prod\nolimits_{i_0 < i_1} \mathcal{O}_{ff_{i_0}f_{i_1}} \to
\ldots \to \mathcal{O}_{ff_1\ldots f_r}), K)
$$
$f = f_i$가 어떤 $i$에 대하여 성립하면 $f_1, \ldots, f_r$는
$\mathcal{O}_f$ 안의 단위 아이디얼을 생성한다. 따라서 확장 교대 Čech 복합체
$$
\mathcal{O}_f \to \prod\nolimits_{i_0} \mathcal{O}_{ff_{i_0}} \to
\prod\nolimits_{i_0 < i_1} \mathcal{O}_{ff_{i_0}f_{i_1}} \to
\ldots \to \mathcal{O}_{ff_1\ldots f_r}
$$
는 영이다(심지어 영과 호모토픽이다). 이로부터 $K^\wedge$가 유도 완비임을
알 수 있다.

\medskip\noindent
$K$가 유도 완비이면 $R\SheafHom_\mathcal{O}(\mathcal{O}_f, K)$는 모든
$f = f_{i_0} \ldots f_{i_p}$, $p \geq 0$에 대하여 영이다. 따라서
$K \to K^\wedge$는 $D(\mathcal{O})$에서 동형사상이다.
\end{proof}

\noindent
다음으로 유도 완비가 왜 완비인지 설명한다.

\begin{lemma}
\label{algebraization-lemma-derived-completion-koszul}
$(\mathcal{C}, \mathcal{O})$를 환 달린 사이트라 하자. $f_1, \ldots, f_r$를
$\mathcal{O}$의 대역 단면이라 하자. $\mathcal{I} \subset \mathcal{O}$를
$f_1, \ldots, f_r$가 생성하는 아이디얼 층이라 하자.
$K \in D(\mathcal{O})$라 하자. 유도 완비 $K^\wedge$는 보조정리
\ref{algebraization-lemma-derived-completion}의 것이며 다음 식으로 주어진다.
$$
K^\wedge = R\lim K \otimes^\mathbf{L}_\mathcal{O} K_n
$$
여기서 $K_n = K(\mathcal{O}, f_1^n, \ldots, f_r^n)$은
$f_1^n, \ldots, f_r^n$에 대한 $\mathcal{O}$ 위의 Koszul 복합체이다.
\end{lemma}

\begin{proof}
대수학 심화의 보조정리
\ref{more-algebra-lemma-extended-alternating-Cech-is-colimit-koszul}에서
확장 교대 Čech 복합체
$$
\mathcal{O} \to \prod\nolimits_{i_0} \mathcal{O}_{f_{i_0}} \to
\prod\nolimits_{i_0 < i_1} \mathcal{O}_{f_{i_0}f_{i_1}} \to
\ldots \to \mathcal{O}_{f_1\ldots f_r}
$$
가 Koszul 복합체 $K^n = K(\mathcal{O}, f_1^n, \ldots, f_r^n)$들의
여극한임을 보았다. 이 복합체들은 차수 $0, \ldots, r$에 놓여 있다.
$K^n$은 유한 자유 $\mathcal{O}$-가군으로 이루어진 유한 연쇄 복합체이고,
그 쌍대는 $\SheafHom_\mathcal{O}(K^n, \mathcal{O}) = K_n$이다. 여기서
$K_n$은 통상대로 차수 $-r, \ldots, 0$에 놓인 Koszul 공사슬 복합체이다.
보조정리 \ref{algebraization-lemma-derived-completion}에 의해 함자 $E \mapsto E^\wedge$는
확장 교대 Čech 복합체에서 $R\SheafHom$을 취하여 얻으며, 그 둘째 인자는
$E$이다.
$$
E^\wedge = R\SheafHom(\colim K^n, E)
$$
이는 $R\lim (E \otimes_\mathcal{O}^\mathbf{L} K_n)$과 같다. 이 등식은
사이트 위의 코호몰로지의 보조정리
\ref{sites-cohomology-lemma-colim-and-lim-of-duals}에서 따른다.
\end{proof}

\begin{lemma}
\label{algebraization-lemma-all-rings}
다음을 구성하는 방법이 존재한다.
\begin{enumerate}
\item 각 쌍 $(A, I)$은 환 $A$와 유한 생성 아이디얼 $I \subset A$로
이루어져 있으며, 이에 복합체 $K(A, I)$를 대응시킨다. 이 복합체는
$A$-가군들로 이루어진다,
\item 사상 $K(A, I) \to A$로서 $A$-가군 복합체들의 사상인 것,
\item 각 환 준동형 $A \to B$와 유한 생성 아이디얼 $I \subset A$에 대하여
복합체의 사상 $K(A, I) \to K(B, IB)$,
\end{enumerate}
이들은 다음을 만족한다.
\begin{enumerate}
\item[(a)] $A \to B$와 유한 생성 아이디얼 $I \subset A$에 대하여 도식
$$
\xymatrix{
K(A, I) \ar[r] \ar[d] & A \ar[d] \\
K(B, IB) \ar[r] & B
}
$$
은 가환한다,
\item[(b)] $A \to B \to C$와 유한 생성 아이디얼 $I \subset A$에 대하여
사상들의 합성 $K(A, I) \to K(B, IB) \to K(C, IC)$은 사상
$K(A, I) \to K(C, IC)$이다,
\item[(c)] $A \to B$와 유한 생성 아이디얼 $I \subset A$에 대하여 유도된
사상 $K(A, I) \otimes_A^\mathbf{L} B \to K(B, IB)$은 $D(B)$에서
동형사상이다, 그리고
\item[(d)] $I = (f_1, \ldots, f_r) \subset A$이면 가환도식
$$
\xymatrix{
(A \to \prod\nolimits_{i_0} A_{f_{i_0}} \to
\prod\nolimits_{i_0 < i_1} A_{f_{i_0}f_{i_1}} \to
\ldots \to A_{f_1\ldots f_r}) \ar[r] \ar[d] &  K(A, I) \ar[d] \\
A \ar[r]^1 & A
}
$$
이 $D(A)$ 안에 존재하며, 수평 화살표들은 동형사상이다.
\end{enumerate}
\end{lemma}

\begin{proof}
$S$를 다음 꼴의 환 $A_0$들의 집합이라 하자.
$A_0 = \mathbf{Z}[x_1, \ldots, x_n]/J$.
모든 유한형 $\mathbf{Z}$-대수는 $S$의 한 원소와 동형이다. $\mathcal{A}_0$를
다음 범주라 하자. 그 대상은 쌍 $(A_0, I_0)$이며, 여기서 $A_0 \in S$이고
$I_0 \subset A_0$는 아이디얼이다. 그 사상
$(A_0, I_0) \to (B_0, J_0)$은 환 준동형 $\varphi : A_0 \to B_0$로서
$J_0 = \varphi(I_0)B_0$를 만족하는 것이다.

\medskip\noindent
$K(A_0, I_0) \to A_0$를 $\mathcal{A}_0$의 대상에 대해 함자적으로 구성하되
성질 (a), (b), (c), (d)를 만족하게 할 수 있다고 가정하자. 그러면
$$
K(A, I) = \colim_{\varphi : (A_0, I_0) \to (A, I)} K(A_0, I_0)
$$
로 놓는다. 여기서 여극한은 환 준동형 $\varphi : A_0 \to A$로서
$\varphi(I_0)A = I$를 만족하는 것들에 대하여 취하고, $(A_0, I_0)$는 $\mathcal{A}_0$에
속한다. $(A_0, I_0) \to (A, I)$과
$(A_0', I_0') \to (A, I)$ 사이의 사상은
$(A_0, I_0) \to (A_0', I_0')$ 꼴의 사상으로 주어진다. 이는 $\mathcal{A}_0$
안에 있고 $A$로 가는 사상들과 가환한다. 이들
$(A_0, I_0) \to (A, I)$의 범주는 여과 범주이다(세부사항은 생략한다).
더욱이 $\colim_{\varphi : (A_0, I_0) \to (A, I)} A_0 = A$이므로
$K(A, I)$은 $A$-가군의 복합체이다. 마지막으로 $\varphi : A \to B$와
$I \subset A$가 주어졌다고 하자. 여극한에 등장하는 각
$(A_0, I_0) \to (A, I)$에 대하여 합성
$(A_0, I_0) \to (B, IB)$은 $(B, IB)$에 대한 여극한에 등장한다. 이로부터
여극한 사이의 사상을 얻는다. 성질 (a), (b), (c), (d)는 이 사실과 복합체
$K(A_0, I_0)$의 대응하는 성질들에서 곧바로 따른다.

\medskip\noindent
$\mathcal{C}_0 = \mathcal{A}_0^{opp}$에 혼돈 위상을 주자. $\mathcal{C}_0$에
환층 $\mathcal{O} : (A, I) \mapsto A$를 주자. 아이디얼 $I$들은 함께
아이디얼 층 $\mathcal{I} \subset \mathcal{O}$을 이룬다. 단사 분해
$\mathcal{O} \to \mathcal{J}^\bullet$을 택하자. 대상
$$
\mathcal{F}^\bullet = \bigcup\nolimits_n \mathcal{J}^\bullet[\mathcal{I}^n]
$$
을 생각하자. $U = (A, I) \in \Ob(\mathcal{C}_0)$라 하자.
$\mathcal{C}_0$의 위상은 혼돈 위상이므로 값 $\mathcal{J}^\bullet(U)$는
단사 $A$-가군들로 이루어진 $A$의 분해이다. 따라서 값
$\mathcal{F}^\bullet(U)$는 $D(A)$의 대상이며, $R\Gamma_I(A)$의
$D(A)$ 안에서의 상을 나타낸다. 쌍대화 복합체의 절
\ref{dualizing-section-local-cohomology}을 보라. $\mathcal{O}$-가군의
복합체 $\mathcal{K}^\bullet$과 가환도식
$$
\xymatrix{
\mathcal{O} \ar[r] & \mathcal{J}^\bullet \\
\mathcal{K}^\bullet \ar[r] \ar[u] & \mathcal{F}^\bullet \ar[u]
}
$$
을 택하되 수평 화살표들이 준동형사상이 되게 하자. 이는 유도 범주
$D(\mathcal{O})$의 구성으로 가능하다. $K(A, I) = \mathcal{K}^\bullet(U)$로
놓되 여기서 $U = (A, I)$이다. 성질 (a), (b)는 명백하고 성질
(c), (d)는 쌍대화 복합체의 보조정리
\ref{dualizing-lemma-compute-local-cohomology-noetherian}와
\ref{dualizing-lemma-local-cohomology-change-rings}에서 따른다.
\end{proof}

\begin{lemma}
\label{algebraization-lemma-global-extended-cech-complex}
$(\mathcal{C}, \mathcal{O})$를 환 달린 사이트라 하자.
$\mathcal{I} \subset \mathcal{O}$를 유한형 아이디얼 층이라 하자.
사상 $K \to \mathcal{O}$가 $D(\mathcal{O})$ 안에 존재하며 다음 성질을
만족한다. $U \in \Ob(\mathcal{C})$이고 $\mathcal{I}|_U$가
$f_1, \ldots, f_r \in \mathcal{I}(U)$에 의해 생성될 때마다 동형사상
$$
(\mathcal{O}_U \to \prod\nolimits_{i_0} \mathcal{O}_{U, f_{i_0}} \to
\prod\nolimits_{i_0 < i_1} \mathcal{O}_{U, f_{i_0}f_{i_1}} \to
\ldots \to \mathcal{O}_{U, f_1\ldots f_r}) \longrightarrow K|_U
$$
이 존재하고, 이는 $\mathcal{O}_U$로 가는 사상들과 양립한다.
\end{lemma}

\begin{proof}
$\mathcal{C}' \subset \mathcal{C}$를 다음 대상 $U$들로 이루어진 충만
부분범주라 하자. 즉 $\mathcal{I}|_U$가 유한 개의 단면으로 생성되는
대상들이다. 그러면 $\mathcal{C}' \to \mathcal{C}$는 특수 쌍연속 함자이다
(사이트의 정의 \ref{sites-definition-special-cocontinuous-functor}).
따라서 $\mathcal{C}'$에서 작업하는 것으로 충분하다. 사이트의 보조정리
\ref{sites-lemma-equivalence}를 보라. 다시 말해 각 대상 $U$가
$\mathcal{C}$에 속할 때 유한 생성 아이디얼 $I \subset \mathcal{I}(U)$가 존재하여
$\mathcal{I}|_U = \Im(I \otimes \mathcal{O}_U \to \mathcal{O}_U)$라고
가정해도 된다. 이때 $I$가 $\mathcal{I}|_U$를 생성한다고 말하겠다.
주의: $\mathcal{I}(U)$가 $\mathcal{O}(U)$ 안의 유한 생성 아이디얼인지는
알 수 없다.

\medskip\noindent
$U$를 대상이라 하고, $I \subset \mathcal{O}(U)$를 $\mathcal{I}|_U$를
생성하는 유한 생성 아이디얼이라 하자. 범주 $\mathcal{C}/U$ 위의 전층
복합체
$$
K_{U, I}^\bullet : U'/U \longmapsto K(\mathcal{O}(U'), I\mathcal{O}(U'))
$$
를 생각하자. 여기서 $K(-, -)$는 보조정리 \ref{algebraization-lemma-all-rings}의 것이다.
이 복합체의 층화는 $I$의 선택과 무관하다고 주장한다. 실제로
$I' \subset \mathcal{O}(U)$가 역시 $\mathcal{I}|_U$를 생성하는 유한 생성
아이디얼이면, 덮개 $\{U_j \to U\}$가 존재하여
$I\mathcal{O}(U_j) = I'\mathcal{O}(U_j)$이다. (힌트: $I$와 $I'$가 모두
유한 생성이고 $\mathcal{I}|_U$를 생성하기 때문에 성립한다.) 따라서
$K_{U, I}^\bullet$과 $K_{U, I'}^\bullet$은 어느 $U_j$ 위에 놓인 대상에
대해서도 {\it 동일하다}. 이제 층화에 대한 주장이 따른다.
그 공통값을 $K_U^\bullet$로 나타내자.

\medskip\noindent
$I$의 선택과 무관하다는 사실은 다음도 보인다.
$K_U^\bullet|_{\mathcal{C}/U'} = K_{U'}^\bullet$이며, 이는 사상
$U' \to U$와 이에 따른 국소화 사상
$\mathcal{C}/U' \to \mathcal{C}/U$가 주어질 때마다 성립한다. 따라서 복합체들
$K_U^\bullet$은 이어 붙어서 잘 정의된 하나의 복합체 $K^\bullet$을 주며,
이는 $\mathcal{O}$-가군들로 이루어진다. 사상 $K^\bullet \to \mathcal{O}$의 존재와 보조정리의
준동형사상은 보조정리 \ref{algebraization-lemma-all-rings}에 있는 복합체 $K(-, -)$의
대응하는 성질들에서 곧바로 따른다.
\end{proof}

\begin{proposition}
\label{algebraization-proposition-derived-completion}
$(\mathcal{C}, \mathcal{O})$를 환 달린 사이트라 하자.
$\mathcal{I} \subset \mathcal{O}$를 유한형 아이디얼 층이라 하자. 포함 함자
$D_{comp}(\mathcal{O}) \to D(\mathcal{O})$의 왼쪽 수반이 존재한다.
\end{proposition}

\begin{proof}
사상 $K \to \mathcal{O}$를 $D(\mathcal{O})$ 안에서 보조정리
\ref{algebraization-lemma-global-extended-cech-complex}에서 구성한 것으로 택하자.
$E \in D(\mathcal{O})$라 하자. 그러면 $E^\wedge = R\SheafHom(K, E)$와
사상 $E \to E^\wedge$가 원하는 것을 준다. 실제로 사이트
$\mathcal{C}$ 위에서 국소적으로 보조정리 \ref{algebraization-lemma-derived-completion}의
수반을 되찾는다. 이로부터 $E^\wedge$는 항상 유도 완비이고, 사상
$E \to E^\wedge$는 $E$가 유도 완비이면 동형사상임을 알 수 있다.
\end{proof}

\begin{remark}[완비와의 비교]
\label{algebraization-remark-compare-with-completion}
$(\mathcal{C}, \mathcal{O})$를 환 달린 사이트라 하자.
$\mathcal{I} \subset \mathcal{O}$를 유한형 아이디얼 층이라 하자.
$K \mapsto K^\wedge$를 명제 \ref{algebraization-proposition-derived-completion}의 유도 완비
함자라 하자. 임의의 $n \geq 1$에 대하여 대상
$K \otimes_\mathcal{O}^\mathbf{L} \mathcal{O}/\mathcal{I}^n$은
$\mathcal{I}$의 국소 단면들의 거듭제곱들에 의해 소멸되므로 유도 완비이다.
따라서 표준 사상의 표준 분해
$$
K \to K^\wedge \to K \otimes_\mathcal{O}^\mathbf{L} \mathcal{O}/\mathcal{I}^n
$$
가 존재한다. 여기서 표준 사상은
$K \to K \otimes_\mathcal{O}^\mathbf{L} \mathcal{O}/\mathcal{I}^n$이다.
이 사상들은 $n$의 변화에 대해 양립하며, 비교 사상
$$
K^\wedge
\longrightarrow
R\lim \left(K \otimes_\mathcal{O}^\mathbf{L} \mathcal{O}/\mathcal{I}^n\right)
$$
을 얻는다. 오른쪽 항은 어떤 종류의 완비로 더 쉽게 알아볼 수 있다.
일반적으로 이 비교 사상은 동형사상이 아니다.
\end{remark}

\begin{remark}[국소화와 유도 완비]
\label{algebraization-remark-localization-and-completion}
$(\mathcal{C}, \mathcal{O})$를 환 달린 사이트라 하자.
$\mathcal{I} \subset \mathcal{O}$를 유한형 아이디얼 층이라 하자.
$K \mapsto K^\wedge$를 명제 \ref{algebraization-proposition-derived-completion}의 유도 완비
함자라 하자. 명제의 증명에 있는 구성으로부터 $K^\wedge|_U$는
$K|_U$의 유도 완비이며, 이는 모든 $U \in \Ob(\mathcal{C})$에 대해
성립함이 따른다. 그러나 다음과 같이 증명할 수도 있다. 유도 완비 대상의 정의로부터
$K^\wedge|_U$가 유도 완비임이 따른다. 따라서 표준 사상
$a : (K|_U)^\wedge \to K^\wedge|_U$를 얻는다. 한편 $E$가
$D(\mathcal{O}_U)$의 유도 완비 대상이면, 보조정리
\ref{algebraization-lemma-pushforward-derived-complete}에 의해 $Rj_*E$는
$D(\mathcal{O})$의 유도 완비 대상이다. 여기서 $j$는 국소화 사상이다
(사이트 위의 가군의 절 \ref{sites-modules-section-localize}). 따라서 표준
사상 $b : K^\wedge \to Rj_*((K|_U)^\wedge)$도 얻는다. $b$의 수반이
$a$의 역임을 확인하는 형식적인 과정은 생략한다.
\end{remark}

\begin{remark}[완비 텐서곱]
\label{algebraization-remark-completed-tensor-product}
$(\mathcal{C}, \mathcal{O})$를 환 달린 사이트라 하자.
$\mathcal{I} \subset \mathcal{O}$를 유한형 아이디얼 층이라 하자.
$K \mapsto K^\wedge$로 명제 \ref{algebraization-proposition-derived-completion}의 수반을
나타내자. 이제
$$
K \otimes^\wedge_\mathcal{O} L = (K \otimes_\mathcal{O}^\mathbf{L} L)^\wedge
$$
로 놓는다. 이 {\it 완비 텐서곱}은 함자
$D_{comp}(\mathcal{O}) \times D_{comp}(\mathcal{O}) \to D_{comp}(\mathcal{O})$
를 정의하며 다음 등식을 만족한다.
$$
\Hom_{D_{comp}(\mathcal{O})}(K, R\SheafHom_\mathcal{O}(L, M))
=
\Hom_{D_{comp}(\mathcal{O})}(K \otimes_\mathcal{O}^\wedge L, M)
$$
$K, L, M \in D_{comp}(\mathcal{O})$에 대하여 성립한다. 보조정리
\ref{algebraization-lemma-derived-complete-internal-hom}에 의해
$R\SheafHom_\mathcal{O}(L, M) \in D_{comp}(\mathcal{O})$임에 유의하자.
\end{remark}

\begin{lemma}
\label{algebraization-lemma-map-identifies-koszul-and-cech-complexes}
$\mathcal{C}$를 사이트라 하자. $\varphi : \mathcal{O} \to \mathcal{O}'$가
환층들의 평탄 준동형사상이라고 가정하자. $f_1, \ldots, f_r$를
$\mathcal{O}$의 대역 단면들로서
$\mathcal{O}/(f_1, \ldots, f_r) \cong
\mathcal{O}'/(f_1, \ldots, f_r)\mathcal{O}'$를 만족하는 것들이라 하자.
그러면 확장 교대 Čech 복합체들의 사상
$$
\xymatrix{
\mathcal{O} \to
\prod_{i_0} \mathcal{O}_{f_{i_0}} \to
\prod_{i_0 < i_1} \mathcal{O}_{f_{i_0}f_{i_1}} \to \ldots \to
\mathcal{O}_{f_1\ldots f_r} \ar[d] \\
\mathcal{O}' \to
\prod_{i_0} \mathcal{O}'_{f_{i_0}} \to
\prod_{i_0 < i_1} \mathcal{O}'_{f_{i_0}f_{i_1}} \to \ldots \to
\mathcal{O}'_{f_1\ldots f_r}
}
$$
은 준동형사상이다.
\end{lemma}

\begin{proof}
둘째 복합체는 첫째 복합체를 $\mathcal{O}'$와 텐서곱한 것임을 관찰하자.
첫째 확장 교대 Čech 복합체는 Koszul 복합체
$K_n = K(\mathcal{O}, f_1^n, \ldots, f_r^n)$들의 여극한으로 쓸 수 있다.
대수학 심화의 보조정리
\ref{more-algebra-lemma-extended-alternating-Cech-is-colimit-koszul}를 보라.
따라서 $K_n \to K_n \otimes_\mathcal{O} \mathcal{O}'$가 준동형사상임을
보이면 충분하다. $\mathcal{O} \to \mathcal{O}'$가 평탄하므로
$H^i \to H^i \otimes_\mathcal{O} \mathcal{O}'$가 동형사상임을 보이면
충분하다. 여기서 $H^i$는 제$i$ 코호몰로지 층
$H^i = H^i(K_n)$이다. 이 층들은 $f_1^n, \ldots, f_r^n$에 의해 소멸된다.
대수학 심화의 보조정리 \ref{more-algebra-lemma-homotopy-koszul}를 보라.
따라서 이 층들은 $(f_1, \ldots, f_r)^m$에 의해 소멸되며, 여기서
$m \gg 0$는 적당히 택한다. 그러므로 사이트 위의 가군의 보조정리
\ref{sites-modules-lemma-neighbourhood-isomorphism}에 의해
$H^i \to H^i \otimes_\mathcal{O} \mathcal{O}'$는 동형사상이다.
\end{proof}

\begin{lemma}
\label{algebraization-lemma-restriction-derived-complete-equivalence}
$\mathcal{C}$를 사이트라 하자. $\mathcal{O} \to \mathcal{O}'$를 환층들의
준동형사상이라 하자. $\mathcal{I} \subset \mathcal{O}$를 유한형 아이디얼
층이라 하자. $\mathcal{O} \to \mathcal{O}'$가 평탄하고
$\mathcal{O}/\mathcal{I} \cong \mathcal{O}'/\mathcal{I}\mathcal{O}'$이면,
제한 함자 $D(\mathcal{O}') \to D(\mathcal{O})$는 동치
$D_{comp}(\mathcal{O}', \mathcal{I}\mathcal{O}') \to
D_{comp}(\mathcal{O}, \mathcal{I})$를 유도한다.
\end{lemma}

\begin{proof}
보조정리 \ref{algebraization-lemma-pushforward-derived-complete}에 의해 제한 함자
$r : D(\mathcal{O}') \to D(\mathcal{O})$는
$D_{comp}(\mathcal{O}', \mathcal{I}\mathcal{O}')$를
$D_{comp}(\mathcal{O}, \mathcal{I})$ 안으로 보낸다. 준역함자
$E \mapsto E'$를 구성하겠다.

\medskip\noindent
사상 $K \to \mathcal{O}$를 보조정리
\ref{algebraization-lemma-global-extended-cech-complex}에서 구성한 $D(\mathcal{O})$의
사상이라 하자. $K' = K \otimes_\mathcal{O}^\mathbf{L} \mathcal{O}'$로
$D(\mathcal{O}')$ 안에서 놓자. 그러면
$K' \to \mathcal{O}'$는 $D(\mathcal{O}')$ 안의 사상이며, 보조정리
\ref{algebraization-lemma-global-extended-cech-complex}의 결론을
$\mathcal{I}' = \mathcal{I}\mathcal{O}'$에 대해 만족한다. 보조정리
\ref{algebraization-lemma-map-identifies-koszul-and-cech-complexes}에 의해 사상
$K \to r(K')$는 준동형사상이다. 이제
$E \in D_{comp}(\mathcal{O}, \mathcal{I})$에 대해
$$
E' = R\SheafHom_\mathcal{O}(r(K'), E)
$$
로 놓고, $D(\mathcal{O}')$의 대상으로 본다. 이는 $\mathcal{O}'$-가군
구조를 사용하며 그 구조는 $K'$ 위의 것이다. $E$는 유도 완비이므로
$E = R\SheafHom_\mathcal{O}(K, E)$이다. 명제
\ref{algebraization-proposition-derived-completion}의 증명을 보라. 한편 사상
$K \to r(K')$는 동형사상이므로, % P07-ERR-0144
$E \to r(E')$가 $D(\mathcal{O})$ 안의 동형사상임을 알 수 있다. 증명을
끝내기 위해 다음을 보이면 된다. $E = r(M')$이고 $M'$가
$D_{comp}(\mathcal{O}', \mathcal{I}')$의 대상이면 $E' \cong M'$이다.
사상을 얻기 위해 다음을 사용한다.
$$
M' = R\SheafHom_{\mathcal{O}'}(\mathcal{O}', M') \to
R\SheafHom_\mathcal{O}(r(\mathcal{O}'), r(M')) \to
R\SheafHom_\mathcal{O}(r(K'), r(M')) = E'
$$
여기서 둘째 화살표는 사상 $K' \to \mathcal{O}'$를 사용한다. 이것이
동형사상임을 보이려면, 이 화살표에 $r$을 적용한 것이 위의 동형사상
$E \to r(E')$와 같음을 보이면 된다. 세부사항은 생략한다.
\end{proof}

\begin{lemma}
\label{algebraization-lemma-pushforward-derived-complete-adjoint}
$f : (\Sh(\mathcal{D}), \mathcal{O}') \to (\Sh(\mathcal{C}), \mathcal{O})$를
환 달린 토포스들의 사상이라 하자. $\mathcal{I} \subset \mathcal{O}$와
$\mathcal{I}' \subset \mathcal{O}'$를 유한형 아이디얼 층들이라 하고,
$f^\sharp$가 $f^{-1}\mathcal{I}$를 $\mathcal{I}'$ 안으로 보낸다고 하자.
그러면 $Rf_*$는 $D_{comp}(\mathcal{O}', \mathcal{I}')$를
$D_{comp}(\mathcal{O}, \mathcal{I})$ 안으로 보내며 왼쪽 수반
$Lf_{comp}^*$를 갖는다. 이 수반은 $Lf^*$를 취한 뒤 유도 완비를 취하는
함자이다.
\end{lemma}

\begin{proof}
첫째 명제는 보조정리 \ref{algebraization-lemma-pushforward-derived-complete}에서 보았다.
명제 \ref{algebraization-proposition-derived-completion}에 의해 유도 완비 함자
$D(\mathcal{O}') \to D_{comp}(\mathcal{O}', \mathcal{I}')$가 있으므로 둘째
명제도 의미가 있다. 이제 $K \in D_{comp}(\mathcal{O}, \mathcal{I})$와
$M \in D_{comp}(\mathcal{O}', \mathcal{I}')$를 택하자. 그러면
$$
\Hom(K, Rf_*M) = \Hom(Lf^*K, M) = \Hom(Lf_{comp}^*K, M)
$$
이며, 이는 유도 완비의 보편성에서 따른다.
\end{proof}

\begin{lemma}
\label{algebraization-lemma-pushforward-commutes-with-derived-completion}
\begin{reference}
\cite[보조정리 6.5.9 (2)]{BS}의 일반화이다. 준연접 가군과 (유도) 대수적
스택의 사상이라는 설정에서는 \cite[정리 6.5]{HL-P}와 비교하라.
\end{reference}
$f : (\Sh(\mathcal{D}), \mathcal{O}') \to (\Sh(\mathcal{C}), \mathcal{O})$를
환 달린 토포스들의 사상이라 하자. $\mathcal{I} \subset \mathcal{O}$를
유한형 아이디얼 층이라 하자. $\mathcal{I}' \subset \mathcal{O}'$를
$f^\sharp(f^{-1}\mathcal{I})$가 생성하는 아이디얼이라 하자. 그러면
$Rf_*$는 유도 완비와 교환한다. 즉,
$Rf_*(K^\wedge) = (Rf_*K)^\wedge$.
\end{lemma}

\begin{proof}
명제 \ref{algebraization-proposition-derived-completion}에 의해 유도 완비 함자들이 존재한다.
보조정리 \ref{algebraization-lemma-pushforward-derived-complete}에 의해 대상
$Rf_*(K^\wedge)$는 유도 완비이다. 따라서 유도 완비의 보편성으로부터
표준 사상 $(Rf_*K)^\wedge \to Rf_*(K^\wedge)$를 얻는다. 이 사상이
동형사상인지는 $\mathcal{C}$ 위에서 국소적으로 확인할 수 있다. 유도 완비는
국소화와 교환하므로(비고 \ref{algebraization-remark-localization-and-completion}),
$\mathcal{I}$가 대역 단면 $f_1, \ldots, f_r$에 의해 생성된다고 가정해도
된다. 그러면 $\mathcal{I}'$는 $g_i = f^\sharp(f_i)$에 의해 생성된다.
보조정리 \ref{algebraization-lemma-derived-completion-koszul}에 의해 다음을 증명하면 된다.
$$
R\lim \left(
Rf_*K \otimes^\mathbf{L}_\mathcal{O} K(\mathcal{O}, f_1^n, \ldots, f_r^n)
\right)
=
Rf_*\left(
R\lim
K \otimes^\mathbf{L}_{\mathcal{O}'} K(\mathcal{O}', g_1^n, \ldots, g_r^n)
\right)
$$
$Rf_*$가 $R\lim$와 교환하므로(사이트 위의 코호몰로지의 보조정리
\ref{sites-cohomology-lemma-Rf-commutes-with-Rlim}), 다음을 증명하면 충분하다.
$$
Rf_*K \otimes^\mathbf{L}_\mathcal{O} K(\mathcal{O}, f_1^n, \ldots, f_r^n) =
Rf_*\left(
K \otimes^\mathbf{L}_{\mathcal{O}'} K(\mathcal{O}', g_1^n, \ldots, g_r^n)
\right)
$$
이는 사영 공식(사이트 위의 코호몰로지의 보조정리
\ref{sites-cohomology-lemma-projection-formula})과 다음 사실에서 따른다.
$Lf^*K(\mathcal{O}, f_1^n, \ldots, f_r^n) =
K(\mathcal{O}', g_1^n, \ldots, g_r^n)$.
\end{proof}

\begin{lemma}
\label{algebraization-lemma-formal-functions-general}
$A$를 환이라 하고 $I \subset A$를 유한 생성 아이디얼이라 하자.
$\mathcal{C}$를 사이트라 하고 $\mathcal{O}$를 $A$-대수층이라 하자.
$\mathcal{F}$를 $\mathcal{O}$-가군층이라 하자. 그러면
$$
R\Gamma(\mathcal{C}, \mathcal{F})^\wedge =
R\Gamma(\mathcal{C}, \mathcal{F}^\wedge)
$$
가 $D(A)$ 안에서 성립한다. 여기서 $\mathcal{F}^\wedge$는
$\mathcal{F}$의 $I\mathcal{O}$에 대한 유도 완비이고, 왼쪽 항에는
$I$에 대한 유도 완비를 취한다. % P07-ERR-0145
이로부터 두 스펙트럼 열
$$
E_2^{i, j} = H^i(H^j(\mathcal{C}, \mathcal{F})^\wedge)
\quad\text{그리고}\quad
E_2^{p, q} = H^p(\mathcal{C}, H^q(\mathcal{F}^\wedge))
$$
을 얻으며, 둘 다
$H^*(R\Gamma(\mathcal{C}, \mathcal{F})^\wedge) =
H^*(\mathcal{C}, \mathcal{F}^\wedge)$
로 수렴한다.
\end{lemma}

\begin{proof}
보조정리 \ref{algebraization-lemma-pushforward-commutes-with-derived-completion}를 환 달린
토포스들의 사상 $(\mathcal{C}, \mathcal{O}) \to (pt, A)$에 적용하고
코호몰로지를 취하면 첫째 명제를 얻는다. 둘째 스펙트럼 열은 유도 범주의
보조정리 \ref{derived-lemma-two-ss-complex-functor}에 나오는 둘째 스펙트럼
열이다. 첫째 스펙트럼 열은 대수학 심화의 예
\ref{more-algebra-example-derived-completion-spectral-sequence}에 나오는
스펙트럼 열을 $R\Gamma(\mathcal{C}, \mathcal{F})^\wedge$에 적용한 것이다.
\end{proof}

\begin{remark}
\label{algebraization-remark-local-calculation-derived-completion}
$(\mathcal{C}, \mathcal{O})$를 환 달린 사이트라 하자.
$\mathcal{I} \subset \mathcal{O}$를 유한형 아이디얼 층이라 하자.
$K \mapsto K^\wedge$를 명제 \ref{algebraization-proposition-derived-completion}의 유도
완비라 하자. $U \in \Ob(\mathcal{C})$를 대상으로서 $\mathcal{I}$가
아이디얼 층으로서 $f_1, \ldots, f_r \in \mathcal{I}(U)$에 의해 생성되는
것이라 하자. $A = \mathcal{O}(U)$와
$I = (f_1, \ldots, f_r) \subset A$로 놓자. 주의:
$I = \mathcal{I}(U)$가 성립하지 않을 수도 있다. 그러면
$$
R\Gamma(U, K^\wedge) = R\Gamma(U, K)^\wedge
$$
이다. 여기서 오른쪽 항은 대상 $R\Gamma(U, K)$를, 이를 $D(A)$의 대상으로
보아, $I$에 대해 유도 완비한 것이다. 이는 유도 완비가 국소화와 교환하고(비고
\ref{algebraization-remark-localization-and-completion}), 보조정리
\ref{algebraization-lemma-formal-functions-general}가 성립하기 때문이다.
\end{remark}








\section{형식함수 정리}
\label{algebraization-section-formal-functions}

\noindent
이 절에서는 설명의 흐름을 잠시 멈추고 스킴 위 준연접 가군의 설정에서
유도 완비를 조금 논의한다. 이를 사용하여 형식함수 정리의 다소 다른
증명을 제시한다. 관련 문헌은 비고 \ref{algebraization-remark-references}에서 안내한다.

\medskip\noindent
보조정리 \ref{algebraization-lemma-pushforward-commutes-with-derived-completion}는 형식함수
정리의 매우 형식적인 유도 버전이다(스킴의 코호몰로지의 정리
\ref{coherent-theorem-formal-functions}). 이를 더 명시적으로 쓰기 위해
$f : X \to S$가 스킴의 사상이고,
$\mathcal{I} \subset \mathcal{O}_S$가 유한형 준연접 아이디얼 층이며,
$\mathcal{F}$가 $X$ 위의 준연접 층이라고 하자. 그러면 이 보조정리는
\begin{equation}
\label{algebraization-equation-formal-functions}
Rf_*(\mathcal{F}^\wedge) = (Rf_*\mathcal{F})^\wedge
\end{equation}
를 말한다. 여기서 $\mathcal{F}^\wedge$는 $\mathcal{F}$의
$f^{-1}\mathcal{I} \cdot \mathcal{O}_X$에 대한 유도 완비이고, 오른쪽 항은
$Rf_*\mathcal{F}$의 $\mathcal{I}$에 대한 유도 완비이다. 이것이
형식함수 정리를 되돌려 준다는 것을 보이려면 약간의 작업이 필요하다.

\begin{lemma}
\label{algebraization-lemma-sections-derived-completion-pseudo-coherent}
$X$를 국소 Noether 스킴이라 하자. $\mathcal{I} \subset \mathcal{O}_X$를
준연접 아이디얼 층이라 하자. $K$를 $D(\mathcal{O}_X)$의 유사연접
대상이라 하고 그 유도 완비를 $K^\wedge$라 하자. 그러면
$$
H^p(U, K^\wedge) = \lim H^p(U, K)/I^nH^p(U, K) =
H^p(U, K)^\wedge
$$
가 임의의 아핀 열린집합 $U \subset X$에 대해 성립한다. 여기서
$I = \mathcal{I}(U)$이고, 오른쪽 항에서는 $I$에 대한 유도 완비를 취한다.
\end{lemma}

\begin{proof}
$U = \Spec(A)$로 쓰자. 환 $A$는 Noether 환이므로 $I \subset A$는 유한
생성이다. 따라서
$$
R\Gamma(U, K^\wedge) = R\Gamma(U, K)^\wedge
$$
이다. 비고 \ref{algebraization-remark-local-calculation-derived-completion}을 보라. 이제
$R\Gamma(U, K)$는 $A$-가군들의 유사연접 복합체이다(스킴의 유도 범주의
보조정리 \ref{perfect-lemma-pseudo-coherent-affine}). 대수학 심화의 보조정리
\ref{more-algebra-lemma-derived-completion-pseudo-coherent}에 의해 제$p$
코호몰로지 가군, 즉 $R\Gamma(U, K^\wedge)$의 그것은 $I$-진 완비로서
$H^p(U, K)$에서 얻어지는 것과 같다. 이것으로 첫째 등식을 증명했다. 둘째 등식은 덜
중요하며, 방금 사용한 보조정리를 다시 한 번 적용하면 곧바로 따른다.
\end{proof}

\begin{lemma}
\label{algebraization-lemma-derived-completion-pseudo-coherent}
$X$를 국소 Noether 스킴이라 하자. $\mathcal{I} \subset \mathcal{O}_X$를
준연접 아이디얼 층이라 하자. $K$를 $D(\mathcal{O}_X)$의 대상이라 하자.
그러면
\begin{enumerate}
\item 유도 완비 $K^\wedge$는 다음과 같다.
$R\lim (K \otimes_{\mathcal{O}_X}^\mathbf{L} \mathcal{O}_X/\mathcal{I}^n)$.
\end{enumerate}
$K$가 $D(\mathcal{O}_X)$의 유사연접 대상이라고 하자. % P07-ERR-0146
그러면
\begin{enumerate}
\item[(2)] 코호몰로지 층 $H^q(K^\wedge)$는 다음과 같다.
$\lim H^q(K)/\mathcal{I}^nH^q(K)$.
\end{enumerate}
$\mathcal{F}$를 연접 $\mathcal{O}_X$-가군이라 하자\footnote{예를 들어
$H^q(K)$이며, 여기서 $K$는 우리가 다루는 국소 Noether 스킴 $X$ 위에서
유사연접이다.}. 그러면
\begin{enumerate}
\item[(3)] 유도 완비 $\mathcal{F}^\wedge$는 다음과 같다.
$\lim \mathcal{F}/\mathcal{I}^n\mathcal{F}$,
\item[(4)]
$\lim \mathcal{F}/\mathcal{I}^n \mathcal{F} =
R\lim \mathcal{F}/\mathcal{I}^n \mathcal{F}$,
\item[(5)] $H^p(U, \mathcal{F}^\wedge) = 0$가 $p \not = 0$일 때 모든
아핀 열린집합 $U \subset X$에 대하여 성립한다.
\end{enumerate}
\end{lemma}

\begin{proof}
(1)을 증명하자. 표준 사상
$$
K \longrightarrow
R\lim (K \otimes_{\mathcal{O}_X}^\mathbf{L} \mathcal{O}_X/\mathcal{I}^n),
$$
이 존재한다. 비고 \ref{algebraization-remark-compare-with-completion}을 보라. 유도 완비는
열린 부분스킴으로의 제한과 교환한다(비고
\ref{algebraization-remark-localization-and-completion}). $R\lim$의 구성도 열린
부분스킴으로의 제한과 교환한다. 따라서 이 사상이 동형사상인지는 국소적으로
확인하면 된다. 그러므로 $X = U = \Spec(A)$라고 가정해도 된다.
$I = (f_1, \ldots, f_r)$라 쓰자. Koszul 복합체를
$K_n = K(A, f_1^n, \ldots, f_r^n)$라 하자. 대수학 심화의 보조정리
\ref{more-algebra-lemma-sequence-Koszul-complexes}에서 전사영계
$\{K_n\}$와 $\{A/I^n\}$가 $D(A)$ 안에서 동형임을 보았다. 스킴의 유도
범주의 보조정리 \ref{perfect-lemma-affine-compare-bounded}에 나오는 동치
$D(A) = D_{\QCoh}(\mathcal{O}_X)$를 사용하면, 전사영계
$\{K(\mathcal{O}_X, f_1^n, \ldots, f_r^n)\}$와
$\{\mathcal{O}_X/\mathcal{I}^n\}$가 $D(\mathcal{O}_X)$ 안에서 동형임을
알 수 있다. 이것은 다음 식의 둘째 등식을 증명한다.
$$
K^\wedge = R\lim \left(
K \otimes_{\mathcal{O}_X}^\mathbf{L} K(\mathcal{O}_X, f_1^n, \ldots, f_r^n)
\right) =
R\lim (K \otimes_{\mathcal{O}_X}^\mathbf{L} \mathcal{O}_X/\mathcal{I}^n)
$$
첫째 등식은 보조정리 \ref{algebraization-lemma-derived-completion-koszul}이다.

\medskip\noindent
$K$가 유사연접이라고 가정하자. 아핀 열린집합 $U \subset X$에 대해
보조정리 \ref{algebraization-lemma-sections-derived-completion-pseudo-coherent}에서
$H^q(U, K^\wedge) = \lim H^q(U, K)/\mathcal{I}^n(U)H^q(U, K)$를 얻는다.
이것은 모든 $U$에 대해 성립하므로 층으로서
$H^q(K^\wedge) = \lim H^q(K)/\mathcal{I}^nH^q(K)$임을 알 수 있다.
이것으로 (2)를 증명했다.

\medskip\noindent
(3)은 (2)의 특수한 경우이다. (4)와 (5)는 스킴의 유도 범주의 보조정리
\ref{perfect-lemma-Rlim-quasi-coherent}에서 따른다.
\end{proof}

\begin{lemma}
\label{algebraization-lemma-formal-functions}
$A$를 Noether 환이라 하고 $I \subset A$를 아이디얼이라 하자. $X$를
$A$ 위의 Noether 스킴이라 하자. $\mathcal{F}$를 연접
$\mathcal{O}_X$-가군이라 하자. $H^p(X, \mathcal{F})$가 유한 $A$-가군이라고
모든 $p$에 대해 가정하자. 그러면 짧은 완전열
$$
0 \to R^1\lim H^{p - 1}(X, \mathcal{F}/I^n\mathcal{F}) \to
H^p(X, \mathcal{F})^\wedge \to \lim H^p(X, \mathcal{F}/I^n\mathcal{F}) \to 0
$$
이 존재하며, 이는 $A$-가군들의 열이다. 여기서
$H^p(X, \mathcal{F})^\wedge$는 통상적인 $I$-진
완비이다. $f$가 고유이면 $R^1\lim$ 항은 영이다. % P07-ERR-0147
\end{lemma}

\begin{proof}
보조정리 \ref{algebraization-lemma-formal-functions-general}의 두 스펙트럼 열을 생각하자.
첫째 것은 대수학 심화의 보조정리
\ref{more-algebra-lemma-derived-completion-pseudo-coherent}에 의해 퇴화한다.
$H^p(X, \mathcal{F})^\wedge$를 차수 $p$에서 얻는다. 여기서
$H^p(X, \mathcal{F})$가 유한 $A$-가군이라는 가정을 사용한다. 둘째 것은
다음 이유로 퇴화한다.
$$
\mathcal{F}^\wedge = \lim \mathcal{F}/I^n\mathcal{F} =
R\lim \mathcal{F}/I^n\mathcal{F}
$$
가 보조정리 \ref{algebraization-lemma-derived-completion-pseudo-coherent}에 의해 층이기
때문이다. $H^p(X, \lim \mathcal{F}/I^n\mathcal{F})$를 차수 $p$에서 얻는다.
$R\Gamma(X, -)$가 유도 극한과 교환하므로(단사 대상의 보조정리
\ref{injectives-lemma-RF-commutes-with-Rlim}), 다음도 얻는다.
$$
R\Gamma(X, \lim \mathcal{F}/I^n\mathcal{F}) =
R\Gamma(X, R\lim \mathcal{F}/I^n\mathcal{F}) =
R\lim R\Gamma(X, \mathcal{F}/I^n\mathcal{F})
$$
대수학 심화의 비고 \ref{more-algebra-remark-how-unique}에 의해 완전열
$$
0 \to
R^1\lim H^{p - 1}(X, \mathcal{F}/I^n\mathcal{F}) \to
H^p(X, \lim \mathcal{F}/I^n\mathcal{F}) \to
\lim H^p(X, \mathcal{F}/I^n\mathcal{F}) \to 0
$$
을 얻는다. 이는 $A$-가군들의 열이다. 위의 결과들을 합치면 보조정리의 첫째
명제를 얻는다. $R^1\lim$ 항의 소멸은 스킴의 코호몰로지의 보조정리
\ref{coherent-lemma-ML-cohomology-powers-ideal}에서 따른다.
\end{proof}

\begin{remark}
\label{algebraization-remark-references}
관련 내용을 논의하는 문헌 몇 가지를 소개한다. 고유 사상에 대한 ``유도
형식함수 정리''는 \cite[정리 6.3.1]{lurie-thesis}까지 거슬러 올라가는 듯하다.
\cite{dag12}에도 논의가 있으며, 특히 제4장은 아핀 경우를 다룬다. 대수학
심화의 절 \ref{more-algebra-section-derived-completion}을 보라.
\cite[절 2.9]{G-R}에서는 (ind)-연접 가군을 사용하여 고유 밑변환과 유도
완비를 논의한다. 준연접 가군들의 복합체에 대한
(\ref{algebraization-equation-formal-functions})의 유사 결과는
\cite[정리 6.5]{HL-P}에서 찾을 수 있다.
\end{remark}







\section{국소 코호몰로지의 대수화 I}
\label{algebraization-section-algebraization-sections-general}

\noindent
$A$를 Noether 환이라 하고 $I$와 $J$를 $A$의 두 아이디얼이라 하자.
$M$을 유한 $A$-가군이라 하자. 이 절에서는 대상
$$
R\Gamma_J(M)^\wedge
\quad\text{이며 다음에 속하는}\quad
D(A)
$$
의 코호몰로지 군들을 연구한다. 여기서 ${}^\wedge$는 유도 $I$-진 완비를
나타낸다. 쌍대화 복합체의 보조정리
\ref{dualizing-lemma-completion-local-H0}에서 $A$가 $I$에 관해 완비이면
동형사상
$$
\colim H^0_Z(M) \longrightarrow H^0(R\Gamma_J(M)^\wedge)
$$
이 존재함을 보였다. 여기서 (유향) 여극한은 $Z = V(J')$ 꼴의 닫힌
부분집합들에 대해 취하며, 이들은 $J' \subset J$ 및
$V(J') \cap V(I) = V(J) \cap V(I)$를 만족한다. 이 닫힌 부분집합들의 합집합은
\begin{equation}
\label{algebraization-equation-associated-subset}
T = \{\mathfrak p \in \Spec(A) :
V(\mathfrak p) \cap V(I) \subset V(J) \cap V(I)\}
\end{equation}
이다. 이는 특수화에 대해 안정적인 $\Spec(A)$의 부분집합이다. 위의 결과는
$$
H^0_T(M) \longrightarrow H^0(R\Gamma_J(M)^\wedge)
$$
가 동형사상이라는 명제가 된다. 단, $A$가 $I$에 관해 완비라고 가정한다.
국소 코호몰로지의 보조정리 \ref{local-cohomology-lemma-adjoint-ext}와 비고
\ref{local-cohomology-remark-upshot}을 보라. 이 동형사상을 더 높은
코호몰로지 군들로 확장하는 방법은 다음 보조정리에 기초한다.

\begin{lemma}
\label{algebraization-lemma-kill-completion-general}
$I, J$를 Noether 환 $A$의 아이디얼들이라 하자. $M$을 유한 $A$-가군이라
하자. $\mathfrak p \subset A$를 소아이디얼이라 하자. $s$와 $d$를 정수라
하자. 다음을 가정하자.
\begin{enumerate}
\item $A$는 쌍대화 복합체를 갖는다,
\item $\mathfrak p \not \in V(J) \cap V(I)$,
\item $\text{cd}(A, I) \leq d$, 그리고
\item 모든 소아이디얼 $\mathfrak p' \subset \mathfrak p$에 대해
$\text{depth}_{A_{\mathfrak p'}}(M_{\mathfrak p'}) +
\dim((A/\mathfrak p')_\mathfrak q) > d + s$
가 모든 $\mathfrak q \in V(\mathfrak p') \cap V(J) \cap V(I)$에 대해
성립한다.
\end{enumerate}
그러면 $f \in A$, $f \not \in \mathfrak p$인 원소가 존재하여
$H^i(R\Gamma_J(M)^\wedge)$를 모든 $i \leq s$에 대해 소멸시킨다. 여기서
${}^\wedge$는 $I$-진 완비를 나타낸다.
\end{lemma}

\begin{proof}
$R\Gamma_J = R\Gamma_{V(J)}$임과 $I + J$에 대한 유사한 사실을 사용하겠다.
쌍대화 복합체의 보조정리
\ref{dualizing-lemma-local-cohomology-noetherian}을 보라. 다음을 관찰하자.
$R\Gamma_J(M)^\wedge = R\Gamma_I(R\Gamma_J(M))^\wedge =
R\Gamma_{I + J}(M)^\wedge$이다. 쌍대화 복합체의 보조정리
\ref{dualizing-lemma-complete-and-local}와
\ref{dualizing-lemma-local-cohomology-ss}를 보라. 따라서 $J$를 $I + J$로
바꾸고 $I \subset J$ 및 $\mathfrak p \not \in V(J)$를 가정해도 된다.
다음을 상기하자.
$$
R\Gamma_J(M)^\wedge = R\Hom_A(R\Gamma_I(A), R\Gamma_J(M))
$$
이는 대수학 심화의 보조정리
\ref{more-algebra-lemma-derived-completion}에 있는 유도 완비의 기술과 쌍대화
복합체의 보조정리
\ref{dualizing-lemma-compute-local-cohomology-noetherian}에 있는 국소
코호몰로지의 기술을 결합하면 얻어진다. 가정 (3)은 $R\Gamma_I(A)$가
차수 $\leq d$에서만 영이 아닌 코호몰로지를 갖는다는 뜻이다.
$R\Gamma_I(A)$의 표준 절단을 사용하면 다음을 보이는 것으로 충분하다.
$$
\text{Ext}^i(N, R\Gamma_J(M))
$$
이 어떤 $f \in A$, $f \not \in \mathfrak p$에 의해 소멸되며, 이는 모든
$i \leq s + d$와 모든 $A$-가군 $N$에 대해 성립함을 보이면 된다. 다시
$R\Gamma_J(M)$의 표준 절단을 사용하면, $H^i_J(M)$이 어떤
$f \in A$, $f \not \in \mathfrak p$에 의해 모든 $i \leq s + d$에 대해
소멸됨을 보이는 것으로 충분하다. 이는 국소 코호몰로지의 보조정리
\ref{local-cohomology-lemma-kill-local-cohomology-at-prime}에서 따른다.
\end{proof}

\begin{lemma}
\label{algebraization-lemma-kill-colimit-weak-general}
% P07-ERR-0149
$I, J$를 Noether 환의 아이디얼들이라 하자. $M$을 유한 $A$-가군이라 하자.
$s$와 $d$를 정수라 하자. (\ref{algebraization-equation-associated-subset})의 $T$에 관해
다음을 가정하자.
\begin{enumerate}
\item $A$는 쌍대화 복합체를 갖는다,
\item $\mathfrak p \in V(I)$이면 조건이 없다,
\item $\mathfrak p \not \in V(I)$이고 $\mathfrak p \in T$이면,
$\dim((A/\mathfrak p)_\mathfrak q) \leq d$인 어떤
$\mathfrak q \in V(\mathfrak p) \cap V(J) \cap V(I)$가 존재한다,
\item $\mathfrak p \not \in V(I)$이고 $\mathfrak p \not \in T$이면,
$$
\text{depth}_{A_\mathfrak p}(M_\mathfrak p) \geq s
\quad\text{또는}\quad
\text{depth}_{A_\mathfrak p}(M_\mathfrak p) +
\dim((A/\mathfrak p)_\mathfrak q) > d + s
$$
가 모든 $\mathfrak q \in V(\mathfrak p) \cap V(J) \cap V(I)$에 대해
성립한다.
\end{enumerate}
그러면 아이디얼 $J_0 \subset J$가 존재하여
$V(J_0) \cap V(I) = V(J) \cap V(I)$이고, 임의의 $J' \subset J_0$가
$V(J') \cap V(I) = V(J) \cap V(I)$를 만족할 때 사상
$$
R\Gamma_{J'}(M) \longrightarrow R\Gamma_{J_0}(M)
$$
은 차수 $\leq s$인 코호몰로지에서 동형사상을 유도한다. 더욱이 이
가군들은 $J_0I$의 어떤 거듭제곱에 의해 소멸된다.
\end{lemma}

\begin{proof}
다음 집합을 생각하자.
$$
B = \{\mathfrak p \not \in V(I),\ \mathfrak p \in T,\text{이고 }
\text{depth}(M_\mathfrak p) \leq s\}
$$
$J_0 \subset J$를 $V(J_0)$가 $B \cup V(J)$의 폐포가 되도록 택하자.

\medskip\noindent
주장 I: $V(J_0) \cap V(I) = V(J) \cap V(I)$이다.

\medskip\noindent
주장 I의 증명. 포함관계 $\supset$는 구성에 의해 성립한다.
$\mathfrak p$를 $V(J_0)$의 극소 소아이디얼이라 하자.
$\mathfrak p \in B \cup V(J)$이면 $\mathfrak p \in T$이거나
$\mathfrak p \in V(J)$이며, 어느 경우든 원하는 대로
$V(\mathfrak p) \cap V(I) \subset V(J) \cap V(I)$이다.
$\mathfrak p \not \in B \cup V(J)$이면 $V(\mathfrak p) \cap B$는
조밀하고 따라서 무한이다. 그러므로 국소 코호몰로지의 보조정리
\ref{local-cohomology-lemma-depth-function}에 의해
$\text{depth}(M_\mathfrak p) < s$이다. 실제로 다음과 같이 쓰자.
$V(\mathfrak p) \cap B = \{\mathfrak p_\lambda\}_{\lambda \in \Lambda}$.
(3)에서와 같이
$\mathfrak q_\lambda \in V(\mathfrak p_\lambda) \cap V(J) \cap V(I)$를
택하자. $\delta : \Spec(A) \to \mathbf{Z}$를 쌍대화 복합체
$\omega_A^\bullet$에 결부된 $A$의 차원함수라 하자. $\Lambda$는 무한이고
$\delta$는 유계이므로, 무한 부분집합 $\Lambda' \subset \Lambda$가
존재하여 그 위에서 $\delta(\mathfrak q_\lambda)$가 상수이다.
$\lambda \in \Lambda'$에 대해 다음이 성립한다.
$$
\text{depth}(M_{\mathfrak p_\lambda}) +
\delta(\mathfrak p_\lambda) - \delta(\mathfrak q_\lambda) =
\text{depth}(M_{\mathfrak p_\lambda}) +
\dim((A/\mathfrak p_\lambda)_{\mathfrak q_\lambda})
\leq d + s
$$
이는 (3)과 $B$의 정의에서 따른다. 스킴의 쌍대성의 보조정리
\ref{duality-lemma-sitting-in-degrees}에서 증명한 함수
$\text{depth} + \delta$의 반연속성에 의해 다음을 얻는다.
$$
\text{depth}(M_\mathfrak p) +
\dim((A/\mathfrak p)_{\mathfrak q_\lambda}) =
\text{depth}(M_\mathfrak p) + \delta(\mathfrak p) - \delta(\mathfrak q_\lambda)
\leq d + s
$$
또한 $\mathfrak p \not \in V(I)$이므로 (4)에서 $\mathfrak p \in T$임을
읽어 낼 수 있다. 즉,
$V(\mathfrak p) \cap V(I) \subset V(J) \cap V(I)$이다. 이로써 주장 I의
증명이 끝난다.

\medskip\noindent
% P07-ERR-0150
주장 II: $H^i_{J_0}(M) \to H^i_J(M)$은 $i \leq s$이고
$J' \subset J_0$가 $V(J') \cap V(I) = V(J) \cap V(I)$를 만족할 때
동형사상이다.

\medskip\noindent
주장 II의 증명. $\mathfrak p \in V(J')$이지만 $V(J_0)$에는 속하지 않는
것을 하나 택하자. $H^i_{\mathfrak pA_\mathfrak p}(M_\mathfrak p) = 0$이
모든 $i \leq s$에 대해 성립함을 보이면 충분하다. 국소 코호몰로지의
보조정리 \ref{local-cohomology-lemma-isomorphism}을 보라.
$\mathfrak p \in T$임을 관찰하자. 따라서 $\mathfrak p$는 $B$에 속하지
않으므로 $\text{depth}(M_\mathfrak p) > s$임을 알 수 있고, 이 군들은
쌍대화 복합체의 보조정리 \ref{dualizing-lemma-depth}에 의해 소멸한다.

\medskip\noindent
주장 III. 보조정리의 마지막 명제가 참이다.

\medskip\noindent
주장 II에 의해 $i \leq s$에 대해 다음이 성립한다.
$$
H^i_T(M) = H^i_{J_0}(M) = H^i_{J'}(M)
$$
여기서 이는 모든 아이디얼 $J' \subset J_0$ 중
$V(J')  \cap V(I) = V(J) \cap V(I)$를 만족하는 것에 대해 성립한다.
국소 코호몰로지의 보조정리
\ref{local-cohomology-lemma-adjoint-ext}를 보라. 부분집합들
$T \subset T \cup V(I)$, 가군 $M$, 정수 $s$에 대해 국소 코호몰로지의
명제 \ref{local-cohomology-proposition-annihilator}의 가정들을 확인하자.
$\mathfrak p \subset \mathfrak q$가 주어지고
$\mathfrak p \not \in T \cup V(I)$이며 $\mathfrak q \in T$일 때 다음을
보여야 한다.
$$
\text{depth}_{A_\mathfrak p}(M_\mathfrak p) +
\dim((A/\mathfrak p)_\mathfrak q) > s
$$
$\text{depth}(M_\mathfrak p) \geq s$이면
$(A/\mathfrak p)_\mathfrak q$의 차원이 적어도 $1$이므로 이는 참이다.
따라서 $\text{depth}(M_\mathfrak p) < s$라고 가정해도 된다.
$\mathfrak q \in V(I)$이면 $\mathfrak q \in V(J) \cap V(I)$이고 부등식은
(4)에 의해 성립한다. $\mathfrak q \not \in V(I)$이면 (3)을 사용하여
$\mathfrak q' \in V(\mathfrak q) \cap V(J) \cap V(I)$를 택하되
$\dim((A/\mathfrak q)_{\mathfrak q'}) \leq d$가 되게 할 수 있다.
그러면 가정 (4)에서 다음을 얻는다.
$$
\text{depth}_{A_\mathfrak p}(M_\mathfrak p) +
\dim((A/\mathfrak p)_{\mathfrak q'}) > s + d
$$
$A$는 사슬적이므로 이로부터 원하는 부등식이 따른다. 국소 코호몰로지의
명제 \ref{local-cohomology-proposition-annihilator}를 적용하면
$J'' \subset A$를 얻는데, 이는 $V(J'') \subset T \cup V(I)$를 만족하고
$J''$가 $H^i_T(M)$을 모든 $i \leq s$에 대해 소멸시킨다. 그러면
$V(J'') \cup V(J_0) \cup V(I) = V(J'I)$로 쓸 수 있으며, 여기서
$J' \subset J_0$이고 $V(J') \cap V(I) = V(J) \cap V(I)$이다. $J_0$를
$J'$로 바꾸면 증명이 끝난다.
\end{proof}

\begin{lemma}
\label{algebraization-lemma-kill-colimit-general}
보조정리 \ref{algebraization-lemma-kill-colimit-weak-general}에서 아무 조건도 없는
조건 (2) 대신 다음을 가정하자.
\begin{enumerate}
\item[(2')] $\mathfrak p \in V(I)$이고
$\mathfrak p \not \in V(J) \cap V(I)$이면
$\text{depth}_{A_\mathfrak p}(M_\mathfrak p) +
\dim((A/\mathfrak p)_\mathfrak q) > s$
가 모든 $\mathfrak q \in V(\mathfrak p) \cap V(J) \cap V(I)$에 대해
성립한다.
\end{enumerate}
그러면 이 조건들은 $H^i_{J_0}(M)$이 유한 $A$-가군임도 함의하며, 이는
$i \leq s$에 대해 성립한다.
\end{lemma}

\begin{proof}
보조정리 \ref{algebraization-lemma-kill-colimit-weak-general}의 증명을 보면
$H^i_{J_0}(M) = H^i_T(M)$임을 상기하자. 따라서
$\mathfrak p \not \in T$이고 $\mathfrak q \in T$이며
$\mathfrak p \subset \mathfrak q$일 때
$\text{depth}_{A_\mathfrak p}(M_\mathfrak p) +
\dim((A/\mathfrak p)_\mathfrak q) > s$가 성립함을 확인하면 충분하다.
국소 코호몰로지의 명제 \ref{local-cohomology-proposition-finiteness}를 보라.
조건 (2')에 따르면 이는 $\mathfrak p \in V(I)$일 때 성립한다.
$H^i_T(M)$이 $IJ_0$의 어떤 거듭제곱에 의해 소멸됨을 알고 있으므로,
국소 코호몰로지의 명제
\ref{local-cohomology-proposition-annihilator}에 의해 이 조건은
$\mathfrak p \not \in V(IJ_0)$일 때도 성립한다. 이로써 모든 경우를
다루었고 증명이 끝난다.
\end{proof}

\begin{lemma}
\label{algebraization-lemma-kill-colimit-support-general}
보조정리 \ref{algebraization-lemma-kill-colimit-weak-general}에서 다음을 추가로 가정하자.
\begin{enumerate}
\item[(6)] $\mathfrak p \not \in V(I)$이고 $\mathfrak p \in T$이면
$\text{depth}_{A_\mathfrak p}(M_\mathfrak p) > s$이다.
\end{enumerate}
그러면 $H^i_{J_0}(M) = H^i_J(M) = H^i_{J + I}(M)$가 $i \leq s$에 대해
성립하며 이 가군들은 $I$의 어떤 거듭제곱에 의해 소멸된다.
\end{lemma}

\begin{proof}
$\mathfrak p \in V(J)$이거나 $\mathfrak p \in V(J_0)$이지만
$\mathfrak p \not \in V(J + I) = V(J_0 + I)$인 것을 택하자.
$H^i_{\mathfrak pA_\mathfrak p}(M_\mathfrak p) = 0$이 $i \leq s$에 대해
성립함을 보이면 충분하다. 국소 코호몰로지의 보조정리
\ref{local-cohomology-lemma-isomorphism}을 보라. 이 군들은 조건 (6)과
쌍대화 복합체의 보조정리 \ref{dualizing-lemma-depth}에 의해 소멸한다.
마지막 명제는 국소 코호몰로지의 명제
\ref{local-cohomology-proposition-annihilator}에서 따른다.
\end{proof}

\begin{lemma}
\label{algebraization-lemma-algebraize-local-cohomology-general}
$I, J$를 Noether 환 $A$의 아이디얼들이라 하자.
$M$을 유한 $A$-가군이라 하자.
$s$와 $d$를 정수라 하자. (\ref{algebraization-equation-associated-subset})의 $T$에 관해
다음을 가정하자.
\begin{enumerate}
\item $A$는 $I$에 관해 진 완비이고 쌍대화 복합체를 갖는다,
\item $\mathfrak p \in V(I)$이면 조건이 없다,
\item $\text{cd}(A, I) \leq d$이다,
\item $\mathfrak p \not \in V(I)$이고 $\mathfrak p \not \in T$이면
$$
\text{depth}_{A_\mathfrak p}(M_\mathfrak p) \geq s
\quad\text{또는}\quad
\text{depth}_{A_\mathfrak p}(M_\mathfrak p) +
\dim((A/\mathfrak p)_\mathfrak q) > d + s
$$
가 모든 $\mathfrak q \in V(\mathfrak p) \cap V(J) \cap V(I)$에 대해
성립한다,
\item $\mathfrak p \not \in V(I)$이고 $\mathfrak p \not \in T$이며,
$V(\mathfrak p) \cap V(J) \cap V(I) \not = \emptyset$이고
$\text{depth}(M_\mathfrak p) < s$이면 다음 중 하나가 성립한다
% P07-ERR-0151
\footnote{우리 방법은 이 추가 조건을 강제한다. 이 점은 나중에 다시 다룰
것이다(향후 참고문헌 삽입).}:
\begin{enumerate}
\item $\dim(\text{Supp}(M_\mathfrak p)) < s + 2$이다\footnote{예를 들어
$M$이 Serre 조건 $(S_s)$를 $V(I) \cup T$의 여집합에서 만족하는
경우이다.}, 또는
\item  $\delta(\mathfrak p) > d + \delta_{max} - 1$
이다. 여기서 $\delta$는 차원함수이고 $\delta_{max}$는 $\delta$가
$V(J) \cap V(I)$ 위에서 갖는 최댓값이다, 또는
\item $\text{depth}_{A_\mathfrak p}(M_\mathfrak p) +
\dim((A/\mathfrak p)_\mathfrak q) > d + s + \delta_{max} - \delta_{min} - 2$
가 모든 $\mathfrak q \in V(\mathfrak p) \cap V(J) \cap V(I)$에 대해
성립한다.
\end{enumerate}
\end{enumerate}
그러면 아이디얼 $J_0 \subset J$가 존재하여
$V(J_0) \cap V(I) = V(J) \cap V(I)$이고, 임의의 $J' \subset J_0$가
$V(J') \cap V(I) = V(J) \cap V(I)$를 만족할 때 사상
$$
R\Gamma_{J'}(M) \longrightarrow R\Gamma_J(M)^\wedge
$$
은 차수 $\leq s$인 코호몰로지에서 동형사상을 유도한다. 여기서
${}^\wedge$는 유도 $I$-진 완비를 나타낸다.
\end{lemma}

\noindent
독자에게 먼저 국소 경우의 증명(보조정리
\ref{algebraization-lemma-algebraize-local-cohomology})을 읽기를 권한다. 모든 부등식을
다루지 않고도 증명의 구조를 설명해 주기 때문이다.

\begin{proof}
아이디얼 $\mathfrak a \subset A$에 대해
$R\Gamma_\mathfrak a = R\Gamma_{V(\mathfrak a)}$이다. 쌍대화 복합체의
보조정리 \ref{dualizing-lemma-local-cohomology-noetherian}을 보라.
다음으로 다음을 관찰하자.
$$
R\Gamma_J(M)^\wedge =
R\Gamma_I(R\Gamma_J(M))^\wedge =
R\Gamma_{I + J}(M)^\wedge =
R\Gamma_{I + J'}(M)^\wedge =
R\Gamma_I(R\Gamma_{J'}(M))^\wedge =
R\Gamma_{J'}(M)^\wedge
$$
이는 쌍대화 복합체의 보조정리
\ref{dualizing-lemma-local-cohomology-ss}와
\ref{dualizing-lemma-complete-and-local}에서 따른다. 이로써 보조정리의
명제에 나오는 화살표를 어떻게 정의하는지 설명된다.

\medskip\noindent
보조정리 \ref{algebraization-lemma-kill-colimit-weak-general}의 가정들은 $M$에 대한 현재
가정들에서 함의된다고 주장한다. 확인할 것은 가정 (3)뿐이다. 따라서
$\mathfrak p \not \in V(I)$이고 $\mathfrak p \in T$라 하자.
$V(\mathfrak p) \cap V(I)$는 공집합이 아닌데, 이는 $I$가 $A$의
Jacobson 근기에 포함되기 때문이다. 대수학의 보조정리
\ref{algebra-lemma-radical-completion}을 보라. $\mathfrak p \in T$이므로
$V(\mathfrak p) \cap V(I) = V(\mathfrak p) \cap V(J) \cap V(I)$이다.
$\mathfrak q \in V(\mathfrak p) \cap V(I)$를 한 기약 성분의 일반점이라
하자. 국소 코호몰로지의 보조정리
\ref{local-cohomology-lemma-cd-local}에 의해
$\text{cd}(A_\mathfrak q, I_\mathfrak q) \leq d$이다.
$V(\mathfrak pA_\mathfrak q) \cap V(I_\mathfrak q) =
\{\mathfrak qA_\mathfrak q\}$인데, 이는 $\mathfrak q$의 선택에 의한다.
국소 코호몰로지의 보조정리
\ref{local-cohomology-lemma-cd-bound-dim-local}에 의해
$\dim((A/\mathfrak p)_\mathfrak q) \leq d$를 얻는다.

\medskip\noindent
$s < 0$일 때 보조정리가 성립함을 관찰하자. 이는 자명한 경우가 아니다.
$H^i(R\Gamma_J(M)^\wedge)$이 $i < 0$일 때 영이라는 것이 선험적으로
분명하지 않기 때문이다. 그러나 이 소멸은 쌍대화 복합체의 보조정리
\ref{dualizing-lemma-completion-local}에서 확립되었다. 이제 $s \geq 0$에
대해 귀납법으로 보조정리를 증명하겠다.

\medskip\noindent
$s = 0$인 보조정리는 보조정리
\ref{algebraization-lemma-kill-colimit-weak-general}의 결론과 쌍대화 복합체의 보조정리
\ref{dualizing-lemma-completion-local-H0}에서 즉시 따른다.

\medskip\noindent
$s > 0$이고 더 작은 $s$의 값들에 대해 보조정리가 증명되었다고 가정하자.
$M' \subset M$을 지지집합이 $V(I) \cup T$에 포함되는 극대 부분가군이라
하자. 그러면 $M'$은 유한 $A$-가군이고 그 지지집합은
$V(J') \cup V(I)$에 포함되는데, 여기서 어떤 아이디얼 $J' \subset J$가
$V(J') \cap V(I) = V(J) \cap V(I)$를 만족한다. 다음을 주장한다.
$$
R\Gamma_{J'}(M') \to R\Gamma_J(M')^\wedge
$$
이는 $J'$를 어떻게 택하든 동형사상이다. 실제로 짧은 완전열
$0 \to M_1 \oplus M_2 \to M' \to N \to 0$을 택할 수 있는데,
$M_1$은 $J'$의 어떤 거듭제곱에 의해 소멸되고, $M_2$는 $I$의 어떤
거듭제곱에 의해 소멸되며, $N$은 $I + J'$의 어떤 거듭제곱에 의해
소멸된다. 따라서 $M_1$, $M_2$, $N$ 각각에 대해 주장이 성립함을 보이면
충분하다. $M_1$의 경우 $R\Gamma_{J'}(M_1) = M_1$이고, $M_1$은 유한
$A$-가군이면서 $I$에 관해 진 완비이므로 $M_1^\wedge = M_1$이다. 증명
첫머리의 관찰에 의해 이는 $M_1$에 대한 주장을 증명한다. $M_2$의 경우
다음을 얻는다.
% P07-ERR-0152
$H^i_J(M_2) = H^i_{I + J}(M) = H^i_{I + J'}(M) = H^i_{J'}(M_2)$
이들은 $I$의 어떤 거듭제곱에 의해 소멸되므로 유도 완비이다.
% P07-ERR-0153
따라서 이 경우에도 주장이 성립한다. $N$에 대해서는 방금 제시한 두
논증 중 어느 것이든 사용할 수 있다. 짧은 완전열
$0 \to M' \to M \to M/M' \to 0$
을 고찰하면 $M/M'$에 대해 보조정리를 증명하는 것으로 충분함을 알 수
있다. 따라서 $\text{Ass}(M) \cap (V(I) \cup T) = \emptyset$라고 가정해도
된다.

\medskip\noindent
$\mathfrak p \in \text{Ass}(M)$가
$V(\mathfrak p) \cap V(J) \cap V(I) = \emptyset$를 만족한다고 하자.
$I$가 $A$의 Jacobson 근기에 포함되므로, 이는
$V(\mathfrak p) \cap V(J') = \emptyset$임을 함의하며, 이 명제는 임의의
$J' \subset J$ 중 $V(J') \cap V(I) = V(J) \cap V(I)$를 만족하는 것에
대해 성립한다. 따라서
$N = H^0_\mathfrak p(M)$으로 놓으면
$R\Gamma_J(N) = R\Gamma_{J'}(N) = 0$이 모든 $J' \subset J$ 중
$V(J') \cap V(I) = V(J) \cap V(I)$를 만족하는 것에 대해 성립함을 알 수
있다. 특히 $R\Gamma_J(N)^\wedge = 0$이다. 그러므로 $M$을 $M/N$으로
바꾸어도 된다. 이 치환은 조건 (4)나 (5)에 관여하지 않는 소아이디얼에서만
$M$의 구조를 바꾸기 때문이다. 이를 반복하면 다음을 가정해도 된다.
$V(\mathfrak p) \cap V(J) \cap V(I) \not = \emptyset$
가 모든 $\mathfrak p \in \text{Ass}(M)$에 대해 성립한다.

\medskip\noindent
다음을 가정하자.
$\text{Ass}(M) \cap (V(I) \cup T) = \emptyset$이고
$V(\mathfrak p) \cap V(J) \cap V(I) \not = \emptyset$
가 모든 $\mathfrak p \in \text{Ass}(M)$에 대해 성립한다.
$\mathfrak p \in \text{Ass}(M)$라 하자. 보조정리
\ref{algebraization-lemma-kill-completion-general}을 적용할 수 있음을 보이고자 한다.
이를 확인하는 과정에서 추가 조건 (5)를 사용한다.
$\mathfrak p' \subset \mathfrak p$를 택하고
% P07-ERR-0154
$\mathfrak q' \subset V(\mathfrak p) \cap V(J) \cap V(I)$를 택하자.
\begin{enumerate}
\item $M_{\mathfrak p'} = 0$이면
$\text{depth}(M_{\mathfrak p'}) = \infty$이고
$\text{depth}(M_{\mathfrak p'}) +
\dim((A/\mathfrak p')_{\mathfrak q'}) > d + s$이다.
\item $\text{depth}(M_{\mathfrak p'}) < s$이면
$\text{depth}(M_{\mathfrak p'}) +
\dim((A/\mathfrak p')_{\mathfrak q'}) > d + s$가 (4)에 의해 성립한다.
\end{enumerate}
나머지 경우에는 $M_{\mathfrak p'} \not = 0$이고
$\text{depth}(M_{\mathfrak p'}) \geq s$이다. 특히 $\mathfrak p'$가 $M$의
지지집합에 속함을 알 수 있고, $\mathfrak p'' \subset \mathfrak p'$이면서
$\mathfrak p'' \in \text{Ass}(M)$인 것을 택할 수 있다.
\begin{enumerate}
\item[(a)] 다음을 관찰하자.
$\dim((A/\mathfrak p'')_{\mathfrak p'}) \geq \text{depth}(M_{\mathfrak p'})$
이는 대수학의 보조정리
\ref{algebra-lemma-depth-dim-associated-primes}에서 따른다. 등호가
성립하면 다음을 얻는다.
$$
\text{depth}(M_{\mathfrak p'}) + \dim((A/\mathfrak p')_{\mathfrak q'}) =
\text{depth}(M_{\mathfrak p''}) + \dim((A/\mathfrak p'')_{\mathfrak q'})
> s + d
$$
(4)를 $\mathfrak p''$에 적용했으며, 이로써 끝난다. 따라서 문제가 될 수
있는 경우는 다음뿐이다.
$\dim((A/\mathfrak p'')_{\mathfrak p'}) > \text{depth}(M_{\mathfrak p'})$.
이는 $\dim(M_\mathfrak p) \geq s + 2$를 함의한다. 따라서 (5)(a)가
성립하면 이 경우는 일어나지 않는다.
\item[(b)] (5)(b)가 성립하면 다음을 얻는다.
$$
\text{depth}(M_{\mathfrak p'}) + \dim((A/\mathfrak p')_{\mathfrak q'})
\geq s + \delta(\mathfrak p') - \delta(\mathfrak q')
\geq s + 1 + \delta(\mathfrak p) - \delta_{max}
> s + d
$$
이는 원하는 바이다.
\item[(c)] (5)(c)가 성립하면 다음을 얻는다.
\begin{align*}
\text{depth}(M_{\mathfrak p'}) + \dim((A/\mathfrak p')_{\mathfrak q'})
& \geq
s + \delta(\mathfrak p') - \delta(\mathfrak q') \\
& \geq
s + 1 + \delta(\mathfrak p) - \delta(\mathfrak q') \\
& =
s + 1 + \delta(\mathfrak p) - \delta(\mathfrak q) +
\delta(\mathfrak q) - \delta(\mathfrak q') \\
& >
s + 1 + (s + d + \delta_{max} - \delta_{min} - 2) +
\delta(\mathfrak q) - \delta(\mathfrak q') \\
& \geq 
2s + d - 1 \geq s + d
\end{align*}
이는 원하는 바이다. 이 논증이 성립하는 까닭은 소아이디얼
$\mathfrak q \in V(\mathfrak p) \cap V(J) \cap V(I)$가 존재함을 알고
있기 때문이다.
\end{enumerate}
이제 귀납 단계를 수행할 준비가 되었다.

\medskip\noindent
$J_0$를 보조정리 \ref{algebraization-lemma-kill-colimit-weak-general}에서와 같이 택하고,
정수 $t > 0$을 $(J_0I)^t$가 $H^s_J(M)$을 소멸시키도록 택하자.
보조정리 \ref{algebraization-lemma-kill-completion-general}의 가정들은 모든
$\mathfrak p \in \text{Ass}(M)$에 대해 만족된다(앞 문단을 보라).
따라서 $\mathfrak a \subset A$를 $H^s(R\Gamma_J(M)^\wedge)$의 소멸자라
하면 이는 $\mathfrak p$에 포함되지 않으며,
이는 $\mathfrak p \in \text{Ass}(M)$에 대해 성립한다. 따라서
$f \in \mathfrak a(J_0I)^t$를 택할 수 있는데, 이는 $M$의 어느 결합
소아이디얼에도 속하지 않으며 $H^s(R\Gamma_J(M)^\wedge)$과
$H^s_J(M)$을 모두 소멸시킨다. 그러면 $f$는 $M$ 위의 영인자가 아니므로
다음 짧은 완전열을 고찰할 수 있다.
$$
0 \to M \xrightarrow{f} M \to M/fM \to 0
$$
$f$의 선택에 의해 다음을 얻는다.
$$
\xymatrix{
H^{s - 1}_{J'}(M) \ar[d] \ar[r] &
H^{s - 1}_{J'}(M/fM) \ar[d] \ar[r] &
H^s_{J'}(M) \ar[d] \ar[r] & 0 \\
H^{s - 1}(R\Gamma_J(M)^\wedge) \ar[r] &
H^{s - 1}(R\Gamma_J(M/fM)^\wedge) \ar[r] &
H^s(R\Gamma_J(M)^\wedge) \ar[r] & 0
}
$$
이는 임의의 $J' \subset J_0$ 중
$V(J') \cap V(I) = V(J) \cap V(I)$를 만족하는 것에 대해 성립한다.
따라서 $J'$를 $M$, $M/fM$, $s - 1$에
대해 작동하도록 택하면(귀납 가정에 의해 가능하다. 다음 문단을 보라),
보조정리가 참이라는 결론을 얻는다.

\medskip\noindent
증명을 끝내려면 가군 $M/fM$이 $s - 1$에 대해 가정 (4)와 (5)를
만족함을 보여야 한다. 따라서 $\mathfrak p$를 $M/fM$의 지지집합에 속하는
소아이디얼이라 하되 $\text{depth}((M/fM)_\mathfrak p) < s - 1$이고
$V(\mathfrak p) \cap V(J) \cap V(I)$가 공집합이 아니라고 하자. 그러면
$\dim(M_\mathfrak p) = \dim((M/fM)_\mathfrak p) + 1$이고
$\text{depth}(M_\mathfrak p) = \text{depth}((M/fM)_\mathfrak p) + 1$이다.
특히 (4)와 (5)가 $\mathfrak p$와 $M$에 대해 원래 값 $s$로 성립함을
알고 있다. 이제 원하는 부등식들은 직접 확인하면 따른다.
\end{proof}

\begin{example}
\label{algebraization-example-no-ML}
보조정리 \ref{algebraization-lemma-algebraize-local-cohomology-general}에서 역계
$H^i_J(M/I^nM)$이 Mittag--Leffler 조건을 만족하는지는 알 수 없다.
예를 들어 $A = \mathbf{Z}_p[[x, y]]$, $I = (p)$, $J = (p, x)$이고
$M = A/(xy - p)$라고 하자. 그러면
$H^0_J(M/p^nM) \to H^0_J(M/pM)$
의 상은 $y^n$으로 생성되는 아이디얼이며, 이는
$M/pM = A/(p, xy)$ 안의 아이디얼이다.
\end{example}





\section{국소 코호몰로지의 대수화 II}
\label{algebraization-section-algebraization-punctured}

\noindent
이 절에서는 $(A, \mathfrak m)$이 국소환이고 국소 코호몰로지
$R\Gamma_\mathfrak m$을 $\mathfrak m$에 대해 취할 때 절
\ref{algebraization-section-algebraization-sections-general}의 논증을 다시 전개한다.
앞에서와 마찬가지로 주된 도구는 다음 보조정리이다.

\begin{lemma}
\label{algebraization-lemma-kill-completion}
$(A, \mathfrak m)$을 Noether 국소환이라 하자. $I \subset A$를 아이디얼이라
하자. $M$을 유한 $A$-가군이라 하고 $\mathfrak p \subset A$를 소아이디얼이라
하자. $s$와 $d$를 정수라 하자. 다음을 가정하자.
\begin{enumerate}
\item $A$는 쌍대화 복합체를 갖는다,
\item $\text{cd}(A, I) \leq d$이다, 그리고
\item
$\text{depth}_{A_\mathfrak p}(M_\mathfrak p) + \dim(A/\mathfrak p) > d + s$이다.
\end{enumerate}
그러면 $f \in A \setminus \mathfrak p$인 원소가 존재하여
$H^i(R\Gamma_\mathfrak m(M)^\wedge)$을 $i \leq s$에 대해 소멸시킨다.
여기서 ${}^\wedge$는 $I$-진 완비를 나타낸다.
\end{lemma}

\begin{proof}
국소 코호몰로지의 보조정리
\ref{local-cohomology-lemma-sitting-in-degrees}에 따르면 함수
$$
\mathfrak p' \longmapsto
\text{depth}_{A_{\mathfrak p'}}(M_{\mathfrak p'}) + \dim(A/\mathfrak p')
$$
는 $\Spec(A)$ 위에서 아래 반연속이다. 따라서
$\mathfrak p' \subset \mathfrak p$일 때 이 함수의 값은 $> s + d$이다.
그러므로 $\mathfrak p \not = \mathfrak m$이면 이 보조정리는 보조정리
\ref{algebraization-lemma-kill-completion-general}의 특수한 경우이다.
$\mathfrak p = \mathfrak m$이면 깊이와 국소 코호몰로지의 관계에 의해
$H^i_\mathfrak m(M) = 0$이 $i \leq s + d$에 대해 성립한다. 쌍대화
복합체의 보조정리 \ref{dualizing-lemma-depth}를 보라. 따라서 보조정리
\ref{algebraization-lemma-kill-completion-general}의 증명에 주어진 논증은 이 퇴화한
경우에 $H^i(R\Gamma_\mathfrak m(M)^\wedge) = 0$이 $i \leq s$에 대해
성립함을 보여 준다.
\end{proof}

\begin{lemma}
\label{algebraization-lemma-kill-colimit-weak}
$(A, \mathfrak m)$을 Noether 국소환이라 하자. $I \subset A$를 아이디얼이라
하자. $M$을 유한 $A$-가군이라 하자. $s$와 $d$를 정수라 하자. 다음을
가정하자.
\begin{enumerate}
\item $A$는 쌍대화 복합체를 갖는다,
\item $\mathfrak p \in V(I)$이면 조건이 없다,
\item $\mathfrak p \not \in V(I)$이고
$V(\mathfrak p) \cap V(I) = \{\mathfrak m\}$이면
$\dim(A/\mathfrak p) \leq d$이다,
\item $\mathfrak p \not \in V(I)$이고
$V(\mathfrak p) \cap V(I) \not = \{\mathfrak m\}$이면
$$
\text{depth}_{A_\mathfrak p}(M_\mathfrak p) \geq s
\quad\text{또는}\quad
\text{depth}_{A_\mathfrak p}(M_\mathfrak p) + \dim(A/\mathfrak p) > d + s
$$
\end{enumerate}
그러면 아이디얼 $J_0 \subset A$가 존재하여
$V(J_0) \cap V(I) = \{\mathfrak m\}$이고, 임의의 $J \subset J_0$가
$V(J) \cap V(I) = \{\mathfrak m\}$를 만족할 때 사상
$$
R\Gamma_J(M) \longrightarrow R\Gamma_{J_0}(M)
$$
은 차수 $\leq s$인 코호몰로지에서 동형사상을 유도한다. 더욱이 이
가군들은 $J_0I$의 어떤 거듭제곱에 의해 소멸된다.
\end{lemma}

\begin{proof}
이는 보조정리 \ref{algebraization-lemma-kill-colimit-weak-general}의 특수한 경우이다.
\end{proof}

\begin{lemma}
\label{algebraization-lemma-kill-colimit}
보조정리 \ref{algebraization-lemma-kill-colimit-weak}에서 아무 조건도 없는 조건 (2) 대신
다음을 가정하자.
\begin{enumerate}
\item[(2')] $\mathfrak p \in V(I)$이고 $\mathfrak p \not = \mathfrak m$이면
$\text{depth}_{A_\mathfrak p}(M_\mathfrak p) + \dim(A/\mathfrak p) > s$이다.
\end{enumerate}
그러면 이 조건들은 $H^i_{J_0}(M)$이 유한 $A$-가군임도 함의하며, 이는
$i \leq s$에 대해 성립한다.
\end{lemma}

\begin{proof}
이는 보조정리 \ref{algebraization-lemma-kill-colimit-general}의 특수한 경우이다.
\end{proof}

\begin{lemma}
\label{algebraization-lemma-kill-colimit-support}
보조정리 \ref{algebraization-lemma-kill-colimit-weak}에서 다음을 추가로 가정하자.
\begin{enumerate}
\item[(6)] $\mathfrak p \not \in V(I)$이고
$V(\mathfrak p) \cap V(I) = \{\mathfrak m\}$이면
$\text{depth}_{A_\mathfrak p}(M_\mathfrak p) > s$이다.
\end{enumerate}
그러면 $H^i_{J_0}(M) = H^i_J(M) = H^i_\mathfrak m(M)$이 $i \leq s$에
대해 성립하고 이 가군들은 $I$의 어떤 거듭제곱에 의해 소멸된다.
\end{lemma}

\begin{proof}
이는 보조정리 \ref{algebraization-lemma-kill-colimit-support-general}의 특수한 경우이다.
\end{proof}

\begin{lemma}
\label{algebraization-lemma-algebraize-local-cohomology}
$(A, \mathfrak m)$을 Noether 국소환이라 하자. $I \subset A$를 아이디얼이라
하자. $M$을 유한 $A$-가군이라 하자. $s$와 $d$를 정수라 하자. 다음을
가정하자.
\begin{enumerate}
\item $A$는 $I$에 관해 진 완비이고 쌍대화 복합체를 갖는다,
\item $\mathfrak p \in V(I)$이면 조건이 없다,
\item $\text{cd}(A, I) \leq d$이다,
\item $\mathfrak p \not \in V(I)$이고
$V(\mathfrak p) \cap V(I) \not = \{\mathfrak m\}$이면
$$
\text{depth}_{A_\mathfrak p}(M_\mathfrak p) \geq s
\quad\text{또는}\quad
\text{depth}_{A_\mathfrak p}(M_\mathfrak p) + \dim(A/\mathfrak p) > d + s
$$
\end{enumerate}
그러면 아이디얼 $J_0 \subset A$가 존재하여
$V(J_0) \cap V(I) = \{\mathfrak m\}$이고, 임의의 $J \subset J_0$가
$V(J) \cap V(I) = \{\mathfrak m\}$를 만족할 때 사상
$$
R\Gamma_J(M) \longrightarrow
R\Gamma_J(M)^\wedge = R\Gamma_\mathfrak m(M)^\wedge
$$
은 차수 $\leq s$인 코호몰로지에서 동형사상을 유도한다. 여기서
${}^\wedge$는 유도 $I$-진 완비를 나타낸다.
\end{lemma}

\begin{proof}
이 보조정리는 보조정리
\ref{algebraization-lemma-algebraize-local-cohomology-general}의 특수한 경우이다. 실제로
$\delta_{max} = \delta_{min} = \delta(\mathfrak m)$이므로 조건 (4)는 조건
(5)(c)를 함의한다. 이 중요한 특수한 경우의 증명은 다소 더 쉽고 확인할
것도 더 적으므로 여기에서 제시한다.

\medskip\noindent
현재 상황에서는 $R\Gamma_\mathfrak a$와
$R\Gamma_{V(\mathfrak a)}$ 사이에 차이가 없다. 쌍대화 복합체의 보조정리
\ref{dualizing-lemma-local-cohomology-noetherian}을 보라. 다음으로 다음을
관찰하자.
$$
R\Gamma_\mathfrak m(M)^\wedge =
R\Gamma_I(R\Gamma_J(M))^\wedge =
R\Gamma_J(M)^\wedge
$$
이는 쌍대화 복합체의 보조정리
\ref{dualizing-lemma-local-cohomology-ss}와
\ref{dualizing-lemma-complete-and-local}
에서 따르며, 보조정리의 명제에 있는 등호를 설명한다.

\medskip\noindent
$s < 0$일 때 보조정리가 성립함을 관찰하자. 이는 자명한 경우가 아니다.
$H^s(R\Gamma_\mathfrak m(M)^\wedge)$이 음의 $s$에 대해 영이라는 것이
선험적으로 분명하지 않기 때문이다. 그러나 이 소멸은 보조정리
\ref{algebraization-lemma-local-cohomology-derived-completion}에서 확립되었다. 이제
$s \geq 0$에 대해 귀납법으로 보조정리를 증명하겠다.

\medskip\noindent
보조정리 \ref{algebraization-lemma-kill-colimit-weak}의 가정들은 국소 코호몰로지의
보조정리 \ref{local-cohomology-lemma-cd-bound-dim-local}에 의해 만족된다.
$s = 0$인 보조정리는 보조정리 \ref{algebraization-lemma-kill-colimit-weak}과 쌍대화
복합체의 보조정리 \ref{dualizing-lemma-completion-local-H0}에서 따른다.

\medskip\noindent
$s > 0$이고 더 작은 $s$의 값들에 대해 보조정리가 성립한다고 가정하자.
$M' \subset M$을 원소들의 지지집합이
% P07-ERR-0155
$V(I) \cup V(J)$에 포함되는 원소들로 이루어진 부분가군이라 하자. 여기서
어떤 아이디얼 $J$가 주어져 있으며
$V(J) \cap V(I) = \{\mathfrak m\}$이다. 그러면 $M'$은 유한 $A$-가군이다.
다음을 주장한다.
$$
R\Gamma_J(M') \to R\Gamma_\mathfrak m(M')^\wedge
$$
이는 $J$를 어떻게 택하든 동형사상이다. 실제로 이러한 임의의 가군에
대해 짧은 완전열 $0 \to M_1 \oplus M_2 \to M' \to N \to 0$이 존재하며,
$M_1$은 $J$의 어떤 거듭제곱에 의해 소멸되고, $M_2$는 $I$의 어떤
거듭제곱에 의해 소멸되며, $N$은 $\mathfrak m$의 어떤 거듭제곱에 의해
소멸된다. $M_1$의 경우 $R\Gamma_J(M_1) = M_1$이고, $M_1$은 유한
$A$-가군이면서 $I$에 관해 진 완비이므로 $M_1^\wedge = M_1$이다.
따라서 $M_1$에 대해 주장이 성립한다. $M_2$의 경우 $H^i_J(M_2)$는
$I$의 어떤 거듭제곱에 의해 소멸되므로 유도 완비임을 알 수 있다.
% P07-ERR-0156
따라서 $M_2$에 대해서도 주장이 성립한다. 같은 논증에 의해 $N$에
대해서도 주장이 성립하므로 주장이 증명된다. 짧은 완전열
$0 \to M' \to M \to M/M' \to 0$을 고찰하면 $M/M'$에 대해 보조정리를
증명하는 것으로 충분함을 알 수 있다.
% P07-ERR-0157
따라서 $\mathfrak p \in \text{Ass}(M)$이면
$V(\mathfrak p) \cap V(I) \not = \{\mathfrak m\}$, 즉 $\mathfrak p$가
(4)에 나오는 소아이디얼이라고 가정해도 된다.

\medskip\noindent
$J_0$를 보조정리 \ref{algebraization-lemma-kill-colimit-weak}에서와 같이 택하고, 정수
$t > 0$을 $(J_0I)^t$가 $H^s_J(M)$을 소멸시키도록 택하자. 여기서 $J$는
임의의 아이디얼 $J \subset J_0$를 나타내며
$V(J) \cap V(I) = \{\mathfrak m\}$를 만족한다. 보조정리
\ref{algebraization-lemma-kill-completion}의 가정들은 모든
$\mathfrak p \in \text{Ass}(M)$에 대해 만족된다(앞 문단을 보라). 따라서
$\mathfrak a \subset A$를 $H^s(R\Gamma_\mathfrak m(M)^\wedge)$의
소멸자라 하면 이는 $\mathfrak p$에 포함되지 않으며, 이는
$\mathfrak p \in \text{Ass}(M)$에 대해 성립한다. 따라서
$f \in \mathfrak a(J_0I)^t$를 택할 수 있는데, 이는 $M$의 어느 결합
소아이디얼에도 속하지 않으며 $H^s(R\Gamma_\mathfrak m(M)^\wedge)$과
$H^s_J(M)$을 모두 소멸시킨다. 그러면 $f$는 $M$ 위의 영인자가 아니므로
다음 짧은 완전열을 고찰할 수 있다.
$$
0 \to M \xrightarrow{f} M \to M/fM \to 0
$$
$f$의 선택에 의해 다음을 얻는다.
$$
\xymatrix{
H^{s - 1}_J(M) \ar[d] \ar[r] &
H^{s - 1}_J(M/fM) \ar[d] \ar[r] &
H^s_J(M) \ar[d] \ar[r] & 0 \\
H^{s - 1}(R\Gamma_\mathfrak m(M)^\wedge) \ar[r] &
H^{s - 1}(R\Gamma_\mathfrak m(M/fM)^\wedge) \ar[r] &
H^s(R\Gamma_\mathfrak m(M)^\wedge) \ar[r] & 0
}
$$
이는 임의의 $J \subset J_0$ 중
$V(J) \cap V(I) = \{\mathfrak m\}$를 만족하는 것에 대해 성립한다.
따라서 $J$를 $M$, $M/fM$, $s - 1$에 대해 작동하도록 택하면(귀납 가정에
의해 가능하다), 보조정리가 참이라는 결론을 얻는다.
\end{proof}








\section{국소 코호몰로지의 대수화 III}
\label{algebraization-section-bootstrap}

\noindent
이 절에서는 절 \ref{algebraization-section-algebraization-sections-general}과
\ref{algebraization-section-algebraization-punctured}의 내용을 부트스트랩하여 다음
상황에서 더 강한 결과를 얻는다.

\begin{situation}
\label{algebraization-situation-bootstrap}
여기서 $A$는 Noether 환이다. 아이디얼 $I \subset A$, 유한 $A$-가군 $M$,
그리고 특수화에 대해 안정적인 부분집합 $T \subset V(I)$가 주어져 있다.
정수 $s$와 $d$가 주어져 있다. 다음을 가정한다.
\begin{enumerate}
\item[(1)] $A$는 쌍대화 복합체를 갖는다,
\item[(3)] $\text{cd}(A, I) \leq d$이다,
\item[(4)] 소아이디얼들 $\mathfrak p \subset \mathfrak r \subset \mathfrak q$가
$\mathfrak p \not \in V(I)$, $\mathfrak r \in V(I) \setminus T$,
$\mathfrak q \in T$를 만족하면 다음이 성립한다.
$$
\text{depth}_{A_\mathfrak p}(M_\mathfrak p) \geq s
\quad\text{또는}\quad
\text{depth}_{A_\mathfrak p}(M_\mathfrak p) +
\dim((A/\mathfrak p)_\mathfrak q) > d + s
$$
\item[(6)] $\mathfrak q \in T$가 주어졌을 때,
% P07-ERR-0159
$A', \mathfrak m', I', M'$로 통상적인 $I$-진 완비들, 즉
$A_\mathfrak q, \mathfrak qA_\mathfrak q, I_\mathfrak q, M_\mathfrak q$의
완비들을 나타내면 다음이 성립한다.
$$
\text{depth}(M'_{\mathfrak p'}) > s
$$
이는 모든 $\mathfrak p' \in \Spec(A') \setminus V(I')$ 중
$V(\mathfrak p') \cap V(I') = \{\mathfrak m'\}$를 만족하는 것에 대해
성립한다.
\end{enumerate}
\end{situation}

\noindent
다음 보조정리는 상황 \ref{algebraization-situation-bootstrap}에서 (4)의 소아이디얼
삼중열 $\mathfrak p \subset \mathfrak r \subset \mathfrak q$만 살펴보면
충분한 이유를 설명한다. 실제 가정에는 $\mathfrak p$와 $\mathfrak q$만
나타난다.

\begin{lemma}
\label{algebraization-lemma-helper-bootstrap}
상황 \ref{algebraization-situation-bootstrap}에서 $\mathfrak p \subset \mathfrak q$를
$A$의 소아이디얼들이라 하고, $\mathfrak p \not \in V(I)$ 및
$\mathfrak q \in T$를 가정하자.
$\mathfrak r \in V(I) \setminus T$ 중
$\mathfrak p \subset \mathfrak r \subset \mathfrak q$를 만족하는 것이
존재하지 않으면
$\text{depth}(M_\mathfrak p) > s$이다.
\end{lemma}

\begin{proof}
$\mathfrak q' \in T$를 택하되
$\mathfrak p \subset \mathfrak q' \subset \mathfrak q$이고 $T$에는
$\mathfrak p$와 $\mathfrak q'$ 사이에 엄밀히 놓이는 소아이디얼이 없도록
하자. 보조정리를 증명할 때 $\mathfrak q$를 $\mathfrak q'$로 바꾸어도
되며 그렇게 하겠다. 다음으로 $\mathfrak p' \subset A_\mathfrak q$를
$\mathfrak p$에 대응하는 소아이디얼이라 하자. 그러면
$V(\mathfrak p') \cap V(IA_\mathfrak q) = \{\mathfrak q A_\mathfrak q\}$
이다. 이는 보조정리에 나오는 $\mathfrak r$가 존재하지 않기 때문이다.
$A', I', \mathfrak m', M'$를 각각 $I$-진 완비들, 즉
$A_\mathfrak q, I_\mathfrak q, \mathfrak qA_\mathfrak q, M_\mathfrak q$의
완비라 하자.
$A_\mathfrak q \to A'$는 충실하게 평탄하므로(대수학의 보조정리
\ref{algebra-lemma-completion-faithfully-flat}), $\mathfrak p'' \subset A'$를
택하되 $\mathfrak p'$ 위에 놓이고
$\dim(A'_{\mathfrak p''}/\mathfrak p' A'_{\mathfrak p''}) = 0$이 되게 할 수
있다. 그러면 다음을 알 수 있다.
$$
\text{depth}(M'_{\mathfrak p''}) =
\text{depth}((M_\mathfrak q \otimes_{A_\mathfrak q} A')_{\mathfrak p''}) =
\text{depth}(M_\mathfrak p \otimes_{A_\mathfrak p} A'_{\mathfrak p''}) =
\text{depth}(M_\mathfrak p)
$$
이는 $A \to A'$의 평탄성과 $\mathfrak p''$의 선택에서 따른다. 대수학의
보조정리 \ref{algebra-lemma-apply-grothendieck-module}를 보라.
$\mathfrak p''$는 $\mathfrak p'$ 위에 놓이므로
$V(\mathfrak p'') \cap V(I') = \{\mathfrak m'\}$이다. 따라서 상황
\ref{algebraization-situation-bootstrap}의 조건 (6)은
$\text{depth}(M'_{\mathfrak p''}) > s$를 함의하고 증명이 끝난다.
\end{proof}

\noindent
다음의 다소 번거로운 보조정리는 부과할 수 있는 여러 조건 묶음 사이의
관계를 설명한다.

\begin{lemma}
\label{algebraization-lemma-bootstrap-inherited}
상황 \ref{algebraization-situation-bootstrap}에서 다음이 성립한다.
\begin{enumerate}
\item[(E)] $T' \subset T$가 특수화에 대해 안정적인 더 작은 부분집합이면
$A, I, T', M$은 상황 \ref{algebraization-situation-bootstrap}의 가정들을 만족한다,
\item[(F)] $S \subset A$가 곱셈적 부분집합이면
$S^{-1}A, S^{-1}I, T', S^{-1}M$
은 상황 \ref{algebraization-situation-bootstrap}의 가정들을 만족한다. 여기서
$T' \subset V(S^{-1}I)$는 $T$의 역상이다,
\item[(G)] 사중항 $A', I', T', M'$은 상황
\ref{algebraization-situation-bootstrap}의 가정들을 만족한다. 여기서 $A', I', M'$는
통상적인 $I$-진 완비들, 즉 $A, I, M$의 완비들이고
$T' \subset V(I')$는 $T$의 역상이다.
\end{enumerate}
$I \subset \mathfrak a \subset A$를 $V(\mathfrak a) \subset T$를 만족하는
아이디얼이라 하자. 그러면 다음이 성립한다.
\begin{enumerate}
\item[(A)] $I$가 $A$의 Jacobson 근기에 포함되면 보조정리
\ref{algebraization-lemma-kill-colimit-weak-general}과
\ref{algebraization-lemma-kill-colimit-support-general}의 모든 가정이
$A, I, \mathfrak a, M$에 대해 만족된다,
\item[(B)] $A$가 $I$에 관해 완비이면 보조정리
\ref{algebraization-lemma-algebraize-local-cohomology-general}의 가정 중, 만족되지 않을
가능성이 있는 (5)를 제외한 모든 가정이 $A, I, \mathfrak a, M$에 대해
만족된다,
\item[(C)] $A$가 극대 아이디얼 $\mathfrak m = \mathfrak a$를 갖는
국소환이면 보조정리 \ref{algebraization-lemma-kill-colimit-weak}과
\ref{algebraization-lemma-kill-colimit-support}의 모든 가정이 $A, \mathfrak m, I, M$에
대해 성립한다,
\item[(D)] $A$가 극대 아이디얼 $\mathfrak m = \mathfrak a$를 갖는
국소환이고 $I$에 관해 진 완비이면 보조정리
\ref{algebraization-lemma-algebraize-local-cohomology}의 모든 가정이
$A, \mathfrak m, I, M$에 대해 성립한다.
\end{enumerate}
\end{lemma}

\begin{proof}
(E)의 증명. 상황 \ref{algebraization-situation-bootstrap}의 가정 (1), (3), (4), (6)이
% P07-ERR-0160
$A, I, T, M$에 대해 성립함을 증명해야 한다. $T$를 $T'$로 줄이면 가정
(6)은 약해지고 가정 (4)는 강해진다. 그러나
$\mathfrak p \subset \mathfrak r \subset \mathfrak q$이고
$\mathfrak p \not \in V(I)$, $\mathfrak r \in V(I) \setminus T'$,
$\mathfrak q \in T'$라고 하자. 이는 $A, I, T', M$에 대한 가정 (4)에
나오는 경우이다. 그러면
$\mathfrak r \in V(I) \setminus T$인 것을 택할 수 있어서 $A, I, T, M$에
대한 조건 (4)를 적용하거나, 그러한 $\mathfrak r$를 찾을 수 없어서
보조정리 \ref{algebraization-lemma-helper-bootstrap}에 의해
$\text{depth}(M_\mathfrak p) > s$를 얻는다. 이로써 원하는 대로 (4)가
$A, I, T', M$에 대해 성립함이 증명된다.

\medskip\noindent
(F)의 증명. 이는 직접적이므로 세부사항을 생략한다.

\medskip\noindent
(G)의 증명. 상황 \ref{algebraization-situation-bootstrap}의 가정 (1), (3), (4), (6)이
$I$-진 완비들 $A', I', T', M'$에 대해 성립함을 증명해야 한다.
$\Spec(A') \to \Spec(A)$가 동형사상 $V(I') \to V(I)$를 유도함을
기억하자.

\medskip\noindent
가정 (1): 환 $A'$는 쌍대화 복합체를 갖는다. 쌍대화 복합체의 보조정리
\ref{dualizing-lemma-ubiquity-dualizing}을 보라.

\medskip\noindent
가정 (3): $I' = IA'$이므로 이는 국소 코호몰로지의 보조정리
\ref{local-cohomology-lemma-cd-change-rings}에서 따른다.

\medskip\noindent
가정 (4): 소아이디얼들
$\mathfrak p' \subset \mathfrak r' \subset \mathfrak q'$이 $A'$ 안에 있고
$\mathfrak p' \not \in V(I')$, $\mathfrak r' \in V(I') \setminus T'$,
$\mathfrak q' \in T'$를 만족한다고 하자. 그러면 그 상들
$\mathfrak p \subset \mathfrak r \subset \mathfrak q$는 $A$의 스펙트럼에서
$\mathfrak p \not \in V(I)$, $\mathfrak r \in V(I) \setminus T$,
$\mathfrak q \in T$를 만족한다. 따라서 다음이 성립한다.
$$
\text{depth}_{A_\mathfrak p}(M_\mathfrak p) \geq s
\quad\text{또는}\quad
\text{depth}_{A_\mathfrak p}(M_\mathfrak p) +
\dim((A/\mathfrak p)_\mathfrak q) > d + s
$$
이는 $A, I, T, M$에 대한 가정 (4)에서 따른다. 또한
$\text{depth}(M'_{\mathfrak p'}) \geq \text{depth}(M_\mathfrak p)$이고
$\text{depth}(M'_{\mathfrak p'}) +
\dim((A'/\mathfrak p')_{\mathfrak q'}) =
\text{depth}(M_\mathfrak p) +
\dim((A/\mathfrak p)_\mathfrak q)$
가 국소 코호몰로지의 보조정리
\ref{local-cohomology-lemma-change-completion}에 의해 성립한다. 따라서
가정 (4)는 $A', I', T', M'$에 대해 성립한다.

\medskip\noindent
가정 (6): $\mathfrak q' \in T'$를 소아이디얼 $\mathfrak q \in T$ 위에
놓이도록 택하자. 그러면 $A'_{\mathfrak q'}$와
$A_\mathfrak q$의 $I$-진 완비들은 서로 동형이고, $M_\mathfrak q$와
$M'_{\mathfrak q'}$에 대해서도 마찬가지이다. 따라서 $A', I', T', M'$에
대한 가정 (6)은 $A, I, T, M$에 대한 가정 (6)과 동치이다.

\medskip\noindent
(A)의 증명. 보조정리 \ref{algebraization-lemma-kill-colimit-weak-general}과
\ref{algebraization-lemma-kill-colimit-support-general}의 조건 (1), (2), (3), (4), (6)을
$(A, I, \mathfrak a, M)$에 대해 확인해야 한다. 주의: 이 보조정리들의
명제에 나오는 집합 $T$는 위의 집합 $T$와 같지 않다.

\medskip\noindent
조건 (1): 상황 \ref{algebraization-situation-bootstrap}에서 $A$가 쌍대화 복합체를
갖는다고 가정했으므로 성립한다.

\medskip\noindent
조건 (2): 아무 조건도 없다.

\medskip\noindent
조건 (3): $\mathfrak p \subset A$가
$V(\mathfrak p) \cap V(I) \subset V(\mathfrak a)$를 만족한다고 하자.
$I$는 $A$의 Jacobson 근기에 포함되므로
$V(\mathfrak p) \cap V(I) \not = \emptyset$임을 알 수 있다.
$\mathfrak q \in V(\mathfrak p) \cap V(I)$를 일반점이라 하자.
$\text{cd}(A_\mathfrak q, I_\mathfrak q) \leq d$이고(국소 코호몰로지의
보조정리 \ref{local-cohomology-lemma-cd-local})
$V(\mathfrak p A_\mathfrak q) \cap V(I_\mathfrak q) =
\{\mathfrak q A_\mathfrak q\}$이므로
$\dim((A/\mathfrak p)_\mathfrak q) \leq d$를 얻는다. 국소 코호몰로지의
보조정리 \ref{local-cohomology-lemma-cd-bound-dim-local}를 보라. 이로써
(3)이 증명된다.

\medskip\noindent
조건 (4): $\mathfrak p \not \in V(I)$이고
$\mathfrak q \in V(\mathfrak p) \cap V(\mathfrak a)$라고 하자. 다음을
보이면 충분하다.
$$
\text{depth}_{A_\mathfrak p}(M_\mathfrak p) \geq s
\quad\text{또는}\quad
\text{depth}_{A_\mathfrak p}(M_\mathfrak p) +
\dim((A/\mathfrak p)_\mathfrak q) > d + s
$$
소아이디얼 $\mathfrak p \subset \mathfrak r \subset \mathfrak q$ 중
$\mathfrak r \in V(I) \setminus T$인 것이 존재하면 이는 상황
\ref{algebraization-situation-bootstrap}의 가정 (4)에서 즉시 따른다. 존재하지 않으면
$\text{depth}(M_\mathfrak p) > s$가 보조정리
\ref{algebraization-lemma-helper-bootstrap}에 의해 성립한다.

\medskip\noindent
조건 (6): $\mathfrak p \not \in V(I)$가
$V(\mathfrak p) \cap V(I) \subset V(\mathfrak a)$를 만족한다고 하자.
$I$는 $A$의 Jacobson 근기에 포함되므로
$V(\mathfrak p) \cap V(I) \not = \emptyset$임을 알 수 있다.
$\mathfrak q \in V(\mathfrak p) \cap V(I) \subset V(\mathfrak a)$를 택하자.
$\mathfrak p \subset \mathfrak r \subset \mathfrak q$를 만족하는
소아이디얼 중 $\mathfrak r \in V(I) \setminus T$인 것은 존재하지 않음이
분명하다. 보조정리 \ref{algebraization-lemma-helper-bootstrap}에 의해
$\text{depth}(M_\mathfrak p) > s$이고, 이로써 (6)이 증명된다.

\medskip\noindent
(B)의 증명. 보조정리 \ref{algebraization-lemma-algebraize-local-cohomology-general}의
조건 (1), (2), (3), (4)를 확인해야 한다. 주의: 이 보조정리의 명제에
나오는 집합 $T$는 위의 집합 $T$와 같지 않다.

\medskip\noindent
조건 (1): $A$는 완비이고 쌍대화 복합체를 가지므로 성립한다.

\medskip\noindent
조건 (2): 아무 조건도 없다.

\medskip\noindent
조건 (3): 이는 상황 \ref{algebraization-situation-bootstrap}의 가정 (3)과 같다.

\medskip\noindent
조건 (4): 이는 (A)에서 증명한 보조정리
\ref{algebraization-lemma-kill-colimit-weak-general}의 가정 (4)와 같다.

\medskip\noindent
(C)의 증명. 국소 경우에 보조정리 \ref{algebraization-lemma-kill-colimit-weak}과
\ref{algebraization-lemma-kill-colimit-support}의 가정들은 보조정리
\ref{algebraization-lemma-kill-colimit-weak-general}과
\ref{algebraization-lemma-kill-colimit-support-general}의 가정들과 같고, 이들이 성립함을
(A)에서 증명했으므로 참이다.

\medskip\noindent
(D)의 증명. 보조정리 \ref{algebraization-lemma-algebraize-local-cohomology}의 가정들은
보조정리 \ref{algebraization-lemma-algebraize-local-cohomology-general}의 가정 (1), (2),
(3), (4)와 같고, 이들이 성립함을 (B)에서 증명했으므로 참이다.
\end{proof}

\begin{lemma}
\label{algebraization-lemma-algebraize-local-cohomology-bis}
상황 \ref{algebraization-situation-bootstrap}에서 $A$가 극대 아이디얼 $\mathfrak m$을 갖는
국소환이고 $T = \{\mathfrak m\}$라고 가정하자. 그러면
$H^i_\mathfrak m(M) \to \lim H^i_\mathfrak m(M/I^nM)$
은 $i \leq s$에 대해 동형사상이고 이 가군들은 $I$의 어떤 거듭제곱에
의해 소멸된다.
\end{lemma}

\begin{proof}
$A', I', \mathfrak m', M'$를 통상적인 $I$-진 완비들, 즉
$A, I, \mathfrak m, M$의 완비들이라 하자. 다음을 상기하자.
$H^i_\mathfrak m(M) \otimes_A A' = H^i_{\mathfrak m'}(M')$
이는 $A \to A'$의 평탄성과 쌍대화 복합체의 보조정리
\ref{dualizing-lemma-torsion-change-rings}에서 따른다.
$H^i_\mathfrak m(M)$은 $\mathfrak m$-멱 꼬임이므로
$H^i_\mathfrak m(M) = H^i_\mathfrak m(M) \otimes_A A'$이다. 대수학 심화의
보조정리 \ref{more-algebra-lemma-neighbourhood-equivalence}를 보라. 따라서
$H^i_\mathfrak m(M) = H^i_{\mathfrak m'}(M')$이다. 완전히 같은 논증은
$H^i_\mathfrak m(M/I^nM) =  H^i_{\mathfrak m'}(M'/(I')^nM')$
가 모든 $n$과 $i$에 대해 성립함을 보여 준다.

\medskip\noindent
보조정리 \ref{algebraization-lemma-algebraize-local-cohomology},
\ref{algebraization-lemma-kill-colimit-weak}, 그리고
\ref{algebraization-lemma-kill-colimit-support}
는 보조정리 \ref{algebraization-lemma-bootstrap-inherited}의 (C), (D)에 의해
$A', \mathfrak m', I', M'$에 적용된다. 따라서 동형사상
$$
H^i_{\mathfrak m'}(M') \longrightarrow H^i(R\Gamma_{\mathfrak m'}(M')^\wedge)
$$
을 $i \leq s$에 대해 얻는다. 여기서 ${}^\wedge$는 유도 $I'$-진 완비이고
이 가군들은 $I'$의 어떤 거듭제곱에 의해 소멸된다. 보조정리
\ref{algebraization-lemma-local-cohomology-derived-completion}에 의해 동형사상들
% P07-ERR-0161
$$
H^i_{\mathfrak m'}(M') \longrightarrow
\lim H^i_{\mathfrak m'}(M'/(I')^nM'))
$$
을 $i \leq s$에 대해 얻는다. 이를 이미 확립한 $A$ 위의 국소 코호몰로지와
비교하면 보조정리가 참이라는 결론을 얻는다.
\end{proof}

\begin{lemma}
\label{algebraization-lemma-bootstrap-bis-bis}
$I \subset \mathfrak a$를 Noether 환 $A$의 아이디얼들이라 하자. $M$을
유한 $A$-가군이라 하자. $s$와 $d$를 정수라 하자. 다음을 가정하자.
\begin{enumerate}
\item[(a)] $A$는 쌍대화 복합체를 갖는다,
\item[(b)] $\text{cd}(A, I) \leq d$이다,
\item[(c)] $\mathfrak p \not \in V(I)$이고
$\mathfrak q \in V(\mathfrak p) \cap V(\mathfrak a)$이면
$\text{depth}_{A_\mathfrak p}(M_\mathfrak p) > s$이거나
$\text{depth}_{A_\mathfrak p}(M_\mathfrak p) +
\dim((A/\mathfrak p)_\mathfrak q) > d + s$이다.
\end{enumerate}
그러면 $A, I, V(\mathfrak a), M, s, d$는 상황
\ref{algebraization-situation-bootstrap}에서와 같다.
\end{lemma}

\begin{proof}
상황 \ref{algebraization-situation-bootstrap}의 가정 (1), (3), (4), (6)이 성립함을 보여야
한다. (a) $\Rightarrow$ (1), (b) $\Rightarrow$ (3), (c) $\Rightarrow$ (4)는
분명하다. 다음 문단에서 (6)이 성립함을 보여 증명을 끝낸다.

\medskip\noindent
$\mathfrak q \in V(\mathfrak a)$라 하자. $A', I', \mathfrak m', M'$로
$I$-진 완비들, 즉
$A_\mathfrak q, I_\mathfrak q, \mathfrak qA_\mathfrak q, M_\mathfrak q$의
완비들을 나타내자. $\mathfrak p' \subset A'$를 비극대
소아이디얼이라 하고 $V(\mathfrak p') \cap V(I') = \{\mathfrak m'\}$라고
하자. 국소 코호몰로지의 보조정리
\ref{local-cohomology-lemma-cd-bound-dim-local}에 의해 이로부터
$\dim(A'/\mathfrak p') \leq d$가 따른다. $\mathfrak p \subset A$로
$\mathfrak p'$의 상을 나타내자. 다음이 성립한다.
$\text{depth}(M'_{\mathfrak p'}) \geq \text{depth}(M_\mathfrak p)$이고
$\text{depth}(M'_{\mathfrak p'}) +
\dim(A'/\mathfrak p') =
\text{depth}(M_\mathfrak p) +
\dim((A/\mathfrak p)_\mathfrak q)$
이는 국소 코호몰로지의 보조정리
\ref{local-cohomology-lemma-change-completion}에서 따른다. 가정 (c)에 의해
다음이 성립하거나
$\text{depth}(M'_{\mathfrak p'}) \geq \text{depth}(M_\mathfrak p) > s$
하여 끝나거나, 다음이 성립한다.
$\text{depth}(M'_{\mathfrak p'}) +
\dim(A'/\mathfrak p') > s + d$. 후자의 경우
$\text{depth}(M'_{\mathfrak p'}) > s$가 따르는데, 이는 이미 보인 부등식
$\dim(A'/\mathfrak p') \leq d$ 때문이다. 어느 경우든 원하는 것을 얻는다.
\end{proof}

\begin{lemma}
\label{algebraization-lemma-bootstrap}
상황 \ref{algebraization-situation-bootstrap}에서 역계
$\{H^i_T(I^nM)\}_{n \geq 0}$은 $i \leq s$에 대해 프로-영이다. 더욱이
정수 $m_0$가 존재하여, 모든 $m \geq m_0$에 대해 정수
$m'(m) \geq m$을 택할 수 있고, $k \geq m'(m)$일 때
$H^{s + 1}_T(I^kM) \to H^{s + 1}_T(I^mM)$
의 상은 $H^{s + 1}_T(I^{m_0}M)$으로 단사적으로 사상된다.
\end{lemma}

\begin{proof}
$m$을 고정하자. $\mathfrak q \in T$라 하자. 보조정리
\ref{algebraization-lemma-bootstrap-inherited}와
\ref{algebraization-lemma-algebraize-local-cohomology-bis}
에 의해 다음을 알 수 있다.
$$
H^i_\mathfrak q(M_\mathfrak q)
\longrightarrow
\lim H^i_\mathfrak q(M_\mathfrak q/I^nM_\mathfrak q)
$$
이는 $i \leq s$에 대해 동형사상이다. 역계
$\{H^i_\mathfrak q(I^nM_\mathfrak q)\}_{n \geq 0}$과
$\{H^i_\mathfrak q(M/I^nM)\}_{n \geq 0}$
은 모든 $i$에 대해 Mittag--Leffler 조건을 만족한다. 보조정리
\ref{algebraization-lemma-ML-local}을 보라. 따라서 긴 완전열의 역계
$$
0 \to H^0_\mathfrak q(I^nM_\mathfrak q) \to
H^0_\mathfrak q(M_\mathfrak q) \to
H^0_\mathfrak q(M_\mathfrak q/I^nM_\mathfrak q) \to
H^1_\mathfrak q(I^nM_\mathfrak q) \to
H^1_\mathfrak q(M_\mathfrak q) \to \ldots
$$
을 살펴보면, 정수 $m'(m, \mathfrak q) \geq m$이 존재하여 모든
$k \geq m'(m, \mathfrak q)$에 대해 사상
$H^i_\mathfrak q(I^kM_\mathfrak q) \to H^i_\mathfrak q(I^mM_\mathfrak q)$
이 $i \leq s$에 대해 영이고, 사상
$H^{s + 1}_\mathfrak q(I^kM_\mathfrak q) \to
H^{s + 1}_\mathfrak q(I^mM_\mathfrak q)$
의 상은 $k \geq m'(m, \mathfrak q)$에 무관하며
$H^{s + 1}_\mathfrak q(M_\mathfrak q)$으로 단사적으로 사상된다는 결론을
얻는다(일부 세부사항은 생략한다).

\medskip\noindent
$m'(m, \mathfrak q)$를 $\mathfrak q \in T$에 무관하게 택할 수 있음을
보인다고 하자. 그러면 보조정리는 국소 코호몰로지의 보조정리
\ref{local-cohomology-lemma-zero}와
\ref{local-cohomology-lemma-essential-image}에서 즉시 따른다.

\medskip\noindent
$\omega_A^\bullet$을 쌍대화 복합체라 하자.
$\delta : \Spec(A) \to \mathbf{Z}$를 이에 대응하는 차원함수라 하자.
$\delta$는 유한 개의 값만 취함을 상기하자. 쌍대화 복합체의 보조정리
\ref{dualizing-lemma-universally-catenary}를 보라. 주장: 각
$d \in \mathbf{Z}$에 대해 정수 $m'(m, \mathfrak q)$를
$\mathfrak q \in T$ 중 $\delta(\mathfrak q) = d$를 만족하는 것에 무관하게
택할 수 있다. 위에서 말한 바에 의해 이 주장은 분명히 보조정리를
함의한다.

\medskip\noindent
$\mathfrak q \in T$를 $\delta(\mathfrak q) = d$가 되도록 택하자. Ext 가군
$$
E(n, j) = \text{Ext}^j_A(I^nM, \omega_A^\bullet)
$$
들을 생각하자. 우리가 사용할 핵심 성질은 이들이 유한 $A$-가군이라는
점이다. 쌍대화 복합체에 결부된 차원함수의 정의에 의해
$(\omega_A^\bullet)_\mathfrak q[-d]$는 $A_\mathfrak q$의 정규화된 쌍대화
복합체이다. 쌍대화 복합체의 절
\ref{dualizing-section-dimension-function}을 보라. 국소 쌍대성 정리,
즉 쌍대화 복합체의 보조정리
\ref{dualizing-lemma-special-case-local-duality}에 따르면
$\mathfrak qA_\mathfrak q$-진 완비를 $E(n, -d - i)_\mathfrak q$에 취하면
$H^i_\mathfrak q(I^nM_\mathfrak q)$의 Matlis 쌍대를 얻는다. 따라서 첫 문단에서
$m'(m, \mathfrak q)$를 택한 결과는 $i \leq s$에 대해 다음을 말한다.
$k \geq m'(m, \mathfrak q)$이고 $j \geq -d - s$일 때 사상
$$
E(m, j)_\mathfrak q \to E(k, j)_\mathfrak q
$$
은 영이다. 이 가군들은 유한이고 가능한 $j$ 중 유한 개에 대해서만 영이
아니므로(작은 세부사항은 생략한다), 열린 근방 $W \subset \Spec(A)$를
$\mathfrak q$의 근방으로 찾아서 다음이 성립하게 할 수 있다.
$$
E(m, j)_{\mathfrak q'} \to E(m'(m, \mathfrak q), j)_{\mathfrak q'}
$$
이는 $j \geq -d - s$이고 모든 $\mathfrak q' \in W$에 대해 영이다. 그러면
물론 사상 $E(m, j)_{\mathfrak q'} \to E(k, j)_{\mathfrak q'}$도
$k \geq m'(m, \mathfrak q)$에 대해 영이다.

\medskip\noindent
$i = s + 1$과 이에 대응하는 $j = - d - s - 1$에 대해서는 국소 쌍대성과
첫 문단의 결과로부터 다음을 얻는다.
$$
K_{k, \mathfrak q} =
\Ker(E(m, -d - s - 1)_\mathfrak q \to E(k, -d - s - 1)_\mathfrak q)
$$
이는 $k \geq m'(m, \mathfrak q)$에 무관하며, 또한
$$
E(0, -d - s - 1)_\mathfrak q \to
E(m, -d - s - 1)_\mathfrak q/K_{m'(m, \mathfrak q), \mathfrak q}
$$
는 전사이다. $k \geq m'(m, \mathfrak q)$에 대해 다음과 같이 놓자.
$$
K_k = \Ker(E(m, -d - s - 1) \to E(k, -d - s - 1))
$$
$K_k$는 유한 가군 $E(m, -d - s - 1)$의 부분가군들로 이루어진 증가열이므로,
$m'(m, \mathfrak q)$를 조금 늘리는 대가로
$K_{m'(m, \mathfrak q)} = K_k$가 $k \geq m'(m, \mathfrak q)$에 대해
성립한다고 가정해도 된다. 필요하면 $W$를 더 줄여서 다음도 가정할 수
있다.
$$
E(0, -d - s - 1)_{\mathfrak q'} \to
E(m, -d - s - 1)_{\mathfrak q'}/K_{m'(m, \mathfrak q), \mathfrak q'}
$$
이는 모든 $\mathfrak q' \in W$에 대해 전사이다. 앞에서처럼 이 가군들이
유한이고 사상이 $\mathfrak q$에서 국소화한 뒤 전사라는 사실을 사용한다.

\medskip\noindent
임의의 부분집합, 특히
$T_d = \{\mathfrak q \in T \text{이고 }\delta(\mathfrak q) = d\}$는 Noether
위상공간 $\Spec(A)$로부터 유도된 위상에 관해 Noether이고 따라서
준콤팩트하다. 위에서 모든
$\mathfrak q \in T_d$에 대해 열린 근방 $W$가 존재하여
$m'(m, \mathfrak q)$가 모든 $\mathfrak q' \in T_d \cap W$에 대해
작동함을 보았다. 따라서 정수 $m'(m, d)$를 찾아서 모든
$\mathfrak q \in T_d$에 대해 다음이 성립하게 할 수 있다.
$$
E(m, j)_\mathfrak q \to E(m'(m, d), j)_\mathfrak q
$$
이는 $j \geq -d - s$에 대해 영이다. 또한
$K_{m'(m, d)} = \Ker(E(m, -d - s - 1) \to E(m'(m, d), -d - s - 1))$
로 놓으면 다음이 성립한다.
$$
K_{m'(m, d), \mathfrak q} =
\Ker(E(m, -d - s - 1)_{\mathfrak q} \to E(k, -d - s - 1)_{\mathfrak q})
$$
이는 모든 $k \geq m'(m, d)$에 대해 성립하고, 사상
$$
E(0, -d - s - 1)_\mathfrak q \to
E(m, -d - s - 1)_\mathfrak q/K_{m'(m, d), \mathfrak q}
$$
은 전사이다. 국소 쌍대성 정리를 반대 방향으로 다시 사용하면 주장이
옳다는 결론을 얻는다. 이로써 증명이 끝난다.
\end{proof}

\begin{lemma}
\label{algebraization-lemma-final-bootstrap}
상황 \ref{algebraization-situation-bootstrap}에서 다음을 만족하는 정수 $m_0 \geq 0$이
존재한다.
\begin{enumerate}
\item $H^i_T(M) \to H^i_T(M/I^nM)$는 모든 $i \leq s$와
$n \geq m_0$에 대해 단사이고,
\item $\{H^i_T(M/I^nM)\}_{n \geq 0}$
$i < s$에 대해 Mittag--Leffler 조건을 만족하고,
\item $\{H^i_T(I^{m_0}M/I^nM)\}_{n \geq m_0}$
$i \leq s$에 대해 Mittag--Leffler 조건을 만족하고,
\item $H^i_T(M) \to \lim H^i_T(M/I^nM)$
$i < s$에 대해 동형이고,
% P07-ERR-0162: 아래 고정된 차수 s의 사상 뒤에 사용되지 않는 i에 대한 범위가 붙어 있다.
\item $H^s_T(I^{m_0}M) \to \lim H^s_T(I^{m_0}M/I^nM)$
$i \leq s$에 대해 동형이고,
\item $H^s_T(M) \to \lim H^s_T(M/I^nM)$는
단사이며 그 여핵은 $I^{m_0}$에 의해 소멸되고,
\item $R^1\lim H^s_T(M/I^nM)$는 $I^{m_0}$에 의해 소멸된다.
\end{enumerate}
\end{lemma}

\begin{proof}
(1)을 보이기 위해 $0 \leq i \leq s$라 하고, 정수 $m > n > 0$을 택하여
$H^i_T(I^mM) \to H^i_T(I^nM)$가 영이 되게 한다.
% P07-ERR-0163: 원문의 ``this is possibly''는 ``this is possible''이어야 한다.
이는 보조정리 \ref{algebraization-lemma-bootstrap}에 의해 가능하다. 행들이 완전한 가환 그림
$$
\xymatrix{
H^i_T(I^nM) \ar[r] &
H^i_T(M) \ar[r] &
H^i_T(M/I^nM) \\
H^i_T(I^mM) \ar[r] \ar[u]^0 &
H^i_T(M) \ar[r] \ar[u]^1 &
H^i_T(M/I^mM) \ar[u] \\
}
$$
은 사상 $H^i_T(M) \to H^i_T(M/I^mM)$가 단사임을 보인다. 따라서 충분히 큰
$m_0$마다 (1)이 성립한다.

\medskip\noindent
긴 완전열
$$
0 \to H^0_T(I^nM) \to H^0_T(M) \to
H^0_T(M/I^nM) \to H^1_T(I^nM) \to
H^1_T(M) \to \ldots
$$
을 생각하자. 이 열과 보조정리 \ref{algebraization-lemma-bootstrap}로부터 (2)와 (4)가 따른다.

\medskip\noindent
$m_0$과 $m'(-)$를 보조정리 \ref{algebraization-lemma-bootstrap}에서와 같이 잡는다.
$m \geq m_0$에 대해 긴 완전열
$$
H^s_T(I^mM) \to H^s_T(I^{m_0}M) \to
H^s_T(I^{m_0}M/I^mM) \to H^{s + 1}_T(I^mM) \to
% P07-ERR-0164: 다음 고정 항의 차수 1은 완전열상 s + 1이어야 하는 것으로 보인다.
H^1_T(I^{m_0}M)
$$
을 생각하자. 그러면 $k \geq m'(m)$에 대해
$H^{s + 1}_T(I^kM) \to H^{s + 1}_T(I^mM)$
의 상은 $H^{s + 1}_T(I^{m_0}M)$으로 단사적으로 사상된다. 따라서
$H^s_T(I^{m_0}M/I^kM) \to H^s_T(I^{m_0}M/I^mM)$
의 상은 $H^{s + 1}_T(I^mM)$에서 영으로 사상되며, 이는 모든
$k \geq m'(m)$에 대해 성립한다.
그러므로 (3)과 (5)가 성립한다.

\medskip\noindent
짧은 완전열
$0 \to I^{m_0}M \to M \to M/I^{m_0} M \to 0$과
$0 \to I^{m_0}M/I^nM \to M/I^nM \to M/I^{m_0} M \to 0$을 생각하자.
그림
$$
\xymatrix{
H^{s - 1}_T(M/I^{m_0}M) \ar[r] &
\lim H^s_T(I^{m_0}M/I^nM) \ar[r] &
\lim H^s_T(M/I^nM) \ar[r] &
H^s_T(M/I^{m_0}M) \\
H^{s - 1}_T(M/I^{m_0}M) \ar[r] \ar@{=}[u] &
H^s_T(I^{m_0}M) \ar[r] \ar[u]_{\cong} &
H^s_T(M) \ar[r] \ar[u] &
H^s_T(M/I^{m_0}M) \ar@{=}[u]
}
$$
을 얻으며 그 아래 행은 완전하다. 위 행도 Homology의 보조정리
\ref{homology-lemma-apply-Mittag-Leffler}에 의해
(가운데 두 자리에서) 완전하다. 이제 (6)이 따른다.

\medskip\noindent
$B_n = H^s_T(M/I^nM)$라 쓰자. $A_n \subset B_n$을
$H^s_T(I^{m_0}M/I^nM) \to H^s_T(M/I^nM)$의 상이라 하자.
그러면 (3)과 Homology의 보조정리 \ref{homology-lemma-Mittag-Leffler}에 의해
$(A_n)$은 Mittag--Leffler 조건을 만족한다. 또한 $C_n = B_n/A_n$은
$I^{m_0}$에 의해 소멸된다. 따라서 $R^1\lim B_n \cong R^1\lim C_n$도
$I^{m_0}$에 의해 소멸되고 (7)을 얻는다.
\end{proof}

\begin{theorem}
\label{algebraization-theorem-final-bootstrap}
상황 \ref{algebraization-situation-bootstrap}에서 $i \leq s$에 대해 역계
$\{H^i_T(M/I^nM)\}_{n \geq 0}$은 값 $H^i_T(M)$을 갖는 본질적 상수계이다.
특히 $\{H^i_T(M/I^nM)\}_{n \geq 0}$은 Mittag--Leffler 조건을 만족하고,
$H^i_T(M)$은 $I$의 어떤 거듭제곱에 의해 소멸되며,
$H^i_T(M) = \lim H^i_T(M/I^nM)$이다.
\end{theorem}

\begin{proof}
$0 \leq i \leq s$라 하자.
$\{H^i_T(M/I^nM)\}_{n \geq 0}$이 Mittag--Leffler 조건을 만족하고
$H^i_T(M) = \lim H^i_T(M/I^nM)$임을 보인다면, 보조정리
\ref{algebraization-lemma-final-bootstrap}에서 보인 $H^i_T(M) \to H^i_T(M/I^nM)$의
$n \gg 0$에서의 단사성으로부터
$\{H^i_T(M/I^nM)\}_{n \geq 0}$이 값 $H^i_T(M)$을 갖는 본질적 상수계임이
따른다. 작은 세부사항은 생략한다.

\medskip\noindent
보조정리 \ref{algebraization-lemma-final-bootstrap}에 의해 이미
$\{H^i_T(M/I^nM)\}_{n \geq 0}$이 Mittag--Leffler 조건을 만족하고
$H^i_T(M) = \lim H^i_T(M/I^nM)$가 $i < s$에 대해 성립함을 안다.
또한 이 명제들이 $i = s$와 가군 $I^{m_0}M$에 대해 성립하도록 하는
어떤 $m_0 > 0$이 있음을 안다.
% P07-ERR-0165: 다음 두 문장의 원문에는 불필요한 ``of''와 수 일치 오류가 있다.
증명을 마치기 위해 이제 이 명제들이 실제로 $i = s$와 $M$에 대해
성립함을 보이겠다.

\medskip\noindent
$M' = H^0_I(M)$ 및 $M'' = M/M'$라 놓으면 짧은 완전열
$$
0 \to M' \to M \to M'' \to 0
$$
을 얻는다. 또한 쌍대화 복합체의 보조정리
\ref{dualizing-lemma-divide-by-torsion}에 의해 $M''$은
% P07-ERR-0166: 아래 원문의 M'은 앞서 정의한 몫가군 M''이어야 한다.
$H^0_I(M') = 0$을 만족한다. Artin--Rees 정리(대수학의 보조정리
\ref{algebra-lemma-Artin-Rees})에 의해 짧은 완전열
$$
0 \to M' \to M/I^n M \to M''/I^n M'' \to 0
$$
을 충분히 큰 $n$에 대해 얻는다. 긴 완전열
$$
H^s_T(M') \to
H^s_T(M/I^nM) \to
H^s_T(M''/I^nM'') \to
H^{s + 1}_T(M')
$$
을 살펴보면, 역계 $\{H^s_T(M''/I^nM'')\}_{n \geq 0}$에 대해
Mittag--Leffler 조건이 성립할 경우 역계
$\{H^s_T(M/I^nM)\}_{n \geq 0}$에 대해서도 이 조건이 성립함을 쉽게 알 수 있다.
(군들의 역계 $G_n$에 대한 Mittag--Leffler 조건은 충분히 큰 $n$에서의
역계 값들에만 의존함에 유의하라.) 더욱이 열
$$
H^s_T(M') \to
\lim H^s_T(M/I^nM) \to
\lim H^s_T(M''/I^nM'') \to
H^{s + 1}_T(M')
$$
은 필요한 자리에서 Mittag--Leffler 조건이 성립하므로 완전하다. 호몰로지의 보조정리
\ref{homology-lemma-apply-Mittag-Leffler}를 보라. 따라서
$H^s_T(M'') \to \lim H^s_T(M''/I^nM'')$가 동형이면
$H^s_T(M) \to \lim H^s_T(M/I^nM)$도 오항 보조정리(호몰로지의 보조정리
\ref{homology-lemma-five-lemma})에 의해 동형이다. 이로써 다음 문단에서
다루는 경우로 환원된다.

\medskip\noindent
$H^0_I(M) = 0$이라 가정하자. $f_1, \ldots, f_r$를 $I^{m_0}$의
생성원으로 택하되, 여기서 $m_0$은 보조정리
\ref{algebraization-lemma-final-bootstrap}에서 $M$에 대해 얻은 정수이다.
이제 완전열
$$
0 \to M \xrightarrow{f_1, \ldots, f_r}
(I^{m_0}M)^{\oplus r} \to Q \to 0
$$
로 $Q$를 정의한다. 몇 가지를 관찰하자. 첫 사상은 바로
$H^0_I(M) = 0$이기 때문에 단사이다. 이 단사의 여핵 $Q$는 유한
$A$-가군이며, 모든 $1 \leq j \leq r$에 대해
$Q_{f_j} \cong (M_{f_j})^{\oplus r - 1}$이다. 특히 소아이디얼
$\mathfrak p \subset A$가 $\mathfrak p \not \in V(I)$를 만족하면
$Q_\mathfrak p \cong (M_\mathfrak p)^{\oplus r - 1}$이다. 마찬가지로
$\mathfrak q \in T$와 $\mathfrak p' \subset A' = (A_\mathfrak q)^\wedge$가
주어지고 후자가 $V(IA')$에 포함되지 않으면
$Q'_{\mathfrak p'} \cong (M'_{\mathfrak p'})^{\oplus r - 1}$이며, 여기서
$Q' = (Q_\mathfrak q)^\wedge$이고 $M' = (M_\mathfrak q)^\wedge$이다.
따라서 상황 \ref{algebraization-situation-bootstrap}의 조건들은 $A, I, T, Q$에 대해
성립한다. ($Q$에 영이 아닌 $H^0_I(Q)$가 있을 수 있지만 이는 문제가 되지 않는다.)

\medskip\noindent
임의의 $n \geq 0$에 대해 $F^nM = M \cap I^n(I^{m_0}M)^{\oplus r}$라 놓으면
짧은 완전열
$$
0 \to F^nM \to I^n(I^{m_0}M)^{\oplus r} \to I^nQ \to 0
$$
을 얻는다. Artin--Rees 정리(대수학의 보조정리
\ref{algebra-lemma-Artin-Rees})에 의해 $c \geq 0$이 존재하여
$I^n M \subset F^nM \subset I^{n - c}M$가 모든 $n \geq c$에 대해 성립한다.
$m_0$을 정수로, $m'(m)$을 $m \geq m_0$에 대해 정의된 함수로 잡되,
이들은 보조정리 \ref{algebraization-lemma-bootstrap}를 $M$에 적용하여 얻은 것이라 하자.
정수 $m_0$은 위에서 택한 정수 $m_0$과 같음에 유의하라(이를 확인할 필요는
없으며, 원한다면 두 정수의 최댓값을 택하면 된다). 마지막으로 보조정리
\ref{algebraization-lemma-bootstrap}를 $Q$에 적용하면, 모든 정수 $m$에 대해 정수
$m''(m) \geq m$이 존재하여 $H^s_T(I^kQ) \to H^s_T(I^mQ)$가 모든
$k \geq m''(m)$에 대해 영이 된다.

\medskip\noindent
$m \geq m_0$을 고정하고 $k \geq m'(m''(m + c))$를 택하자.
$\xi \in H^{s + 1}_T(I^kM)$를 택하되, 그 상이 $H^{s + 1}_T(M)$에서
영이 되게 한다. $\xi$가 $H^{s + 1}_T(I^mM)$에서 영으로 사상됨을
보이려 한다. 실제로 이는 보조정리 \ref{algebraization-lemma-final-bootstrap}의 증명에서와
꼭 같이 $\{H^s_T(M/I^nM)\}_{n \geq 0}$이 Mittag--Leffler 조건을
만족함을 보일 것이다. 논증을 시각화하는 데 도움이 되는 그림은 다음과 같다.
$$
\xymatrix{
&
H^{s + 1}_T(I^kM) \ar[r] \ar[d] &
H^{s + 1}_T(I^k(I^{m_0}M)^{\oplus r}) \ar[d] &
\\
H^s_T(I^{m''(m + c)}Q) \ar[r]_-\delta \ar[d] &
H^{s + 1}_T(F^{m''(m + c)}M) \ar[r] \ar[d] &
H^{s + 1}_T(I^{m''(m + c)}(I^{m_0}M)^{\oplus r}) \\
H^s_T(I^{m + c}Q) \ar[r] &
H^{s + 1}_T(F^{m + c}M) \ar[d] &
\\
&
H^{s + 1}_T(I^mM)
}
$$
$\xi$의 $H^{s + 1}_T(I^k(I^{m_0}M)^{\oplus r})$에서의 상은
$H^{s + 1}_T((I^{m_0}M)^{\oplus r})$에서 영으로 사상되므로,
$H^{s + 1}_T(I^{m''(m + c)}(I^{m_0}M)^{\oplus r})$에서도 영으로
사상된다. 이는 $m'(-)$의 선택에 따른다. 따라서 상
$\xi' \in H^{s + 1}_T(F^{m''(m + c)}M)$은
$H^{s + 1}_T(I^{m''(m + c)}(I^{m_0}M)^{\oplus r})$에서 영으로 사상되고,
그러므로 $\xi' = \delta(\eta)$이다. 여기서
$\eta \in H^s_T(I^{m''(m + c)}Q)$이다. $m''(-)$의 선택에 의해 $\eta$는
$H^s_T(I^{m + c}Q)$에서 영으로 사상된다. 이는 다시 $\xi'$가
$H^{s + 1}_T(F^{m + c}M)$에서 영으로 사상됨을 뜻한다.
$F^{m + c}M \subset I^mM$이므로 결론을 얻는다.

\medskip\noindent
마지막으로 극한에 관한 명제를 증명하자. 짧은 완전열
$$
0 \to M/F^nM \to (I^{m_0}M)^{\oplus r}/I^n (I^{m_0}M)^{\oplus r}
\to Q/I^nQ \to 0
$$
을 생각하자. 해당 역계들이 프로-동형이므로
$\lim H^s_T(M/I^nM) = \lim H^s_T(M/F^nM)$이다. 가환 그림
$$
\xymatrix{
H^{s - 1}_T(Q) \ar[r] \ar[d] &
\lim H^{s - 1}_T(Q/I^nQ) \ar[d] \\
H^s_T(M) \ar[r] \ar[d] &
\lim H^s_T(M/I^nM) \ar[d] \\
H^s_T((I^{m_0}M)^{\oplus r}) \ar[r] \ar[d] &
\lim H^s_T((I^{m_0}M)^{\oplus r}/I^n(I^{m_0}M)^{\oplus r}) \ar[d] \\
H^s_T(Q) \ar[r] &
\lim H^s_T(Q/I^nQ)
}
$$
을 얻는다. 필요한 자리에서 Mittag--Leffler 조건이 성립하므로 오른쪽 열은
완전하다. 호몰로지의 보조정리
\ref{homology-lemma-apply-Mittag-Leffler}를 보라. 가장 아래 가로 화살표는
보조정리 \ref{algebraization-lemma-final-bootstrap}의 (6)에 의해 단사이다(!). 그 위의
가로 화살표는 보조정리 \ref{algebraization-lemma-final-bootstrap}의 (5)에 의해 전단사이다.
코호몰로지 차수 $\leq s - 1$에서의 화살표들은 동형이다. 따라서
$H^s_T(M) \to \lim H^s_T(M/I^nM)$가 오항 보조정리
(호몰로지의 보조정리 \ref{homology-lemma-five-lemma})에 의해 동형이라는
결론을 얻는다.
이로써 정리의 증명이 끝난다.
\end{proof}

\begin{lemma}
\label{algebraization-lemma-combine-two}
$I \subset \mathfrak a \subset A$를 Noether 환 $A$의 아이디얼들이라 하고,
$M$을 유한 $A$-가군이라 하자. $s$와 $d$를 정수라 하자. 다음을 가정한다.
\begin{enumerate}
\item $A, I, V(\mathfrak a), M$은
$s$와 $d$에 대해 상황 \ref{algebraization-situation-bootstrap}의 가정들을 만족하고,
\item $A, I, \mathfrak a, M$은
$s + 1$과 $d$에 대해, 그리고 $J = \mathfrak a$로 두었을 때 보조정리
\ref{algebraization-lemma-algebraize-local-cohomology-general}의 조건들을 만족한다.
\end{enumerate}
그러면 아이디얼 $J_0 \subset \mathfrak a$가 존재하여
$V(J_0) \cap V(I) = V(\mathfrak a)$를 만족하고, 임의의 $J \subset J_0$가
$V(J) \cap V(I) = V(\mathfrak a)$를 만족하면 사상
$$
H^{s + 1}_J(M) \longrightarrow \lim H^{s + 1}_\mathfrak a(M/I^nM)
$$
은 동형이다.
\end{lemma}

\begin{proof}
실제로 $J_0$가 존재하고 동형
$H^{s + 1}_J(M) = H^{s + 1}(R\Gamma_\mathfrak a(M)^\wedge)$
을 얻으며, 이는 보조정리 \ref{algebraization-lemma-algebraize-local-cohomology-general}에
따른다. 또한 짧은 완전열
$$
0 \to R^1\lim H^s_\mathfrak a(M/I^nM) \to
H^{s + 1}(R\Gamma_\mathfrak a(M)^\wedge) \to
\lim H^{s + 1}_\mathfrak a(M/I^nM) \to 0
$$
은 쌍대화 복합체의 보조정리 \ref{dualizing-lemma-completion-local}에 의해
성립한다. 한편 가군 $R^1\lim H^s_\mathfrak a(M/I^nM)$은 영이다. 실제로
$\{H^s_\mathfrak a(M/I^nM)\}_{n \geq 0}$은 정리
\ref{algebraization-theorem-final-bootstrap}에 의해 Mittag--Leffler 조건을 만족한다.
\end{proof}








\section{형식적 단면의 대수화, I}
\label{algebraization-section-algebraization-sections}

\noindent
이 절에서는 국소적인 경우 형식적 단면의 대수화 문제를 연구한다.
$(A, \mathfrak m)$을 Noether 국소환이라 하자. $I \subset A$를 아이디얼이라 하자.
$$
X = \Spec(A) \supset U = \Spec(A) \setminus \{\mathfrak m\}
$$
% P07-ERR-0167: 원문의 표기 도입 구문 ``denote Y''에는 전치사 ``by''가 빠져 있다.
그리고 $Y = V(I)$로 $I$에 대응하는 닫힌 부분스킴을 나타내자.
$\mathcal{F}$를 연접 $\mathcal{O}_U$-가군이라 하자. 이 절에서는 극한
$$
\lim_n H^i(U, \mathcal{F}/I^n\mathcal{F})
$$
들을 생각한다. 이는 $\mathcal{F}$를 $U$의 $Y$를 따른 형식적 완비화로
당겨 얻는 층의 코호몰로지와 밀접하게 관련된다. 그러나 아직 형식 스킴을
도입하지 않았으므로 여기서는 이 용어를 사용할 수 없다.

\begin{lemma}
\label{algebraization-lemma-compare-with-derived-completion}
$U$를 Noether 국소환 $A$의 천공 스펙트럼이라 하자.
$\mathcal{F}$를 연접 $\mathcal{O}_U$-가군이라 하고, $I \subset A$를
아이디얼이라 하자. 그러면
$$
H^i(R\Gamma(U, \mathcal{F})^\wedge) =
\lim H^i(U, \mathcal{F}/I^n\mathcal{F})
$$
가 모든 $i$에 대해 성립한다. 여기서 $R\Gamma(U, \mathcal{F})^\wedge$는
유도 $I$-진 완비를 나타낸다.
\end{lemma}

\begin{proof}
보조정리 \ref{algebraization-lemma-formal-functions-general}와
\ref{algebraization-lemma-derived-completion-pseudo-coherent}에 의해
$$
R\Gamma(U, \mathcal{F})^\wedge =
R\Gamma(U, \mathcal{F}^\wedge) =
R\Gamma(U, R\lim \mathcal{F}/I^n\mathcal{F})
$$
이다. 따라서 짧은 완전열
$$
0 \to R^1\lim H^{i - 1}(U, \mathcal{F}/I^n\mathcal{F}) \to
H^i(R\Gamma(U, \mathcal{F})^\wedge) \to
\lim H^i(U, \mathcal{F}/I^n\mathcal{F}) \to 0
$$
을 얻는다. 이는 코호몰로지의 보조정리
\ref{cohomology-lemma-RGamma-commutes-with-Rlim}에 따른다.
$R^1\lim$ 항들은 영이다. 실제로 군들의 역계
$H^i(U, \mathcal{F}/I^n\mathcal{F})$는 보조정리 \ref{algebraization-lemma-ML-local}에 의해
Mittag--Leffler 조건을 만족한다.
\end{proof}

\begin{theorem}
\label{algebraization-theorem-algebraization-formal-sections}
\begin{reference}
증명 방법은 대략 \cite[정리 1]{Faltings-algebraisation}과
\cite[정리 2]{Faltings-uber}의 증명 방법을 따른다. 이 결과는
\cite[정리 1.1]{MRaynaud-paper}(여집합이 아핀인 경우) 및
\cite[정리 3.9]{MRaynaud-book}(여집합이 소수의 아핀들의 합집합인 경우)와
거의 같다.
\end{reference}
$(A, \mathfrak m)$을 쌍대화 복합체를 가지며 아이디얼 $I$에 관해 완비인
Noether 국소환이라 하자. $X = \Spec(A)$, $Y = V(I)$,
$U = X \setminus \{\mathfrak m\}$로 놓자. $\mathcal{F}$를 $U$ 위의
연접층이라 하자. 다음을 가정한다.
\begin{enumerate}
\item $\text{cd}(A, I) \leq d$, 즉 $H^i(X \setminus Y, \mathcal{G}) = 0$가
$i \geq d$이고 $\mathcal{G}$가 $X$ 위의 준연접층일 때 성립한다.
\item 임의의 $x \in X \setminus Y$에 대해, 그 폐포 $\overline{\{x\}}$가
$X$ 안에서 $U \cap Y$와 만나면 다음이 성립한다.
$$
\text{depth}_{\mathcal{O}_{X, x}}(\mathcal{F}_x) \geq s
\quad\text{또는}\quad
\text{depth}_{\mathcal{O}_{X, x}}(\mathcal{F}_x)
+ \dim(\overline{\{x\}}) > d + s
$$
\end{enumerate}
그러면 열린집합 $V_0 \subset U$가 존재하여 $U \cap Y$를 포함하고,
$V \subset V_0$가 $U \cap Y$를 포함하는 임의의 열린집합이면 사상
$$
H^i(V, \mathcal{F}) \to \lim H^i(U, \mathcal{F}/I^n\mathcal{F})
$$
은 $i < s$에 대해 동형이다. 더욱이
$
\text{depth}_{\mathcal{O}_{X, x}}(\mathcal{F}_x) +
\dim(\overline{\{x\}}) > s
$
가 모든 $x \in U \cap Y$에 대해 성립하면 이 코호몰로지 군들은 유한
$A$-가군이다.
\end{theorem}

\begin{proof}
유한 $A$-가군 $M$을 택하여 $\mathcal{F}$가 $U$로 제한한 연접
$\mathcal{O}_X$-가군, 곧 $M$에 결부된 가군층이 되게 한다. 국소 코호몰로지의
보조정리 \ref{local-cohomology-lemma-finiteness-pushforwards-and-H1-local}를 보라.
그러면 보조정리 \ref{algebraization-lemma-algebraize-local-cohomology}의 가정들이 만족된다.
그 보조정리에서와 같이 $J_0$를 택하고 $V_0 = X \setminus V(J_0)$로 놓자.
그러면 열린집합 $V \subset V_0$ 중 $U \cap Y$를 포함하는 것들은
$1$-대-$1$로 아이디얼 $J \subset J_0$ 중
$V(J) \cap V(I) = \{\mathfrak m\}$를 만족하는 것들과 대응한다.
더욱이 이러한 선택에 대해 구별 삼각형
$$
R\Gamma_J(M) \to M \to R\Gamma(V, \mathcal{F}) \to
R\Gamma_J(M)[1]
$$
을 얻는다. 마찬가지로 구별 삼각형
$$
R\Gamma_\mathfrak m(M)^\wedge \to
M \to
R\Gamma(U, \mathcal{F})^\wedge \to
R\Gamma_\mathfrak m(M)^\wedge[1]
$$
을 얻으며, 여기에는 유도 $I$-진 완비가 나타난다. 보조정리
\ref{algebraization-lemma-compare-with-derived-completion}에 의해
$R\Gamma(U, \mathcal{F})^\wedge$의 코호몰로지 군들은 정리의 명제에 나온
극한들과 같다. 이 두 삼각형 사이의 표준 사상과 몇 가지 쉬운 논증으로부터
우리 정리가 주 보조정리 \ref{algebraization-lemma-algebraize-local-cohomology}에서 따름을
알 수 있다(여기서는 $i < s$인 반면 그 보조정리에서는 $i \leq s$이다.
이는 밀림 때문이다). 추가 가정하에서 코호몰로지 군들의 유한성은 보조정리
\ref{algebraization-lemma-kill-colimit}에서 따른다.
\end{proof}

\begin{lemma}
\label{algebraization-lemma-application-theorem}
$(A, \mathfrak m)$을 쌍대화 복합체를 가지며 아이디얼 $I$에 관해 완비인
Noether 국소환이라 하자. $X = \Spec(A)$, $Y = V(I)$,
$U = X \setminus \{\mathfrak m\}$로 놓자. $\mathcal{F}$를 $U$ 위의
연접층이라 하자. 임의의 결부점 $x \in U$가 $\mathcal{F}$의 결부점이면
$\dim(\overline{\{x\}}) > \text{cd}(A, I) + 1$이라 가정하자. 여기서
$\overline{\{x\}}$는 $X$에서의 폐포이다. 그러면 사상
$$
\colim H^0(V, \mathcal{F})
\longrightarrow
\lim H^0(U, \mathcal{F}/I^n\mathcal{F})
$$
은 유한 $A$-가군들의 동형이다. 여기서 여극한은 열린집합 $V \subset U$ 중
$U \cap Y$를 포함하는 것들 전체에 걸쳐 취한다.
\end{lemma}

\begin{proof}
정리 \ref{algebraization-theorem-algebraization-formal-sections}를 $s = 1$로 적용한다
(유한성도 함께 얻는다).
\end{proof}




\section{형식적 단면의 대수화, II}
\label{algebraization-section-algebraization-sections-coherent}

\noindent
절 \ref{algebraization-section-algebraization-sections-general}과 \ref{algebraization-section-bootstrap}의
결과가 준아핀 스킴 위 연접층의 코호몰로지 및 아이디얼에 관한 그 완비에
미치는 가능한 귀결들을 모두 간결하게 서술하기는 다소 어렵다. 이 절에서는
$H^0$에 대한 응용들의 비망라적 목록을 제시한다. 다음 절에는 고차
코호몰로지에 대한 응용들이 들어 있다.

\begin{lemma}
\label{algebraization-lemma-ses-H0}
$I \subset \mathfrak a$를 Noether 환 $A$의 아이디얼들이라 하자.
$0 \to \mathcal{F}' \to \mathcal{F} \to \mathcal{F}'' \to 0$을
$U = \Spec(A) \setminus V(\mathfrak a)$ 위 연접 가군들의 짧은 완전열이라 하자.
$\mathcal{V}$를 열린 부분스킴 $V \subset U$ 중 $U \cap V(I)$를 포함하는 것들의
집합에 역포함관계로 순서를 준 것이라 하자. 가환 그림
$$
\xymatrix{
\colim_\mathcal{V} H^0(V, \mathcal{F}') \ar[d] \ar[r] &
\colim_\mathcal{V} H^0(V, \mathcal{F}) \ar[d] \ar[r] &
\colim_\mathcal{V} H^0(V, \mathcal{F}'') \ar[d] \\
\lim H^0(U, \mathcal{F}'/I^n\mathcal{F}') \ar[r] &
% P07-ERR-0168: 아래 가운데 항의 분자는 완전열상 F여야 하나 고정 원문은 F'을 쓴다.
\lim H^0(U, \mathcal{F}'/I^n\mathcal{F}) \ar[r] &
\lim H^0(U, \mathcal{F}'/I^n\mathcal{F}'')
}
$$
을 생각하자. 왼쪽과 오른쪽 아래 방향 화살표가 동형이면 가운데 것도 동형이다.
% P07-ERR-0169: 다음 결론은 가정한 왼쪽 화살표를 반복하며, 논리상 오른쪽이어야 한다.
가운데와 왼쪽 아래 방향 화살표가 동형이면 왼쪽 것도 동형이다.
\end{lemma}

\begin{proof}
그림의 열들은 가운데에서 완전하고 첫 화살표는 단사이다. 따라서 마지막
명제는 쉬운 그림 추적으로부터 따른다. 나머지 증명에서는 왼쪽과 오른쪽
아래 방향 화살표들이 동형이라고 가정한다. 그림 추적에 의해 가운데 아래
방향 화살표는 단사이다. 이제 전사임만 보이면 된다.

\medskip\noindent
유한 $A$-가군 $M$과 $M'$을 택하여 $\mathcal{F}$와 $\mathcal{F}'$가 각각
$\widetilde{M}$과 $\widetilde{M}'$을 $U$에 제한한 것이 되게 할 수 있다.
국소 코호몰로지의 보조정리
\ref{local-cohomology-lemma-finiteness-pushforwards-and-H1-local}를 보라.
$M'$을 $\mathfrak a^n M'$으로, 어떤 $n \geq 0$에 대해 바꾼 뒤에는
$\mathcal{F}' \to \mathcal{F}$가 가군 사상 $M' \to M$에 대응한다고
가정할 수 있다. 스킴의 코호몰로지의 보조정리
\ref{coherent-lemma-homs-over-open}를 보라. $M'$을 $M' \to M$의 상으로
바꾸고 $M'' = M/M'$로 놓으면, 주어진 짧은 완전열은 짧은 완전열
$0 \to M' \to M \to M'' \to 0$을 이루는 $A$-가군들에 결부된 연접 가군들의 짧은 완전열을
제한한 것에 대응함을 알 수 있다.

\medskip\noindent
$\hat s \in \lim H^0(U, \mathcal{F}/I^n\mathcal{F})$의 상을
$\hat s'' \in \lim H^0(U, \mathcal{F}''/I^n\mathcal{F}'')$라 하자.
가정에 의해 어떤 $V \in \mathcal{V}$와 단면 $s'' \in \mathcal{F}''(V)$를
찾되, 후자가 $\hat s''$로 사상되게 한다. $J \subset A$를
$V(J) = \Spec(A) \setminus V$를 만족하는 아이디얼이라 하자. 스킴의
코호몰로지의 보조정리 \ref{coherent-lemma-homs-over-open}에 의해 $J$를
한 거듭제곱으로 바꾼 뒤에는 $A$-선형 사상 $\varphi : J \to M''$이 있어
$s''$에 대응한다고 가정할 수 있다. 이 $J$의 선택을 고정한다.
나머지 증명에서는 $V$를 더 작은 $V$, 곧 $\mathcal{V}$에 속하는 것으로 바꿀 것이다.
즉 $V \cap V(J) = \emptyset$가 성립한다.

\medskip\noindent
표현 $A^{\oplus m} \to A^{\oplus n} \to J \to 0$을 택하자.
$g_1, \ldots, g_n \in J$로 $A^{\oplus n}$의 기저벡터들의 상을 나타내면
$J = (g_1, \ldots, g_n)$이다. $A^{\oplus m} \to A^{\oplus n}$이 행렬
$(a_{ji})$로 주어져 $\sum a_{ji} g_i = 0$, $j = 1, \ldots, m$이 되게 하자.
$M \to M''$가 전사이므로 각 $i$에 대해 $m_i \in M$을 택하여
$\varphi(g_i) \in M''$로 사상되게 할 수 있다. 그러면 $g_i \hat s - m_i$는
$\lim H^0(U, \mathcal{F}/I^n\mathcal{F})$의 원소로서 부분가군
$\lim H^0(U, \mathcal{F}'/I^n\mathcal{F}')$에 속한다. 가정에 의해 $V$를
축소한 뒤에는 $s'_i \in \mathcal{F}'(V)$, $i = 1, \ldots, n$이 존재하여
$s'_i$가 $g_i \hat s - m_i$로 사상된다고 가정할 수 있다.
$s_i = s'_i + m_i$로 $\mathcal{F}(V)$ 안에서 놓자.
$\sum a_{ji} s_i$가 $\sum a_{ji}g_i\hat s = 0$으로 사상됨에 유의하라. 이때 사상은
$$
\colim_\mathcal{V} \mathcal{F}(V')
\longrightarrow
\lim H^0(U, \mathcal{F}/I^n\mathcal{F})
$$
이다. 이 사상은 단사이므로(위 참조), $V$를 축소한 뒤에는
$\sum a_{ji}s_i = 0$가 $\mathcal{F}(V)$ 안에서 모든
$j = 1, \ldots, m$에 대해 성립한다고 가정할 수 있다. 따라서 $A$-가군 사상
$J \to \mathcal{F}(V)$을 얻으며, 이는 $g_i$를 $s_i$로 보낸다.
$\widetilde{J}$의 보편 성질에 의해 이 $A$-가군 사상은 $\mathcal{O}_V$-가군 사상
$\widetilde{J}|_V \to \mathcal{F}$에 대응한다. 그러나
$V(J) \cap V = \emptyset$이므로 $\widetilde{J}|_V = \mathcal{O}_V$이다.
따라서 단면 $s \in \mathcal{F}(V)$를 구성하였다. $s$가 위에 표시한 사상에
의해 $\hat s$로 사상됨을 보이는 계산은 생략한다.
\end{proof}

\noindent
다음 보조정리는 명제 \ref{algebraization-proposition-application-H0}로 대체될 것이다.

\begin{lemma}
\label{algebraization-lemma-application-H0-pre}
$I \subset \mathfrak a$를 Noether 환 $A$의 아이디얼들이라 하자.
$\mathcal{F}$를 $U = \Spec(A) \setminus V(\mathfrak a)$ 위의 연접 가군이라 하자.
다음을 가정한다.
\begin{enumerate}
\item $A$는 $I$-진 완비이고 쌍대화 복합체를 가진다.
\item $x \in \text{Ass}(\mathcal{F})$, $x \not \in V(I)$,
$\overline{\{x\}} \cap V(I) \not \subset V(\mathfrak a)$이고
$z \in \overline{\{x\}} \cap V(\mathfrak a)$이면
$\dim(\mathcal{O}_{\overline{\{x\}}, z}) > \text{cd}(A, I) + 1$이다.
\item 다음 중 하나가 성립한다.
\begin{enumerate}
\item $\mathcal{F}$를 $U \setminus V(I)$에 제한한 것은 $(S_1)$을 만족한다.
\item $V(\mathfrak a)$의 차원은 많아야 $2$이다.\footnote{이는
차원함수가 $V(\mathfrak a)$에서 취하는 최댓값과 최솟값의 차가, 그 함수가
$\Spec(A)$ 위에 정의되어 있을 때 많아야 $2$라는 뜻이다.}
\end{enumerate}
\end{enumerate}
그러면 동형
$$
\colim H^0(V, \mathcal{F})
\longrightarrow
\lim H^0(U, \mathcal{F}/I^n\mathcal{F})
$$
을 얻는다. 여기서 여극한은 열린집합 $V \subset U$ 중 $U \cap V(I)$를
포함하는 것들 전체에 걸쳐 취한다.
\end{lemma}

\begin{proof}
유한 $A$-가군 $M$을 택하여 $\mathcal{F}$가 $U$로 제한한 연접 가군, 곧
$M$에 결부된 가군층이 되게 한다. 국소 코호몰로지의 보조정리
\ref{local-cohomology-lemma-finiteness-pushforwards-and-H1-local}를 보라.
$d = \text{cd}(A, I)$로 놓자. $\mathfrak p$를 $A$의 소아이디얼 중
$V(I)$에 포함되지 않는 것이라 하고,
$\mathfrak q \in V(\mathfrak p) \cap V(\mathfrak a)$라 하자. 그러면
$\mathfrak p$가 $M$의 결부 소아이디얼이 아니어서
$\text{depth}(M_\mathfrak p) \geq 1$이거나, 아니면 (2)에 의해
$\dim((A/\mathfrak p)_\mathfrak q) > d + 1$이다. 따라서 보조정리
\ref{algebraization-lemma-algebraize-local-cohomology-general}의 가정들은 $s = 1$과
$d$에 대해 만족된다. 여기서 조건 (3)을 사용한다. 그러므로 아이디얼
$J_0 \subset \mathfrak a$가 존재하여 $V(J_0) \cap V(I) = V(\mathfrak a)$이고,
임의의 $J \subset J_0$가 $V(J) \cap V(I) = V(\mathfrak a)$를 만족하면 사상
$$
H^i_J(M) \longrightarrow H^i(R\Gamma_\mathfrak a(M)^\wedge)
$$
들은 $i = 0, 1$에 대해 동형이다. 완전 삼각형들의 사상
$$
\xymatrix{
R\Gamma_J(M)  \ar[d] \ar[r] &
M \ar[r] \ar[d] &
R\Gamma(V, \mathcal{F}) \ar[d] \ar[r] &
R\Gamma_J(M)[1]  \ar[d] \\
R\Gamma_J(M)^\wedge \ar[r] &
M \ar[r] &
R\Gamma(V, \mathcal{F})^\wedge \ar[r] &
R\Gamma_J(M)^\wedge[1] \\
R\Gamma_\mathfrak a(M)^\wedge \ar[r] \ar[u] &
M \ar[r] \ar[u] &
R\Gamma(U, \mathcal{F})^\wedge \ar[r] \ar[u] &
R\Gamma_\mathfrak a(M)^\wedge[1] \ar[u]
}
$$
을 생각하자. 여기서 $V = \Spec(A) \setminus V(J)$이다.
$R\Gamma_\mathfrak a(M)^\wedge \to R\Gamma_J(M)^\wedge$가 동형임을 상기하라.
실제로 $\mathfrak a$, $\mathfrak a + I$, $J + I$는 같은 닫힌 부분스킴을
정의한다. 예를 들어 보조정리
\ref{algebraization-lemma-algebraize-local-cohomology-general}의 증명을 보라. 따라서
$R\Gamma(U, \mathcal{F})^\wedge = R\Gamma(V, \mathcal{F})^\wedge$이다.
이로부터 가환 그림
$$
\xymatrix{
0 \ar[r] &
H^0_J(M) \ar[r] \ar[d] &
M \ar[r] \ar[d] \ar[r] &
\Gamma(V, \mathcal{F}) \ar[d] \ar[r] &
H^1_J(M) \ar[d] \ar[r] &
0 \\
0 \ar[r] &
H^0(R\Gamma_J(M)^\wedge) \ar[r] &
M \ar[r] &
H^0(R\Gamma(V, \mathcal{F})^\wedge) \ar[r] &
H^1(R\Gamma_J(M)^\wedge) \ar[r] &
0 \\
0 \ar[r] &
H^0(R\Gamma_\mathfrak a(M)^\wedge) \ar[r] \ar[u] &
M \ar[r] \ar[u] &
H^0(R\Gamma(U, \mathcal{F})^\wedge) \ar[r] \ar[u] &
H^1(R\Gamma_\mathfrak a(M)^\wedge) \ar[r] \ar[u] &
0
}
$$
을 얻는다. 각 행은 완전하고 아래쪽 수직 화살표들은 동형이다. 따라서 동형
$\Gamma(V, \mathcal{F}) \to H^0(R\Gamma(U, \mathcal{F})^\wedge)$을 얻는다.
보조정리 \ref{algebraization-lemma-formal-functions-general}와
\ref{algebraization-lemma-derived-completion-pseudo-coherent}에 의해
$$
R\Gamma(U, \mathcal{F})^\wedge =
R\Gamma(U, \mathcal{F}^\wedge) =
R\Gamma(U, R\lim \mathcal{F}/I^n\mathcal{F})
$$
이고, 코호몰로지의 보조정리
\ref{cohomology-lemma-RGamma-commutes-with-Rlim}에 의해
$H^0(R\Gamma(U, \mathcal{F})^\wedge) =
\lim H^0(U, \mathcal{F}/I^n\mathcal{F})$임을 얻는다.
\end{proof}

\noindent
이제 앞의 보조정리를 부트스트랩하여 조건 (3)을 없앤다.

\begin{proposition}
\label{algebraization-proposition-application-H0}
$I \subset \mathfrak a$를 Noether 환 $A$의 아이디얼들이라 하자.
$\mathcal{F}$를 $U = \Spec(A) \setminus V(\mathfrak a)$ 위의 연접 가군이라 하자.
다음을 가정한다.
\begin{enumerate}
\item $A$는 $I$-진 완비이고 쌍대화 복합체를 가진다.
\item $x \in \text{Ass}(\mathcal{F})$, $x \not \in V(I)$,
$\overline{\{x\}} \cap V(I) \not \subset V(\mathfrak a)$,
그리고 $z \in \overline{\{x\}} \cap V(\mathfrak a)$이면
$\dim(\mathcal{O}_{\overline{\{x\}}, z}) > \text{cd}(A, I) + 1$이다.
\end{enumerate}
그러면 동형
$$
\colim H^0(V, \mathcal{F})
\longrightarrow
\lim H^0(U, \mathcal{F}/I^n\mathcal{F})
$$
을 얻는다. 여기서 여극한은 열린집합 $V \subset U$ 중 $U \cap V(I)$를
포함하는 것들 전체에 걸쳐 취한다.
\end{proposition}

\begin{proof}
$T \subset U$를 점 $x$들 중
$\overline{\{x\}} \cap V(I) \subset V(\mathfrak a)$를 만족하는 것들의
집합이라 하자. 연접 가군들의 전사 $\mathcal{F} \to \mathcal{F}'$을 택하자.
이는 $U$ 위에서 국소 코호몰로지의 보조정리
\ref{local-cohomology-lemma-get-depth-1-along-Z}에 따라 구성한 사상이다. 이 사상
$\mathcal{F} \to \mathcal{F}'$은 열린집합 $V \subset U$ 중
$U \cap V(I)$를 포함하는 것 위에서 동형이므로, $\mathcal{F}$를 $\mathcal{F}'$로 바꾼
보조정리를 증명하면 충분하다. 따라서 $x \in U$가
$\overline{\{x\}} \cap V(I) \subset V(\mathfrak a)$를 만족하면
$\text{depth}(\mathcal{F}_x) \geq 1$이라고 가정해도 좋고 그렇게 한다.

\medskip\noindent
$\mathcal{V}$를 열린 부분스킴 $V \subset U$ 중 $U \cap V(I)$를 포함하는
것들의 집합에 역포함관계로 순서를 준 것이라 하자. 이는 유향집합이다.
먼저 사상
$$
\mathcal{F}(V)
\longrightarrow
\lim H^0(U, \mathcal{F}/I^n\mathcal{F})
$$
% P07-ERR-0170: 다음 고정 원문의 V \in F는 열린집합의 지표집합 V-script를 써야 한다.
이 임의의 $V \in \mathcal{F}$에 대해 단사라고 주장한다(특히 보조정리의
사상은 단사이다). 실제로 결부점 $x$가 $\mathcal{F}$의 결부점이면 앞 문단에 의해
% P07-ERR-0171: 다음 두 교집합의 Y는 이 명제에서 정의되지 않았으며 V(I)가 의도된 것으로 보인다.
$\overline{\{x\}} \cap U \cap Y \not = \emptyset$를 만족해야 한다.
$y \in \overline{\{x\}} \cap U \cap Y$이면 $\mathcal{F}_x$는
$\mathcal{F}_y$의 국소화이고, Krull 교차 정리(대수학의 보조정리
\ref{algebra-lemma-intersect-powers-ideal-module-zero})에 의해
$\mathcal{F}_y \subset \lim \mathcal{F}_y/I^n \mathcal{F}_y$이다.
이로써 주장을 증명하였다. 실제로 핵에 속하는 단면
$s \in \mathcal{F}(V)$의 지지집합은 공집합이어야 하므로 그 단면은 영이다.

\medskip\noindent
유한 $A$-가군 $M$을 택하여 $\mathcal{F}$가 $\widetilde{M}$을 $U$에 제한한
것이 되게 한다. 국소 코호몰로지의 보조정리
\ref{local-cohomology-lemma-finiteness-pushforwards-and-H1-local}를 보라.
$H^0_\mathfrak a(M) = 0$이라고 가정해도 좋고 그렇게 한다.
$\text{Ass}(M) \setminus V(I) = \{\mathfrak p_1, \ldots, \mathfrak p_n\}$라 하자.
$n$에 대한 귀납법으로 보조정리를 증명하겠다. 순서를 바꾸어
$\mathfrak p_n$이 집합 $\{\mathfrak p_1, \ldots, \mathfrak p_n\}$의 포함관계에
대한 극소원소라고 가정할 수 있다. 즉 $\mathfrak p_n$은 $M$의 지지집합의
일반점이다. 다음과 같이 놓자.
$$
M' = H^0_{\mathfrak p_1 \ldots \mathfrak p_{n - 1} I}(M)
$$
그리고 $M'' = M/M'$라 하자. $\mathcal{F}'$과 $\mathcal{F}''$을 연접
$\mathcal{O}_U$-가군으로서 각각 $M'$과 $M''$에 대응하는 것이라 하자. 쌍대화 복합체의
보조정리 \ref{dualizing-lemma-divide-by-torsion}에 의해 $M''$은 단 하나의
결부 소아이디얼, 즉 $\mathfrak p_n$만을 가진다. 따라서 $\mathcal{F}''$도
결부점이 하나뿐이고, 보조정리 \ref{algebraization-lemma-application-H0-pre}의 조건
(3)(a)가 성립한다. 그러므로 사상
$\colim H^0(V, \mathcal{F}'')
\to \lim H^0(U, \mathcal{F}''/I^n\mathcal{F}'')$는 동형이다. 한편
$\mathfrak p_n \not \in V(\mathfrak p_1 \ldots \mathfrak p_{n - 1} I)$이므로
$\mathfrak p_n$은 $M'$의 결부 소아이디얼이 아니다. 따라서 귀납 가정을
$M'$에 적용할 수 있다. 또한 $\mathcal{F}' \subset \mathcal{F}$이므로, 조건
$\text{depth}(\mathcal{F}'_x) \geq 1$이 점 $x$ 중
$\overline{\{x\}} \cap V(I) \subset V(\mathfrak a)$를 만족하는 것들에서
성립함에 유의하라.
대수학의 보조정리 \ref{algebra-lemma-depth-in-ses}를 보라. 따라서 사상
$\colim H^0(V, \mathcal{F}')
\to \lim H^0(U, \mathcal{F}'/I^n\mathcal{F}')$도 동형이다.
보조정리 \ref{algebraization-lemma-ses-H0}에 의해 결론을 얻는다.
\end{proof}

\begin{lemma}
\label{algebraization-lemma-application-H0}
$I \subset \mathfrak a$를 Noether 환 $A$의 아이디얼들이라 하자.
$\mathcal{F}$를 $U = \Spec(A) \setminus V(\mathfrak a)$ 위의 연접 가군이라 하자.
다음을 가정한다.
\begin{enumerate}
\item $A$는 $I$-진 완비이고 쌍대화 복합체를 가진다.
\item $x \in \text{Ass}(\mathcal{F})$, $x \not \in V(I)$,
$\overline{\{x\}} \cap V(I) \not \subset V(\mathfrak a)$이고
$z \in V(\mathfrak a) \cap \overline{\{x\}}$이면
$\dim(\mathcal{O}_{\overline{\{x\}}, z}) > \text{cd}(A, I) + 1$이다.
\item $x \in U$가 $\overline{\{x\}} \cap V(I) \subset V(\mathfrak a)$를
만족하면 $\text{depth}(\mathcal{F}_x) \geq 2$이다.
\end{enumerate}
그러면 동형
$$
H^0(U, \mathcal{F})
\longrightarrow
\lim H^0(U, \mathcal{F}/I^n\mathcal{F})
$$
을 얻는다.
\end{lemma}

\begin{proof}
$\hat s \in \lim H^0(U, \mathcal{F}/I^n\mathcal{F})$라 하자. 명제
\ref{algebraization-proposition-application-H0}에 의해 $\hat s$가 원소
$s \in \mathcal{F}(V)$의 상임을 알 수 있다. 여기서 $V \subset U$는
$U \cap V(I)$를 포함하는 어떤 열린집합이다. 그러나 조건 (3)에 의해 모든
$\text{depth}(\mathcal{F}_x) \geq 2$이고, 이는 모든 $x \in U \setminus V$에
대해 성립한다. 따라서 인자의 보조정리
\ref{divisors-lemma-depth-2-hartog}에 의해
$\mathcal{F}(V) = \mathcal{F}(U)$이고 증명이 끝난다.
\end{proof}

\begin{lemma}
\label{algebraization-lemma-alternative-colim-H0}
$A$를 Noether 환이라 하자. $f \in \mathfrak a \subset A$를 $A$의
아이디얼의 원소라 하고, $M$을 유한 $A$-가군이라 하자. 다음을 가정한다.
\begin{enumerate}
\item $A$는 $f$-진 완비이다.
\item $f$는 $M$ 위의 영인자가 아니다.
\item $H^1_\mathfrak a(M/fM)$은 유한 $A$-가군이다.
\end{enumerate}
그러면 $U = \Spec(A) \setminus V(\mathfrak a)$로 놓을 때 사상
$$
\colim_V \Gamma(V, \widetilde{M})
\longrightarrow
\lim \Gamma(U, \widetilde{M/f^nM})
$$
은 동형이다. 여기서 여극한은 열린집합 $V \subset U$ 중 $U \cap V(f)$를
포함하는 것들 전체에 걸쳐 취한다.
\end{lemma}

\begin{proof}
$\mathcal{F} = \widetilde{M}|_U$로 놓자. $H^1_\mathfrak a(M/fM)$의 유한성은
$H^0(U, \mathcal{F}/f\mathcal{F})$의 유한성을 함의한다. 국소 코호몰로지의
보조정리 \ref{local-cohomology-lemma-finiteness-pushforwards-and-H1-local}를
보라. 코호몰로지의 보조정리 \ref{cohomology-lemma-limit-finite}에 의해
($f$가 $\mathcal{F}$ 위의 영인자가 아니므로 적용할 수 있다)
$N = \lim H^0(U, \mathcal{F}/f^n\mathcal{F})$은 유한 $A$-가군이고,
$f$-꼬임이 없으며, $N/fN \subset H^0(U, \mathcal{F}/f\mathcal{F})$이다.
한편 사상 $M \to N$과 이와 양립하는 사상
$$
M/fM \longrightarrow H^0(U, \mathcal{F}/f\mathcal{F})
$$
을 가진다. $g \in \mathfrak a$에 대해 $(M/fM)_g$는
% P07-ERR-0172: 다음 고정 원문의 D(f)는 국소화 변수에 맞추면 D(g)여야 한다.
$H^0(U \cap D(f), \mathcal{F}/f\mathcal{F})$로 동형적으로 사상된다.
실제로 $\mathcal{F}/f\mathcal{F}$는 $\widetilde{M/fM}$을 $U$에 제한한 것이다.
따라서 $M_g \to N_g$는 동형
$$
M_g/fM_g = (M/fM)_g \to (N/fN)_g = N_g/fN_g
$$
을 유도한다. $f$는 $N$과 $M$ 모두 위에서 영인자가 아니므로
$M_g \to N_g$가 $f$-진 완비들 사이의 동형을 유도한다. 이는 다시
$M_g \to N_g$가 $V(f) \cap D(g)$의 어떤 열린 근방에서 동형임을 함의한다.
$g \in \mathfrak a$가 임의였으므로, $M$과 $N$은 보조정리의 명제에서와
같은 어떤 열린집합 $V$ 위에서 서로 동형인 연접 가군들을 정한다.
이로써 증명이 끝난다.
\end{proof}

\begin{proposition}
\label{algebraization-proposition-alternative-colim-H0}
$A$를 Noether 환이라 하자. $f \in \mathfrak a \subset A$를 $A$의 아이디얼의
원소라 하자. $\mathcal{F}$를 $U = \Spec(A) \setminus V(\mathfrak a)$ 위의
연접 가군이라 하자. 다음을 가정한다.
\begin{enumerate}
\item $A$는 $f$-진 완비이고 쌍대화 복합체를 가진다.
\item $x \in \text{Ass}(\mathcal{F})$, $x \not \in V(f)$,
$\overline{\{x\}} \cap V(f) \not \subset V(\mathfrak a)$,
그리고 $z \in \overline{\{x\}} \cap V(\mathfrak a)$이면
$\dim(\mathcal{O}_{\overline{\{x\}}, z}) > 2$이다.
\end{enumerate}
그러면 사상
$$
\colim_V \Gamma(V, \mathcal{F})
\longrightarrow
\lim \Gamma(U, \mathcal{F}/f^n\mathcal{F})
$$
은 동형이다. 여기서 여극한은 열린집합 $V \subset U$ 중 $U \cap V(f)$를
포함하는 것들 전체에 걸쳐 취한다.
\end{proposition}

\begin{proof}[첫째 증명]
% P07-ERR-0173: 원문은 두 번째 서술어의 동사를 빠뜨린다.
$A$는 보편적 사슬환이고 Gorenstein 형식적 섬유를 가짐을 상기하라. 쌍대화
복합체의 보조정리들
\ref{dualizing-lemma-dualizing-gorenstein-formal-fibres} 및
\ref{dualizing-lemma-universally-catenary}를 보라. 따라서 사상
$\mathcal{F} \to \mathcal{F}'$을 생각할 수 있다. 이는 국소 코호몰로지의
보조정리 \ref{local-cohomology-lemma-make-S2-along-Z}에서 닫힌 부분집합
$V(f) \cap U$, 곧 $U$의 닫힌 부분집합에 대해 구성한 사상이다. 다음을 관찰하라.
\begin{enumerate}
\item $\mathcal{F} \to \mathcal{F}'$의 핵과 여핵은 $V(f) \cap U$ 위에
지지되어 있다.
\item $\mathcal{F}'$은 $f$-꼬임이 없다. 실제로 그 줄기들은 깊이 $\geq 1$이고,
이는 $V(f) \cap U$의 모든 점에서 성립한다. 즉 $\mathcal{F}'$은
$V(f) \cap U$ 안에 결부점을 갖지 않는다.
\item $y \in V(f) \cap U$가 $\mathcal{F}'/f\mathcal{F}'$의 결부점이면
$\text{depth}(\mathcal{F}'_y) = 1$이고, 따라서 ($\mathcal{F}'$의 구성에 의해)
즉시 특수화 $x \leadsto y$가 존재하며 $x \not \in V(f)$는
$\mathcal{F}$의 결부점이다. 그러므로 가정 (2)에 의해 $y$는
$\Spec(A)$ 안에서 점 $z \in V(\mathfrak a)$로의 즉시 특수화를 가질 수 없다.
\item (3)으로부터 $H^0(U, \mathcal{F}'/f\mathcal{F}')$은 유한
$A$-가군이다. 국소 코호몰로지의 보조정리
\ref{local-cohomology-lemma-finiteness-Rjstar}를 보라.
\end{enumerate}
이 관찰들을 이용하여 증명을 마칠 수 있다.

\medskip\noindent
먼저 보조정리가 $\mathcal{F}'$에 대해 성립한다고 주장한다. 실제로 유한
$A$-가군 $M'$을 택하여 $\mathcal{F}'$이 $U$에 제한한 연접 가군, 곧
$M'$에 결부된 것이 되게 한다. 국소 코호몰로지의 보조정리
\ref{local-cohomology-lemma-finiteness-pushforwards-and-H1-local}를 보라.
$\mathcal{F}'$은 $f$-꼬임이 없으므로 $M'$도 $f$-꼬임이 없다고 가정할 수 있다.
% P07-ERR-0174: 다음 고정 원문의 H^1_a(M')에는 보조정리에 필요한 /fM'가 빠져 있다.
위 관찰 (4)에 의해 $H^1_\mathfrak a(M')$은 유한 $A$-가군이다. 국소
코호몰로지의 보조정리
\ref{local-cohomology-lemma-finiteness-pushforwards-and-H1-local}를 보라.
% P07-ERR-0175: 원문 ``Thus the claim by''에는 추론 동사가 빠져 있다.
따라서 보조정리 \ref{algebraization-lemma-alternative-colim-H0}에 의해 주장이 따른다.

\medskip\noindent
둘째, 임의의 연접 $\mathcal{O}_U$-가군 중 $V(f) \cap U$ 위에 지지된 것에 대해
보조정리가 자명하게 성립함을 관찰한다. $\mathcal{K}$, $\mathcal{G}$,
$\mathcal{Q}$를 각각 사상 $\mathcal{F} \to \mathcal{F}'$의 핵, 상, 여핵이라
하자. 짧은 완전열
$0 \to \mathcal{G} \to \mathcal{F}' \to \mathcal{Q} \to 0$과 보조정리
\ref{algebraization-lemma-ses-H0}에 의해 결과는 $\mathcal{G}$에 대해 성립한다. 이제 짧은
완전열 $0 \to \mathcal{K} \to \mathcal{F} \to \mathcal{G} \to 0$에 대해
같은 논증을 다시 적용하면 증명이 끝난다.
\end{proof}

\begin{proof}[둘째 증명]
이 명제는 명제 \ref{algebraization-proposition-application-H0}의 특수한 경우이다.
\end{proof}

\begin{lemma}
\label{algebraization-lemma-alternative-H0}
$A$를 Noether 환이라 하자. $f \in \mathfrak a \subset A$를 $A$의 아이디얼의
원소라 하고, $M$을 유한 $A$-가군이라 하자. 다음을 가정한다.
\begin{enumerate}
\item $A$는 $f$-진 완비이다.
\item $H^1_\mathfrak a(M)$과 $H^2_\mathfrak a(M)$은 $f$의 어떤 거듭제곱에
의해 소멸된다.
\end{enumerate}
그러면 $U = \Spec(A) \setminus V(\mathfrak a)$로 놓을 때 사상
$$
\Gamma(U, \widetilde{M})
\longrightarrow
\lim \Gamma(U, \widetilde{M/f^nM})
$$
은 동형이다.
\end{lemma}

\begin{proof}
보조정리 \ref{algebraization-lemma-formal-functions-principal}를 $U$와
$\mathcal{F} = \widetilde{M}|_U$에 적용할 수 있다. 실제로 $\mathcal{F}$은
연접 $\mathcal{O}_U$-가군들의 범주에서 Noether 대상이다.
$H^1(U, \mathcal{F}) = H^2_\mathfrak a(M)$은 (국소 코호몰로지의 보조정리
\ref{local-cohomology-lemma-finiteness-pushforwards-and-H1-local})
$f$의 어떤 거듭제곱에 의해 소멸되므로, 그 $f$-진 Tate 가군은 영이다.
따라서 이 보조정리에 의해 $\lim H^0(U, \mathcal{F}/f^n \mathcal{F})$은
통상적인 $f$-진 완비, 곧 $H^0(U, \mathcal{F})$의 완비와 같다. 짧은 완전열
$$
0 \to M/H^0_\mathfrak a(M) \to H^0(U, \mathcal{F}) \to H^1_\mathfrak a(M) \to 0
$$
을 생각하자. 이는 국소 코호몰로지의 보조정리
\ref{local-cohomology-lemma-finiteness-pushforwards-and-H1-local}에 나온다.
$M/H^0_\mathfrak a(M)$은 유한 $A$-가군이므로 완비이다.
대수학의 보조정리 \ref{algebra-lemma-completion-tensor}를 보라.
$H^1_\mathfrak a(M)$은 $f$의 어떤 거듭제곱에 의해 소멸되므로, 대수학의
보조정리 \ref{algebra-lemma-completion-differ-by-torsion}로부터
$H^0(U, \mathcal{F})$도 완비라는 결론을 얻는다. 이로써 증명이 끝난다.
\end{proof}







\section{형식적 단면의 대수화, III}
\label{algebraization-section-algebraization-sections-coherent-III}

\noindent
% P07-ERR-0176: 다음 절이라는 원문 지시는 실제 내용이 이 절에서 시작하므로 의심스럽다.
다음 절에는 국소 코호몰로지의 완비에 관한 자료를, 준아핀 스킴 위 연접
가군들의 고차 코호몰로지와 아이디얼에 관한 그 완비에 적용한 결과들의
비망라적 목록이 들어 있다.

\begin{proposition}
\label{algebraization-proposition-application-higher}
$I \subset \mathfrak a$를 Noether 환 $A$의 아이디얼들이라 하자.
$\mathcal{F}$를 $U = \Spec(A) \setminus V(\mathfrak a)$ 위의 연접 가군이라 하자.
$s \geq 0$이라 하자. 다음을 가정한다.
\begin{enumerate}
\item $A$는 $I$-진 완비이고 쌍대화 복합체를 가진다.
\item $x \in U \setminus V(I)$이면 $\text{depth}(\mathcal{F}_x) > s$이거나
$$
\text{depth}(\mathcal{F}_x) +
\dim(\mathcal{O}_{\overline{\{x\}}, z}) > \text{cd}(A, I) + s + 1
$$
가 모든 $z \in V(\mathfrak a) \cap \overline{\{x\}}$에 대해 성립한다.
\item 다음 조건 중 하나가 성립한다.
\begin{enumerate}
\item $\mathcal{F}$를 $U \setminus V(I)$에 제한한 것은 $(S_{s + 1})$을 만족한다.
\item $V(\mathfrak a)$의 차원은 많아야 $2$이다.\footnote{이는 차원함수가
$V(\mathfrak a)$에서 취하는 최댓값과 최솟값의 차가, 그 함수가
$\Spec(A)$ 위에 정의되어 있을 때 많아야 $2$라는 뜻이다.}
\end{enumerate}
\end{enumerate}
그러면 사상들
$$
H^i(U, \mathcal{F})
\longrightarrow
\lim H^i(U, \mathcal{F}/I^n\mathcal{F})
$$
은 $i < s$에 대해 동형이다. 더욱이 동형
$$
\colim H^s(V, \mathcal{F})
\longrightarrow
\lim H^s(U, \mathcal{F}/I^n\mathcal{F})
$$
을 얻는다. 여기서 여극한은 열린집합 $V \subset U$ 중 $U \cap V(I)$를
포함하는 것들 전체에 걸쳐 취한다.
\end{proposition}

\begin{proof}
$s > 0$이라 가정해도 좋다. 실제로 $s = 0$인 경우는 명제
\ref{algebraization-proposition-application-H0}에서 다루었다.

\medskip\noindent
유한 $A$-가군 $M$을 택하여 $\mathcal{F}$가 $U$로 제한한 연접 가군, 곧
$M$에 결부된 것이 되게 한다. 국소 코호몰로지의 보조정리
\ref{local-cohomology-lemma-finiteness-pushforwards-and-H1-local}를 보라.
$d = \text{cd}(A, I)$로 놓자. $\mathfrak p$를 $A$의 소아이디얼 중
% P07-ERR-0177: 원문의 점 p가 닫힌집합 V(I)에 ``포함되지 않는다''는 표현은 ``속하지 않는다''여야 한다.
$V(I)$에 포함되지 않는 것이라 하고,
$\mathfrak q \in V(\mathfrak p) \cap V(\mathfrak a)$라 하자. 그러면
$\text{depth}(M_\mathfrak p) \geq s + 1 > s$이거나, 아니면 (2)에 의해
$\dim((A/\mathfrak p)_\mathfrak q) > d + s + 1$이다. 보조정리
\ref{algebraization-lemma-bootstrap-bis-bis}에 의해 상황 \ref{algebraization-situation-bootstrap}의 가정들이
$A, I, V(\mathfrak a), M, s, d$에 대해 만족된다는 결론을 얻는다. 한편 보조정리
\ref{algebraization-lemma-algebraize-local-cohomology-general}의 가정들은 $s + 1$과 $d$에
대해 만족된다. 바로 여기서 조건 (3)을 사용한다.

\medskip\noindent
보조정리 \ref{algebraization-lemma-algebraize-local-cohomology-general}를 적용하면 아이디얼
$J_0 \subset \mathfrak a$가 존재하여 $V(J_0) \cap V(I) = V(\mathfrak a)$이고,
임의의 $J \subset J_0$가 $V(J) \cap V(I) = V(\mathfrak a)$를 만족하면 사상들
$$
H^i_J(M) \longrightarrow H^i(R\Gamma_\mathfrak a(M)^\wedge)
$$
은 $i \leq s + 1$에 대해 동형임을 알 수 있다.

\medskip\noindent
$i \leq s$에 대해 사상 $H^i_\mathfrak a(M) \to H^i_J(M)$은 보조정리들
\ref{algebraization-lemma-bootstrap-inherited}와 \ref{algebraization-lemma-kill-colimit-support-general}에
의해 동형이다. 코호몰로지와 국소 코호몰로지의 비교(국소 코호몰로지의
보조정리 \ref{local-cohomology-lemma-local-cohomology})를 사용하면
$H^i(U, \mathcal{F}) \to H^i(V,\mathcal{F})$
가 $V = \Spec(A) \setminus V(J)$ 및 $i < s$에 대해 동형임을 얻는다.

\medskip\noindent
정리 \ref{algebraization-theorem-final-bootstrap}에 의해
$H^i_\mathfrak a(M) = \lim H^i_\mathfrak a(M/I^nM)$가 $i \leq s$에 대해
성립한다. 보조정리 \ref{algebraization-lemma-combine-two}에 의해
$H^{s + 1}_\mathfrak a(M) = \lim H^{s + 1}_\mathfrak a(M/I^nM)$이다.

\medskip\noindent
$H^0(U, \mathcal{F}) = H^0(V, \mathcal{F}) =
\lim H^0(U, \mathcal{F}/I^n\mathcal{F})$라는 동형은 위 결과와 명제
\ref{algebraization-proposition-application-H0}에서 따른다. $0 < i < s$에 대해서도
원하는 동형들
$H^i(U, \mathcal{F}) = H^i(V, \mathcal{F}) =
\lim H^i(U, \mathcal{F}/I^n\mathcal{F})$을 국소 코호몰로지와 코호몰로지의
관계를 사용하여 같은 방식으로 얻는다. 이는 $i = 0$인 경우보다 쉽다.
실제로 $i > 0$이면
$$
H^i(U, \mathcal{F}) = H^{i + 1}_\mathfrak a(M),
\quad
H^i(V, \mathcal{F}) = H^{i + 1}_J(M),
\quad
H^i(R\Gamma(U, \mathcal{F})^\wedge) = 
H^{i + 1}(R\Gamma_\mathfrak a(M)^\wedge)
$$
% P07-ERR-0179: 원문의 마지막 문장은 주어와 동사가 없는 단편이다.
마지막 명제에 대한 논증도 마찬가지이다.
\end{proof}

\begin{lemma}
\label{algebraization-lemma-alternative-higher}
$A$를 Noether 환이라 하자. $f \in \mathfrak a \subset A$를 $A$의 아이디얼의
원소라 하고, $M$을 유한 $A$-가군이라 하자. $s \geq 0$이라 하자. 다음을 가정한다.
\begin{enumerate}
\item $A$는 $f$-진 완비이다.
\item $H^i_\mathfrak a(M)$은 $f$의 어떤 거듭제곱에 의해 소멸되며, 이는
$i \leq s + 1$에 대해 성립한다.
\end{enumerate}
그러면 $U = \Spec(A) \setminus V(\mathfrak a)$로 놓을 때 사상
$$
H^i(U, \widetilde{M})
\longrightarrow
\lim H^i(U, \widetilde{M/f^nM})
$$
은 $i < s$에 대해 동형이다.
\end{lemma}

\begin{proof}
$s$에 대한 귀납법을 사용한다. $s = 0$이면 명제는 공허하다. $s = 1$이면
결과는 보조정리 \ref{algebraization-lemma-alternative-H0}이다. $s > 1$이라 가정하자.
귀납법에 의해 $i = s - 1 \geq 1$인 경우를 증명하면 충분하다.
보조정리 \ref{algebraization-lemma-formal-functions-principal}를 $U$와
$\mathcal{F} = \widetilde{M}|_U$에 적용할 수 있다. 실제로 $\mathcal{F}$은
연접 $\mathcal{O}_U$-가군들의 범주에서 Noether 대상이다.
$H^j(U, \mathcal{F}) = H^{j + 1}_\mathfrak a(M)$가 모든 $j$에 대해
성립함을 관찰하라. 이는 국소 코호몰로지의 보조정리
\ref{local-cohomology-lemma-finiteness-pushforwards-and-H1-local}에 따른다.
따라서 $j = s = (s - 1) + 1$이면 이는 가정에 의해 $f$의 어떤
거듭제곱에 의해 소멸된다. 그러므로 보조정리
\ref{algebraization-lemma-formal-functions-principal}에 의해
$\lim H^{s - 1}(U, \mathcal{F}/f^n\mathcal{F})$은 통상적인 $f$-진 완비,
곧 $H^{s - 1}(U, \mathcal{F})$의 완비이다. 다시 이 가군이
$f$의 어떤 거듭제곱에 의해 소멸됨을 사용하면, 그 완비는 단순히
$H^{s - 1}(U, \mathcal{F})$와 같으며 이것이 원하는 바이다.
\end{proof}





\section{연결성에의 응용}
\label{algebraization-section-connected}

\noindent
이 절에서는 Grothendieck의 연결성 정리와 그 변형들을 논한다. 원래 판본은
% P07-ERR-0181: 다음 인용 위치의 ``Exposee''는 프랑스어 Exposé의 오기이다.
\cite[Exposee XIII, 정리 2.1]{SGA2}에서 찾을 수 있다. 문헌에는 Faltings의
연결성 정리라고 불리는 판본이 있으며, 이는 \cite[정리 6]{Faltings-some}을
가리키는 것으로 추측한다. \cite[정리 1.6]{Varbaro}에 주어진 완비 국소환에
대한 최적 판본을 진술하고 증명하자.

\begin{lemma}
\label{algebraization-lemma-punctured-still-connected}
\begin{reference}
\cite[정리 1.6]{Varbaro}
\end{reference}
$(A, \mathfrak m)$을 Noether 완비 국소환이라 하자. $I$를 $A$의 진아이디얼이라
하자. $X = \Spec(A)$ 및 $Y = V(I)$로 놓자. 다음과 같이 나타내자.
\begin{enumerate}
\item $d$는 $X$의 기약 성분들의 차원 중 최솟값이다.
\item $c$는 닫힌 부분집합 $Z \subset X$들 중 $X \setminus Z$가
비연결이 되는 것들의 차원 중 최솟값이다.
\end{enumerate}
그러면 닫힌 부분집합 $Z \subset Y$에 대해 $Y \setminus Z$는
$\dim(Z) < \min(c, d - 1) - \text{cd}(A, I)$이면 연결이다.
특히 $A/I$의 천공 스펙트럼은
$\text{cd}(A, I) < \min(c, d - 1)$이면 연결이다.
\end{lemma}

\begin{proof}
먼저 마지막 주장을 증명하자. 첫째 경우로 $A/I$의 천공 스펙트럼이 공집합이면,
국소 코호몰로지의 보조정리 \ref{local-cohomology-lemma-cd-bound-dim-local}에
의해 $X$의 모든 기약 성분의 차원은 $\leq \text{cd}(A, I)$이다. 따라서
$\min(c, d - 1) - \text{cd}(A, I) < 0$이고, 이 경우 보조정리가 성립한다.
그러므로 $U \cap Y$가 공집합이 아니라고 가정해도 좋다. 여기서
$U = X \setminus \{\mathfrak m\}$는 $A$의 천공 스펙트럼이다. $A$를 그 축약으로
바꾸어도 좋다. $A$는 쌍대화 복합체를 가지고(쌍대화 복합체의 보조정리
\ref{dualizing-lemma-ubiquity-dualizing}), 또한 $A$는 $I$에 관해 완비임을 관찰하라
(대수학의 보조정리 \ref{algebra-lemma-complete-by-sub}).
$d - 1 > \text{cd}(A, I)$라 가정하면 보조정리
\ref{algebraization-lemma-application-theorem}를 적용하여 사상
$$
\colim H^0(V, \mathcal{O}_V)
\longrightarrow
\lim H^0(U, \mathcal{O}_U/I^n\mathcal{O}_U)
$$
이 동형임을 알 수 있다. 여기서 여극한은 열린집합 $V \subset U$ 중
$U \cap Y$를 포함하는 것들 전체에 걸쳐 취한다. $U \cap Y$가 비연결이면
그 $n$차 무한소 근방은 $U$ 안에서 모든 $n$에 대해 비연결이고, 따라서 오른쪽
항에 비자명 멱등원이 있음을 알 수 있다(여기서 $U \cap Y$가 공집합이 아님을
사용한다). 그러므로 비연결인 어떤 $V$를 찾을 수 있다.
$Z = X \setminus V$로 놓자. 국소 코호몰로지의 보조정리
\ref{local-cohomology-lemma-cd-bound-dim-local}에 의해 $Z$의 모든 기약 성분의
차원은 $\leq \text{cd}(A, I)$이다. 따라서 $c \leq \text{cd}(A, I)$이고,
이로써 마지막 명제를 증명하였다.

\medskip\noindent
방금 증명한 것으로부터 다음과 같이 보조정리의 명제를 도출할 수 있다.
$Z \subset Y$가 닫혀 있고 $Y \setminus Z$가 비연결이며
$\dim(Z) = e$라고 가정하자. 관례상 연결공간은 공집합이 아님을 상기하라.
따라서 (a) $Y = Z$이거나, (b)
$Y \setminus Z = W_1 \amalg W_2$이고 $W_i$들이 $Y \setminus Z$ 안에서
공집합이 아닌 열린닫힌집합이다. (b)의 경우 $w_i \in W_i$인 점들을 택할 수
있고, 이 점들은 $U$에서 닫혀 있다. 사상의 보조정리
\ref{morphisms-lemma-ubiquity-Jacobson-schemes}를 보라. 이제
$f_1, \ldots, f_e \in \mathfrak m$을 찾아
$V(f_1, \ldots, f_e) \cap Z = \{\mathfrak m\}$가 되게 할 수 있고, (b)의
경우에는 $w_i \in V(f_1, \ldots, f_e)$라고 가정할 수 있다. 실제로 소아이디얼
회피를 사용하여 귀납적으로 $f_i$를 택하되
$\dim V(f_1, \ldots, f_i) \cap Z = e - i$가 되고, (b)의 경우에는
$w_1, w_2 \in V(f_i)$가 되게 할 수 있다. 따라서
% P07-ERR-0183: 다음 고정 원문의 몫환에는 I + (...)를 묶는 괄호가 빠져 있다.
$A/I + (f_1, \ldots, f_e)$의 천공 스펙트럼은 비연결이다(작은 세부사항은
생략한다). 부등식
$\text{cd}(A, I + (f_1, \ldots, f_e)) \leq \text{cd}(A, I) + e$은
국소 코호몰로지의 보조정리들 \ref{local-cohomology-lemma-cd-sum} 및
\ref{local-cohomology-lemma-bound-cd}에 의해 성립하므로
$$
\text{cd}(A, I) + e \geq \min(c, d - 1)
$$
라는 결론을 증명의 첫 부분으로부터 얻는다. 이는
$e \geq \min(c, d - 1) - \text{cd}(A, I)$를 함의하며, 이것이 보일 바였다.
\end{proof}

\begin{lemma}
\label{algebraization-lemma-connected}
$I \subset \mathfrak a$를 Noether 환 $A$의 아이디얼들이라 하자. 다음을 가정한다.
\begin{enumerate}
\item $A$는 $I$-진 완비이고 쌍대화 복합체를 가진다.
\item $\mathfrak p \subset A$가 $V(I)$에 포함되지 않는 극소 소아이디얼이고
$\mathfrak q \in V(\mathfrak p) \cap V(\mathfrak a)$이면
$\dim((A/\mathfrak p)_\mathfrak q) > \text{cd}(A, I) + 1$이다.
% P07-ERR-0184: 앞의 ``V(I)에 포함되지 않는''은 점과 부분집합을 혼동한 고정 원문 표현이다.
\item 공집합 아닌 임의의 열린집합 $V \subset \Spec(A)$가
$V(I) \setminus V(\mathfrak a)$를 포함하면 연결이다.\footnote{예를 들어
$A$가 정역인 경우이다.}
\end{enumerate}
그러면 $V(I) \setminus V(\mathfrak a)$는 공집합이거나 연결이다.
\end{lemma}

\begin{proof}
$A$를 그 축약으로 바꾸어도 좋다. 그러면 (2)의 부등식은 $A$의 모든 결부
소아이디얼에 대해 성립한다. 명제 \ref{algebraization-proposition-application-H0}에 의해
$$
\colim H^0(V, \mathcal{O}_V) = \lim H^0(T_n, \mathcal{O}_{T_n})
$$
임을 알 수 있다. 여기서 여극한은 (3)과 같은 열린집합 $V$ 전체에 걸쳐
취하고, $T_n$은 $n$차 무한소 근방, 곧 $T = V(I) \setminus V(\mathfrak a)$의
$U = \Spec(A) \setminus V(\mathfrak a)$ 안에서의 근방이다. 따라서 $T$는
공집합이거나 연결이다. 그렇지 않다면 오른쪽 항은 비자명 멱등원을 가질
것이나, 왼쪽 항은 그렇지 않다고 가정했기 때문이다. 일부 세부사항은 생략한다.
\end{proof}

\begin{lemma}
\label{algebraization-lemma-H0}
$A$를 쌍대화 복합체를 가지며 영이 아닌 $f \in A$에 관해 완비인 Noether
정역이라 하자. $f \in \mathfrak a \subset A$를 아이디얼이라 하자.
$Z = V(\mathfrak a)$의 모든 기약 성분이 여차원 $> 2$로
$X = \Spec(A)$ 안에 놓인다고 가정하자. 즉 $Z$의 모든 기약 성분이 여차원
$> 1$로 $Y = V(f)$ 안에 놓인다고 가정하자. 그러면 $Y \setminus Z$는 연결이다.
\end{lemma}

\begin{proof}
이는 보조정리 \ref{algebraization-lemma-connected}의 특수한 경우이다(그 증명은 명제
\ref{algebraization-proposition-application-H0}에 의존한다). 아래에서는 더 쉬운 명제
\ref{algebraization-proposition-alternative-colim-H0}를 사용하여 증명한다.

\medskip\noindent
$U = X \setminus Z$로 놓자. 명제 \ref{algebraization-proposition-alternative-colim-H0}에 의해
동형
$$
\colim \Gamma(V, \mathcal{O}_V) \to
\lim_n \Gamma(U, \mathcal{O}_U/f^n \mathcal{O}_U)
$$
을 얻는다. 여기서 여극한은 열린집합 $V \subset U$ 중 $U \cap Y$를 포함하는
것들 전체에 걸쳐 취한다. 따라서 $U \cap Y$가 비연결이면 어떤 $V$에 대해
$\Gamma(V, \mathcal{O}_V)$ 안에 비자명 멱등원이 존재한다. 그러나 $V$는
$X$가 정역의 스펙트럼이므로 정역 스킴이고, 이는 불가능하다.
\end{proof}






\section{완비화 함자}
\label{algebraization-section-completion}

\noindent
$X$를 Noether 스킴이라 하자. $Y \subset X$를 준연접 아이디얼층
$\mathcal{I} \subset \mathcal{O}_X$로 주어지는 닫힌 부분스킴이라 하자.
이 절에서는 연접 $\mathcal{O}_X$-가군들의 역계 $(\mathcal{F}_n)$ 중
$\mathcal{F}_n$이 $I^n$에 의해 소멸되고 전이사상이 동형
$\mathcal{F}_{n + 1}/I^n\mathcal{F}_{n + 1} \to \mathcal{F}_n$을 유도하는
것들을 고찰한다. 이러한 역계들의 범주는 스킴의 코호몰로지의 절
$$
\textit{Coh}(X, \mathcal{I})
$$
\ref{coherent-section-coherent-formal}에서 위와 같이 표기하였다.
이 범주는 $Y$를 따른 $X$의 형식적 완비 위의 연접 가군 범주와 동치이다.
그러나 아직 형식 스킴이나 그 위의 연접 가군을 도입하지 않았으므로 여기서는
이 용어를 사용할 수 없다. 특히 다음 완비화 함자에 관심을 둔다.
$$
\textit{Coh}(\mathcal{O}_X)
\longrightarrow
\textit{Coh}(X, \mathcal{I}),\quad
\mathcal{F} \longmapsto \mathcal{F}^\wedge
$$
스킴의 코호몰로지의 식 (\ref{coherent-equation-completion-functor})을 보라.

\begin{lemma}
\label{algebraization-lemma-completion-fully-faithful}
$X$를 Noether 스킴, $Y \subset X$를 닫힌 부분스킴이라 하자.
$Y_n \subset X$를 $X$ 안에서 $Y$의 $n$차 무한소 근방이라 하자.
다음 조건들을 고찰하자.
\begin{enumerate}
\item $X$는 준아핀이며
$\Gamma(X, \mathcal{O}_X) \to \lim \Gamma(Y_n, \mathcal{O}_{Y_n})$
은 동형이다.
\item $X$는 풍부한 가역 가군 $\mathcal{L}$을 가지며
$\Gamma(X, \mathcal{L}^{\otimes m}) \to
\lim \Gamma(Y_n, \mathcal{L}^{\otimes m}|_{Y_n})$
은 모든 $m \gg 0$에 대해 동형이다.
\item 모든 유한 국소 자유 $\mathcal{O}_X$-가군 $\mathcal{E}$에 대해 사상
$\Gamma(X, \mathcal{E}) \to \lim \Gamma(Y_n, \mathcal{E}|_{Y_n})$
은 동형이다.
\item 완비화 함자
$\textit{Coh}(\mathcal{O}_X) \to \textit{Coh}(X, \mathcal{I})$
는 유한 국소 자유 대상들의 충만한 부분범주 위에서 완전충실하다.
\end{enumerate}
그러면 (1) $\Rightarrow$ (2) $\Rightarrow$ (3) $\Rightarrow$ (4)이고
(4) $\Rightarrow$ (3)이다.
\end{lemma}

\begin{proof}
먼저 (3) $\Rightarrow$ (4)를 증명하자. $\mathcal{F}$와 $\mathcal{G}$가
$X$ 위에서 유한 국소 자유이면
$\mathcal{H} = \SheafHom_{\mathcal{O}_X}(\mathcal{G}, \mathcal{F})$
를 고찰하고 스킴의 코호몰로지의 보조정리
\ref{coherent-lemma-completion-internal-hom}
를 사용하면 (3)에서 (4)가 따름을 알 수 있다.

\medskip\noindent
(2) $\rightarrow$ (3)을 증명하자. $\mathcal{L}$이 $X$ 위에서 풍부하고
$\mathcal{E}$가 유한 국소 자유 $\mathcal{O}_X$-가군이라고 하자.
어떤 $r, s \geq 0$ 및 $0 \ll p \ll q$에 대해 보편적으로 완전한 열
$$
0 \to \mathcal{E} \to
(\mathcal{L}^{\otimes p})^{\oplus r} \to
(\mathcal{L}^{\otimes q})^{\oplus s}
$$
을 찾을 수 있다고 주장한다. 이것이 성립하면 완전열
$$
0 \to \lim \Gamma(\mathcal{E}|_{Y_n}) \to
\lim \Gamma((\mathcal{L}^{\otimes p})^{\oplus r}|_{Y_n}) \to
\lim \Gamma((\mathcal{L}^{\otimes q})^{\oplus s}|_{Y_n})
$$
과 (2)의 동형들을 사용하여 (3)의 동형을 얻는다. 주장을 증명하기 위해
쌍대 국소 자유 가군
$\SheafHom_{\mathcal{O}_X}(\mathcal{E}, \mathcal{O}_X)$
을 고찰하고 성질의 명제 \ref{properties-proposition-characterize-ample}를
적용하여 전사
$$
(\mathcal{L}^{\otimes -p})^{\oplus r}
\longrightarrow
\SheafHom_{\mathcal{O}_X}(\mathcal{E}, \mathcal{O}_X)
$$
를 찾는다. 쌍대를 취하면 위 완전열의 첫째 사상을 얻는다(전사라는 성질이
보편적이므로 이 사상은 보편적으로 단사이다). 여핵에 대해 이를 반복하면
둘째 사상을 얻는다. 일부 세부사항은 생략한다.

\medskip\noindent
(1) $\Rightarrow$ (2)를 증명하자. $X$가 준아핀이면 $\mathcal{O}_X$가
풍부한 가역 가군이기 때문이다. 성질의 보조정리
\ref{properties-lemma-quasi-affine-O-ample}를 보라.

\medskip\noindent
(4) $\Rightarrow$ (3)의 증명은 생략한다.
\end{proof}

\noindent
Noether 스킴과 준연접 아이디얼층 $\mathcal{I} \subset \mathcal{O}_X$가
주어졌다고 하자. $\textit{Coh}(X, \mathcal{I})$의 대상
$(\mathcal{F}_n)$에서 각 $\mathcal{F}_n$이 유한 국소 자유
$\mathcal{O}_X/\mathcal{I}^n$-가군이면 이를 {\it 유한 국소 자유}라고 한다.

\begin{lemma}
\label{algebraization-lemma-completion-fully-faithful-general}
$X$를 Noether 스킴, $Y \subset X$를 아이디얼층
$\mathcal{I} \subset \mathcal{O}_X$로 주어진 닫힌 부분스킴이라 하자.
$Y_n \subset X$를 $X$ 안에서 $Y$의 $n$차 무한소 근방이라 하자.
$\mathcal{V}$를 $Y$를 포함하는 열린 부분스킴 $V \subset X$들의 집합에
포함관계의 역순으로 순서를 준 것으로 놓자.
\begin{enumerate}
\item $X$는 준아핀이며
$$
\colim_\mathcal{V} \Gamma(V, \mathcal{O}_V)
\longrightarrow
\lim \Gamma(Y_n, \mathcal{O}_{Y_n})
$$
은 동형이다.
\item $X$는 풍부한 가역 가군 $\mathcal{L}$을 가지며
$$
\colim_\mathcal{V} \Gamma(V, \mathcal{L}^{\otimes m})
\longrightarrow
\lim \Gamma(Y_n, \mathcal{L}^{\otimes m}|_{Y_n})
$$
은 모든 $m \gg 0$에 대해 동형이다.
\item 모든 $V \in \mathcal{V}$와 모든 유한 국소 자유
$\mathcal{O}_V$-가군 $\mathcal{E}$에 대해 사상
$$
\colim_{V' \geq V} \Gamma(V', \mathcal{E}|_{V'})
\longrightarrow
\lim \Gamma(Y_n, \mathcal{E}|_{Y_n})
$$
은 동형이다.
\item 완비화 함자
$$
\colim_\mathcal{V} \textit{Coh}(\mathcal{O}_V)
\longrightarrow
\textit{Coh}(X, \mathcal{I}),
\quad
\mathcal{F} \longmapsto \mathcal{F}^\wedge
$$
는 유한 국소 자유 대상들의 충만한 부분범주 위에서 완전충실하다
(위 설명을 보라).
\end{enumerate}
그러면 (1) $\Rightarrow$ (2) $\Rightarrow$ (3) $\Rightarrow$ (4)이고
(4) $\Rightarrow$ (3)이다.
\end{lemma}

\begin{proof}
$\mathcal{V}$는 유향 집합이므로 여극한들은 범주론의 절
\ref{categories-section-directed-colimits}에서 정의한 것과 같다.
나머지 논증은 보조정리 \ref{algebraization-lemma-completion-fully-faithful}의 증명에
나온 논증과 거의 정확히 같으므로 독자는 건너뛰어도 좋다.

\medskip\noindent
먼저 (3) $\Rightarrow$ (4)를 증명하자. $\mathcal{F}$와 $\mathcal{G}$가
$V \in \mathcal{V}$ 위에서 유한 국소 자유이면
$\mathcal{H} = \SheafHom_{\mathcal{O}_V}(\mathcal{G}, \mathcal{F})$
를 고찰하고 스킴의 코호몰로지의 보조정리
\ref{coherent-lemma-completion-internal-hom}
를 사용하면 (3)에서 (4)가 따름을 알 수 있다.

\medskip\noindent
(2) $\Rightarrow$ (3)을 증명하자. $\mathcal{L}$이 $X$ 위에서 풍부하고
어떤 $V \in \mathcal{V}$에 대해 $\mathcal{E}$가 유한 국소 자유
$\mathcal{O}_V$-가군이라고 하자. 어떤 $r, s \geq 0$ 및
$0 \ll p \ll q$에 대해 보편적으로 완전한 열
$$
0 \to \mathcal{E} \to
(\mathcal{L}^{\otimes p})^{\oplus r}|_{V} \to
(\mathcal{L}^{\otimes q})^{\oplus s}|_{V}
$$
을 찾을 수 있다고 주장한다. 이것이 참이면 (2)의 동형에서 (3)의 동형이
따른다. 주장을 증명하기 위해 $\mathcal{L}|_V$가 풍부함을 상기하자.
성질의 보조정리 \ref{properties-lemma-ample-on-locally-closed}를 보라.
쌍대 국소 자유 가군
$\SheafHom_{\mathcal{O}_V}(\mathcal{E}, \mathcal{O}_V)$
을 고찰하고 성질의 명제 \ref{properties-proposition-characterize-ample}를
적용하여 전사
$$
(\mathcal{L}^{\otimes -p})^{\oplus r}|_V \longrightarrow
\SheafHom_{\mathcal{O}_V}(\mathcal{E}, \mathcal{O}_V)
$$
(전사라는 성질이 보편적이므로 이 사상은 보편적으로 단사이다). 쌍대를
취하면 완전열의 첫째 사상을 얻는다. 여핵에 대해 이를 반복하면 둘째 사상을
얻는다. 일부 세부사항은 생략한다.

\medskip\noindent
(1) $\Rightarrow$ (2)를 증명하자. $X$가 준아핀이면 $\mathcal{O}_X$가
풍부한 가역 가군이기 때문이다. 성질의 보조정리
\ref{properties-lemma-quasi-affine-O-ample}를 보라.

\medskip\noindent
(4) $\Rightarrow$ (3)의 증명은 생략한다.
\end{proof}

\begin{lemma}
\label{algebraization-lemma-recognize-formal-coherent-modules}
$X$를 Noether 스킴이라 하자. $\mathcal{I} \subset \mathcal{O}_X$를
준연접 아이디얼층이라 하자. 함자
$$
\textit{Coh}(X, \mathcal{I}) \longrightarrow \text{Pro-}\QCoh(\mathcal{O}_X)
$$
는 완전충실하다. 범주론의 비고 \ref{categories-remark-pro-category}를 보라.
\end{lemma}

\begin{proof}
$(\mathcal{F}_n)$과 $(\mathcal{G}_n)$을 $\textit{Coh}(X, \mathcal{I})$의
대상이라 하자. $(\mathcal{F}_n)$에서 $(\mathcal{G}_n)$으로 가는 전사영 대상의
사상 $\alpha$는 전이사상과 양립하는 사상계
$\alpha_n : \mathcal{F}_{m(n)} \to \mathcal{G}_n$으로 주어진다. 여기서
$\mathbf{N} \to \mathbf{N}$, $n \mapsto m(n)$은 증가함수이다(특히
$m(n) \geq n$이다). 이제
$\mathcal{F}_n = \mathcal{F}_{m(n)}/\mathcal{I}^n\mathcal{F}_{m(n)}$이고
$\mathcal{G}_n$은 $\mathcal{I}^n$에 의해 소멸되므로 $\alpha_n$은 사상
$\mathcal{F}_n \to \mathcal{G}_n$을 유도한다.
\end{proof}

\noindent
다음으로 이 장 앞부분의 결과를 사용하여 증명할 수 있는 완전충실성 결과의
몇 가지 예를 덧붙인다.

\begin{lemma}
\label{algebraization-lemma-fully-faithful}
$I \subset \mathfrak a$를 Noether 환 $A$의 아이디얼들이라 하자.
$U = \Spec(A) \setminus V(\mathfrak a)$라 하자. 다음을 가정한다.
\begin{enumerate}
\item $A$는 $I$-진 완비이고 쌍대화 복합체를 가진다.
\item 다음을 만족하는 임의의 동반 소아이디얼 $\mathfrak p \subset A$와
$\mathfrak p \not \in V(I)$이고
$V(\mathfrak p) \cap V(I) \not \subset V(\mathfrak a)$이며
$\mathfrak q \in V(\mathfrak p) \cap V(\mathfrak a)$에 대해
$\dim((A/\mathfrak p)_\mathfrak q) > \text{cd}(A, I) + 1$,
이다.
\item $\mathfrak p \not \in V(I)$이고
$V(\mathfrak p) \cap V(I) \subset V(\mathfrak a)$인
$\mathfrak p \subset A$에 대해 $\text{depth}(A_\mathfrak p) \geq 2$이다.
\end{enumerate}
그러면 완비화 함자
$$
\textit{Coh}(\mathcal{O}_U)
\longrightarrow
\textit{Coh}(U, I\mathcal{O}_U),
\quad
\mathcal{F} \longmapsto \mathcal{F}^\wedge
$$
는 유한 국소 자유 대상들의 충만한 부분범주 위에서 완전충실하다.
\end{lemma}

\begin{proof}
보조정리 \ref{algebraization-lemma-completion-fully-faithful}에 의해
$$
\Gamma(U, \mathcal{O}_U) =
\lim \Gamma(U, \mathcal{O}_U/I^n\mathcal{O}_U)
$$
를 보이면 충분하다. 이는 보조정리 \ref{algebraization-lemma-application-H0}에서 곧바로
따른다.
\end{proof}

\begin{lemma}
\label{algebraization-lemma-fully-faithful-alternative}
$A$를 Noether 환이라 하자. $f \in \mathfrak a \subset A$를 $A$의
아이디얼의 원소라 하자. $U = \Spec(A) \setminus V(\mathfrak a)$라 하자.
다음을 가정한다.
\begin{enumerate}
\item $A$는 $f$-진 완비이다.
\item $H^1_\mathfrak a(A)$와 $H^2_\mathfrak a(A)$는 $f$의 어떤 거듭제곱에
의해 소멸된다.
\end{enumerate}
그러면 완비화 함자
$$
\textit{Coh}(\mathcal{O}_U)
\longrightarrow
\textit{Coh}(U, I\mathcal{O}_U),
\quad
\mathcal{F} \longmapsto \mathcal{F}^\wedge
$$
는 유한 국소 자유 대상들의 충만한 부분범주 위에서 완전충실하다.
\end{lemma}

\begin{proof}
보조정리 \ref{algebraization-lemma-completion-fully-faithful}에 의해
$$
\Gamma(U, \mathcal{O}_U) =
\lim \Gamma(U, \mathcal{O}_U/I^n\mathcal{O}_U)
$$
를 보이면 충분하다. 이는 보조정리 \ref{algebraization-lemma-alternative-H0}에서 곧바로
따른다.
\end{proof}

\begin{lemma}
\label{algebraization-lemma-fully-faithful-simple-two}
$A$를 Noether 환이라 하자. $f \in \mathfrak a$를 $A$의 아이디얼의
원소라 하자. $U = \Spec(A) \setminus V(\mathfrak a)$라 하자. 다음을 가정한다.
\begin{enumerate}
\item $A$는 쌍대화 복합체를 가지며 $f$에 관해 완비이다.
\item 모든 소아이디얼 $\mathfrak p \subset A$에 대해,
$f \not \in \mathfrak p$이고
$\mathfrak q \in V(\mathfrak p) \cap V(\mathfrak a)$이면
$\text{depth}(A_\mathfrak p) + \dim((A/\mathfrak p)_\mathfrak q) > 2$.
\end{enumerate}
그러면 완비화 함자
$$
\textit{Coh}(\mathcal{O}_U)
\longrightarrow
\textit{Coh}(U, I\mathcal{O}_U),
\quad
\mathcal{F} \longmapsto \mathcal{F}^\wedge
$$
는 유한 국소 자유 대상들의 충만한 부분범주 위에서 완전충실하다.
\end{lemma}

\begin{proof}
보조정리 \ref{algebraization-lemma-fully-faithful-alternative}와 국소 코호몰로지의 명제
\ref{local-cohomology-proposition-annihilator}에서 따른다.
\end{proof}

\begin{lemma}
\label{algebraization-lemma-fully-faithful-general}
$I \subset \mathfrak a \subset A$를 Noether 환 $A$의 아이디얼들이라 하자.
$U = \Spec(A) \setminus V(\mathfrak a)$라 하자. $\mathcal{V}$를
$U \cap V(I)$를 포함하는 $U$의 열린 부분스킴들의 집합에 포함관계의
역순으로 순서를 준 것으로 놓자. 다음을 가정한다.
\begin{enumerate}
\item $A$는 $I$-진 완비이고 쌍대화 복합체를 가진다.
\item 다음을 만족하는 임의의 동반 소아이디얼
$\mathfrak p \subset A$와
$I \not \subset \mathfrak p$이고
$V(\mathfrak p) \cap V(I) \not \subset V(\mathfrak a)$
및 $\mathfrak q \in V(\mathfrak p) \cap V(\mathfrak a)$에 대해
$\dim((A/\mathfrak p)_\mathfrak q) > \text{cd}(A, I) + 1$.
\end{enumerate}
그러면 완비화 함자
$$
\colim_\mathcal{V} \textit{Coh}(\mathcal{O}_V)
\longrightarrow
\textit{Coh}(U, I\mathcal{O}_U),
\quad
\mathcal{F} \longmapsto \mathcal{F}^\wedge
$$
는 유한 국소 자유 대상들의 충만한 부분범주 위에서 완전충실하다.
\end{lemma}

\begin{proof}
보조정리 \ref{algebraization-lemma-completion-fully-faithful-general}에 의해
$$
\colim_\mathcal{V} \Gamma(V, \mathcal{O}_V) =
\lim \Gamma(U, \mathcal{O}_U/I^n\mathcal{O}_U)
$$
를 보이면 충분하다. 이는 명제 \ref{algebraization-proposition-application-H0}에서 곧바로
따른다.
\end{proof}

\begin{lemma}
\label{algebraization-lemma-fully-faithful-general-alternative}
$A$를 Noether 환이라 하자. $f \in \mathfrak a \subset A$를 $A$의
아이디얼의 원소라 하자. $U = \Spec(A) \setminus V(\mathfrak a)$라 하자.
$\mathcal{V}$를 $U \cap V(f)$를 포함하는 $U$의 열린 부분스킴들의 집합에
포함관계의 역순으로 순서를 준 것으로 놓자. 다음을 가정한다.
\begin{enumerate}
\item $A$는 $f$-진 완비이다.
\item $f$는 영인자가 아니다.
\item $H^1_\mathfrak a(A/fA)$는 유한 $A$-가군이다.
\end{enumerate}
그러면 완비화 함자
$$
\colim_\mathcal{V} \textit{Coh}(\mathcal{O}_V)
\longrightarrow
\textit{Coh}(U, f\mathcal{O}_U),
\quad
\mathcal{F} \longmapsto \mathcal{F}^\wedge
$$
는 유한 국소 자유 대상들의 충만한 부분범주 위에서 완전충실하다.
\end{lemma}

\begin{proof}
보조정리 \ref{algebraization-lemma-completion-fully-faithful-general}에 의해
$$
\colim_\mathcal{V} \Gamma(V, \mathcal{O}_V) =
\lim \Gamma(U, \mathcal{O}_U/I^n\mathcal{O}_U)
$$
를 보이면 충분하다. 이는 보조정리 \ref{algebraization-lemma-alternative-colim-H0}에서
곧바로 따른다.
\end{proof}

\begin{lemma}
\label{algebraization-lemma-fully-faithful-very-general}
$I \subset \mathfrak a \subset A$를 Noether 환 $A$의 아이디얼들이라 하자.
$U = \Spec(A) \setminus V(\mathfrak a)$라 하자. $\mathcal{V}$를
$U \cap V(I)$를 포함하는 $U$의 열린 부분스킴들의 집합에 포함관계의
역순으로 순서를 준 것으로 놓자. 어떤 $V \in \mathcal{V}$에 대해
$\mathcal{F}$와 $\mathcal{G}$를 연접 $\mathcal{O}_V$-가군이라 하자. 사상
$$
\colim_{V' \geq V} \Hom_V(\mathcal{G}|_{V'}, \mathcal{F}|_{V'})
\longrightarrow
\Hom_{\textit{Coh}(U, I\mathcal{O}_U)}(\mathcal{G}^\wedge, \mathcal{F}^\wedge)
$$
은 다음 가정들이 성립하면 전단사이다.
\begin{enumerate}
\item $A$는 $I$-진 완비이고 쌍대화 복합체를 가진다.
\item $x \in \text{Ass}(\mathcal{F})$, $x \not \in V(I)$이고
$\overline{\{x\}} \cap V(I) \not \subset V(\mathfrak a)$
또 $z \in \overline{\{x\}} \cap V(\mathfrak a)$이면
$\dim(\mathcal{O}_{\overline{\{x\}}, z}) > \text{cd}(A, I) + 1$.
\end{enumerate}
\end{lemma}

\begin{proof}
$V$로의 제한이 $\mathcal{F}$와 $\mathcal{G}$인 연접
$\mathcal{O}_U$-가군 $\mathcal{F}'$와 $\mathcal{G}'$을 택할 수 있다.
성질의 보조정리 \ref{properties-lemma-lift-finite-presentation}를 보라.
$\text{Ass}(\mathcal{F}') \subset V$가 되도록 $\mathcal{F}'$의 선택을
수정할 수 있다. 예를 들어 국소 코호몰로지의 보조정리
\ref{local-cohomology-lemma-get-depth-1-along-Z}를 보라. 따라서 $V$를
$U$로, $\mathcal{F}$와 $\mathcal{G}$를 $\mathcal{F}'$와 $\mathcal{G}'$으로
바꾸어도 되며 그렇게 한다.
$\mathcal{H} = \SheafHom_{\mathcal{O}_U}(\mathcal{G}, \mathcal{F})$로
놓자. 이는 연접 $\mathcal{O}_U$-가군이다. 다음이 성립한다.
$$
\Hom_V(\mathcal{G}|_V, \mathcal{F}|_V) =
H^0(V, \mathcal{H})
\quad\text{이고}\quad
\lim H^0(U, \mathcal{H}/\mathcal{I}^n\mathcal{H}) =
\Mor_{\textit{Coh}(U, I\mathcal{O}_U)}
(\mathcal{G}^\wedge, \mathcal{F}^\wedge)
$$
스킴의 코호몰로지의 보조정리
\ref{coherent-lemma-completion-internal-hom}를 보라. 따라서
$\mathcal{H}$에 대해 명제 \ref{algebraization-proposition-application-H0}의 가정들이
성립함을 보이면 증명이 끝난다. 이는
$\text{Ass}(\mathcal{H}) \subset \text{Ass}(\mathcal{F})$이기 때문이다.
스킴의 코호몰로지의 보조정리 \ref{coherent-lemma-hom-into-depth}를 보라.
\end{proof}








\section{연접 형식 가군의 대수화 I}
\label{algebraization-section-algebraization-modules}

\noindent
완비화 함자의 본질적 전사성(아래를 보라)은 \cite{SGA2},
\cite{MRaynaud-book}, \cite{MRaynaud-paper}에서 체계적으로 연구되었다.
다음 아핀 상황에서 작업한다.

\begin{situation}
\label{algebraization-situation-algebraize}
여기서 $A$는 Noether 환이고 $I \subset \mathfrak a \subset A$는
아이디얼들이다. $X = \Spec(A)$, $Y = V(I) = \Spec(A/I)$ 및
$Z = V(\mathfrak a) = \Spec(A/\mathfrak a)$로 놓는다. 또한
$U = X \setminus Z$로 놓는다.
\end{situation}

\noindent
이 절에서는 $\textit{Coh}(U, I\mathcal{O}_U)$의 대상이 완비화 함자
$\textit{Coh}(\mathcal{O}_U) \to \textit{Coh}(U, I\mathcal{O}_U)$의 상에
속함을 보장하는 조건들을 찾는다. 스킴의 코호몰로지의 절
\ref{coherent-section-coherent-formal}과 절 \ref{algebraization-section-completion}을 보라.

\begin{lemma}
\label{algebraization-lemma-system-of-modules}
상황 \ref{algebraization-situation-algebraize}에서 다음을 만족하는 $A$-가군의 역계
$(M_n)$을 고찰하자.
\begin{enumerate}
\item $M_n$은 유한 $A$-가군이다.
\item $M_n$은 $I^n$에 의해 소멸된다.
\item $M_{n + 1}/I^nM_{n + 1} \to M_n$의 핵과 여핵은
$\mathfrak a$-거듭제곱 꼬임이다.
\end{enumerate}
그러면 $(\widetilde{M}_n|_U)$는 $\textit{Coh}(U, I\mathcal{O}_U)$에 속한다.
역으로 $\textit{Coh}(U, I\mathcal{O}_U)$의 모든 대상은 이 방식으로 얻어진다.
\end{lemma}

\begin{proof}
$(\widetilde{M}_n|_U)$가 $\textit{Coh}(U, I\mathcal{O}_U)$에 속한다는
확인은 생략한다. $(\mathcal{F}_n)$을 $\textit{Coh}(U, I\mathcal{O}_U)$의
대상이라 하자. 국소 코호몰로지의 보조정리
\ref{local-cohomology-lemma-finiteness-pushforwards-and-H1-local}
에 의해 어떤 유한 $A/I^n$-가군 $M_n$에 대해
$\mathcal{F}_n = \widetilde{M_n}$임을 알 수 있다. $M_n$을
$H^0_\mathfrak a(M_n)$으로 나누면 $M_n \subset H^0(U, \mathcal{F}_n)$이라고
가정할 수 있다. 쌍대화 복합체의 보조정리
\ref{dualizing-lemma-divide-by-torsion}와 이미 인용한 보조정리를 보라.
$M_{n + 1} \to H^0(U, \mathcal{F}_{n + 1}) \to H^0(U, \mathcal{F}_n)$ 아래
$M_n$의 역상으로 $M_{n + 1}$을 귀납적으로 바꾸면, $M_{n + 1}$이
$M_n$으로 사상된다고 가정할 수 있다. 이로써 (1), (2)를 만족하고
$\mathcal{F}_n = \widetilde{M_n}$인 역계 $(M_n)$을 얻는다. (3)을 보이려면
$M_{n + 1}/I^nM_{n + 1} \to M_n$이 유한 $A$-가군의 사상이며
$\widetilde{\ }$를 적용하고 $U$로 제한하면 동형을 유도한다는 사실을
사용한다(여기서 첫째로 인용한 보조정리를 다시 한 번 사용한다).
\end{proof}

\noindent
상황 \ref{algebraization-situation-algebraize}에서는 스킴의 코호몰로지의 식
(\ref{coherent-equation-completion-functor})에 나오는 완비화 함자
\begin{equation}
\label{algebraization-equation-completion}
\textit{Coh}(\mathcal{O}_U)
\longrightarrow
\textit{Coh}(U, I\mathcal{O}_U),\quad
\mathcal{F} \longmapsto \mathcal{F}^\wedge
\end{equation}
를 연구할 수 있다. $A$가 $I$-진 완비이면 형식적 단면의 대수화에 관한
앞선 결과들에 의해 이 함자는 적절한 부분범주들 위에서 완전충실하다.
몇 가지 예는 절 \ref{algebraization-section-completion}과 보조정리
\ref{algebraization-lemma-fully-faithful-inequalities}를 보라. 이제 $(\mathcal{F}_n)$을
$\textit{Coh}(U, I\mathcal{O}_U)$의 대상이라 하자. 여전히 $A$가 $I$-진
완비라고 가정할 때, $(\mathcal{F}_n)$이 언제 위 완비화 함자의 본질적 상에
속하는지 물을 수 있다.

\begin{lemma}
\label{algebraization-lemma-essential-image-completion}
상황 \ref{algebraization-situation-algebraize}에서 $(\mathcal{F}_n)$을
$\textit{Coh}(U, I\mathcal{O}_U)$의 대상이라 하자. 다음 조건들을 고찰하자.
\begin{enumerate}
\item $(\mathcal{F}_n)$은 함자 (\ref{algebraization-equation-completion})의 본질적 상에 속한다.
\item $(\mathcal{F}_n)$은 연접 $\mathcal{O}_U$-가군의 완비이다.
\item 어떤 열린 $U \cap Y \subset V \subset U$에 대해 $(\mathcal{F}_n)$은
연접 $\mathcal{O}_V$-가군의 완비이다.
\item $(\mathcal{F}_n)$은 연접 $\mathcal{O}_X$-가군을 $U$로 제한한 것의
완비이다.
\item $(\mathcal{F}_n)$은 연접 $\mathcal{O}_X$-가군의 완비를 $U$로 제한한
것이다.
\item $U$로의 제한이 $(\mathcal{F}_n)$인
$\textit{Coh}(X, I\mathcal{O}_X)$의 대상 $(\mathcal{G}_n)$이 존재한다.
\end{enumerate}
그러면 조건 (1), (2), (3), (4), (5)는 동치이며 (6)을 함의한다.
$A$가 $I$-진 완비이면 조건 (6)은 나머지 조건들을 함의한다.
\end{lemma}

\begin{proof}
(1)과 (2)는 동치이다. 연접 $\mathcal{O}_U$-가군 $\mathcal{F}$의 완비는
정의상 함자 (\ref{algebraization-equation-completion}) 아래 $\mathcal{F}$의 상이기 때문이다.
$V \subset U$가 $U \cap Y$를 포함하는 열린 부분스킴이면
$$
\textit{Coh}(V, I\mathcal{O}_V) =
\textit{Coh}(U, I\mathcal{O}_U)
$$
이다. $V \cap Y$ 위에 지지되는 연접 $\mathcal{O}_V$-가군의 범주와
$U \cap Y$ 위에 지지되는 연접 $\mathcal{O}_U$-가군의 범주가 같기 때문이다.
따라서 연접 $\mathcal{O}_V$-가군의 완비는
$\textit{Coh}(U, I\mathcal{O}_U)$의 대상이다. 이제 함자들
$\textit{Coh}(\mathcal{O}_X) \to \textit{Coh}(\mathcal{O}_U) \to
\textit{Coh}(\mathcal{O}_V)$는 본질적으로 전사이므로 (2), (3), (4), (5)는
동치이다. 성질의 보조정리 \ref{properties-lemma-lift-finite-presentation}를 보라.

\medskip\noindent
(5)는 언제나 (6)을 함의한다. $A$가 $I$-진 완비라고 가정하자. 그러면
스킴의 코호몰로지의 보조정리
\ref{coherent-lemma-inverse-systems-affine}.
에 의해 $\textit{Coh}(X, I\mathcal{O}_X)$의 모든 대상은 유한 $A$-가군에
대응한다. 따라서 이 경우 (6)은 (5)를 함의한다.
\end{proof}

\begin{example}
\label{algebraization-example-not-algebraizable}
$k$를 체라 하자. $I = (t)$ 및 $\mathfrak a = (x, y, t)$와 함께
$A = k[x, y][[t]]$로 놓고 상황 \ref{algebraization-situation-algebraize}의 표기를 사용하자.
$U \cap Y = (D(x) \cap Y) \cup (D(y) \cap Y)$는 아핀 열린 덮개이다.
$n \geq 1$에 대해 $A_x/t^nA_x$와 $A_y/t^nA_y$를
$A_{xy}/t^nA_{xy}$의 가역원, 즉 다음 꼴의 임의의 거듭제곱급수의 상으로
붙여서 얻는 $\mathcal{O}_U/t^n\mathcal{O}_U$의 가역 가군
$\mathcal{L}_n$을 고찰하자.
$$
u = 1 + \frac{t}{xy} + \sum_{n \geq 2} a_n \frac{t^n}{(xy)^{\varphi(n)}}
$$
여기서 $a_n \in k[x, y]$이고 $\varphi(n) \in \mathbf{N}$이다.
그러면 $(\mathcal{L}_n)$은 $\textit{Coh}(U, I\mathcal{O}_U)$의 가역 대상이지만
연접 $\mathcal{O}_U$-가군 $\mathcal{L}$의 완비는 아니다. 논증의 개요만
제시하고 세부사항 대부분은 생략한다. $y \in U \cap Y$라 하자. 줄기
$\mathcal{L}_y$의 완비는 가역 가군일 것이므로 $\mathcal{L}_y$도 가역이다.
따라서 $U \cap Y$를 포함하는 어떤 열린 $V \subset U$가 존재하여
$\mathcal{L}|_V$가 가역이다. 인자의 보조정리
\ref{divisors-lemma-extend-invertible-module}에 의해
$\widetilde{M}|_V \cong \mathcal{L}|_V$인 가역 $A$-가군 $M$을 찾는다.
그러나 환 $A$는 유일인수분해정역이므로 $M \cong A$이고, 이로부터
$\mathcal{L}_n \cong \mathcal{O}_U/I^n\mathcal{O}_U$가 따를 것이다.
구성상 $\mathcal{L}_2 \not \cong \mathcal{O}_U/I^2\mathcal{O}_U$이므로
원하는 모순을 얻는다.

\medskip\noindent
$n \geq 2$에 대해 $a_n = 0$으로 택하면
$\lim H^0(U, \mathcal{L}_n)$은 영이 아님을 알 수 있다. 이 경우 $D(x)$ 위의
함수 $x$와 $D(y)$ 위의 함수 $x + t/y$가 붙는다. 한편 $a_n = 1$ 및
$\varphi(n) = 2^n$, 나아가 $\varphi(n) = n^2$으로 택하면 독자는
$\lim H^0(U, \mathcal{L}_n)$이 영임을 보일 수 있다. 이는 이 경우
$(\mathcal{L}_n)$이 대수화 가능하지 않다는 또 다른 증명을 준다.
\end{example}

\noindent
상황 \ref{algebraization-situation-algebraize}에서 환 $A$가 $I$-진 완비가 아니면,
보조정리 \ref{algebraization-lemma-essential-image-completion}는 $(\mathcal{F}_n)$이 제한 함자
$$
\textit{Coh}(X, I\mathcal{O}_X)
\longrightarrow
\textit{Coh}(U, I\mathcal{O}_U)
$$
의 본질적 상에 속하는지를 묻는 것이 올바름을 시사한다. 그러나 이제 이것이
$(\mathcal{F}_n)$의 대수화 가능성을 뜻한다고 말할 수는 없다. 따라서 다음
용어를 도입한다.

\begin{definition}
\label{algebraization-definition-algebraizable}
상황 \ref{algebraization-situation-algebraize}에서 $(\mathcal{F}_n)$을
$\textit{Coh}(U, I\mathcal{O}_U)$의 대상이라 하자. $U$로의 제한이
$(\mathcal{F}_n)$과 동형인 $\textit{Coh}(X, I\mathcal{O}_X)$의 대상
$(\mathcal{G}_n)$이 존재하면 {\it $(\mathcal{F}_n)$은 $X$로 연장된다}고 한다.
\end{definition}

\noindent
이 개념은 완비 위에서 대수화 가능하다는 것과 동치이다.

\begin{lemma}
\label{algebraization-lemma-algebraizable}
상황 \ref{algebraization-situation-algebraize}에서 $(\mathcal{F}_n)$을
$\textit{Coh}(U, I\mathcal{O}_U)$의 대상이라 하자. $A', I', \mathfrak a'$을
$A, I, \mathfrak a$의 $I$-진 완비라 하자. $X' = \Spec(A')$ 및
$U' = X' \setminus V(\mathfrak a')$로 놓자. 다음은 동치이다.
\begin{enumerate}
\item $(\mathcal{F}_n)$은 $X$로 연장된다.
\item $(\mathcal{F}_n)$을 $U'$로 당긴 것은 연접 $\mathcal{O}_{U'}$-가군의
완비이다.
\end{enumerate}
\end{lemma}

\begin{proof}
$A \to A'$은 동형 $A/I \to A'/I'$을 유도하는 평탄한 환 준동형임을
상기하자. 대수학의 보조정리 \ref{algebra-lemma-completion-flat}과
\ref{algebra-lemma-completion-complete}를 보라. 따라서 $X' \to X$는
동형 $Y' \to Y$를 유도하는 평탄 사상이다. 그러므로 $U' \to U$는 동형
$U' \cap Y' \to U \cap Y$를 유도하는 평탄 사상이다. 이로부터 가환 그림
$$
\xymatrix{
\textit{Coh}(X', I\mathcal{O}_{X'}) \ar[r] &
\textit{Coh}(U', I\mathcal{O}_{U'}) \\
\textit{Coh}(X, I\mathcal{O}_X) \ar[u] \ar[r] &
\textit{Coh}(U, I\mathcal{O}_U) \ar[u]
}
$$
에서 수직 함자들은 동치이다. 스킴의 코호몰로지의 보조정리
\ref{coherent-lemma-inverse-systems-pullback-equivalence}.
를 보라. 이 사실과 보조정리 \ref{algebraization-lemma-essential-image-completion}의 결과에서
보조정리가 형식적으로 따른다.
\end{proof}

\noindent
상황 \ref{algebraization-situation-algebraize}에서 $(\mathcal{F}_n)$을
$\textit{Coh}(U, I\mathcal{O}_U)$의 대상이라 하자. $(\mathcal{F}_n)$이
$X$로 연장되는지를 알아보기 위해 $A$-가군
\begin{equation}
\label{algebraization-equation-guess}
M = \lim H^0(U, \mathcal{F}_n)
\end{equation}
을 살펴보는 것이 자연스럽다. $M \to H^0(U, \mathcal{F}_n)$은 $I^nM$을
소멸시키므로 $M$은 $I$-진 위상보다 (선험적으로) 더 거친 극한 위상을 가진다.
표준 사상들이 있다.
$$
\widetilde{M}|_U \to \widetilde{M/I^nM}|_U \to
\widetilde{H^0(U, \mathcal{F}_n)}|_U \to
\mathcal{F}_n
$$
$\widetilde{M}$을 $U$로 제한하면 연접 가군이 되고 $(\mathcal{F}_n)$이 이
가군의 완비가 되기를 기대할 수 있다. 그러나 $A$가 완비가 아니면 이것이
참일 가능성이 거의 없으므로 이는 순진한 기대이다. 더구나 $A$가 $I$-진
완비인 경우에도 이 개념은 다루기 매우 어렵다. 좀 더 적절한 방법은 다음
요건을 고찰하는 것이다.

\begin{definition}
\label{algebraization-definition-canonically-algebraizable}
상황 \ref{algebraization-situation-algebraize}에서 $(\mathcal{F}_n)$을
$\textit{Coh}(U, I\mathcal{O}_U)$의 대상이라 하자. $\QCoh(\mathcal{O}_X)$
안의 역계
$$
\{\widetilde{H^0(U, \mathcal{F}_n)}\}_{n \geq 1}
$$
가 $\textit{Coh}(X, I\mathcal{O}_X)$의 어떤 대상 $(\mathcal{G}_n)$과
전사영 동형이면 {\it $(\mathcal{F}_n)$은 $X$로 표준적으로 연장된다}고 한다.
\end{definition}

\noindent
보조정리 \ref{algebraization-lemma-canonically-algebraizable}에서 정의
\ref{algebraization-definition-canonically-algebraizable}의 조건이 정의
\ref{algebraization-definition-algebraizable}의 조건보다 강함을 보게 될 것이다.

\begin{lemma}
\label{algebraization-lemma-canonically-algebraizable}
상황 \ref{algebraization-situation-algebraize}에서 $(\mathcal{F}_n)$을
$\textit{Coh}(U, I\mathcal{O}_U)$의 대상이라 하자. $(\mathcal{F}_n)$이
$X$로 표준적으로 연장되면 다음이 성립한다.
\begin{enumerate}
\item $(\widetilde{H^0(U, \mathcal{F}_n)})$은 유일한 동형까지 유일한
$\textit{Coh}(X, I \mathcal{O}_X)$의 대상 $(\mathcal{G}_n)$과 전사영 동형이다.
\item $(\mathcal{G}_n)$을 $U$로 제한한 것은 $(\mathcal{F}_n)$과 동형이다.
즉 $(\mathcal{F}_n)$은 $X$로 연장된다.
\item 역계 $\{H^0(U, \mathcal{F}_n)\}$은 Mittag-Leffler 조건을 만족한다.
\item (\ref{algebraization-equation-guess})의 가군 $M$은 $A$의 $I$-진 완비 위에서 유한이고
$M$의 극한 위상은 $I$-진 위상이다.
\end{enumerate}
\end{lemma}

\begin{proof}
(1)의 $(\mathcal{G}_n)$의 존재는 정의
\ref{algebraization-definition-canonically-algebraizable}에서 따른다. (1)의
$(\mathcal{G}_n)$의 유일성은 보조정리
\ref{algebraization-lemma-recognize-formal-coherent-modules}에서 따른다.
$\mathcal{G}_n = \widetilde{M_n}$으로 쓰자. 그러면 $\{M_n\}$은
$M_n = M_{n + 1}/I^n M_{n + 1}$인 유한 $A$-가군의 역계이다. 정의
\ref{algebraization-definition-canonically-algebraizable}에 의해 역계
$\{H^0(U, \mathcal{F}_n)\}$은 $\{M_n\}$과 전사영 동형이다. 따라서 역계
$\{H^0(U, \mathcal{F}_n)\}$은 Mittag-Leffler 조건을 만족하고
$M = \lim M_n$이다(위상 가군으로서). 이제 (4)에 나오는 $M$의 성질들은
대수학의 보조정리 \ref{algebra-lemma-limit-complete},
\ref{algebra-lemma-finite-over-complete-ring}, 그리고
\ref{algebra-lemma-hathat-finitely-generated}.
$U$가 준아핀이므로 표준 사상
$$
\widetilde{H^0(U, \mathcal{F}_n)}|_U \to \mathcal{F}_n
$$
은 동형이다(성질의 보조정리
\ref{properties-lemma-quasi-coherent-quasi-affine}). 따라서
$(\mathcal{G}_n|_U)$와 $(\mathcal{F}_n)$은 전사영 동형이며, 보조정리
\ref{algebraization-lemma-recognize-formal-coherent-modules}에 의해 동형이라고 결론내린다.
\end{proof}

\begin{lemma}
\label{algebraization-lemma-canonically-extend-base-change}
상황 \ref{algebraization-situation-algebraize}에서 $(\mathcal{F}_n)$을
$\textit{Coh}(U, I\mathcal{O}_U)$의 대상이라 하자. $A \to A'$을 평탄한
환 준동형이라 하자. $X' = \Spec(A')$로 놓고, $U' \subset X'$를 $U$의
역상이라 하며, 유도된 사상을 $g : U' \to U$로 표기하자.
$(\mathcal{F}'_n) = (g^*\mathcal{F}_n)$로 놓는다. 스킴의 코호몰로지의
보조정리 \ref{coherent-lemma-inverse-systems-pullback}을 보라.
$(\mathcal{F}_n)$이 $X$로 표준적으로 연장되면 $(\mathcal{F}'_n)$은
$X'$로 표준적으로 연장된다. 또한 보조정리
\ref{algebraization-lemma-canonically-algebraizable}에서 $(\mathcal{F}_n)$에 대해 얻은
연장을 당기면 $(\mathcal{F}'_n)$의 연장이 된다.
\end{lemma}

\begin{proof}
$f : X' \to X$를 유도된 사상이라 하자. 평탄 밑변환에 의해
$H^0(U', \mathcal{F}'_n) = H^0(U, \mathcal{F}_n) \otimes_A A'$이다.
스킴의 코호몰로지의 보조정리
\ref{coherent-lemma-flat-base-change-cohomology}.
를 보라. 따라서 $\textit{Coh}(X, I\mathcal{O}_X)$의 $(\mathcal{G}_n)$이
$(\widetilde{H^0(U, \mathcal{F}_n)})$과 전사영 동형이면
$(f^*\mathcal{G}_n)$은 다음과 전사영 동형이다.
$$
(f^*\widetilde{H^0(U, \mathcal{F}_n)}) =
(\widetilde{H^0(U, \mathcal{F}_n) \otimes_A A'}) =
(\widetilde{H^0(U', \mathcal{F}'_n)})
$$
이로써 증명이 끝난다.
\end{proof}

\begin{lemma}
\label{algebraization-lemma-when-done}
상황 \ref{algebraization-situation-algebraize}에서 $(\mathcal{F}_n)$을
$\textit{Coh}(U, I\mathcal{O}_U)$의 대상이라 하자. $M$을
(\ref{algebraization-equation-guess})에서와 같이 놓자. 다음을 가정한다.
\begin{enumerate}
\item[(a)] 역계 $H^0(U, \mathcal{F}_n)$은 Mittag-Leffler 조건을 만족한다.
\item[(b)] $M$의 극한 위상은 $I$-진 위상과 일치한다.
\item[(c)] 모든 $n$에 대해 $M \to H^0(U, \mathcal{F}_n)$의 상은 유한
$A$-가군이다.
\end{enumerate}
그러면 $(\mathcal{F}_n)$은 $X$로 표준적으로 연장된다. 특히 $A$가
$I$-진 완비이면 $(\mathcal{F}_n)$은 연접 $\mathcal{O}_U$-가군의 완비이다.
\end{lemma}

\begin{proof}
$H^0(U, \mathcal{F}_n)$이 Mittag-Leffler 조건을 만족하고 $M$의 극한 위상이
$I$-진 위상이므로, $\{M/I^nM\}$과 $\{H^0(U, \mathcal{F}_n)\}$은 서로
전사영 동형인 $A$-가군의 역계이다. 따라서
$$
\mathcal{G}_n = \widetilde{M/I^n M}
$$
로 놓으면 정의 \ref{algebraization-definition-canonically-algebraizable}의 조건을 확인하기
위해 $M$이 $A$의 $I$-진 완비 위에서 유한 가군임을 보이는 것으로 충분하다.
이는 조건 (c)와 위의 사실에 의해 $M/I^n M$이 유한이라는 점 및 대수학의
보조정리 \ref{algebra-lemma-finite-over-complete-ring}에서 따른다.
\end{proof}

\noindent
다음은 어떤 의미에서 위 보조정리 \ref{algebraization-lemma-when-done}의 가장 직접적인
응용이다.

\begin{lemma}
\label{algebraization-lemma-algebraization-principal-variant}
상황 \ref{algebraization-situation-algebraize}에서 $(\mathcal{F}_n)$을
$\textit{Coh}(U, I\mathcal{O}_U)$의 대상이라 하자. 다음을 가정한다.
\begin{enumerate}
\item 영인자가 아닌 $f \in \mathfrak a$에 대해 $I = (f)$는 주 아이디얼이다.
\item $\mathcal{F}_n$은 유한 국소 자유
$\mathcal{O}_U/f^n\mathcal{O}_U$-가군이다.
\item $H^1_\mathfrak a(A/fA)$와 $H^2_\mathfrak a(A/fA)$는 유한
$A$-가군이다.
\end{enumerate}
그러면 $(\mathcal{F}_n)$은 $X$로 표준적으로 연장된다. 특히 $A$가
완비이면 $(\mathcal{F}_n)$은 연접 $\mathcal{O}_U$-가군의 완비이다.
\end{lemma}

\begin{proof}
보조정리 \ref{algebraization-lemma-when-done}의 가정 (a), (b), (c)를 확인하여 증명하겠다.

\medskip\noindent
$\mathcal{F}_n$은 $\mathcal{O}_U/f^n\mathcal{O}_U$ 위에서 국소 자유이므로
모든 $n$에 대해 짧은 완전열
$0 \to \mathcal{F}_n \to \mathcal{F}_{n + 1} \to \mathcal{F}_1 \to 0$을
얻는다. 따라서 코호몰로지의 보조정리
\ref{cohomology-lemma-topology-I-adic-f}에 의해 조건 (b)가 성립한다.

\medskip\noindent
$f$가 영인자가 아니므로 짧은 완전열
$$
0 \to A/f^nA \xrightarrow{f} A/f^{n + 1}A \to A/fA \to 0
$$
과 이에 대응하는 짧은 완전열
$0 \to \mathcal{F}_n \to \mathcal{F}_{n + 1} \to \mathcal{F}_1 \to 0$을
얻는다. 국소 코호몰로지의 보조정리
\ref{local-cohomology-lemma-finiteness-pushforwards-and-H1-local}
를 더 언급하지 않고 사용하겠다. 가정들로부터
$H^0(U, \mathcal{O}_U/f\mathcal{O}_U)$와
$H^1(U, \mathcal{O}_U/f\mathcal{O}_U)$
은 유한 $A$-가군이다. 따라서 $\mathcal{F}_1$에 대해서도 같은 사실이
성립한다. 국소 코호몰로지의 보조정리
\ref{local-cohomology-lemma-finiteness-for-finite-locally-free}.
를 보라. 귀납법과 짧은 완전열들을 사용하면 모든 $n$에 대해
$H^0(U, \mathcal{F}_n)$이 유한 $A$-가군임을 알 수 있다. 이로써 가정 (c)가
성립한다.

\medskip\noindent
마지막으로 $H^1(U, \mathcal{F}_1)$은 유한 $A$-가군이므로 코호몰로지의
보조정리 \ref{cohomology-lemma-ML}을 적용하여 가정 (a)가 성립함을 안다.
\end{proof}

\begin{remark}
\label{algebraization-remark-interesting-case-variant}
보조정리 \ref{algebraization-lemma-algebraization-principal-variant}에서 $A$가
Cohen--Macaulay 형식 올을 갖는 보편적 사슬환이면(예를 들어 $A$가
쌍대화 복합체를 가지면),
$H^1_\mathfrak a(A/fA)$와 $H^2_\mathfrak a(A/fA)$
이 유한 $A$-가군이라는 조건은
$$
\text{depth}((A/f)_\mathfrak p) + \dim((A/\mathfrak p)_\mathfrak q) > 2
$$
가 모든 $\mathfrak p \in V(f) \setminus V(\mathfrak a)$와
$\mathfrak q \in V(\mathfrak p) \cap V(\mathfrak a)$에 대해 성립하는 것과
동치이다. 국소 코호몰로지의 정리
\ref{local-cohomology-theorem-finiteness}를 보라.

\medskip\noindent
예를 들어 $A/fA$가 $(S_2)$이고 $Z = V(\mathfrak a)$의 모든 기약 성분이
$Y = \Spec(A/fA)$ 안에서 여차원 $\geq 3$이면,
$H^1_\mathfrak a(A/fA)$와 $H^2_\mathfrak a(A/fA)$의 유한성을 얻는다.
이는 보조정리 \ref{algebraization-lemma-algebraization-principal}에 나오는 조금 더 약한
조건들과 대조해야 한다(비고 \ref{algebraization-remark-interesting-case}도 보라).
\end{remark}






\section{연접 형식 가군의 대수화 II}
\label{algebraization-section-algebraization-modules-general}

\noindent
절 \ref{algebraization-section-algebraization-modules}에서 시작한 논의를 계속한다.
처음 읽을 때에는 이 절을 건너뛰어도 된다.

\begin{lemma}
\label{algebraization-lemma-map-kernel-cokernel-on-closed}
상황 \ref{algebraization-situation-algebraize}에서 핵과 여핵이 $I$의 어떤 거듭제곱에 의해
소멸되는 $\textit{Coh}(U, I\mathcal{O}_U)$의 사상
$(\mathcal{F}_n) \to (\mathcal{F}'_n)$을 취하자. 그러면
\begin{enumerate}
\item $(\mathcal{F}_n)$이 $X$로 연장될 필요충분조건은 $(\mathcal{F}'_n)$이
$X$로 연장되는 것이다.
\item $(\mathcal{F}_n)$이 연접 $\mathcal{O}_U$-가군의 완비일 필요충분조건은
$(\mathcal{F}'_n)$이 그러한 것이다.
\end{enumerate}
\end{lemma}

\begin{proof}
(2)는 스킴의 코호몰로지의 보조정리
\ref{coherent-lemma-existence-easy}에서 곧바로 따른다. (1)을 보기 위해 먼저
보조정리 \ref{algebraization-lemma-algebraizable}를 사용하여 $A$가 $I$-진 완비인 경우로
환원한다. 이 경우에는 보조정리 \ref{algebraization-lemma-essential-image-completion}에 의해
(1)이 (2)로 환원된다.
\end{proof}

\noindent
다음 두 보조정리는 본래 보조정리 \ref{algebraization-lemma-when-done}의 증명에서
사용되었다. 어떤 중간 결과들을 얻을 수 있는지 알고 싶은 독자를 위해
여기에 남겨 둔다.

\begin{lemma}
\label{algebraization-lemma-when-ML}
상황 \ref{algebraization-situation-algebraize}에서 $(\mathcal{F}_n)$을
$\textit{Coh}(U, I\mathcal{O}_U)$의 대상이라 하자. 역계
$H^0(U, \mathcal{F}_n)$이 Mittag-Leffler 조건을 만족하면 표준 사상
$$
\widetilde{M/I^nM}|_U \to \mathcal{F}_n
$$
은 모든 $n$에 대해 전사이다. 여기서 $M$은 (\ref{algebraization-equation-guess})에서와 같다.
\end{lemma}

\begin{proof}
전사성은 어떤 점 $y \in Y \setminus Z$에서의 줄기에서 확인할 수 있다.
$y$가 소아이디얼 $\mathfrak q \subset A$에 대응하면
$f \in \mathfrak a$, $f \not \in \mathfrak q$를 택할 수 있다. 그러면
$$
M_f \longrightarrow H^0(U, \mathcal{F}_n)_f = H^0(D(f), \mathcal{F}_n)
$$
가 전사임을 보이면 충분하다. $D(f)$가 아핀이기 때문이다(등식은 성질의
보조정리 \ref{properties-lemma-invert-f-sections}에 의해 성립한다).
Mittag-Leffler 성질에 의해 어떤 $m \geq n$에 대해
$$
\Im(M \to H^0(U, \mathcal{F}_n)) =
\Im(H^0(U, \mathcal{F}_m) \to H^0(U, \mathcal{F}_n))
$$
이다. 코호몰로지 긴 완전열을 사용하면
$$
\Coker(H^0(U, \mathcal{F}_m) \to H^0(U, \mathcal{F}_n))
\subset
H^1(U, \Ker(\mathcal{F}_m \to \mathcal{F}_n))
$$
$U = X \setminus V(\mathfrak a)$이므로 이 $H^1$은
$\mathfrak a$-거듭제곱 꼬임이다. 따라서 $f$를 가역화하면 여핵은 영이 된다.
\end{proof}

\begin{lemma}
\label{algebraization-lemma-when-topology}
상황 \ref{algebraization-situation-algebraize}에서 $(\mathcal{F}_n)$을
$\textit{Coh}(U, I\mathcal{O}_U)$의 대상이라 하자. $M$을
(\ref{algebraization-equation-guess})에서와 같이 놓고
$$
\mathcal{G}_n = \widetilde{M/I^nM}.
$$
로 놓자. $M$의 극한 위상이 $I$-진 위상과 일치하면
$\mathcal{G}_n|_U$는 연접 $\mathcal{O}_U$-가군이고 역계의 사상
$$
(\mathcal{G}_n|_U) \longrightarrow (\mathcal{F}_n)
$$
은 아벨 범주 $\textit{Coh}(U, I\mathcal{O}_U)$에서 단사이다.
\end{lemma}

\begin{proof}
$\mathcal{G}_n$은 $I^n$에 의해 소멸되는 준연접
$\mathcal{O}_X$-가군이고
$\mathcal{G}_{n + 1}/I^n\mathcal{G}_{n + 1} = \mathcal{G}_n$.
이다. 다음을 고찰하자.
$$
M_n = \Im(M \longrightarrow H^0(U, \mathcal{F}_n))
$$
가정에 따르면 역계 $(M_n)$과 $(M/I^nM)$은 $\text{Mod}_A$의 전사영 대상으로서
동형이다. $D(f) \subset U$가 아핀 열린집합이 되도록
$f \in \mathfrak a$를 택하자. 그러면
$$
(M_n)_f \subset H^0(U, \mathcal{F}_n)_f = H^0(D(f), \mathcal{F}_n)
$$
이다. 등식은 성질의 보조정리 \ref{properties-lemma-invert-f-sections}에 의해
성립한다. 따라서 $\widetilde{M_n}|_U \to \mathcal{F}_n$은 단사이고,
$\widetilde{M_n}|_U$는 연접 가군이다(스킴의 코호몰로지의 보조정리
\ref{coherent-lemma-coherent-Noetherian-quasi-coherent-sub-quotient}).
). $M \to M/I^nM$은 전사이고 어떤 $n' \geq n$에 대해
$M_{n'} \to M/I^nM$을 통해 분해되므로, $\mathcal{G}_n|_U$는 연접 가군의
몫으로서 연접이다. 이를 증명 첫머리의 관찰과 결합하면
$(\mathcal{G}_n|_U)$가 실제로 $\textit{Coh}(U, I\mathcal{O}_U)$의 대상을
이룸을 얻는다. 마지막으로 사상의 단사성을 보이기 위해
$$
\lim (M/I^nM)_f = \lim H^0(D(f), \mathcal{G}_n) \to
\lim H^0(D(f), \mathcal{F}_n)
$$
가 단사임을 보이면 충분하다. 스킴의 코호몰로지의 보조정리
\ref{coherent-lemma-inverse-systems-abelian}과
\ref{coherent-lemma-inverse-systems-affine}.
들을 보라. $\lim (M_n)_f \to \lim H^0(D(f), \mathcal{F}_n)$의 단사성은
명백하고(위를 보라), 전사영계에 관한 관찰에 의해
$\lim (M_n)_f = \lim (M/I^nM)_f$이다. 이로써 증명이 끝난다.
\end{proof}





\section{거리 함수}
\label{algebraization-section-distance}

\noindent
$Y$를 Noether 스킴, $Z \subset Y$를 닫힌 부분집합이라 하자. $Y$의 점에서
$Z$까지의 ``거리''를 측정하는 함수
\begin{equation}
\label{algebraization-equation-delta-Z}
\delta^Y_Z = \delta_Z : Y \longrightarrow \mathbf{Z}_{\geq 0} \cup \{\infty\}
\end{equation}
를 정의한다. 직관적인 논의는 비고 \ref{algebraization-remark-discussion}을 보라.
$y \in Y$라 하자. $y$가 $Z$와 만나지 않는 $Y$의 연결 성분에 속하면
$\delta_Z(y) = \infty$로 놓는다. $y$가 $Z$와 만나는 $Y$의 연결 성분에
속하면 어떤 $k \geq 0$와 $Y$의 정수적 닫힌 부분스킴들의 계
$$
V_0 \subset W_0 \supset V_1 \subset W_1 \supset \ldots \supset
V_k \subset W_k
$$
를 찾을 수 있는데, 여기서 $V_0 \subset Z$이고 $y \in W_k$는 일반점이다.
$i = 0, \ldots, k$에 대해 $c_i = \text{codim}(V_i, W_i)$로,
$i = 0, \ldots, k - 1$에 대해
$b_i = \text{codim}(V_{i + 1}, W_i)$로 놓는다. 이러한 계에 대해
$$
\delta(V_0, W_0, V_1, \ldots, W_k) =
k +
\max_{i = 0, 1, \ldots, k}
(c_i + c_{i + 1} + \ldots + c_k - b_i - b_{i + 1} - \ldots - b_{k - 1})
$$
$i = k$를 택할 수 있고 $c_k \geq 0$이므로 이는 $\geq k$이다. 마지막으로
$$
\delta_Z(y) = \min \delta(V_0, W_0, V_1, \ldots, W_k)
$$
로 놓는다. 여기서 최솟값은 위와 같은 $Y$의 정수적 닫힌 부분스킴들의 모든
계에 걸쳐 취한다.

\begin{lemma}
\label{algebraization-lemma-discussion}
$Y$를 Noether 스킴, $Z \subset Y$를 닫힌 부분집합이라 하자.
\begin{enumerate}
\item $y \in Y$에 대해 $\delta_Z(y) = 0 \Leftrightarrow y \in Z$이다.
\item 부분집합 $\{y \in Y \mid \delta_Z(y) \leq k\}$는 특수화에 대해
안정적이다.
\item $y \in Y$ 및 $z \in \overline{\{y\}} \cap Z$에 대해
$\dim(\mathcal{O}_{\overline{\{y\}}, z}) \geq \delta_Z(y)$.
\item $\delta$가 $Y$ 위의 차원 함수이면
$\delta_Z(y) \geq \delta(y) - \delta_{max}$이다. 여기서 $\delta_{max}$는
$Z$ 위에서 $\delta$의 최댓값이다.
\item $Y = \Spec(A)$가 극대 아이디얼 $\mathfrak m$을 갖는 사슬적 Noether
국소환의 스펙트럼이고 $Z = \{\mathfrak m\}$이면
$\delta_Z(y) = \dim(\overline{\{y\}})$.
\item $Y' \subset Y$가 열린 부분스킴이면
$y' \in Y'$에 대해
$\delta^{Y'}_{Y' \cap Z}(y') \geq \delta^Y_Z(y')$이다.
\end{enumerate}
$Y$가 사슬적이라고 가정하자. 그러면
\begin{enumerate}
\setcounter{enumi}{6}
\item $y' \leadsto y$를 $Y$의 점들의 바로 다음 특수화라 하자.
$Y$가 사슬적이면 $\delta_Z(y') \leq \delta_Z(y) + 1$이다.
\item 다음 특수화들의 도식이 주어졌다고 하자.
$$
\xymatrix{
& y'_0 \ar@{~>}[ld] \ar@{~>}[rd] &
& y'_1 \ar@{~>}[ld] & \ldots
& y'_{k - 1} \ar@{~>}[rd] &
\\
y_0 & &
y_1 & &
\ldots & &
y_k = y
}
$$
여기서 점들은 $Y$에 속하고 $y_0 \in Z$이며 $y_i' \leadsto y_i$는
바로 다음 특수화이다. 그러면 $\delta_Z(y_k) \leq k$이다.
\end{enumerate}
\end{lemma}

\begin{proof}
(1)을 증명하자. $y \in Z$이면 $k = 0$ 및
$V_0 = W_0 = \overline{\{y\}}$로 택할 수 있고
$\delta(V_0, W_0) = 0$을 얻으므로 $\delta_Z(y) = 0$이다.
$y \not \in Z$이면 $y$에 대한 모든 계
$V_0 \subset W_0 \supset V_1 \subset W_1 \supset \ldots \subset W_k$
에서 $k = 0$이고 $V_0 \not = W_0$이거나 $k > 0$이다. 어느 경우든
$\delta(V_0, W_0, \ldots, W_k) > 0$이다. 따라서 $\delta_Z(y) > 0$이다.

\medskip\noindent
(2)를 증명하자. $y \leadsto y'$를 비자명 특수화라 하고
$V_0 \subset W_0 \supset V_1 \subset W_1 \supset \ldots \subset W_k$
를 $y$에 대한 계라 하자. 두 경우가 있다.
경우 I: $V_k = W_k$, 즉 $c_k = 0$이다. 이 경우
$V'_k = W'_k = \overline{\{y'\}}$로 놓을 수 있다. 간단한 계산으로
$\delta(V_0, W_0, \ldots, V'_k, W'_k) \leq
\delta(V_0, W_0, \ldots, V_k, W_k)$
임을 안다. $b_{k - 1}$만 더 큰 정수로 바뀌기 때문이다.
경우 II: $V_k \not = W_k$, 즉 $c_k > 0$이다. 이 경우
$V_{k + 1} = W_{k + 1} = \overline{\{y'\}}$로 놓으면
$V_0 \subset W_0 \supset \ldots \subset W_k \supset V_{k + 1}
\subset W_{k + 1}$은 $y'$에 대한 계이다. 이제 $c_{k + 1} = 0$이고
$b_k > 0$이므로 다음을 얻는다.
\begin{align*}
& \delta(V_0, \ldots, W_{k + 1}) \\
& =
k + 1 +
\max_{i = 0, 1, \ldots, k + 1}
(c_i + c_{i + 1} + \ldots + c_k + c_{k + 1}
- b_i - b_{i + 1} - \ldots - b_{k - 1} - b_k) \\
& =
k + 1 +
\max_{i = 0, 1, \ldots, k + 1}
(c_i + c_{i + 1} + \ldots + c_k
- b_i - b_{i + 1} - \ldots - b_{k - 1} - b_k) \\
& \leq
k + \max_{i = 0, 1, \ldots, k}
(c_i + c_{i + 1} + \ldots + c_k - b_i - b_{i + 1} - \ldots - b_{k - 1}) \\
& = \delta(V_0, \ldots, W_k)
\end{align*}
부등식이 성립하는 까닭은 $c_k > 0$이고 $b_k > 0$이어서 특히
$\delta(V_0, \ldots, W_k) \geq k + c_k \geq k + 1$이기 때문이다.

\medskip\noindent
(3)을 증명하자. $y \in Y$와 $z \in \overline{\{y\}} \cap Z$가 주어지면
계
$$
V_0 = \overline{\{z\}} \subset W_0 = \overline{\{y\}}
$$
를 얻고, 성질의 보조정리 \ref{properties-lemma-codimension-local-ring}에 의해
$c_0 = \text{codim}(V_0, W_0) = \dim(\mathcal{O}_{\overline{\{y\}}, z})$이다.
따라서 $\delta(V_0, W_0) = 0 + c_0 = c_0$이고 원하는 결론이 따른다.

\medskip\noindent
(4)를 증명하자. $\delta$를 $Y$ 위의 차원 함수라 하자.
$V_0 \subset W_0 \supset V_1 \subset W_1 \supset \ldots \subset W_k$
를 $y$에 대한 계라 하자. $y'_i \in W_i$와 $y_i \in V_i$를 일반점들이라
하자. 그러면 $y_0 \in Z$이고 $y_k = y$이며
$$
\delta(y_i) - \delta(y_{i - 1}) =
\delta(y'_{i - 1}) - \delta(y_{i - 1}) - \delta(y'_{i - 1}) + \delta(y_i) =
c_{i - 1} - b_{i - 1}
$$
마지막으로 $\delta(y'_k) - \delta(y_{k - 1}) = c_k$이다. 따라서
$$
\delta(y) - \delta(y_0) =
c_0 + \ldots + c_k - b_0 - \ldots - b_{k - 1}
$$
를 얻는다. 이로부터
$\delta(V_0, W_0, \ldots, W_k) \geq k + \delta(y) - \delta(y_0)$이고,
원하는 결론이 따른다.

\medskip\noindent
(5)를 증명하자. 함수 $\delta(y) = \dim(\overline{\{y\}})$는 차원 함수이다.
따라서 (4)에 의해 $\delta(y) \leq \delta_Z(y)$이다. (3)에 의해
$\delta_Z(y) \leq \delta(y)$이므로 끝난다.

\medskip\noindent
(6)을 증명하자. $\delta^{Y'}_{Y' \cap Z}$의 계산에서 고찰할 계가 더
적으므로 명백하다.

\medskip\noindent
(7)을 증명하자.
$V_0 \subset W_0 \supset V_1 \subset W_1 \supset \ldots \subset W_k$
를 $y$에 대한 계라 하자. $W'_k = \overline{\{y'\}}$로 놓는다.
$Y$가 사슬적이므로
$\text{codim}(V_k, W'_k) = \text{codim}(V_k, W_k) + 1$.
이로부터 $\delta(V_0, \ldots, V_k, W'_k) \leq
\delta(V_0, \ldots, V_k, W_k) + 1$이 쉽게 따른다. 원하는 결론이다.

\medskip\noindent
(8)은 (7)과 (2)를 결합하면 된다.
\end{proof}

\begin{lemma}
\label{algebraization-lemma-change-distance-function}
$Y$를 보편적 사슬 Noether 스킴, $Z \subset Y$를 닫힌 부분스킴이라 하자.
$f : Y' \to Y$를 모든 올의 차원이 $\leq e$인 유한형 사상이라 하자.
$Z' = f^{-1}(Z)$로 놓자. 그러면
$$
\delta_Z(y) \leq \delta_{Z'}(y') + e - \text{trdeg}_{\kappa(y)}(\kappa(y'))
$$
이다. 여기서 $y' \in Y'$의 상은 $y \in Y$이다.
\end{lemma}

\begin{proof}
$\delta_{Z'}(y') = \infty$이면 증명할 것이 없다.
$\delta_{Z'}(y') < \infty$이면 정수적 닫힌 부분스킴들의 계
$$
V'_0 \subset W'_0 \supset V'_1 \subset W'_1 \supset \ldots \subset W'_k
$$
$Y'$의 계를 택하되 $V'_0 \subset Z'$이고 $y'$가 $W'_k$의 일반점이며
$\delta_{Z'}(y') = \delta(V'_0, W'_0, \ldots, W'_k)$가 되게 한다. 위
스킴들의 $Y$ 안에서의 스킴론적 상을
$$
V_0 \subset W_0 \supset V_1 \subset W_1 \supset \ldots \subset W_k
$$
로 표기하자. $y$는 $W_k$의 일반점이고 $V_0 \subset Z$임을 주목하라.
각 $i$에 대해 그림
$$
\xymatrix{
V'_i \ar[r] \ar[d] & W'_i \ar[d] & V'_{i + 1} \ar[l] \ar[d] \\
V_i \ar[r] & W_i & V_{i + 1} \ar[l]
}
$$
$V'_i/V_i$의 상대차원을 $n_i$, $W'_i/W_i$의 상대차원을 $m_i$로
표기하자. 더 정확히 말해 이들은 대응하는 함수체 확대의 초월차수이다.
$c_i = \text{codim}(V_i, W_i)$,
$c'_i = \text{codim}(V'_i, W'_i)$,
$b_i = \text{codim}(V_{i + 1}, W_i)$ 및
$b'_i = \text{codim}(V'_{i + 1}, W'_i)$로 놓자. 차원 공식에 의해
$$
c_i = c'_i + n_i - m_i
\quad\text{이고}\quad
b_i = b'_i + n_{i + 1} - m_i
$$
이다. 사상의 보조정리 \ref{morphisms-lemma-dimension-formula}를 보라.
따라서 $c_i - b_i = c'_i - b'_i + n_i - n_{i + 1}$이고 다음을 얻는다.
\begin{align*}
& c_i + c_{i + 1} + \ldots + c_k - b_i - b_{i + 1} - \ldots - b_{k - 1} \\
& =
c'_i + c'_{i + 1} + \ldots + c'_k - b'_i - b'_{i + 1} - \ldots - b'_{k - 1}
+ n_i - n_k + c_k - c'_k \\
& =
c'_i + c'_{i + 1} + \ldots + c'_k - b'_i - b'_{i + 1} - \ldots - b'_{k - 1}
+ n_i - m_k
\end{align*}
따라서
\begin{align*}
\max_{i = 0, \ldots, k}
& (c_i + c_{i + 1} + \ldots + c_k - b_i - b_{i + 1} - \ldots - b_{k - 1}) \\
& =
\max_{i = 0, \ldots, k}
(c'_i + c'_{i + 1} + \ldots + c'_k - b'_i - b'_{i + 1} - \ldots - b'_{k - 1}
+ n_i - m_k) \\
& =
\max_{i = 0, \ldots, k}
(c'_i + c'_{i + 1} + \ldots + c'_k - b'_i - b'_{i + 1} - \ldots - b'_{k - 1}
+ n_i) - m_k \\
& \leq
\max_{i = 0, \ldots, k}
(c'_i + c'_{i + 1} + \ldots + c'_k - b'_i - b'_{i + 1} - \ldots - b'_{k - 1})
+ e - m_k
\end{align*}
$m_k = \text{trdeg}_{\kappa(y)}(\kappa(y'))$이므로
$$
\delta(V_0, W_0, \ldots, W_k) \leq
\delta(V'_0, W'_0, \ldots, W'_k) + e - \text{trdeg}_{\kappa(y)}(\kappa(y'))
$$
를 얻으며, 이것이 원하는 부등식이다.
\end{proof}

\begin{remark}
\label{algebraization-remark-discussion}
$Y$를 사슬적 Noether 스킴, $Z \subset Y$를 닫힌 부분집합이라 하자.
보조정리 \ref{algebraization-lemma-discussion}에 의해
$$
\delta_Z(y) \leq \min
\left\{ k \middle|
\begin{matrix}
\text{ 다음과 같은 }Y\text{ 안의 특수화들이 존재한다} \\
y_0 \leftarrow y'_0 \rightarrow y_1 \leftarrow y'_1 \rightarrow \ldots
\leftarrow y'_{k - 1} \rightarrow y_k = y \\
\text{ 여기서 }y_0 \in Z\text{이고 }y_i' \leadsto y_i
\text{는 바로 다음 특수화이다}
\end{matrix}
\right\}
$$
이다. $Y$가 어떤 체 위에서 유한형이면 등식이 성립한다고 주장한다.
이 결과가 필요해지면 정확한 명제를 서술하고 여기서 증명할 것이다.
그러나 일반적으로 이 부등식의 우변으로 $\delta_Z$를 정의하면 보조정리
\ref{algebraization-lemma-change-distance-function}이 여전히 참인지는 알 수 없다.
\end{remark}

\begin{example}
\label{algebraization-example-distance}
$k$를 체라 하고 $Y = \mathbf{A}^n_k$라 하자. 통상적인 차원 함수를
$\delta : Y \to \mathbf{Z}_{\geq 0}$로 표기하자.
\begin{enumerate}
\item 어떤 닫힌점 $z$에 대해 $Z = \{z\}$이면
\begin{enumerate}
\item $y \leadsto z$이면 $\delta_Z(y) = \delta(y)$이고,
\item $y \not \leadsto z$이면 $\delta_Z(y) = \delta(y) + 1$이다.
\end{enumerate}
\item $Z$가 닫힌 부분다양체이고 $W = \overline{\{y\}}$이면
\begin{enumerate}
\item $W \subset Z$이면 $\delta_Z(y) = 0$이다.
\item $Z$가 $W$에 포함되면 $\delta_Z(y) = \dim(W) - \dim(Z)$이다.
\item $\dim(W) \leq \dim(Z)$이고 $W \not \subset Z$이면
$\delta_Z(y) = 1$이다.
\item $\dim(W) > \dim(Z)$이고 $Z \not \subset W$이면
$\delta_Z(y) = \dim(W) - \dim(Z) + 1$이다.
\end{enumerate}
\end{enumerate}
경우 (1)의 일반화로, $Y$가 어떤 체 위에서 유한형이고
$Z = \{z\}$가 닫힌점이라고 하자. 그러면
$\delta_Z(y) = \delta(y) + t$이다. 여기서 $t$는 $z$와
$\overline{\{y\}}$의 닫힌점을 잇는 곡선 사슬의 최소 길이이다.
\end{example}







\section{연접 형식 가군의 대수화 III}
\label{algebraization-section-uniqueness}

\noindent
절 \ref{algebraization-section-algebraization-modules}과
\ref{algebraization-section-algebraization-modules-general}에서 시작한 논의를 계속한다.
절 \ref{algebraization-section-distance}의 거리 함수를 사용하여 상황
\ref{algebraization-situation-algebraize}의 연접 형식 가군에 관한 몇 가지 자연스러운
조건을 서술할 것이다.

\medskip\noindent
상황 \ref{algebraization-situation-algebraize}에서 점 $y \in U \cap Y$가 주어지면
$I$-진 완비
$$
\mathcal{O}_{X, y}^\wedge = \lim \mathcal{O}_{X, y}/I^n\mathcal{O}_{X, y}
$$
를 고찰할 수 있다. 이는 극대 아이디얼 $\mathfrak m_y^\wedge$을 가지며
$I\mathcal{O}_{X, y}^\wedge$에 관해 완비인 Noether 국소환이다. 대수학의 절
\ref{algebra-section-completion-noetherian}을 보라. $(\mathcal{F}_n)$을
$\textit{Coh}(U, I\mathcal{O}_U)$의 대상이라 하자. $y$에서
$(\mathcal{F}_n)$의 ``줄기''를 다음 식으로 정의한다.
$$
\mathcal{F}_y^\wedge = \lim \mathcal{F}_{n, y}
$$
이는 $\mathcal{O}_{X, y}^\wedge$ 위의 유한 가군이다. 대수학의 보조정리
\ref{algebra-lemma-limit-complete}과
\ref{algebra-lemma-finite-over-complete-ring}을 보라.

\begin{definition}
\label{algebraization-definition-s-d-inequalities}
상황 \ref{algebraization-situation-algebraize}에서 $(\mathcal{F}_n)$을
$\textit{Coh}(U, I\mathcal{O}_U)$의 대상이라 하고 $a, b$를 정수라 하자.
$\delta^Y_Z$는 (\ref{algebraization-equation-delta-Z})에서와 같이 놓는다.
$y \in U \cap Y$와
$\mathfrak p \not \in V(I\mathcal{O}_{X, y}^\wedge)$인 소아이디얼
$\mathfrak p \subset \mathcal{O}_{X, y}^\wedge$에 대해 다음 조건들이
성립하면 {\it $(\mathcal{F}_n)$은 $(a, b)$-부등식들을 만족한다}고 한다.
\begin{enumerate}
\item $V(\mathfrak p) \cap V(I\mathcal{O}_{X, y}^\wedge) \not =
\{\mathfrak m_y^\wedge\}$이면
$$
\text{depth}((\mathcal{F}^\wedge_y)_\mathfrak p) + \delta^Y_Z(y) \geq a
\quad\text{또는}\quad
\text{depth}((\mathcal{F}^\wedge_y)_\mathfrak p) +
\dim(\mathcal{O}_{X, y}^\wedge/\mathfrak p) + \delta^Y_Z(y) > b
$$
\item $V(\mathfrak p) \cap V(I\mathcal{O}_{X, y}^\wedge) =
\{\mathfrak m_y^\wedge\}$이면
$$
\text{depth}((\mathcal{F}^\wedge_y)_\mathfrak p) + \delta^Y_Z(y) > a
$$
\end{enumerate}
다음 조건이 성립하면 {\it $(\mathcal{F}_n)$은 엄격한
$(a, b)$-부등식들을 만족한다}고 한다.
$y \in U \cap Y$와
$\mathfrak p \not \in V(I\mathcal{O}_{X, y}^\wedge)$인 소아이디얼
$\mathfrak p \subset \mathcal{O}_{X, y}^\wedge$에 대해
$$
\text{depth}((\mathcal{F}^\wedge_y)_\mathfrak p) + \delta^Y_Z(y) > a
\quad\text{또는}\quad
\text{depth}((\mathcal{F}^\wedge_y)_\mathfrak p) +
\dim(\mathcal{O}_{X, y}^\wedge/\mathfrak p) + \delta^Y_Z(y) > b
$$
\end{definition}

\noindent
몇 가지 초등적인 관찰을 제시한다.

\begin{lemma}
\label{algebraization-lemma-elementary}
상황 \ref{algebraization-situation-algebraize}에서 $(\mathcal{F}_n)$을
$\textit{Coh}(U, I\mathcal{O}_U)$의 대상이라 하고 $a, b$를 정수라 하자.
\begin{enumerate}
\item $(\mathcal{F}_n)$이 $I$의 어떤 거듭제곱에 의해 소멸되면, 임의의
$a, b$에 대해 $(\mathcal{F}_n)$은 $(a, b)$-부등식들을 만족한다.
\item $(\mathcal{F}_n)$이 $(a + 1, b)$-부등식들을 만족하면 엄격한
$(\mathcal{F}_n)$은 $(a, b)$-부등식들을 만족한다.
\end{enumerate}
$\text{cd}(A, I) \leq d$이고 $A$가 쌍대화 복합체를 가지면 다음이 성립한다.
\begin{enumerate}
\item[(3)] $(\mathcal{F}_n)$이 $(s, s + d)$-부등식들을 만족할 필요충분조건은
모든 $y \in U \cap Y$에 대해 튜플
$\mathcal{O}_{X, y}^\wedge, I\mathcal{O}_{X, y}^\wedge,
\{\mathfrak m_y^\wedge\}, \mathcal{F}_y^\wedge, s - \delta^Y_Z(y), d$
이 상황 \ref{algebraization-situation-bootstrap}에서와 같은 것이다.
\item[(4)]
$(\mathcal{F}_n)$이 엄격한 $(s, s + d)$-부등식들을 만족하면
$(\mathcal{F}_n)$은 $(s, s + d)$-부등식들을 만족한다.
\end{enumerate}
\end{lemma}

\begin{proof}
(4)를 제외하면 즉시 알 수 있다. (4)는 보조정리
\ref{algebraization-lemma-bootstrap-bis-bis}와 (3)의 번역에서 따른다.
\end{proof}

\begin{lemma}
\label{algebraization-lemma-explain-2-3-cd-1}
상황 \ref{algebraization-situation-algebraize}에서 $(\mathcal{F}_n)$을
$\textit{Coh}(U, I\mathcal{O}_U)$의 대상이라 하자.
$\text{cd}(A, I) = 1$이면 $\mathcal{F}$가 $(2, 3)$-부등식들을 만족할
필요충분조건은
$$
\text{depth}((\mathcal{F}^\wedge_y)_\mathfrak p) +
\dim(\mathcal{O}_{X, y}^\wedge/\mathfrak p) + \delta^Y_Z(y) > 3
$$
가 모든 $y \in U \cap Y$와
$\mathfrak p \not \in V(I\mathcal{O}_{X, y}^\wedge)$인
$\mathfrak p \subset \mathcal{O}_{X, y}^\wedge$에 대해 성립하는 것이다.
\end{lemma}

\begin{proof}
소아이디얼 $\mathfrak p \subset \mathcal{O}_{X, y}^\wedge$에 대해
$\mathfrak p \not \in V(I\mathcal{O}_{X, y}^\wedge)$이면
$V(\mathfrak p) \cap V(I\mathcal{O}_{X, y}^\wedge) =
\{\mathfrak m_y^\wedge\}
\Leftrightarrow \dim(\mathcal{O}_{X, y}^\wedge/\mathfrak p) = 1$
임을 주목하자. $\text{cd}(A, I) = 1$이기 때문이다. 국소 코호몰로지의
보조정리
\ref{local-cohomology-lemma-cd-change-rings}과
\ref{local-cohomology-lemma-cd-bound-dim-local}.
이제 세 수
$\alpha = \text{depth}((\mathcal{F}^\wedge_y)_\mathfrak p) \geq 0$,
$\beta = \dim(\mathcal{O}_{X, y}^\wedge/\mathfrak p) \geq 1$, 그리고
$\gamma = \delta^Y_Z(y) \geq 1$.
를 고찰하자. 정의 \ref{algebraization-definition-s-d-inequalities}가 요구하는 것은
\begin{enumerate}
\item $\beta > 1$이면 $\alpha + \gamma \geq 2$ 또는
$\alpha + \beta + \gamma > 3$이고,
\item $\beta = 1$이면 $\alpha + \gamma > 2$인 것이다.
\end{enumerate}
이는 $\alpha + \beta + \gamma > 3$과 동치임이 자명하다.
\end{proof}

\noindent
처음 읽을 때 건너뛰기를 권하는 이 절의 나머지 부분에서는 $A$가 $I$-진
완비일 때, $X$로 연장되고 엄격한
$(1, 1 + \text{cd}(A, I))$-부등식들을 만족하는
$\textit{Coh}(U, I\mathcal{O}_U)$의 $(\mathcal{F}_n)$들의 범주가 연접
$\mathcal{O}_U$-가군 범주의 어떤 충만한 부분범주와 동치임을 보일 것이다.

\begin{lemma}
\label{algebraization-lemma-sanity}
상황 \ref{algebraization-situation-algebraize}에서 $\mathcal{F}$를 연접
$\mathcal{O}_U$-가군이라 하고 $d \geq 1$이라 하자. 다음을 가정한다.
\begin{enumerate}
\item $A$는 $I$-진 완비이고 쌍대화 복합체를 가지며
$\text{cd}(A, I) \leq d$이다.
\item $\mathcal{F}$의 완비 $\mathcal{F}^\wedge$은 엄격한
$(1, 1 + d)$-부등식들을 만족한다.
\end{enumerate}
$x \in X$를 점이라 하고 $W = \overline{\{x\}}$로 놓자. $W \cap Y$가
$Z$에 포함되는 기약 성분과 포함되지 않는 기약 성분을 모두 가지면
$\text{depth}(\mathcal{F}_x) \geq 1$이다.
\end{lemma}

\begin{proof}
$W \cap Y = W_1 \cup \ldots \cup W_n$을 기약 성분 분해라 하자. 가정에
따라 번호를 다시 붙이면 어떤 $0 < m < n$에 대해
$W_1, \ldots, W_m  \subset Z$이고
$W_{m + 1}, \ldots, W_n \not \subset Z$이 되게 할 수 있다. 그러면
$$
W \cap Y \setminus
\left((W_1 \cup \ldots \cup W_m) \cap (W_{m + 1} \cup \ldots \cup W_n)\right)
$$
는 비연결이다. 보조정리 \ref{algebraization-lemma-connected}에 의해
$1 \leq i \leq m < j \leq n$ 및 $z \in W_i \cap W_j$를 택하여
$\dim(\mathcal{O}_{W, z}) \leq d + 1$이 되게 할 수 있다.
$y \in W_j$, $y \not \in Z$인 바로 다음 특수화 $y \leadsto z$를 택하자.
$y$의 존재는 성질의 보조정리
\ref{properties-lemma-complement-closed-point-Jacobson}에서 따른다.
$\delta^Y_Z(y) = 1$이고 $\dim(\mathcal{O}_{W, y}) \leq d$임을 주목하라.
$\mathfrak p \subset \mathcal{O}_{X, y}$를 $x$에 대응하는 소아이디얼이라
하자. $\mathfrak p' \subset \mathcal{O}_{X, y}^\wedge$을
$\mathfrak p\mathcal{O}_{X, y}^\wedge$ 위의 극소 소아이디얼이라 하자. 그러면
$$
\text{depth}(\mathcal{F}_x) =
\text{depth}((\mathcal{F}^\wedge_y)_{\mathfrak p'})
\quad\text{이고}\quad
\dim(\mathcal{O}_{W, y}) = \dim(\mathcal{O}_{X, y}^\wedge/\mathfrak p')
$$
이다. 대수학의 보조정리 \ref{algebra-lemma-apply-grothendieck-module}과
국소 코호몰로지의 보조정리
\ref{local-cohomology-lemma-change-completion}을 보라. 이제 주어진
부등식들에서 결론을 읽어낸다.
\end{proof}

\begin{lemma}
\label{algebraization-lemma-recover}
상황 \ref{algebraization-situation-algebraize}에서 $\mathcal{F}$를 연접
$\mathcal{O}_U$-가군이라 하고 $d \geq 1$이라 하자. 다음을 가정한다.
\begin{enumerate}
\item $A$는 $I$-진 완비이고 쌍대화 복합체를 가지며
$\text{cd}(A, I) \leq d$이다.
\item $\mathcal{F}$의 완비 $\mathcal{F}^\wedge$은 엄격한
$(1, 1+ d)$-부등식들을 만족한다.
\item $\overline{\{x\}} \cap Y \subset Z$인 $x \in U$에 대해
$\text{depth}(\mathcal{F}_x) \geq 2$이다.
\end{enumerate}
그러면 $H^0(U, \mathcal{F}) \to \lim H^0(U, \mathcal{F}/I^n\mathcal{F})$은
동형이다.
\end{lemma}

\begin{proof}
보조정리 \ref{algebraization-lemma-application-H0}를 적용할 수 있음을 보여 증명하겠다.
$x \not \in Y$인 $x \in \text{Ass}(\mathcal{F})$를 취하고
$W = \overline{\{x\}}$로 놓자. 조건 (3)에 의해
$W \cap Y \not \subset Z$이다. 보조정리 \ref{algebraization-lemma-sanity}에 의해
$W \cap Y$의 어떤 기약 성분도 $Z$에 포함되지 않는다. 따라서
$z \in W \cap Z$이면 $y \in W \cap Y$, $y \not \in Z$인 바로 다음 특수화
$y \leadsto z$가 있다. $y$의 존재에는 성질의 보조정리
\ref{properties-lemma-complement-closed-point-Jacobson}를 사용한다. 그러면
$\delta^Y_Z(y) = 1$이고 가정으로부터
$\dim(\mathcal{O}_{W, y}) > d$이다. 따라서
$\dim(\mathcal{O}_{W, z}) > 1 + d$이고 증명이 끝난다.
\end{proof}

\begin{lemma}
\label{algebraization-lemma-fully-faithful-inequalities}
상황 \ref{algebraization-situation-algebraize}에서 $\mathcal{F}$를 연접
$\mathcal{O}_U$-가군이라 하고 $d \geq 1$이라 하자. 다음을 가정한다.
\begin{enumerate}
\item $A$는 $I$-진 완비이고 쌍대화 복합체를 가지며
$\text{cd}(A, I) \leq d$이다.
\item $\mathcal{F}$의 완비 $\mathcal{F}^\wedge$은 엄격한
$(1, 1 + d)$-부등식들을 만족한다.
\item $\overline{\{x\}} \cap Y \subset Z$인 $x \in U$에 대해
$\text{depth}(\mathcal{F}_x) \geq 2$이다.
\end{enumerate}
그러면 사상
$$
\Hom_U(\mathcal{G}, \mathcal{F})
\longrightarrow
\Hom_{\textit{Coh}(U, I\mathcal{O}_U)}(\mathcal{G}^\wedge, \mathcal{F}^\wedge)
$$
은 모든 연접 $\mathcal{O}_U$-가군 $\mathcal{G}$에 대해 전단사이다.
\end{lemma}

\begin{proof}
$\mathcal{H} = \SheafHom_{\mathcal{O}_U}(\mathcal{G}, \mathcal{F})$로 놓자.
스킴의 코호몰로지의 보조정리
\ref{coherent-lemma-hom-into-depth} 또는 대수학 더 보기의 보조정리
\ref{more-algebra-lemma-hom-into-depth}를 사용하면 $\mathcal{H}$의 완비가
엄격한 $(1, 1 + d)$-부등식들을 만족하고,
$\overline{\{x\}} \cap Y \subset Z$인 $x \in U$에 대해
$\text{depth}(\mathcal{H}_x) \geq 2$임을 알 수 있다. 세부사항은 생략한다.
따라서 보조정리 \ref{algebraization-lemma-recover}에 의해
$$
\Hom_U(\mathcal{G}, \mathcal{F}) =
H^0(U, \mathcal{H}) =
\lim H^0(U, \mathcal{H}/\mathcal{I}^n\mathcal{H}) =
\Mor_{\textit{Coh}(U, I\mathcal{O}_U)}
(\mathcal{G}^\wedge, \mathcal{F}^\wedge)
$$
이다. 마지막 등식은 스킴의 코호몰로지의 보조정리
\ref{coherent-lemma-completion-internal-hom}를 보라.
\end{proof}

\begin{lemma}
\label{algebraization-lemma-construct-unique}
상황 \ref{algebraization-situation-algebraize}에서 $(\mathcal{F}_n)$을
$\textit{Coh}(U, I\mathcal{O}_U)$의 대상이라 하고 $d \geq 1$이라 하자.
다음을 가정한다.
\begin{enumerate}
\item $A$는 $I$-진 완비이고 쌍대화 복합체를 가지며
$\text{cd}(A, I) \leq d$이다.
\item $(\mathcal{F}_n)$은 연접 $\mathcal{O}_U$-가군의 완비이다.
\item $(\mathcal{F}_n)$은 엄격한 $(1, 1 + d)$-부등식들을 만족한다.
\end{enumerate}
그러면 완비가 $(\mathcal{F}_n)$이고, $\overline{\{x\}} \cap Y \subset Z$인
$x \in U$에 대해 $\text{depth}(\mathcal{F}_x) \geq 2$인 연접
$\mathcal{O}_U$-가군 $\mathcal{F}$가 유일하게 존재한다.
\end{lemma}

\begin{proof}
완비가 $(\mathcal{F}_n)$인 연접 $\mathcal{O}_U$-가군 $\mathcal{F}$를 택하자.
$T = \{x \in U \mid \overline{\{x\}} \cap Y \subset Z\}$로 놓자.
$\mathcal{F}$와 $T$에 국소 코호몰로지의 보조정리
\ref{local-cohomology-lemma-make-S2-along-T-simple}를 적용하여
$\mathcal{F}$를 구성하겠다. 그러면 유일성은 보조정리
\ref{algebraization-lemma-fully-faithful-inequalities}의 사상 성질에서 따른다.

\medskip\noindent
$T$는 $U$ 안의 특수화에 대해 안정적이므로 다음만 확인하면 된다.
$x \in T$, $x' \not \in T$인 $U$의 점들의 바로 다음 특수화를
$x' \leadsto x$라 하면 $\text{depth}(\mathcal{F}_{x'}) \geq 1$이다.
$W = \overline{\{x\}}$ 및 $W' = \overline{\{x'\}}$로 놓자.
$x' \not \in T$이므로 $W' \cap Y$는 $Z$에 포함되지 않는다.
$W' \cap Y$가 $Z$에 포함되는 기약 성분을 가지면 보조정리
\ref{algebraization-lemma-sanity}에 의해 끝난다. 그렇지 않으면 $W \cap Y$의 기약 성분
$W_1$과 $W_1 \subset W'_1$인 $W' \cap Y$의 기약 성분 $W'_1$을 택한다.
$z \in W_1$을 일반점이라 하자. $y \not \in Z$인 바로 다음 특수화
$y \leadsto z$, $y \in W'_1$을 택하자. $y$의 존재는
$W'_1 \not \subset Z$(위를 보라)와 성질의 보조정리
\ref{properties-lemma-complement-closed-point-Jacobson}에서 따른다. 이제
$z \in Z$, $x \leadsto z$, $x' \leadsto y \leadsto z$,
$y \in Y \setminus Z$, $\delta^Y_Z(y) = 1$이다. 국소 코호몰로지의 보조정리
\ref{local-cohomology-lemma-cd-bound-dim-local}와 $z$가 $W \cap Y$의
일반점이라는 사실에 의해 $\dim(\mathcal{O}_{W, z}) \leq d$이다.
$x' \leadsto x$가 바로 다음 특수화이므로
$\dim(\mathcal{O}_{W', z}) \leq d + 1$이다. $y \not = z$이므로
$\dim(\mathcal{O}_{W', y}) \leq d$를 얻는다.
$\text{depth}(\mathcal{F}_{x'}) = 0$이면 가정 (3)과 모순이다.
$\mathcal{O}_{X, y}$에서 그 완비로 옮기는 세부사항은 생략한다.
이로써 증명이 끝난다.
\end{proof}








\section{연접 형식 가군의 대수화 IV}
\label{algebraization-section-algebraization-modules-local}

\noindent
이 절에서는 국소적인 경우에 보조정리
\ref{algebraization-lemma-algebraization-principal-variant}의 더 강한 두 버전, 즉 보조정리
\ref{algebraization-lemma-algebraization-principal}과
\ref{algebraization-lemma-algebraization-principal-bis}를 증명한다. 이 보조정리들은 더
일반적인 명제 \ref{algebraization-proposition-cd-1}에 의해 대체되지만 증명은 훨씬 쉽다.

\begin{lemma}
\label{algebraization-lemma-algebraization-principal}
상황 \ref{algebraization-situation-algebraize}에서 $(\mathcal{F}_n)$을
$\textit{Coh}(U, I\mathcal{O}_U)$의 대상이라 하자. 다음을 가정한다.
\begin{enumerate}
\item $A$는 국소환이고 $\mathfrak a = \mathfrak m$은 극대 아이디얼이다.
\item $A$는 쌍대화 복합체를 가진다.
\item 영인자가 아닌 $f \in \mathfrak m$에 대해 $I = (f)$는 주 아이디얼이다.
\item $\mathcal{F}_n$은 유한 국소 자유
$\mathcal{O}_U/f^n\mathcal{O}_U$-가군이다.
\item $\mathfrak p \in V(f) \setminus \{\mathfrak m\}$이면
$\text{depth}((A/f)_\mathfrak p) + \dim(A/\mathfrak p) > 1$이다.
\item $\mathfrak p \not \in V(f)$이고
$V(\mathfrak p) \cap V(f) \not = \{\mathfrak m\}$이면
$\text{depth}(A_\mathfrak p) + \dim(A/\mathfrak p) > 3$.
\end{enumerate}
그러면 $(\mathcal{F}_n)$은 $X$로 표준적으로 연장된다. 특히 $A$가
완비이면 $(\mathcal{F}_n)$은 연접 $\mathcal{O}_U$-가군의 완비이다.
\end{lemma}

\begin{proof}
보조정리 \ref{algebraization-lemma-when-done}의 가정 (a), (b), (c)를 확인하여 증명하겠다.

\medskip\noindent
$\mathcal{F}_n$은 $\mathcal{O}_U/f^n\mathcal{O}_U$ 위에서 국소 자유이므로
모든 $n$에 대해 짧은 완전열
$0 \to \mathcal{F}_n \to \mathcal{F}_{n + 1} \to \mathcal{F}_1 \to 0$을
얻는다. 따라서 코호몰로지의 보조정리
\ref{cohomology-lemma-topology-I-adic-f}에 의해 조건 (b)가 성립한다.

\medskip\noindent
$n$에 관한 귀납법과 짧은 완전열
$0 \to A/f^n \to A/f^{n + 1} \to A/f \to 0$을 사용하면
$A/f^nA$의 동반 소아이디얼들은 $A/fA$의 것들과 일치한다. 가정 (3)에
의해 $\mathcal{F}_n$의 동반점들은 $\mathfrak m$과 다른 $A/f^nA$의 동반
소아이디얼들에 대응하므로, (5)와 국소 코호몰로지의 명제
\ref{local-cohomology-proposition-kollar}에 의해
$M_n = H^0(U, \mathcal{F}_n)$은 유한 $A$-가군이다. 따라서 가정 (c)가
성립한다.

\medskip\noindent
증명을 끝내려면 어떤 $n > 1$에 대해
$$
H^1(U, \mathcal{F}_n) \longrightarrow H^1(U, \mathcal{F}_1)
$$
의 상이 $A$-가군으로서 유한 길이임을 보이면 충분하다. 그러면 코호몰로지의
보조정리 \ref{cohomology-lemma-ML-better}에 의해 가정 (a)가 따르기 때문이다.
보조정리 \ref{algebraization-lemma-ML-local}에 의해 충분히 큰 $n$에 대해서는 이 상이
$n$에 무관하다. $\omega_A^\bullet$을 $A$의 정규화된 쌍대화 복합체라 하자.
국소 쌍대성 정리와 Matlis 쌍대성(쌍대화 복합체의 보조정리
\ref{dualizing-lemma-special-case-local-duality}와 명제
\ref{dualizing-proposition-matlis})에 의해 주장은 다음 사상의 상이
$$
\text{Ext}^{-2}_A(M_1, \omega_A^\bullet) \to
\text{Ext}^{-2}_A(M_n, \omega_A^\bullet)
$$
$n \gg 1$에 대해 유한 길이를 갖는다는 것과 동치이다. 문제의 가군들은
$V(f)$ 위에 지지되는 유한 $A$-가군이다. 따라서 $f$를 포함하고
$\mathfrak m$과 다른 소아이디얼 $\mathfrak q$에서 국소화한 뒤 이 사상이
영임을 보이면 충분하다. $\omega_{A_\mathfrak q}^\bullet$을
$A_\mathfrak q$ 위의 정규화된 쌍대화 복합체라 하자. 쌍대화 복합체의
보조정리 \ref{dualizing-lemma-dimension-function}에 의해
$\omega_{A_\mathfrak q}^\bullet =
(\omega_A^\bullet)_\mathfrak q[\dim(A/\mathfrak q)]$임을 상기하자.
(4)에 주어진 $\mathcal{F}_n$의 국소 구조를 사용하면 충분히 큰 $n$에 대해
$$
\text{Ext}^{-2 + \dim(A/\mathfrak q)}_{A_\mathfrak q}(
A_\mathfrak q/f, \omega_{A_\mathfrak q}^\bullet)
\to
\text{Ext}^{-2 + \dim(A/\mathfrak q)}_{A_\mathfrak q}(
A_\mathfrak q/f^n, \omega_{A_\mathfrak q}^\bullet)
$$
가 영임을 보이는 것으로 충분하다. $\dim(A/\mathfrak q) > 3$이면 이는
국소 코호몰로지의 보조정리
\ref{local-cohomology-lemma-sitting-in-degrees}에서 즉시 따른다. 다른 경우에는
긴 완전열
$$
\ldots
\xrightarrow{f^n}
H^{-1}(\omega_{A_\mathfrak q}^\bullet)
\to
\text{Ext}^{-1}_{A_\mathfrak q}(
A_\mathfrak q/f^n, \omega_{A_\mathfrak q}^\bullet) \to
H^0(\omega_{A_\mathfrak q}^\bullet)
\xrightarrow{f^n}
H^0(\omega_{A_\mathfrak q}^\bullet)
\to
\text{Ext}^0_{A_\mathfrak q}(
A_\mathfrak q/f^n, \omega_{A_\mathfrak q}^\bullet) \to 0
$$
을 사용한다. $\dim(A/\mathfrak q) = 2$이면
$H^0(\omega_{A_\mathfrak q}^\bullet) = 0$
이다. $f$가 영인자가 아니므로 $\text{depth}(A_\mathfrak q) \geq 1$이기
때문이다. 따라서 긴 완전열로부터 필요한 조건은
$$
f^{n - 1} :
H^{-1}(\omega_{A_\mathfrak q}^\bullet)/f \to
H^{-1}(\omega_{A_\mathfrak q}^\bullet)/f^n
$$
가 영인 것이다. 이제 $H^{-1}(\omega^\bullet_\mathfrak q)$은 다음을
만족하는 소아이디얼 $\mathfrak p \subset A_\mathfrak q$들 위에 지지되는
유한 가군이다.
$\text{depth}(A_\mathfrak p) + \dim((A/\mathfrak p)_\mathfrak q) \leq 1$.
$\dim((A/\mathfrak p)_\mathfrak q) = \dim(A/\mathfrak p) - 2$이므로 조건
(6)에 의해 이 소아이디얼들은 $V(f)$에 포함된다. 따라서 충분히 큰 $n$에
대해 원하는 소멸을 얻는다. 마지막으로 $\dim(A/\mathfrak q) = 1$이면
조건 (5)와 $f$가 영인자가 아니라는 사실로부터 $A_\mathfrak q$의 깊이가
적어도 $2$임을 안다. 따라서
$H^0(\omega_{A_\mathfrak q}^\bullet) =
H^{-1}(\omega_{A_\mathfrak q}^\bullet) = 0$
이고, 긴 완전열에 의해 주장은
$$
f^{n - 1} :
H^{-2}(\omega_{A_\mathfrak q}^\bullet)/f \to
H^{-2}(\omega_{A_\mathfrak q}^\bullet)/f^n
$$
의 소멸과 동치이다. 이제 $H^{-2}(\omega^\bullet_\mathfrak q)$은
$\text{depth}(A_\mathfrak p) + \dim((A/\mathfrak p)_\mathfrak q)
\leq 2$인 소아이디얼 $\mathfrak p \subset A_\mathfrak q$들 위에 지지되는
유한 가군이다. 조건 (6)에 의해 이 소아이디얼들은 모두 $V(f)$에 포함된다.
따라서 충분히 큰 $n$에 대해 원하는 소멸을 얻는다.
\end{proof}

\begin{remark}
\label{algebraization-remark-interesting-case}
$(A, \mathfrak m)$을 차원 $\geq 4$인 완비 Noether 정규 국소 정역이라 하고
$f \in \mathfrak m$은 영이 아니라고 하자. 그러면 보조정리
\ref{algebraization-lemma-algebraization-principal}의 가정 (1), (2), (3), (5), (6)이
성립한다. 따라서 $U \cap V(f)$를 따른 $U$의 형식적 완비 위의 벡터다발은
대수화할 수 있다. 보조정리 \ref{algebraization-lemma-algebraization-principal-bis}에서 이를
더 일반적인 연접 형식 가군으로 일반화할 것이다. 비고
\ref{algebraization-remark-interesting-case-bis}도 비교하라.
\end{remark}

\begin{lemma}
\label{algebraization-lemma-helper-algebraize}
상황 \ref{algebraization-situation-algebraize}에서 $(M_n)$을 보조정리
\ref{algebraization-lemma-system-of-modules}에서와 같은 $A$-가군의 역계라 하고,
$(\mathcal{F}_n)$을 이에 대응하는 $\textit{Coh}(U, I\mathcal{O}_U)$의
대상이라 하자. $d \geq \text{cd}(A, I)$와 $s \geq 0$을 정수라 하자.
위의 표기 아래 다음을 가정한다.
\begin{enumerate}
\item $A$는 극대 아이디얼 $\mathfrak m = \mathfrak a$을 갖는 국소환이다.
\item $A$는 쌍대화 복합체를 가진다.
\item $(\mathcal{F}_n)$은 $(s, s + d)$-부등식들을 만족한다
(정의 \ref{algebraization-definition-s-d-inequalities}).
\end{enumerate}
$E$를 $A$의 잉여체의 단사 포락이라 하자. 그러면 $i \leq s$에 대해,
$I$의 어떤 거듭제곱에 의해 소멸되는 유한 $A$-가군 $N$과
$n \gg 0$에 대해 서로 양립하는 사상
$$
H^i_\mathfrak m(M_n) \to \Hom_A(N, E)
$$
들이 존재한다. 그 여핵들은 유한 길이 $A$-가군이고, 핵들 $K_n$은 역계를
이루며 $n'' \gg n' \gg 0$에 대해 $\Im(K_{n''} \to K_{n'})$는 유한
길이를 갖는다.
\end{lemma}

\begin{proof}
$\omega_A^\bullet$을 정규화된 쌍대화 복합체라 하자. 그러면
$\delta^Y_Z = \delta$는 이 쌍대화 복합체에 대응하는 차원 함수이다.
$\Ext^{-i}_A(M_n, \omega_A^\bullet)$은 $I^n$에 의해 소멸되는 유한
$A$-가군임을 주목하라. $0 \leq i \leq s$를 고정하자. 아래에서
$n_1 > n_0 > 0$을 찾아, 다음과 같이 놓으면
$$
N = \Im(\Ext^{-i}_A(M_{n_0}, \omega_A^\bullet) \to
\Ext^{-i}_A(M_{n_1}, \omega_A^\bullet))
$$
사상
$$
N \to \Ext^{-i}_A(M_n, \omega_A^\bullet),\quad n \geq n_1
$$
들의 핵은 유한 길이 $A$-가군이고, 여핵들 $Q_n$은
$n'' \gg n' \gg n_1$에 대해 $\Im(Q_{n'} \to Q_{n''})$가 유한 길이를
갖는 계를 이루게 할 것이다. 이는 계
$\{\Ext^{-i}_A(M_n, \omega_A^\bullet)\}_{n \geq 1}$이 유한 $A$-가군의
범주를 유한 길이 $A$-가군들의 Serre 부분범주로 나눈 몫에서 본질적으로
상수라는 것과 동치이다. 국소 쌍대성 정리(쌍대화 복합체의 보조정리
\ref{dualizing-lemma-special-case-local-duality})와 Matlis 쌍대성(쌍대화
복합체의 명제 \ref{dualizing-proposition-matlis})에 의해
$$
H^i_\mathfrak m(M_n) \to \Hom_A(N, E),\quad n \geq n_1
$$
와 같은 사상들을 얻는다.

\medskip\noindent
$f \in \mathfrak m$을 택하자. $B = A_f^\wedge$를 국소화 $A_f$의
$I$-진 완비라 하자.
$\omega_{A_f}^\bullet = \omega_A^\bullet \otimes_A A_f$
및 $\omega_B^\bullet = \omega_A^\bullet \otimes_A B$가 쌍대화 복합체임을
상기하자(쌍대화 복합체의 보조정리
\ref{dualizing-lemma-dualizing-localize}와
\ref{dualizing-lemma-completion-henselization-dualizing}). $M$을 유한
$B$-가군 $\lim M_{n, f}$라 하자(스킴의 코호몰로지의 보조정리
\ref{coherent-lemma-inverse-systems-affine}). 그러면
$$
\Ext^{-i}_A(M_n, \omega_A^\bullet)_f =
\Ext^{-i}_{A_f}(M_{n, f}, \omega_{A_f}^\bullet) =
\Ext^{-i}_B(M/I^n M, \omega_B^\bullet)
$$
$\mathfrak m$은 유한 개의 $f \in \mathfrak m$으로 생성되므로 각 $f$에 대해 계
$$
\{\Ext^{-i}_B(M/I^n M, \omega_B^\bullet)\}_{n \geq 1}
$$
가 본질적으로 상수임을 보이면 충분하다. 일부 세부사항은 생략한다.

\medskip\noindent
$\mathfrak q \subset IB$를 소아이디얼이라 하자. 그러면 $\mathfrak q$는
점 $y \in U \cap Y$에 대응한다.
$\delta(\mathfrak q) = \dim(\overline{\{y\}})$
는 $\omega_B^\bullet$에 대응하는 차원 함수의 값이기도 하다(세부사항은
생략한다. $\omega_B^\bullet$은 $\omega_A^\bullet$을 $B$로 텐서하여
얻는다는 사실을 사용하라). 가정 (3)과 보조정리 \ref{algebraization-lemma-elementary}에
의해 보조정리 \ref{algebraization-lemma-algebraize-local-cohomology-bis}를
$B_\mathfrak q, IB_\mathfrak q, \mathfrak qB_\mathfrak q, M_\mathfrak q$
에 적용할 수 있는데, 여기서 $s$를 $s - \delta(y)$로 바꾼다. 그러면
$$
H^{i - \delta(\mathfrak q)}_{\mathfrak qB_\mathfrak q}(M_\mathfrak q) =
\lim H^{i - \delta(\mathfrak q)}_{\mathfrak qB_\mathfrak q}(
(M/I^nM)_\mathfrak q)
$$
이고 이 가군은 $I$의 어떤 거듭제곱에 의해 소멸된다. 보조정리
\ref{algebraization-lemma-terrific}에 의해 역계
$H^{i - \delta(\mathfrak q)}_{\mathfrak qB_\mathfrak q}((M/I^nM)_\mathfrak q)$
는 값
$H^{i - \delta(\mathfrak q)}_{\mathfrak qB_\mathfrak q}(M_\mathfrak q)$.
을 갖는 본질적으로 상수인 계이다.
$(\omega_B^\bullet)_\mathfrak q[-\delta(\mathfrak q)]$가 $B_\mathfrak q$
위의 정규화된 쌍대화 복합체이므로 국소 쌍대성 정리에 의해 계
$$
\Ext^{-i}_B(M/I^n M, \omega_B^\bullet)_\mathfrak q
$$
는 값 $\Ext^{-i}_B(M, \omega_B^\bullet)_\mathfrak q$을 갖는 본질적으로
상수인 계이다.

\medskip\noindent
증명을 끝내기 위해 보조정리 \ref{algebraization-lemma-bootstrap}의 증명에서와 같이
대역화한다. 여기서는 계의 값을 이미 알고 있으므로 논증이 더 쉽다.
즉 변화하는 $n$에 대해 사상
$$
\alpha_n :
\Ext^{-i}_B(M/I^n M, \omega_B^\bullet)
\longrightarrow
\Ext^{-i}_B(M, \omega_B^\bullet)
$$
를 고찰하자. 위에서 보았듯 모든 $\mathfrak q$에 대해 어떤 $n$을 택하여
$\mathfrak q$에서 국소화한 뒤 $\alpha_n$이 전사가 되게 할 수 있다.
$B$는 Noether 환이고 $\Ext^{-i}_B(M, \omega_B^\bullet)$은 유한 가군이므로
어떤 $n$에 대해 $\alpha_n$이 전사가 되게 할 수 있다. $\alpha_n$이
전사인 임의의 $n$과 소아이디얼 $\mathfrak q \in V(IB)$가 주어지면,
어떤 $n' > n$을 찾아 적어도 $\mathfrak q$에서 국소화한 뒤
$\Ker(\alpha_n)$이 $\Ext^{-i}(M/I^{n'}M, \omega_B^\bullet)$ 안에서 영으로
사상되게 할 수 있다. $\Ker(\alpha_n)$은 유한 $A$-가군이고 단면의 지지는
준콤팩트이므로, 어떤 $n'$에 대해 $\Ker(\alpha_n)$이
$\Ext^{-i}(M/I^{n'}M, \omega_B^\bullet)$ 안에서 영으로 사상되게 할 수 있다.
이로써 $\Ext^{-i}(M/I^n M, \omega_B^\bullet)$은 값
$\Ext^{-i}(M, \omega_B^\bullet)$을 갖는 본질적으로 상수인 계이다.
증명이 끝났다.
\end{proof}

\noindent
다음은 보조정리 \ref{algebraization-lemma-algebraization-principal}의 더 일반적인 버전이다.

\begin{lemma}
\label{algebraization-lemma-algebraization-principal-bis}
상황 \ref{algebraization-situation-algebraize}에서 $(\mathcal{F}_n)$을
$\textit{Coh}(U, I\mathcal{O}_U)$의 대상이라 하자. 다음을 가정한다.
\begin{enumerate}
\item $A$는 국소환이고 $\mathfrak a = \mathfrak m$은 극대 아이디얼이다.
\item $A$는 쌍대화 복합체를 가진다.
\item $I = (f)$는 주 아이디얼이다.
\item $(\mathcal{F}_n)$은 $(2, 3)$-부등식들을 만족한다.
\end{enumerate}
그러면 $(\mathcal{F}_n)$은 $X$로 연장된다. 특히 $A$가 $I$-진 완비이면
$(\mathcal{F}_n)$은 연접 $\mathcal{O}_U$-가군의 완비이다.
\end{lemma}

\begin{proof}
$\textit{Coh}(U, I\mathcal{O}_U)$가 아벨 범주임을 상기하자. 스킴의
코호몰로지의 보조정리 \ref{coherent-lemma-inverse-systems-abelian}을 보라.
$U$의 아핀 열린집합 위에서 대상 $(\mathcal{F}_n)$은 Noether 환 위의 유한
가군에 대응한다(스킴의 코호몰로지의 보조정리
\ref{coherent-lemma-inverse-systems-affine}). 따라서 충분히 큰 $N$에 대해
사상 $f^N : (\mathcal{F}_n) \to (\mathcal{F}_n)$의 핵은 안정화한다.
보조정리 \ref{algebraization-lemma-map-kernel-cokernel-on-closed}와
\ref{algebraization-lemma-essential-image-completion}에 의해 보조정리를 증명할 때
$(\mathcal{F}_n)$을 이러한 사상의 상으로 바꾸어도 된다. 따라서 $f$가
$(\mathcal{F}_n)$ 위에서 단사라고 가정할 수 있다. 이렇게 바꾼 뒤에는
$U$ 위의 역계 $(\mathcal{F}_n)$에 대해 보조정리
\ref{algebraization-lemma-equivalent-f-good}의 동치 조건들이 성립한다. 증명의 나머지에서는
이를 더 언급하지 않고 사용한다.

\medskip\noindent
보조정리 \ref{algebraization-lemma-when-done}의 가정 (a), (b), (c)를 확인하겠다.
가정 (b)는 코호몰로지의 보조정리
\ref{cohomology-lemma-topology-I-adic-f}에 의해 성립한다.

\medskip\noindent
보조정리 \ref{algebraization-lemma-system-of-modules}에서와 같은 가군의 역계
$\{M_n\}$을 택하자. 필요하다면 $M_n$을 $M_n/H^0_\mathfrak m(M_n)$으로
바꾸어 $H^0_\mathfrak m(M_n) = 0$이라고 가정할 수 있다. 그러면 모든
$n$에 대해 짧은 완전열
$$
0 \to M_n \to H^0(U, \mathcal{F}_n) \to H^1_\mathfrak m(M_n) \to 0
$$
을 얻는다. $E$를 $A$의 잉여체의 단사 포락이라 하자. 보조정리
\ref{algebraization-lemma-helper-algebraize}와 현재 가정 (4)에 의해, 정수 $m \geq 0$,
어떤 $c \geq 0$에 대해 $f^c$에 의해 소멸되는 유한 $A$-가군 $N_1$, $N_2$,
그리고 서로 양립하는 사상계
$$
H^i_\mathfrak m(M_n) \to \Hom_A(N_i, E), \quad i = 1, 2
$$
를 $n \geq m$에 대해 택하되 보조정리에서 말한 성질들을 만족하게 할 수 있다.

\medskip\noindent
$M = \lim H^0(U, \mathcal{F}_n)$은 극한 위상이 $f$-진 위상인
$A$-가군이다. 따라서 $n$이 주어지면 어떤 $N \gg n$에 대해 가군
$M/f^nM$은 $H^0(U, \mathcal{F}_N)$의 부분몫이다. 위에서 얻은 정보로부터
$f^cM/f^nM$은 유한 $A$-가군이다. $f$가 $M$ 위에서 영인자가 아니므로
$M/f^{n - c}M$은 유한 $A$-가군이다. 이로써 보조정리
\ref{algebraization-lemma-when-done}의 가정 (c)가 성립한다.

\medskip\noindent
다음으로 가군
$$
Ob = \lim H^1(U, \mathcal{F}_n) = \lim H^2_\mathfrak m(M_n)
$$
을 연구한다. $n \geq m$에 대해 $K_n$을 사상
$H^2_\mathfrak m(M_n) \to \Hom_A(N_2, E)$의 핵이라 하고
$K = \lim K_n$으로 놓자. 완전열
$$
0 \to K \to Ob \to \Hom_A(N_2, E)
$$
을 얻는다. 위에서 보았듯 $Ob = \lim H^2_\mathfrak m(M_n)$의 극한 위상은
$f$-진 위상이다. $N_2$가 $f^c$에 의해 소멸되므로 $K = \lim K_n$의 극한
위상도 마찬가지이다. 따라서 $n \gg 1$에 대해 $K/fK$는 $K_n$의
부분몫이다. 그러나 보조정리 \ref{algebraization-lemma-helper-algebraize}의 결론에 의해
$\{K_n\}$은 유한 길이 $A$-가군의 역계와 전사영 동형이므로, $K/fK$는
어떤 유한 길이 $A$-가군의 부분몫이다. 따라서 $K$는 유한 $A$-가군이다.
대수학의 보조정리 \ref{algebra-lemma-finite-over-complete-ring}을 보라.
(실제로 $\dim(\text{Supp}(K)) = 1$까지 알 수 있지만 이는 필요하지 않다.)

\medskip\noindent
$n \geq 1$이 주어졌을 때 경계사상
$$
\delta_n :
H^0(U, \mathcal{F}_n)
\longrightarrow
\lim_N H^1(U, f^n\mathcal{F}_N) \xrightarrow{f^{-n}} Ob
$$
을 고찰하자(둘째 사상은 동형이다). 이는 짧은 완전열
$$
0 \to f^n\mathcal{F}_N \to \mathcal{F}_N \to \mathcal{F}_n \to 0
$$
에서 온다. 각 $n$에 대해
$$
P_n = \Im(H^0(U, \mathcal{F}_{n + m}) \to H^0(U, \mathcal{F}_n))
$$
로 놓자. 여기서 $m$은 위에서와 같다. $\{P_n\}$은 역계이고, 전역 단면
위에서 사상 $f : \mathcal{F}_n \to \mathcal{F}_{n + 1}$은 $P_n$을
$P_{n + 1}$로 보낸다. $p \in P_n$이면 $\delta_n(p) \in K \subset Ob$이다.
실제로 $\delta_n(p)$은
$H^1(U, f^n\mathcal{F}_{n + m}) = H^2_\mathfrak m(M_m)$
안에서 영으로 사상되고, $m$의 선택에 의해 $\delta_n$과
$Ob \to \Hom_A(N_2, E)$의 합성은 $H^2_\mathfrak m(M_m)$을 통해 분해된다.
따라서
$$
\bigoplus\nolimits_{n \geq 0} \Im(P_n \to Ob)
$$
은 유한 등급 $A[T]$-가군이다. 여기서 $T$는 $f$를 곱하는 것으로 작용한다.
실제로 이는 $K[T]$의 등급 부분가군이고 $K$는 $A$ 위에서 유한이다.
코호몰로지의 보조정리 \ref{cohomology-lemma-ML-general}의 증명에서와 같이
논증하면\footnote{표시된 가군의 동차 생성원을 $\delta_{n_j}(p_j)$ 꼴로
택하자. $k = \max(n_j)$라 하면 $n \geq k$ 및 임의의 $p \in P_n$에 대해
어떤 $a_j \in A$들이 존재하여
$p - \sum a_j f^{n - n_j} p_j$가 $\delta_n$의 핵에 속하고 따라서 모든
$n' \geq n$에 대해 $P_{n'}$의 상에 속한다. 그러므로 모든 $n' \geq n$에
대해 $\Im(P_n \to P_{n - k}) = \Im(P_{n'} \to P_{n - k})$이다.}
역계 $\{P_n\}$은 Mittag-Leffler 조건을 만족한다. $\{P_n\}$은
$\{H^0(U, \mathcal{F}_n)\}$과 전사영 동형이므로
$\{H^0(U, \mathcal{F}_n)\}$도 Mittag-Leffler 조건을 만족한다. 따라서
보조정리 \ref{algebraization-lemma-when-done}의 가정 (a)가 성립하고 증명이 끝난다.
\end{proof}

\noindent
보조정리 \ref{algebraization-lemma-algebraization-principal-bis}의 조건은 다음과 같이
풀어 쓸 수 있다.

\begin{lemma}
\label{algebraization-lemma-unwinding-conditions}
상황 \ref{algebraization-situation-algebraize}에서 $(\mathcal{F}_n)$을
$\textit{Coh}(U, I\mathcal{O}_U)$의 대상이라 하자. 다음을 가정한다.
\begin{enumerate}
\item $A$는 극대 아이디얼 $\mathfrak a = \mathfrak m$을 갖는 국소환이다.
\item $\text{cd}(A, I) = 1$.
\end{enumerate}
그러면 $(\mathcal{F}_n)$이 $(2, 3)$-부등식들을 만족할 필요충분조건은
$\dim(\{y\}) = 1$인 모든 $y \in U \cap Y$와
$\mathfrak p \not \in V(I\mathcal{O}_{X, y}^\wedge)$인 모든 소아이디얼
$\mathfrak p \subset \mathcal{O}_{X, y}^\wedge$에 대해
$$
\text{depth}((\mathcal{F}_y^\wedge)_\mathfrak p) +
\dim(\mathcal{O}_{X, y}^\wedge/\mathfrak p) > 2
$$
\end{lemma}

\begin{proof}
보조정리 \ref{algebraization-lemma-explain-2-3-cd-1}을 더 언급하지 않고 사용하겠다.
특히 조건이 필요함을 알 수 있다. 역으로 조건이 참이라고 하자.
보조정리 \ref{algebraization-lemma-discussion}에 의해
$\delta^Y_Z(y) = \dim(\overline{\{y\}})$임을 주목하라. 이 함수를
$\delta$로 쓰자. $y \in U \cap Y$라 하자. $\delta(y) > 2$이면 보조정리
\ref{algebraization-lemma-explain-2-3-cd-1}의 부등식이 성립한다. 마지막으로
$\delta(y) = 2$라고 하자. 다음을 보여야 한다.
$$
\text{depth}((\mathcal{F}_y^\wedge)_\mathfrak p) +
\dim(\mathcal{O}_{X, y}^\wedge/\mathfrak p) > 1
$$
$\delta(y') = 1$인 특수화 $y \leadsto y'$를 택하자. 그러면 환 준동형
$\mathcal{O}_{X, y'}^\wedge \to \mathcal{O}_{X, y}^\wedge$이 존재하고,
이는 표적을 $\dim(\mathcal{O}_{X, y'}^\wedge/\mathfrak q) = 1$인
소아이디얼 $\mathfrak q$에서 $\mathcal{O}_{X, y'}^\wedge$을 국소화한 것의
완비와 동일시한다. 또한
$$
\mathcal{F}_y^\wedge =
\mathcal{F}_{y'}^\wedge
\otimes_{\mathcal{O}_{X, y'}^\wedge}
\mathcal{O}_{X, y}^\wedge
$$
$\mathfrak p' \subset \mathcal{O}_{X, y'}^\wedge$을 $\mathfrak p$의 상이라
하자. 국소 코호몰로지의 보조정리
\ref{local-cohomology-lemma-change-completion}에 의해
\begin{align*}
\text{depth}((\mathcal{F}_y^\wedge)_\mathfrak p) +
\dim(\mathcal{O}_{X, y}^\wedge/\mathfrak p)
& =
\text{depth}((\mathcal{F}_{y'}^\wedge)_{\mathfrak p'}) +
\dim((\mathcal{O}_{X, y}^\wedge/\mathfrak p)_{\mathfrak p'}) \\
& =
\text{depth}((\mathcal{F}_{y'}^\wedge)_{\mathfrak p'}) +
\dim(\mathcal{O}_{X, y}^\wedge/\mathfrak p') - 1
\end{align*}
이다. 마지막 등식은 특수화가 바로 다음 특수화이기 때문이다. 따라서
$y', \mathfrak p'$에 대해 가정한 부등식으로 보조정리가 증명된다.
\end{proof}

\begin{lemma}
\label{algebraization-lemma-unwinding-conditions-bis}
상황 \ref{algebraization-situation-algebraize}에서 $(\mathcal{F}_n)$을
$\textit{Coh}(U, I\mathcal{O}_U)$의 대상이라 하자. 다음을 가정한다.
\begin{enumerate}
\item $A$는 극대 아이디얼 $\mathfrak a = \mathfrak m$을 갖는 국소환이다.
\item $A$는 쌍대화 복합체를 가진다.
\item $\text{cd}(A, I) = 1$,
\item $y \in U \cap Y$에 대해 가군 $\mathcal{F}_y^\wedge$은
$V(I\mathcal{O}_{X, y}^\wedge)$ 밖에서 유한 국소 자유이다. 예를 들어
$\mathcal{F}_n$이 유한 국소 자유 $\mathcal{O}_U/I^n\mathcal{O}_U$-가군인
경우가 이에 해당한다.
\item 다음 중 하나가 참이다.
\begin{enumerate}
\item $A_f$는 $(S_2)$이고 $Y$에 포함되지 않는 $X$의 모든 기약 성분은
차원 $\geq 4$이다.
\item $\mathfrak p \not \in V(f)$이고
$V(\mathfrak p) \cap V(f) \not = \{\mathfrak m\}$이면
$\text{depth}(A_\mathfrak p) + \dim(A/\mathfrak p) > 3$.
\end{enumerate}
\end{enumerate}
그러면 $(\mathcal{F}_n)$은 $(2, 3)$-부등식들을 만족한다.
\end{lemma}

\begin{proof}
보조정리 \ref{algebraization-lemma-unwinding-conditions}의 판정법을 사용하겠다.
$\dim(\overline{\{y\}} = 1$인 $y \in U \cap Y$와
$\mathfrak p \not \in V(I\mathcal{O}_{X, y}^\wedge)$인 소아이디얼
$\mathfrak p$, 즉 $\mathfrak p \subset \mathcal{O}_{X, y}^\wedge$을 취하자.
조건 (4)에 의해
$\text{depth}((\mathcal{F}_y^\wedge)_\mathfrak p) =
\text{depth}((\mathcal{O}_{X, y}^\wedge)_\mathfrak p)$.
따라서 다음을 증명해야 한다.
$$
\text{depth}((\mathcal{O}_{X, y}^\wedge)_\mathfrak p) +
\dim(\mathcal{O}_{X, y}^\wedge/\mathfrak p) > 2
$$
$\mathfrak p_0 \subset A$를 $\mathfrak p$의 상이라 하고,
$\mathfrak q \subset A$를 $y$에 대응하는 소아이디얼이라 하자.
국소 코호몰로지의 보조정리
\ref{local-cohomology-lemma-change-completion}
다음을 얻는다.
\begin{align*}
\text{depth}((\mathcal{O}_{X, y}^\wedge)_\mathfrak p) +
\dim(\mathcal{O}_{X, y}^\wedge/\mathfrak p)
& =
\text{depth}(A_{\mathfrak p_0}) + \dim((A/\mathfrak p_0)_\mathfrak q) \\
& =
\text{depth}(A_{\mathfrak p_0}) + \dim(A/\mathfrak p_0) - 1
\end{align*}
(5)(a)가 성립하면 이는
$$
\geq \min(2, \dim(A_{\mathfrak p_0})) + \dim(A/\mathfrak p_0) - 1
$$
이상이다. 어느 경우든 $\dim(A/\mathfrak p_0) \geq 2$임을 주목하라.
따라서 최솟값이 $2$이면 끝난다. 그렇지 않으면
$$
\dim(A_{\mathfrak p_0}) + \dim(A/\mathfrak p_0) - 1 \geq 4 - 1
$$
이다. $\mathfrak p_0$을 지나는 $\Spec(A)$의 모든 성분이 차원
$\geq 4$이기 때문이다. (5)(b)가 성립하면 즉시 끝난다.
\end{proof}

\begin{remark}
\label{algebraization-remark-interesting-case-bis}
$(A, \mathfrak m)$을 쌍대화 복합체를 가지며 $f \in \mathfrak m$에 관해
완비인 Noether 국소환이라 하자. $U$를 $A$의 천공 스펙트럼이라 하고,
$(\mathcal{F}_n)$을 $\textit{Coh}(U, f\mathcal{O}_U)$의 대상이라 하자.
$Y = V(f) \subset X = \Spec(A)$로 놓자. $U$에서 닫힌
$y \in U \cap V(f)$, 즉 $\dim(\overline{\{y\}}) = 1$인 점에 대해
$\mathcal{O}_{X, y}^\wedge$-가군 $\mathcal{F}_y^\wedge$이 다음 두 조건을
만족한다고 가정하자.
\begin{enumerate}
\item $\mathcal{F}_y^\wedge[1/f]$는 $\mathcal{O}_{X, y}^\wedge[1/f]$-가군으로서
$(S_2)$이다.
\item $\mathfrak p \in \text{Ass}(\mathcal{F}_y^\wedge[1/f])$에 대해
$\dim(\mathcal{O}_{X, y}^\wedge/\mathfrak p) \geq 3$이다.
\end{enumerate}
그러면 $(\mathcal{F}_n)$은 $U$ 위의 연접 가군의 완비이다. 이는 보조정리
\ref{algebraization-lemma-algebraization-principal-bis}와
\ref{algebraization-lemma-unwinding-conditions}에서 따른다.
\end{remark}
































\section{연접 형식 가군의 개선}
\label{algebraization-section-improving-formal-modules}

\noindent
$X$를 Noether 스킴, $Y \subset X$를 준연접 아이디얼층
$\mathcal{I} \subset \mathcal{O}_X$로 주어지는 닫힌 부분스킴이라 하자.
$(\mathcal{F}_n)$을 $\textit{Coh}(X, \mathcal{I})$의 대상이라 하자.
이 절에서는 국소 코호몰로지의 절 \ref{local-cohomology-section-improve}에서
연접 가군에 대해 구성한 것과 비슷한 사상
$(\mathcal{F}_n) \to (\mathcal{F}'_n)$을 구성한다. 점 $y \in Y$에 대해
$$
\mathcal{O}_{X, y}^\wedge = \lim \mathcal{O}_{X, y}/\mathcal{I}^n_y, \quad
\mathcal{I}_y^\wedge = \lim \mathcal{I}_y/\mathcal{I}^n_y
\quad\text{이고}\quad
\mathfrak m_y^\wedge = \lim \mathfrak m_y/\mathcal{I}_y^n
$$
로 놓는다. 그러면 $\mathcal{O}_{X, y}^\wedge$은 극대 아이디얼
$\mathfrak m_y^\wedge$을 가지며
$\mathcal{I}_y^\wedge = \mathcal{I}_y\mathcal{O}_{X, y}^\wedge$에 관해
완비인 Noether 국소환이다. 또한
$$
\mathcal{F}_y^\wedge = \lim \mathcal{F}_{n, y}
$$
로 놓는다. 그러면 $\mathcal{F}_y^\wedge$은 $\mathcal{O}_{X, y}^\wedge$
위의 유한 가군이고 모든 $n$에 대해
$\mathcal{F}_y^\wedge/(\mathcal{I}_y^\wedge)^n\mathcal{F}_y^\wedge =
\mathcal{F}_{n, y}$이다. 대수학의 보조정리
\ref{algebra-lemma-limit-complete}과
\ref{algebra-lemma-finite-over-complete-ring}.

\begin{lemma}
\label{algebraization-lemma-divide-torsion-formal-coherent-module}
위 상황에서 $X$가 국소적으로 쌍대화 복합체를 가진다고 가정하자.
$T \subset Y$를 특수화에 대해 안정적인 부분집합이라 하자.
$y \in T$와 다음을 만족하는 극대가 아닌 소아이디얼
$\mathfrak p \subset \mathcal{O}_{X, y}^\wedge$에 대해
$V(\mathfrak p) \cap V(\mathcal{I}^\wedge_y) = \{\mathfrak m_y^\wedge\}$
다음이 성립한다고 가정하자.
$$
\text{depth}_{(\mathcal{O}_{X, y})_\mathfrak p}
((\mathcal{F}^\wedge_y)_\mathfrak p) > 0
$$
그러면 다음 성질들을 갖는 연접 $\mathcal{O}_X$-가군의 역계의 표준 사상
$(\mathcal{F}_n) \to (\mathcal{F}_n')$이 존재한다.
\begin{enumerate}
\item $y \in T$에 대해 $\text{depth}(\mathcal{F}'_{n, y}) \geq 1$이다.
\item $(\mathcal{F}'_n)$은 전사영계로서 $\textit{Coh}(X, \mathcal{I})$의
대상 $(\mathcal{G}_n)$과 동형이다.
\item 유도된 $\textit{Coh}(X, \mathcal{I})$의 사상
$(\mathcal{F}_n) \to (\mathcal{G}_n)$은 전사이고 그 핵은
$\mathcal{I}$의 어떤 거듭제곱에 의해 소멸된다.
\end{enumerate}
\end{lemma}

\begin{proof}
모든 $n$에 대해 $\mathcal{F}_n \to \mathcal{F}'_n$을 국소 코호몰로지의
보조정리 \ref{local-cohomology-lemma-get-depth-1-along-Z}에서 구성한
전사라 하자. 이는 $\mathcal{F}_n$을 $T$ 위에 지지되는 단면들의 부분층으로
나눈 몫이므로, 표준 사상 $\mathcal{F}'_{n + 1} \to \mathcal{F}'_n$을 얻어
연접 $\mathcal{O}_X$-가군의 역계의 사상
$(\mathcal{F}_n) \to (\mathcal{F}_n')$을 이룬다. 성질 (1)은 구성상 성립한다.

\medskip\noindent
성질 (2), (3)을 증명할 때 $X = \Spec(A_0)$가 아핀이고 $A_0$가 쌍대화
복합체를 가진다고 가정할 수 있다. $I_0 \subset A_0$를 $Y$에 대응하는
아이디얼이라 하자. $A, I$를 $A_0, I_0$의 $I_0$-진 완비라 하자. 나중에
사용하기 위해 $A$가 쌍대화 복합체를 가짐을 주목한다(쌍대화 복합체의
보조정리 \ref{dualizing-lemma-ubiquity-dualizing}). $M$을 $(\mathcal{F}_n)$에
대응하는 유한 $A$-가군이라 하자. 스킴의 코호몰로지의 보조정리
\ref{coherent-lemma-inverse-systems-affine}을 보라. 그러면
$\mathcal{F}_n$은 $M_n = M/I^nM$에 대응한다. $\mathcal{F}'_n$은 몫
$M'_n = M_n / H^0_T(M_n)$에 대응함을 상기하자. 국소 코호몰로지의 보조정리
\ref{local-cohomology-lemma-get-depth-1-along-Z}와 그 증명을 보라.

\medskip\noindent
$s = 0$, $d = \text{cd}(A, I)$로 놓자. $A, I, T, M, s, d$가 상황
\ref{algebraization-situation-bootstrap}의 가정 (1), (3), (4), (6)을 만족한다고 주장한다.
(1), (3)은 위에서 즉시 따르고, $s = 0$이므로 (4)는 공허한 조건이며,
(6)은 보조정리의 명제에서 둔 가정이다.

\medskip\noindent
정리 \ref{algebraization-theorem-final-bootstrap}에 의해 $\{H^0_T(M_n)\}$은 값
$H^0_T(M)$을 갖는 본질적으로 상수인 계이다. 따라서 짧은 완전열
$0 \to H^0_T(M_n) \to M_n \to M'_n \to 0$들의 극한은 짧은 완전열
$$
0 \to H^0_T(M) \to M \to \lim M'_n \to 0
$$
이다. $M' = \lim M'_n = M/H^0_T(M)$으로 놓자. 이는 유한 $A$-가군이다.
$(\mathcal{G}_n)$을 $M'$에 대응하는 $\textit{Coh}(X, \mathcal{I})$의
대상이라 하자. 증명을 끝내려면 표준 사상
$\{M'/I^nM'\} \to \{M'_n\}$이 전사영 동형임을 보여야 한다. 이는
$\{H^0_T(M) + I^nM\} \to \{\Ker(M \to M'_n)\}$이 전사영 동형이라는 것과
동치이다. $I^nM$으로 나누면
$\{H^0_T(M)/H^0_T(M) \cap I^nM\} \to \{H^0_T(M_n)\}$이 전사영 동형임을
보이면 충분하다. $H^0_T(M)$은 $I$의 어떤 거듭제곱에 의해 소멸되므로
Artin--Rees에 의해 모든 $n \gg 0$에 대해
$H^0_T(M) \cap I^nM = 0$이다. 위에서 $\{H^0_T(M_n)\}$이 값
$H^0_T(M)$을 갖는 본질적으로 상수인 계임을 보았으므로 원하는 전사영
동형을 얻는다.
\end{proof}

\begin{lemma}
\label{algebraization-lemma-improvement-formal-coherent-module-better}
위 상황에서 $X$가 국소적으로 쌍대화 복합체를 가진다고 가정하자.
$T' \subset T \subset Y$를 특수화에 대해 안정적인 부분집합들이라 하고,
$d \geq 0$을 정수라 하자. 다음을 가정한다.
\begin{enumerate}
\item[(a)] 아핀 국소적으로 $X = \Spec(A_0)$ 및 $Y = V(I_0)$이고
$\text{cd}(A_0, I_0) \leq d$이다.
\item[(b)] $y \in T$와 다음을 만족하는 극대가 아닌 소아이디얼
$\mathfrak p \subset \mathcal{O}_{X, y}^\wedge$로서
$V(\mathfrak p) \cap V(\mathcal{I}_y^\wedge) = \{\mathfrak m_y^\wedge\}$
에 대해
$$
\text{depth}_{(\mathcal{O}_{X, y})_\mathfrak p}
((\mathcal{F}^\wedge_y)_\mathfrak p) > 0
$$
이다.
\item[(c)] $y \in T'$와 다음을 만족하는 소아이디얼
$\mathfrak p \subset \mathcal{O}_{X, y}^\wedge$로서
$\mathfrak p \not \in V(\mathcal{I}_y^\wedge)$
및 $V(\mathfrak p) \cap V(\mathcal{I}_y^\wedge) \not =
\{\mathfrak m_y^\wedge\}$에 대해
$$
\text{depth}_{(\mathcal{O}_{X, y})_\mathfrak p}
((\mathcal{F}^\wedge_y)_\mathfrak p) \geq 1
\quad\text{또는}\quad
\text{depth}_{(\mathcal{O}_{X, y})_\mathfrak p}
((\mathcal{F}^\wedge_y)_\mathfrak p) +
\dim(\mathcal{O}_{X, y}^\wedge/\mathfrak p) > 1 + d
$$
\item[(d)] $y \in T'$와 다음을 만족하는 극대가 아닌 소아이디얼
$\mathfrak p \subset \mathcal{O}_{X, y}^\wedge$로서
$V(\mathfrak p) \cap V(\mathcal{I}_y^\wedge) = \{\mathfrak m_y^\wedge\}$
에 대해
$$
\text{depth}_{(\mathcal{O}_{X, y})_\mathfrak p}
((\mathcal{F}^\wedge_y)_\mathfrak p) > 1
$$
\item[(e)] $y \leadsto y'$가 바로 다음 특수화이고 $y' \in T'$이면
$y \in T$이다.
\end{enumerate}
그러면 다음 성질들을 갖는 연접 $\mathcal{O}_X$-가군의 역계의 표준 사상
$(\mathcal{F}_n) \to (\mathcal{F}_n'')$이 존재한다.
\begin{enumerate}
\item $y \in T$에 대해 $\text{depth}(\mathcal{F}''_{n, y}) \geq 1$이다.
\item $y' \in T'$에 대해 $\text{depth}(\mathcal{F}''_{n, y'}) \geq 2$이다.
\item $(\mathcal{F}''_n)$은 전사영계로서 $\textit{Coh}(X, \mathcal{I})$의
대상 $(\mathcal{H}_n)$과 동형이다.
\item 유도된 $\textit{Coh}(X, \mathcal{I})$의 사상
$(\mathcal{F}_n) \to (\mathcal{H}_n)$의 핵과 여핵은 $\mathcal{I}$의 어떤
거듭제곱에 의해 소멸된다.
\end{enumerate}
\end{lemma}

\begin{proof}
보조정리 \ref{algebraization-lemma-divide-torsion-formal-coherent-module}와 그 증명에서와
같이, 모든 $n$에 대해 $\mathcal{F}_n \to \mathcal{F}'_n$을 $T$를 사용하여
국소 코호몰로지의 보조정리
\ref{local-cohomology-lemma-get-depth-1-along-Z}에서 구성한 전사라 하자.
다음으로 $\mathcal{F}'_n \to \mathcal{F}''_n$을 국소 코호몰로지의 보조정리
\ref{local-cohomology-lemma-make-S2-along-T}와 그 증명에서 구성한 단사라
하자. 구성으로부터 표준 사상 $\mathcal{F}''_{n + 1} \to \mathcal{F}''_n$을
얻으며, 따라서 연접 $\mathcal{O}_X$-가군의 역계의 사상
$$
(\mathcal{F}_n) \longrightarrow (\mathcal{F}_n') \longrightarrow
(\mathcal{F}''_n)
$$
을 얻는다. 성질 (1), (2)는 구성상 성립한다.

\medskip\noindent
성질 (3), (4)를 증명할 때 $X = \Spec(A_0)$가 아핀이고 $A_0$가 쌍대화
복합체를 가진다고 가정할 수 있다. $I_0 \subset A_0$를 $Y$에 대응하는
아이디얼이라 하자. $A, I$를 $A_0, I_0$의 $I_0$-진 완비라 하자. 나중에
사용하기 위해 $A$가 쌍대화 복합체를 가짐을 주목한다(쌍대화 복합체의
보조정리 \ref{dualizing-lemma-ubiquity-dualizing}). $M$을 $(\mathcal{F}_n)$에
대응하는 유한 $A$-가군이라 하자. 스킴의 코호몰로지의 보조정리
\ref{coherent-lemma-inverse-systems-affine}을 보라. 그러면
$\mathcal{F}_n$은 $M_n = M/I^nM$에 대응한다. $\mathcal{F}'_n$은 몫
$M'_n = M_n / H^0_T(M_n)$에 대응함을 상기하자. 또한
$M' = \lim M'_n$은 $M$을 $H^0_T(M)$으로 나눈 몫이고, $H^0_T(M)$은
$I$-거듭제곱 꼬임이며, $\{M'_n\}$과 $\{M'/I^nM'\}$은 전사영계로서
동형임을 상기하자. 마지막으로 $\mathcal{F}''_n$은 확대
$$
0 \to M'_n \to M''_n \to H^1_{T'}(M'_n) \to 0
$$
에 대응한다. 국소 코호몰로지의 보조정리
\ref{local-cohomology-lemma-make-S2-along-T}의 증명을 보라.

\medskip\noindent
$s = 1$로 놓자. $A, I, T', M', s, d$가 상황
\ref{algebraization-situation-bootstrap}의 가정 (1), (3), (4), (6)을 만족한다고 주장한다.
(1), (3)은 즉시 따르고, (c)는 (4)를 함의하며, (d)에서 (6)이 따른다.
(c) $\Rightarrow$ (4)와 (d) $\Rightarrow$ (6)의 증명은 생략한다.

\medskip\noindent
정리 \ref{algebraization-theorem-final-bootstrap}에 의해 $\{H^1_{T'}(M'/I^nM')\}$은 값
$H^1_{T'}(M')$을 갖는 본질적으로 상수인 계이다. 따라서
$\{H^1_{T'}(M'_n)\}$도 값 $H^1_{T'}(M')$을 갖는 본질적으로 상수인 계이다.
그러므로 위에 표시한 짧은 완전열들의 극한은 짧은 완전열
$$
0 \to M' \to \lim M''_n \to H^1_{T'}(M') \to 0
$$
이다. $M'' = \lim M''_n$으로 놓자. 국소 코호몰로지의 명제
\ref{local-cohomology-proposition-finiteness}에 의해 $H^1_{T'}(M')$과
따라서 $M''$은 유한 $A$-가군이다(조건을 확인할 때 $M'$이 $T - T'$의
소아이디얼들에서 깊이 적어도 $1$임을 사용하라). $(\mathcal{H}_n)$을
유한 $A$-가군 $M''$에 대응하는 $\textit{Coh}(X, \mathcal{I})$의 대상이라
하자. 증명을 끝내려면 표준 사상
$\{M''/I^nM''\} \to \{M''_n\}$이 전사영 동형임을 보여야 한다.
$\{M'/I^nM'\}$이 $\{M'_n\}$과 전사영 동형임을 이미 알고 있으므로,
이는 $\{H^1_{T'}(M')/I^nH^1_{T'}(M')\} \to \{H^1_{T'}(M'_n)\}$이 전사영
동형이라는 것과 동치임을 독자가 확인할 수 있다(생략한다). 충분히 큰
$n$에 대해 $H^1_{T'}(M')/I^nH^1_{T'}(M') = H^1_{T'}(M')$이고, 위에서
$\{H^1_{T'}(M'_n)\}$이 값 $H^1_{T'}(M')$을 갖는 본질적으로 상수인 계임을
보았으므로 결론이 따른다.
\end{proof}

\begin{lemma}
\label{algebraization-lemma-improvement-application}
상황 \ref{algebraization-situation-algebraize}에서 $A$가 쌍대화 복합체를 가진다고
가정하자. $d \geq \text{cd}(A, I)$라 하자. $(\mathcal{F}_n)$을
$\textit{Coh}(U, I\mathcal{O}_U)$의 대상이라 하자. $(\mathcal{F}_n)$이
$(2, 2 + d)$-부등식들을 만족한다고 가정하자. 정의
\ref{algebraization-definition-s-d-inequalities}를 보라. 그러면 다음 성질들을 갖는 연접
$\mathcal{O}_U$-가군의 역계의 표준 사상
$(\mathcal{F}_n) \to (\mathcal{F}_n'')$이 존재한다.
\begin{enumerate}
\item 모든 $y \in U \cap Y$에 대해
$\text{depth}(\mathcal{F}''_{n, y}) + \delta^Y_Z(y) \geq 3$이다.
\item $(\mathcal{F}''_n)$은 전사영계로서
$\textit{Coh}(U, I\mathcal{O}_U)$의 대상 $(\mathcal{H}_n)$과 동형이다.
\item 유도된 $\textit{Coh}(U, I\mathcal{O}_U)$의 사상
$(\mathcal{F}_n) \to (\mathcal{H}_n)$의 핵과 여핵은 $I$의 어떤
거듭제곱에 의해 소멸된다.
\item 모든 $n$에 대해 가군 $H^0(U, \mathcal{F}''_n)$과
$H^1(U, \mathcal{F}''_n)$은 유한 $A$-가군이다.
\end{enumerate}
\end{lemma}

\begin{proof}
존재와 성질 (1), (2), (3)은 $U$, $U \cap Y$,
$T = \{y \in U \cap Y : \delta^Y_Z(y) \leq 2\}$,
$T' = \{y \in U \cap Y : \delta^Y_Z(y) \leq 1\}$ 및 $(\mathcal{F}_n)$에
보조정리 \ref{algebraization-lemma-improvement-formal-coherent-module-better}를 적용하면
따른다. 가군 $H^0(U, \mathcal{F}''_n)$과 $H^1(U, \mathcal{F}''_n)$의
유한성은 국소 코호몰로지의 보조정리
\ref{local-cohomology-lemma-finiteness-Rjstar}와 보조정리
\ref{algebraization-lemma-discussion}에서 증명한 함수 $\delta^Y_Z(-)$의 초등적 성질들에서
따른다.
\end{proof}











\section{연접 형식 가군의 대수화 V}
\label{algebraization-section-algebraization-modules-conclusion}

\noindent
이 절에서는 연접 형식 가군의 대수화에 관한 가장 일반적인 결과들을
증명한다. 먼저 아이디얼의 코호몰로지 차원이 $1$인 경우를 증명한 뒤,
이를 부풀림에 적용하여 더 일반적인 결과를 증명한다.

\begin{lemma}
\label{algebraization-lemma-cd-1-canonical}
상황 \ref{algebraization-situation-algebraize}에서 $(\mathcal{F}_n)$을
$\textit{Coh}(U, I\mathcal{O}_U)$의 대상이라 하자. 다음을 가정한다.
\begin{enumerate}
\item $A$는 쌍대화 복합체를 가지며 $\text{cd}(A, I) = 1$이다.
\item $(\mathcal{F}_n)$은 다음을 만족하는 연접 $\mathcal{O}_U$-가군의
역계 $(\mathcal{F}_n'')$과 전사영 동형이다.
$\text{depth}(\mathcal{F}''_{n, y}) + \delta^Y_Z(y) \geq 3$
이는 모든 $y \in U \cap Y$에 대해 성립한다.
\end{enumerate}
그러면 $(\mathcal{F}_n)$은 $X$로 표준적으로 연장된다. 정의
\ref{algebraization-definition-canonically-algebraizable}를 보라.
\end{lemma}

\begin{proof}
보조정리 \ref{algebraization-lemma-when-done}의 가정 (a), (b), (c)를 확인하겠다.
시작하기 전에 국소 코호몰로지의 보조정리
\ref{local-cohomology-lemma-finiteness-Rjstar}에 의해 모든 $n$에 대해
가군 $H^0(U, \mathcal{F}''_n)$과 $H^1(U, \mathcal{F}''_n)$이 유한
$A$-가군임을 지적해 둔다.

\medskip\noindent
각 $p \geq 0$에 대해 보조정리 \ref{algebraization-lemma-topology-I-adic}에 의해
$\lim H^p(U, \mathcal{F}_n)$의 극한 위상은 $I$-진 위상이다. 특히
가정 (b)가 성립한다.

\medskip\noindent
$M = \lim H^0(U, \mathcal{F}_n)$은 극한 위상이 $I$-진 위상인
$A$-가군이다. 따라서 $n$이 주어지면 어떤 $N \gg n$에 대해 $M/I^nM$은
$H^0(U, \mathcal{F}_N)$의 부분몫이다. 역계
$\{H^0(U, \mathcal{F}_N)\}$은 유한 $A$-가군의 역계
$\{H^0(U, \mathcal{F}''_N)\}$과 전사영 동형이므로 $M/I^nM$은 유한이다.
따라서 $M$도 유한이다. 대수학의 보조정리
\ref{algebra-lemma-finite-over-complete-ring}을 보라. 특히 가정 (c)가
성립한다.

\medskip\noindent
각 $n \geq 0$에 대해 $Ob_n = \lim_N H^1(U, I^n\mathcal{F}_N)$로 쓰자.
특히 $Ob = Ob_0 = \lim_N H^1(U, \mathcal{F}_N)$이다. 바로 앞 문단과
똑같이 논증하면 $Ob$가 유한 $A$-가군임을 얻는다. (실제로 $Ob/I Ob$가
$\mathfrak a$의 어떤 거듭제곱에 의해 소멸된다는 사실도 알지만, 이를
사용하기는 다소 어려워 보인다.)

\medskip\noindent
$\mathcal{F} = \lim \mathcal{F}_n$으로 놓고, $I$의 생성원
$f_1, \ldots, f_r$을 택하며, $c \geq 1$을 택하고, 보조정리
\ref{algebraization-lemma-cd-is-one-for-system}에서와 같이 $\Phi_\mathcal{F}$를 택한다.
보조정리 \ref{algebraization-lemma-properties-system}의 결과들을 더 언급하지 않고
사용하겠다. 특히 각 $n \geq 1$에 대해 사상
$$
\delta_n :
H^0(U, \mathcal{F}_n)
\longrightarrow
H^1(U, I^n\mathcal{F})
\longrightarrow
Ob_n
$$
들이 있다. 첫째 것은 짧은 완전열
$0 \to I^n\mathcal{F} \to \mathcal{F} \to \mathcal{F}_n \to 0$에서,
둘째 것은 $I^n\mathcal{F} = \lim I^n\mathcal{F}_N$에서 온다. 나중에
$s \in H^0(U, \mathcal{F}_n)$에 대해 $\delta_n(s) = 0$이면 각
$n' \geq n$마다 $s$로 사상되는 $s' \in H^0(U, \mathcal{F}_{n'})$을 찾을
수 있다는 사실을 사용하겠다. 다음 가환 그림이 있다.
$$
\xymatrix{
H^0(U, \mathcal{F}_{nc}) \ar[r] \ar[dd] &
H^1(U, I^{nc}\mathcal{F}) \ar[dd] \ar[rd]^{\Phi_\mathcal{F}} \\
& &
\bigoplus_{e_1 + \ldots + e_r = n}
H^1(U, \mathcal{F}) \cdot T_1^{e_1} \ldots T_r^{e_r} \ar[ld] \\
H^0(U, \mathcal{F}_n) \ar[r] &
H^1(U, I^n\mathcal{F})
}
$$
따라서 장애사상 $H^0(U, \mathcal{F}_n) \to Ob_n$은
$H^0(U, \mathcal{F}_{nc}) \to H^0(U, \mathcal{F}_n)$의 상을 부분가군
$$
Ob'_n =
\Im\left(
\bigoplus\nolimits_{e_1 + \ldots + e_r = n}
Ob \cdot T_1^{e_1} \ldots T_r^{e_r} \to Ob_n
\right)
$$
으로 보낸다. 여기서 직합항 $Ob \cdot T_1^{e_1} \ldots T_r^{e_r}$ 위에서는
곱셈 사상 $f_1^{e_1} \ldots f_r^{e_r} : \mathcal{F} \to I^n\mathcal{F}$을
$I$의 거듭제곱들로 나눈 환원에서 오는 코호몰로지 사상을 사용한다.
구성상
$$
\bigoplus\nolimits_{n \geq 0} Ob'_n
$$
은 Rees 대수 $\bigoplus_{n \geq 0} I^n$ 위의 유한 등급 가군이다.
각 $n$에 대해
$$
M_n = \{s \in H^0(U, \mathcal{F}_n) \mid \delta_n(s) \in Ob'_n\}
$$
로 놓자. $\{M_n\}$은 역계이고 전역 단면 위에서
$f_j : \mathcal{F}_n \to \mathcal{F}_{n + 1}$은 $M_n$을 $M_{n + 1}$로
보낸다. 코호몰로지의 보조정리 \ref{cohomology-lemma-ML-general}의
증명에서와 정확히 같은 논증으로 $\{M_n\}$이 Mittag-Leffler 조건을
만족함을 얻는다. 실제로 Rees 대수가 Noether이므로 등급 부분가군
$\bigoplus_{n \geq 0} \Im(M_n \to Ob'_n)$에 대해 $z_j \in M_{n_j}$인
$\delta_{n_j}(z_j)$ 꼴의 유한 개 동차 생성원을 택할 수 있다.
$k = \max(n_j)$라 하면, $n \geq k$ 및 임의의 $z \in M_n$에 대해
$a_j \in I^{n - n_j}$를 택하여 $z - \sum a_j z_j$가 $\delta_n$의 핵에
속하게 할 수 있고, 따라서 이는 모든 $n' \geq n$에 대해 $M_{n'}$의
상에 속한다($\delta_n$의 소멸은 모든 $n' \ge n$에 대해
$z - \sum a_j z_j$를 어떤 $z' \in H^0(U, \mathcal{F}_{n'c})$로 올릴 수
있음을 뜻하고, 위에서 증명한 바에 따라 $z'$의
$H^0(U, \mathcal{F}_{n'})$ 안의 상은 $M_{n'}$에 속한다). 따라서 모든
$n' \geq n$에 대해
$\Im(M_n \to M_{n - k}) = \Im(M_{n'} \to M_{n - k})$이다.

\medskip\noindent
$n$을 택하자. 방금 증명한 $\{M_n\}$의 Mittag-Leffler 성질에 의해 어떤
$n' \geq n$을 택하여 $M_{n'} \to M_n$의 상이 $M' \to M_n$의 상과 같게
할 수 있다. 위에서 보았듯 $M' \to M_n$의 상은
$H^0(U, \mathcal{F}_{n'c}) \to H^0(U, \mathcal{F}_n)$의 상을 포함한다.
따라서 $\{M_n\}$과 $\{H^0(U, \mathcal{F}_n)\}$은 전사영 동형이다.
그러므로 $\{H^0(U, \mathcal{F}_n)\}$은 Mittag-Leffler 조건을 만족하고,
마침내 가정 (a)가 성립한다. 증명이 끝났다.
\end{proof}

\begin{proposition}[코호몰로지 차원 1에서의 대수화]
\label{algebraization-proposition-cd-1}
\begin{reference}
이 결과의 국소적인 경우는 \cite[IV 따름정리 2.9]{MRaynaud-book}이다.
\end{reference}
상황 \ref{algebraization-situation-algebraize}에서 $(\mathcal{F}_n)$을
$\textit{Coh}(U, I\mathcal{O}_U)$의 대상이라 하자. 다음을 가정한다.
\begin{enumerate}
\item $A$는 쌍대화 복합체를 가지며 $\text{cd}(A, I) = 1$이다.
\item $(\mathcal{F}_n)$은 $(2, 3)$-부등식들을 만족한다. 정의
\ref{algebraization-definition-s-d-inequalities}를 보라.
\end{enumerate}
그러면 $(\mathcal{F}_n)$은 $X$로 연장된다. 특히 $A$가 $I$-진 완비이면
$(\mathcal{F}_n)$은 연접 $\mathcal{O}_U$-가군의 완비이다.
\end{proposition}

\begin{proof}
보조정리 \ref{algebraization-lemma-map-kernel-cokernel-on-closed}에 의해
$(\mathcal{F}_n)$을 보조정리 \ref{algebraization-lemma-improvement-application}에서 얻은
$\textit{Coh}(U, I\mathcal{O}_U)$의 대상 $(\mathcal{H}_n)$으로 바꾸어도
된다. 따라서 $(\mathcal{F}_n)$이 보조정리
\ref{algebraization-lemma-improvement-application}에 나온 성질들을 갖는 역계
$(\mathcal{F}_n'')$과 전사영 동형이라고 가정할 수 있다. 보조정리
\ref{algebraization-lemma-cd-1-canonical}에서 $(\mathcal{F}_n)$이 $X$로 표준적으로
연장됨을 증명하였다. 마지막 명제는 보조정리
\ref{algebraization-lemma-canonically-algebraizable}에서 따른다.
\end{proof}

\begin{lemma}
\label{algebraization-lemma-blowup}
상황 \ref{algebraization-situation-algebraize}에서 $(\mathcal{F}_n)$을
$\textit{Coh}(U, I\mathcal{O}_U)$의 대상이라 하자. 다음을 가정한다.
\begin{enumerate}
\item $A$는 쌍대화 복합체를 가진다.
\item $I$에 대한 부풀림 $b : X' \to X$의 모든 올은 차원
$\leq d - 1$이다.
\item 다음 중 하나가 참이다.
\begin{enumerate}
\item $(\mathcal{F}_n)$은 $(d + 1, d + 2)$-부등식들을 만족한다
(정의 \ref{algebraization-definition-s-d-inequalities}).
\item $y \in U \cap Y$와 다음을 만족하는 소아이디얼
$\mathfrak p \subset \mathcal{O}_{X, y}^\wedge$로서
$\mathfrak p \not \in V(I\mathcal{O}_{X, y}^\wedge)$
에 대해
$$
\text{depth}((\mathcal{F}^\wedge_y)_\mathfrak p) +
\dim(\mathcal{O}_{X, y}^\wedge/\mathfrak p) + \delta^Y_Z(y) > d + 2
$$
\end{enumerate}
\end{enumerate}
그러면 $(\mathcal{F}_n)$은 $X$로 연장된다.
\end{lemma}

\begin{proof}
$Y' \subset X'$를 예외 인자라 하고 $Z' \subset Y'$를 $Z \subset Y$의
역상이라 하자. 그러면 $U' = X' \setminus Z'$는 $U$의 역상이다.
$\delta^{Y'}_{Z'}$를 (\ref{algebraization-equation-delta-Z})에서와 같이 놓고
$$
T' = \{y' \in Y' \mid \delta^{Y'}_{Z'}(y') = 1\}
\subset
T = \{y' \in Y' \mid \delta^{Y'}_{Z'}(y') = 1\text{ 또는 }2\}
$$
로 놓자. 보조정리 \ref{algebraization-lemma-discussion}의 (1), (2)에 의해 이들은
$U' \cap Y' = Y' \setminus Z'$에서 특수화에 대해 안정적인 부분집합이다.
$\textit{Coh}(U', I\mathcal{O}_{U'})$의 대상
$(b|_{U'}^*\mathcal{F}_n)$을 고찰하자. 스킴의 코호몰로지의 보조정리
\ref{coherent-lemma-inverse-systems-pullback}을 보라.
$y' \in U' \cap Y'$에 대해 이 당김의 $y'$에서의 ``줄기''를
$$
\mathcal{F}_{y'}^\wedge = \lim (b|_{U'}^*\mathcal{F}_n)_{y'}
$$
로 표기하자. $d$를 $1$로 바꾸고 부분집합
$T' \subset T \subset U' \cap Y'$을 취했을 때, $U'$ 위의 대상
$(b|_{U'}^*\mathcal{F}_n)$에 대해 보조정리
\ref{algebraization-lemma-improvement-formal-coherent-module-better}의 조건
(a), (b), (c), (d), (e)가 성립한다고 주장한다. $Y'$가 유효 Cartier
인자이고 따라서 국소적으로 방정식 $1$개로 잘리므로 조건 (a)가 성립한다.
조건 (e)는 보조정리 \ref{algebraization-lemma-discussion}의 (7)에 의해 성립한다.
(b), (c), (d)를 증명하려면 약간의 준비가 필요하다.

\medskip\noindent
$y' \in U' \cap Y'$와 $V(I\mathcal{O}_{X', y'}^\wedge)$에 포함되지 않는
소아이디얼 $\mathfrak p' \subset \mathcal{O}_{X', y'}^\wedge$을 취하자.
$y = b(y') \in U \cap Y$로 표기하자. $y'$가 아핀 부풀림 대수
$A[\frac{I}{f}]$의 스펙트럼에 포함되도록 $f \in I$를 택하자. 인자의
보조정리 \ref{divisors-lemma-blowing-up-affine}를 보라. 임의의 $A$-대수
$B$에 대해 대응하는 아핀 부풀림 대수를 $B' = B[\frac{IB}{f}]$로 표기하자.
$I$-진 완비는 ${\ }^\wedge$으로 표기한다. $f$의 선택으로부터 환 준동형
$(\mathcal{O}_{X, y}^\wedge)' \to \mathcal{O}_{X', y'}^\wedge$을 얻는다.
$\mathfrak q' \subset (\mathcal{O}_{X, y}^\wedge)'$을
$\mathfrak m_{y'}^\wedge$의 역상이라 하면
$((\mathcal{O}_{X, y}^\wedge)'_{\mathfrak q'})^\wedge =
\mathcal{O}_{X', y'}^\wedge$.
이다. $\mathfrak p \subset \mathcal{O}_{X, y}^\wedge$을 대응하는
소아이디얼이라 하자. 이제 다음 가환 그림이 있다.
$$
\xymatrix{
\mathcal{O}_{X, y}^\wedge \ar[d] \ar[r] &
(\mathcal{O}_{X, y}^\wedge)' \ar[d]_\alpha \ar[r] &
(\mathcal{O}_{X, y}^\wedge)'_{\mathfrak q'} \ar[d] \ar[r]_\beta &
\mathcal{O}_{X', y'}^\wedge \ar[d] \\
\mathcal{O}_{X, y}^\wedge/\mathfrak p \ar[r] &
(\mathcal{O}_{X, y}^\wedge/\mathfrak p)' \ar[r] &
(\mathcal{O}_{X, y}^\wedge/\mathfrak p)'_{\mathfrak q'} \ar[r]^\gamma &
((\mathcal{O}_{X, y}^\wedge/\mathfrak p)'_{\mathfrak q'})^\wedge \ar[d] \\
 & & &
\mathcal{O}_{X', y'}^\wedge/\mathfrak p'
}
$$
여기서 수직 화살표들은 전사이다. 대수학 더 보기의 보조정리
\ref{more-algebra-lemma-completion-dimension}와 차원 공식(대수학의 보조정리
\ref{algebra-lemma-dimension-formula})에 의해
$$
\dim(((\mathcal{O}_{X, y}^\wedge/\mathfrak p)'_{\mathfrak q'})^\wedge) =
\dim((\mathcal{O}_{X, y}^\wedge/\mathfrak p)'_{\mathfrak q'}) =
\dim(\mathcal{O}_{X, y}^\wedge/\mathfrak p)
- \text{trdeg}(\kappa(y')/\kappa(y))
$$
당김, 줄기, 국소화, 완비의 정의를 따라가면
$$
(\mathcal{F}_y^\wedge)_{\mathfrak p}
\otimes_{(\mathcal{O}_{X, y}^\wedge)_\mathfrak p}
(\mathcal{O}_{X', y'}^\wedge)_{\mathfrak p'}
=
(\mathcal{F}_{y'}^\wedge)_{\mathfrak p'}
$$
이다. 세부사항은 생략한다. 그림의 환 준동형 $\beta$와 $\gamma$는
Gorenstein(따라서 Cohen--Macaulay) 올을 갖는 평탄 사상이다. 이들은
쌍대화 복합체를 갖는 환들의 완비이기 때문이다. 쌍대화 복합체의 보조정리
\ref{dualizing-lemma-formal-fibres-gorenstein}과
\ref{dualizing-lemma-dualizing-gorenstein-formal-fibres}
들과 대수학 더 보기의 절
\ref{more-algebra-section-properties-formal-fibres}의 논의를 보라.
$(\mathcal{O}_{X, y}^\wedge)_\mathfrak p =
(\mathcal{O}_{X, y}^\wedge)'_{\tilde{\mathfrak p}}$임을 주목하라. 여기서
$\tilde{\mathfrak p}$는 그림의 $\alpha$의 핵이다. 한편
$(\mathcal{O}_{X, y}^\wedge)'_{\tilde{\mathfrak p}}
\to (\mathcal{O}_{X', y'}^\wedge)_{\mathfrak p'}$
은 위의 사실에 의해 Cohen--Macaulay 올을 갖는 평탄 사상이다. 따라서
$(\mathcal{O}_{X, y}^\wedge)_\mathfrak p  \to
(\mathcal{O}_{X', y'}^\wedge)_{\mathfrak p'}$
도 Cohen--Macaulay 올을 갖는 평탄 사상이다. 대수학의 보조정리
\ref{algebra-lemma-apply-grothendieck-module}을 사용하면
$$
\text{depth}((\mathcal{F}_{y'}^\wedge)_{\mathfrak p'}) =
\text{depth}((\mathcal{F}_y^\wedge)_{\mathfrak p}) +
\dim(F_\mathfrak r)
$$
이다. 여기서 $F$는
$(\mathcal{O}_{X, y}^\wedge/\mathfrak p)'_{\mathfrak q'}$
의 일반 형식 올이고 $\mathfrak r$은 $\mathfrak p'$에 대응하는
소아이디얼이다. $(\mathcal{O}_{X, y}^\wedge/\mathfrak p)'_{\mathfrak q'}$는
보편적 사슬 국소 정역이므로, Ratliff 정리(대수학 더 보기의 명제
\ref{more-algebra-proposition-ratliff})에 의해 그 $I$-진 완비는 등차원이고
(보편적으로) 사슬적이다. 따라서
$$
\dim(((\mathcal{O}_{X, y}^\wedge/\mathfrak p)'_{\mathfrak q'})^\wedge) =
\dim(F_\mathfrak r) + \dim(\mathcal{O}_{X', y'}^\wedge/\mathfrak p')
$$
이다. 이를 보조정리 \ref{algebraization-lemma-change-distance-function}과 결합하면
\begin{equation}
\label{algebraization-equation-one}
\begin{aligned}
&
\text{depth}((\mathcal{F}_{y'}^\wedge)_{\mathfrak p'}) +
\delta^{Y'}_{Z'}(y') \\
& =
\text{depth}((\mathcal{F}_y^\wedge)_{\mathfrak p}) +
\dim(F_\mathfrak r) + \delta^{Y'}_{Z'}(y') \\
& \geq
\text{depth}((\mathcal{F}_y^\wedge)_{\mathfrak p}) + \delta^Y_Z(y) +
\dim(F_\mathfrak r) + \text{trdeg}(\kappa(y')/\kappa(y)) - (d - 1) \\
& =
\text{depth}((\mathcal{F}_y^\wedge)_{\mathfrak p}) + \delta^Y_Z(y) - (d - 1)
+ \dim(\mathcal{O}_{X, y}^\wedge/\mathfrak p) -
\dim(\mathcal{O}_{X', y'}^\wedge/\mathfrak p')
\end{aligned}
\end{equation}
다음을 유념하라.
$\dim(\mathcal{O}_{X, y}^\wedge/\mathfrak p) \geq
\dim(\mathcal{O}_{X', y'}^\wedge/\mathfrak p')$이다. 이를 다시 쓰면
\begin{equation}
\label{algebraization-equation-two}
\begin{aligned}
&
\text{depth}((\mathcal{F}_{y'}^\wedge)_{\mathfrak p'}) +
\dim(\mathcal{O}_{X', y'}^\wedge/\mathfrak p') +
\delta^{Y'}_{Z'}(y') \\
& \geq
\text{depth}((\mathcal{F}_y^\wedge)_{\mathfrak p}) +
\dim(\mathcal{O}_{X, y}^\wedge/\mathfrak p) +
\delta^Y_Z(y) - (d - 1)
\end{aligned}
\end{equation}
이 부등식으로 나머지 조건들을 확인할 수 있다.

\medskip\noindent
보조정리 \ref{algebraization-lemma-improvement-formal-coherent-module-better}의 조건
(b), (d)를 보자. 다음을 가정하자.
$V(\mathfrak p') \cap V(I\mathcal{O}_{X', y'}^\wedge) =
\{\mathfrak m_{y'}^\wedge\}$.
$Z'$가 유효 Cartier 인자이므로 이는
$\dim(\mathcal{O}_{X', y'}^\wedge/\mathfrak p') = 1$을 함의한다.
(b), (d)를 합친 조건은
$$
\text{depth}((\mathcal{F}_{y'}^\wedge)_{\mathfrak p'}) + \delta^{Y'}_{Z'}(y')
> 2
$$
와 동치이다. $(\mathcal{F}_n)$이 (3)(b)의 부등식을 만족하면
(\ref{algebraization-equation-two})를 적용하여 이것이 참임을 즉시 얻는다.
$(\mathcal{F}_n)$이 (3)(a), 즉 $(d + 1, d + 2)$-부등식들을 만족하면
어느 경우든
$$
\text{depth}((\mathcal{F}_y^\wedge)_{\mathfrak p}) + \delta^Y_Z(y)
\geq d + 1
\quad\text{또는}\quad
\text{depth}((\mathcal{F}_y^\wedge)_{\mathfrak p}) +
\dim(\mathcal{O}_{X, y}^\wedge/\mathfrak p) +
\delta^Y_Z(y) > d + 2
$$
이다. 위의 (\ref{algebraization-equation-one})과 (\ref{algebraization-equation-two})를 보면
$\dim(\mathcal{O}_{X, y}^\wedge/\mathfrak p) = 1$일 가능성을 제외하고
원하는 결론을 얻는다. 그러나
$\dim(\mathcal{O}_{X, y}^\wedge/\mathfrak p) = 1$이면
$V(\mathfrak p) \cap V(I\mathcal{O}_{X, y}^\wedge) = \{\mathfrak m_y^\wedge\}$이고
실제로
$$
\text{depth}((\mathcal{F}_y^\wedge)_{\mathfrak p}) + \delta^Y_Z(y) > d + 1
$$
이다. $(\mathcal{F}_n)$이 $(d + 1, d + 2)$-부등식들을 만족하기 때문이다.
따라서 다시 결론을 얻는다.

\medskip\noindent
보조정리 \ref{algebraization-lemma-improvement-formal-coherent-module-better}의 조건
(c)를 보자. 다음을 가정하자.
$V(\mathfrak p') \cap V(I\mathcal{O}_{X', y'}^\wedge) \not =
\{\mathfrak m_{y'}^\wedge\}$이다. 그러면 조건 (c)는
$$
\text{depth}((\mathcal{F}_{y'}^\wedge)_{\mathfrak p'}) + \delta^{Y'}_{Z'}(y')
\geq 2
\quad\text{또는}\quad
\text{depth}((\mathcal{F}_{y'}^\wedge)_{\mathfrak p'}) +
\dim(\mathcal{O}_{X', y'}^\wedge/\mathfrak p') +
\delta^{Y'}_{Z'}(y') > 3
$$
와 동치이다. $(\mathcal{F}_n)$이 (3)(b)의 부등식을 만족하면
(\ref{algebraization-equation-two})를 적용하여 표시된 두 부등식 중 둘째 것이 성립함을
안다. $(\mathcal{F}_n)$이 (3)(a), 즉 $(d + 1, d + 2)$-부등식들을 만족하면
(\ref{algebraization-equation-one})과 (\ref{algebraization-equation-two})에서 즉시 따른다. 이로써 주장의
증명이 끝난다.

\medskip\noindent
보조정리 \ref{algebraization-lemma-improvement-formal-coherent-module-better}에서와 같이
$(b|_{U'}^*\mathcal{F}_n) \to (\mathcal{F}_n'')$과
$\textit{Coh}(U', I\mathcal{O}_{U'})$의 $(\mathcal{H}_n)$을 택하자.
임의의 아핀 열린집합 $W \subset X'$에 대해 보조정리
\ref{algebraization-lemma-discussion}의 (6)에 의해
$\delta^{W \cap Y'}_{W \cap Z'}(y') \geq \delta^{Y'}_{Z'}(y')$이다.
따라서 $(\mathcal{H}_n|_W)$은 보조정리 \ref{algebraization-lemma-cd-1-canonical}의
가정들을 만족한다. 그러므로 $(\mathcal{H}_n|_W)$은 $W$로 표준적으로
연장된다. $(\mathcal{G}_{W, n})$을
보조정리 \ref{algebraization-lemma-canonically-algebraizable}에서와 같은
$\textit{Coh}(W, I\mathcal{O}_W)$의 표준 연장이라 하자. 보조정리
\ref{algebraization-lemma-canonically-extend-base-change}에 의해 $W' \subset W$에 대해
유일한 동형
$$
(\mathcal{G}_{W, n}|_{W'}) \longrightarrow
(\mathcal{G}_{W', n})
$$
이 존재하고, 이는 주어진 동형
$(\mathcal{G}_{W, n}|_{W \cap U}) \cong (\mathcal{H}_n|_{W \cap U})$과
양립한다. 따라서 $U$로의 제한이 $(\mathcal{H}_n)$과 동형인
$\textit{Coh}(X', I\mathcal{O}_{X'})$의 대상 $(\mathcal{G}_n)$이 존재한다.

\medskip\noindent
$A$가 $I$-진 완비이면 다음과 같이 증명을 끝낼 수 있다. Grothendieck 존재
정리(스킴의 코호몰로지의 보조정리
\ref{coherent-lemma-existence-projective})에 의해 $(\mathcal{G}_n)$은 연접
$\mathcal{O}_{X'}$-가군의 완비이다. 이어서 스킴의 코호몰로지의 보조정리
\ref{coherent-lemma-existence-easy}에 의해 $(b|_{U'}^*\mathcal{F}_n)$은 연접
$\mathcal{O}_{U'}$-가군 $\mathcal{F}'$의 완비이다. 스킴의 코호몰로지의
보조정리 \ref{coherent-lemma-inverse-systems-push-pull}에 의해 사상
$$
(\mathcal{F}_n) \longrightarrow ((b|_{U'})_*\mathcal{F}')^\wedge
$$
이 있고 그 핵과 여핵은 $I$의 어떤 거듭제곱에 의해 소멸된다. 마지막으로
보조정리 \ref{algebraization-lemma-map-kernel-cokernel-on-closed}를 적용하면 끝난다.

\medskip\noindent
$A$가 완비가 아니면 증명을 시작하기 전에 $A$를 그 완비로 바꾸어도 된다.
보조정리 \ref{algebraization-lemma-algebraizable}를 보라. 완비한 뒤에도 가정들은 여전히
성립한다. 조건 (3)은 즉시 알 수 있고, 조건 (1)은 쌍대화 복합체의 보조정리
\ref{dualizing-lemma-ubiquity-dualizing}에서, 조건 (2)는 인자의 보조정리
\ref{divisors-lemma-flat-base-change-blowing-up}에서 따른다. 따라서 완비인
경우에서 일반적인 경우가 따른다.
\end{proof}

\begin{proposition}[생성원이 적은 아이디얼에 대한 대수화]
\label{algebraization-proposition-d-generators}
상황 \ref{algebraization-situation-algebraize}에서 $(\mathcal{F}_n)$을
$\textit{Coh}(U, I\mathcal{O}_U)$의 대상이라 하자. 다음을 가정한다.
\begin{enumerate}
\item $A$는 쌍대화 복합체를 가진다.
\item 어떤 $d \geq 1$ 및 $f_1, \ldots, f_d \in A$에 대해
$V(I) = V(f_1, \ldots, f_d)$이다.
\item 다음 중 하나가 참이다.
\begin{enumerate}
\item $(\mathcal{F}_n)$은 $(d + 1, d + 2)$-부등식들을 만족한다
(정의 \ref{algebraization-definition-s-d-inequalities}).
\item $y \in U \cap Y$와 다음을 만족하는 소아이디얼
$\mathfrak p \subset \mathcal{O}_{X, y}^\wedge$로서
$\mathfrak p \not \in V(I\mathcal{O}_{X, y}^\wedge)$
에 대해
$$
\text{depth}((\mathcal{F}^\wedge_y)_\mathfrak p) +
\dim(\mathcal{O}_{X, y}^\wedge/\mathfrak p) + \delta^Y_Z(y) > d + 2
$$
\end{enumerate}
\end{enumerate}
그러면 $(\mathcal{F}_n)$은 $X$로 연장된다. 특히 $A$가 $I$-진 완비이면
$(\mathcal{F}_n)$은 연접 $\mathcal{O}_U$-가군의 완비이다.
\end{proposition}

\begin{proof}
$I = (f_1, \ldots, f_d)$라고 가정할 수 있다. 스킴의 코호몰로지의 보조정리
\ref{coherent-lemma-inverse-systems-ideals-equivalence}.
를 보라. 그러면 $I$에서 $X$를 부풀린 것의 모든 올은 차원
많아야 $d - 1$이다. 따라서 보조정리 \ref{algebraization-lemma-blowup}에서 연장을 얻는다.
마지막 명제는 보조정리 \ref{algebraization-lemma-essential-image-completion}에서 따른다.
\end{proof}

\noindent
다음 보조정리를 비고 \ref{algebraization-remark-interesting-case-variant},
\ref{algebraization-remark-interesting-case},
\ref{algebraization-remark-interesting-case-bis}, 그리고
\ref{algebraization-remark-interesting-case-ter}와 비교하라.

\begin{lemma}
\label{algebraization-lemma-interesting-case-final}
상황 \ref{algebraization-situation-algebraize}에서 $(\mathcal{F}_n)$을
$\textit{Coh}(U, I\mathcal{O}_U)$의 대상이라 하자. 다음을 가정한다.
\begin{enumerate}
\item $A$는 쌍대화 복합체를 갖는 국소환이다.
\item $X$의 모든 기약 성분은 같은 차원을 갖는다.
\item 스킴 $X \setminus Y$는 Cohen--Macaulay이다.
\item $I$는 $d$개의 원소로 생성된다.
\item $\dim(X) - \dim(Z) > d + 2$이다.
\item $y \in U \cap Y$에 대해 가군 $\mathcal{F}_y^\wedge$은
$V(I\mathcal{O}_{X, y}^\wedge)$ 밖에서 유한 국소 자유이다. 예를 들어
$\mathcal{F}_n$이 유한 국소 자유 $\mathcal{O}_U/I^n\mathcal{O}_U$-가군인
경우가 이에 해당한다.
\end{enumerate}
그러면 $(\mathcal{F}_n)$은 $X$로 연장된다. 특히 $A$가 $I$-진 완비이면
$(\mathcal{F}_n)$은 연접 $\mathcal{O}_U$-가군의 완비이다.
\end{lemma}

\begin{proof}
명제 \ref{algebraization-proposition-d-generators}의 가정 (1), (2), (3)(b)가 성립함을
보이겠다. (1), (2)는 명백하다.

\medskip\noindent
점 $y \in U \cap Y$와 소아이디얼 $\mathfrak p$, 즉
$\mathfrak p \subset \mathcal{O}_{X, y}^\wedge$이면서
$\mathfrak p \not \in V(I\mathcal{O}_{X, y}^\wedge)$.
을 취하자. 마지막 조건에 의해
$\text{depth}((\mathcal{F}_y^\wedge)_\mathfrak p) =
\text{depth}((\mathcal{O}_{X, y}^\wedge)_\mathfrak p)$.
$X \setminus Y$가 Cohen--Macaulay이므로
$(\mathcal{O}_{X, y}^\wedge)_\mathfrak p$도 Cohen--Macaulay이다. 따라서
\begin{align*}
& \text{depth}((\mathcal{F}^\wedge_y)_\mathfrak p) +
\dim(\mathcal{O}_{X, y}^\wedge/\mathfrak p) + \delta^Y_Z(y) \\
& =
\dim((\mathcal{O}_{X, y}^\wedge)_\mathfrak p) +
\dim(\mathcal{O}_{X, y}^\wedge/\mathfrak p) + \delta^Y_Z(y) \\
& =
\dim(\mathcal{O}_{X, y}^\wedge) + \delta^Y_Z(y)
\end{align*}
이다. 마지막 등식은 둘째 조건에 의해 $\mathcal{O}_{X, y}$가
등차원이기 때문이다. $\delta(y) = \dim(\overline{\{y\}})$로 놓자.
$A$가 사슬적 국소환이므로 이는 차원 함수이다. 보조정리
\ref{algebraization-lemma-discussion}에 의해
$\delta^Y_Z(y) \geq \delta(y) - \dim(Z)$이다. $X$가 등차원이므로
$$
\dim(\mathcal{O}_{X, y}^\wedge) + \delta^Y_Z(y)
\geq \dim(\mathcal{O}_{X, y}^\wedge) + \delta(y) - \dim(Z)
= \dim(X) - \dim(Z)
$$
따라서 원하는 부등식을 얻고 증명이 끝난다.
\end{proof}

\begin{remark}
\label{algebraization-remark-question}
명제 \ref{algebraization-proposition-d-generators}에서 $I$가 $d$개의 생성원을 갖는다는
가정을 $\text{cd}(A, I) \leq d$라는 가정으로 바꾼 유사 명제를 증명하거나
반증하지 못하였다. 증명을 알거나 반례가 있다면
\href{mailto:stacks.project@gmail.com}{stacks.project@gmail.com}.
로 전자우편을 보내기 바란다. 또 하나의 자연스러운 질문은 명제
\ref{algebraization-proposition-d-generators}의 조건들이 어느 정도까지 필요한가이다.
\end{remark}




\section{연접 형식 가군의 대수화 VI}
\label{algebraization-section-algebraization-modules-yet-more}

\noindent
이 절에서는 증명이 더 쉬운 몇 가지 경우를 덧붙인다.

\begin{proposition}
\label{algebraization-proposition-algebraization-regular-sequence}
상황 \ref{algebraization-situation-algebraize}에서 $(\mathcal{F}_n)$을
$\textit{Coh}(U, I\mathcal{O}_U)$의 대상이라 하자. 다음을 가정한다.
\begin{enumerate}
\item 어떤 $f_1, \ldots, f_d \in I$가 존재하여, $y \in U \cap Y$에 대해
아이디얼 $I\mathcal{O}_{X, y}$은 $f_1, \ldots, f_d$로 생성되고
$f_1, \ldots, f_d$는 $\mathcal{F}_y^\wedge$-정칙열을 이룬다.
\item $H^0(U, \mathcal{F}_1)$과 $H^1(U, \mathcal{F}_1)$은 유한
$A$-가군이다.
\end{enumerate}
그러면 $(\mathcal{F}_n)$은 $X$로 표준적으로 연장된다. 특히 $A$가
완비이면 $(\mathcal{F}_n)$은 연접 $\mathcal{O}_U$-가군의 완비이다.
\end{proposition}

\begin{proof}
보조정리 \ref{algebraization-lemma-when-done}의 가정 (a), (b), (c)를 확인하여 증명하겠다.
모든 $n$에 대해 짧은 완전열
$$
0 \to I^n\mathcal{F}_{n + 1} \to \mathcal{F}_{n + 1} \to \mathcal{F}_n \to 0
$$
이 있다. $f_1, \ldots, f_d$는 각 ``줄기'' $\mathcal{F}_y^\wedge$ 위에서
정칙열(따라서 준정칙열, 대수학의 보조정리
\ref{algebra-lemma-regular-quasi-regular}를 보라)을 이루고, 모든 $n$에 대해
$I\mathcal{F}_n = (f_1, \ldots, f_d)\mathcal{F}_n$이므로, 줄기에서 확인하면
$$
I^n\mathcal{F}_{n + 1} =
\bigoplus\nolimits_{e_1 + \ldots + e_d = n} \mathcal{F}_1 \cdot
f_1^{e_1} \ldots f_d^{e_d}
$$
이다. $H^0(U, \mathcal{F}_1)$의 유한성 가정과 귀납법을 사용하면 먼저
모든 $n$에 대해 $M_n = H^0(U, \mathcal{F}_n)$이 유한 $A$-가군임을
얻는다. 이로써 보조정리 \ref{algebraization-lemma-when-done}의 조건 (c)가 성립한다.
또한
$$
\bigoplus\nolimits_{n \geq 0} H^1(U, I^n\mathcal{F}_{n + 1})
$$
은 유한 등급 $R = \bigoplus I^n/I^{n +1}$-가군이다. 코호몰로지의
보조정리 \ref{cohomology-lemma-ML-general}에 의해 보조정리
\ref{algebraization-lemma-when-done}의 조건 (a)가 성립한다. 마지막으로
$\bigoplus H^0(U, I^n\mathcal{F}_{n + 1})$가 유한 등급 $R$-가군이고
보조정리 \ref{algebraization-lemma-when-done}의 조건 (b)도 코호몰로지의 보조정리
\ref{cohomology-lemma-topology-I-adic-general}을 적용할 수 있으므로 성립한다.
\end{proof}

\begin{remark}
\label{algebraization-remark-interesting-case-ter}
명제 \ref{algebraization-proposition-algebraization-regular-sequence}의 상황에서 $A$가
쌍대화 복합체를 가진다고 가정하면, $H^0(U, \mathcal{F}_1)$과
$H^1(U, \mathcal{F}_1)$이 유한이라는 조건은
$$
\text{depth}(\mathcal{F}_{1, y}) +
\dim(\mathcal{O}_{\overline{\{y\}}, z}) > 2
$$
가 모든 $y \in U \cap Y$와 $z \in Z \cap \overline{\{y\}}$에 대해
성립하는 것과 동치이다. 국소 코호몰로지의 보조정리
\ref{local-cohomology-lemma-finiteness-Rjstar}를 보라. 예를 들어
$\mathcal{F}_1$이 유한 국소 자유 $\mathcal{O}_{U \cap Y}$-가군이고,
$Y$가 $(S_2)$이며, $Z' \subset Y'$인 $Y$의 기약 성분 $Y'$과 $Z$의
기약 성분 $Z'$의 모든 쌍에 대해 $\text{codim}(Z', Y') \geq 3$이면
이 조건이 성립한다.
\end{remark}

\begin{proposition}
\label{algebraization-proposition-algebraization-flat}
상황 \ref{algebraization-situation-algebraize}에서 $(\mathcal{F}_n)$을
$\textit{Coh}(U, I\mathcal{O}_U)$의 대상이라 하자. 다음을 만족하는
Noether 국소환 $(R, \mathfrak m)$과 환 준동형 $R \to A$가 존재한다고
가정하자.
\begin{enumerate}
\item $I = \mathfrak m A$,
\item $y \in U \cap Y$에 대해 줄기 $\mathcal{F}_y^\wedge$은 $R$-평탄이다.
\item $H^0(U, \mathcal{F}_1)$과 $H^1(U, \mathcal{F}_1)$은 유한
$A$-가군이다.
\end{enumerate}
그러면 $(\mathcal{F}_n)$은 $X$로 표준적으로 연장된다. 특히 $A$가
완비이면 $(\mathcal{F}_n)$은 연접 $\mathcal{O}_U$-가군의 완비이다.
\end{proposition}

\begin{proof}
증명은 명제 \ref{algebraization-proposition-algebraization-regular-sequence}의 증명과
정확히 같다. 즉 $\kappa = R/\mathfrak m$이라 하면 $n \geq 0$에 대해 동형
$$
I^n \mathcal{F}_{n + 1} \cong
\mathcal{F}_1 \otimes_\kappa \mathfrak m^n/\mathfrak m^{n + 1}
$$
이 있고, 우변은 $\mathcal{F}_1$의 유한 직합이다. 이는 줄기에서 확인할 수
있다. 나머지는 모두 정확히 같다.
\end{proof}

\begin{remark}
\label{algebraization-remark-interesting-case-quater}
명제 \ref{algebraization-proposition-algebraization-flat}은
\cite[정리 2.10 (i)]{Baranovsky}의 국소적 형태이다. 국소 결과에서
대역 결과를 유도하는 것은 직접적이며, 그 논증을 개략적으로 설명하겠다.
즉 $(R, \mathfrak m)$을 완비 Noether 국소환이라 하고
$X \to \Spec(R)$을 고유 사상이라 하자. $n \geq 1$에 대해
$X_n = X \times_{\Spec(R)} \Spec(R/\mathfrak m^n)$으로 놓는다.
$Z \subset X_1$을 특수 올의 닫힌 부분집합이라 하자.
$U = X \setminus Z$로 놓고 포함 사상을 $j : U \to X$로 나타낸다.
다음 대상이 주어졌다고 하자.
$$
(\mathcal{F}_n) \text{는 } \textit{Coh}(U, \mathfrak m\mathcal{O}_U)\text{의 대상}
$$
여기서 모든 $n$에 대해 $\mathcal{F}_n$이 $R/\mathfrak m^n$ 위에서
평탄하다는 의미로 이 계는 $R$ 위에서 평탄하다.
$j_*\mathcal{F}_1$과 $R^1j_*\mathcal{F}_1$이 연접 가군이라고 가정하자.
그러면 명제 \ref{algebraization-proposition-algebraization-flat}에 의해 $X$ 위에서
아핀 국소적으로 $(\mathcal{F}_n)$의 표준적 연장을 얻으며, 이 연장의
구성은 국소화와 가환한다(보조정리
\ref{algebraization-lemma-algebraization-principal-variant}). 따라서 그 $U$로의 제한이
$(\mathcal{F}_n)$인 $\textit{Coh}(X, \mathfrak m\mathcal{O}_X)$의 표준적
대역 대상 $(\mathcal{G}_n)$을 얻는다. Grothendieck 존재 정리
(스킴의 코호몰로지, 명제 \ref{coherent-proposition-existence-proper})에 의해
완비가 $(\mathcal{G}_n)$인 연접 $\mathcal{O}_X$-가군 $\mathcal{G}$가
존재한다. 이로부터 $(\mathcal{F}_n)$이 대수화 가능함, 즉 어떤 연접
$\mathcal{O}_U$-가군의 완비임을 알 수 있다.

\medskip\noindent
$j_*\mathcal{F}_1$과 $R^1j_*\mathcal{F}_1$의 연접성은 특수 올 위의
조건임도 덧붙인다. 실제로 $j : U \to X$의 특수 올을
$j_1 : U_1 \to X_1$로 나타내면 $\mathcal{F}_1$을 $U_1$ 위의 연접 층으로
볼 수 있고,
$j_*\mathcal{F}_1 = j_{1, *}\mathcal{F}_1$ 및
$R^1j_*\mathcal{F}_1 = R^1j_{1, *}\mathcal{F}_1$이다.
따라서 예를 들어 $X_1$이 $(S_2)$이고 기약이며,
$\dim(X_1) - \dim(Z) \geq 3$이고, $\mathcal{F}_1$이 국소 자유
$\mathcal{O}_{U_1}$-가군이면 $j_{1, *}\mathcal{F}_1$과
$R^1j_{1, *}\mathcal{F}_1$은 연접 가군이다.
\end{remark}








\section{완비화 함자에의 응용}
\label{algebraization-section-completion-application}

\noindent
이 절에서는 편리하게 인용할 수 있도록 이미 얻은 몇 가지 결과를 단순히
결합한다. 여기서 진술할 수 있는 (더 강한) 결과는 많이 있다.

\begin{lemma}
\label{algebraization-lemma-equivalence-better}
상황 \ref{algebraization-situation-algebraize}에서 다음을 가정하자.
\begin{enumerate}
\item $A$는 쌍대화 복합체를 가지며 $I$-진 완비이다.
\item $I = (f)$는 하나의 원소로 생성된다.
\item $A$는 극대 아이디얼 $\mathfrak a = \mathfrak m$을 갖는 국소환이다.
\item 다음 중 하나가 성립한다.
\begin{enumerate}
\item $A_f$는 $(S_2)$이고, $f \not \in \mathfrak p$인 극소
$\mathfrak p \subset A$에 대해 $\dim(A/\mathfrak p) \geq 4$이다. 또는
\item $\mathfrak p \not \in V(f)$이고
$V(\mathfrak p) \cap V(f) \not = \{\mathfrak m\}$이면
$\text{depth}(A_\mathfrak p) + \dim(A/\mathfrak p) > 3$.
\end{enumerate}
\end{enumerate}
그러면 $U_0 = U \cap V(f)$로 놓을 때 완비화 함자
$$
\colim_{U_0 \subset U' \subset U\text{ 열린 부분스킴}}
\textit{Coh}(\mathcal{O}_{U'})
\longrightarrow
\textit{Coh}(U, f\mathcal{O}_U)
$$
는 유한 국소 자유 대상들로 이루어진 충만한 부분범주 사이의 동치이다.
\end{lemma}

\begin{proof}
보조정리 \ref{algebraization-lemma-fully-faithful-general}에서 이 함자가 완전충실임이
따른다(세부사항은 생략한다). 본질적 전사성을 증명하자.
$(\mathcal{F}_n)$을 $\textit{Coh}(U, f\mathcal{O}_U)$의 유한 국소 자유
대상이라 하자. 보조정리 \ref{algebraization-lemma-algebraization-principal-bis} 또는
명제 \ref{algebraization-proposition-cd-1}에 의해 $(\mathcal{F}_n)$이 $\mathcal{F}$의
완비가 되는 연접 $\mathcal{O}_U$-가군 $\mathcal{F}$가 존재한다. 실제로
두 결과 중 어느 것을 적용하든 확인할 것은 $(\mathcal{F}_n)$이
$(2, 3)$-부등식을 만족한다는 것뿐이며, 이는 보조정리
\ref{algebraization-lemma-unwinding-conditions-bis}에서 확인했다. $y \in U_0$이면 줄기
$\mathcal{F}_y$의 $f$-진 완비는 $\mathcal{O}_{U, y}$의 $f$-진 완비 위의
유한 자유 가군과 동형이다. 따라서 $U_0$의 어떤 열린 근방 $U'$에서
$\mathcal{F}$는 유한 국소 자유이다. 이로써 증명이 끝난다.
\end{proof}

\begin{lemma}
\label{algebraization-lemma-equivalence}
상황 \ref{algebraization-situation-algebraize}에서 다음을 가정하자.
\begin{enumerate}
\item $I = (f)$는 주 아이디얼이다.
\item $A$는 $f$-진 완비이다.
\item $f$는 영인자가 아니다.
\item $H^1_\mathfrak a(A/fA)$와 $H^2_\mathfrak a(A/fA)$
는 유한 $A$-가군이다.
\end{enumerate}
그러면 $U_0 = U \cap V(f)$로 놓을 때 완비화 함자
$$
\colim_{U_0 \subset U' \subset U\text{ 열린 부분스킴}}
\textit{Coh}(\mathcal{O}_{U'})
\longrightarrow
\textit{Coh}(U, f\mathcal{O}_U)
$$
는 유한 국소 자유 대상들로 이루어진 충만한 부분범주 사이의 동치이다.
\end{lemma}

\begin{proof}
보조정리 \ref{algebraization-lemma-fully-faithful-general-alternative}에 의해 이 함자는
완전충실하다. 본질적 전사성은 보조정리
\ref{algebraization-lemma-algebraization-principal-variant}에서 따른다.
\end{proof}








\section{연접 삼중항}
\label{algebraization-section-coherent-triples}

\noindent
$(A, \mathfrak m)$을 Noether 국소환이라 하자.
$f \in \mathfrak m$을 영인자가 아닌 원소라 하자. 다음과 같이 놓는다.
$X = \Spec(A)$, $X_0 = \Spec(A/fA)$, $U = X \setminus V(\mathfrak m)$, 그리고
$U_0 = U \cap X_0$.
다음을 만족할 때 $(\mathcal{F}, \mathcal{F}_0, \alpha)$를
{\it 연접 삼중항}이라 한다.
\begin{enumerate}
\item $\mathcal{F}$는 $f : \mathcal{F} \to \mathcal{F}$가 단사인 연접
$\mathcal{O}_U$-가군이다.
\item $\mathcal{F}_0$는 연접 $\mathcal{O}_{X_0}$-가군이다.
\item $\alpha : \mathcal{F}/f\mathcal{F} \to \mathcal{F}_0|_{U_0}$
는 동형이다.
\end{enumerate}
모든 연접 삼중항의 모임을 범주로 만드는 자명한 {\it 연접 삼중항의 사상}
개념이 있다.

\medskip\noindent
연접 삼중항의 범주는 가법 범주이지만 Abel 범주는 아니다. 그러나 연접
삼중항의 짧은 완전열이 무엇인지는 명백하다.

\medskip\noindent
두 연접 삼중항 $(\mathcal{F}, \mathcal{F}_0, \alpha)$와
$(\mathcal{G}, \mathcal{G}_0, \beta)$가 주어져도
$(\mathcal{F} \otimes_{\mathcal{O}_U} \mathcal{G},
\mathcal{F}_0  \otimes_{\mathcal{O}_{X_0}} \mathcal{G}_0,
\alpha \otimes \beta)$가 연접 삼중항일 필요는 없다\footnote{즉
$\mathcal{F} \otimes_{\mathcal{O}_U} \mathcal{G}$ 위에서 $f$가 단사일
필요는 없다.}. 그러나 모든 $x \in U_0$에 대해 줄기 $\mathcal{G}_x$가
자유이면 이는 성립한다.

\medskip\noindent
연접 삼중항 $(\mathcal{G}, \mathcal{G}_0, \beta)$에서
$\mathcal{G}$와 $\mathcal{G}_0$가 각각 국소 자유 가군이면 이를
{\it 국소 자유}라 하고, 각각 가역 가군이면 {\it 가역}이라 한다. 이 경우
$(\mathcal{G}, \mathcal{G}_0, \beta)$와의 텐서곱은 의미가 있으며(위의
논의를 보라), 연접 삼중항의 짧은 완전열을 연접 삼중항의 짧은 완전열로
보낸다.

\begin{lemma}
\label{algebraization-lemma-prepare-chi-triple}
임의의 연접 삼중항 $(\mathcal{F}, \mathcal{F}_0, \alpha)$에 대해
$f : \mathcal{F}' \to \mathcal{F}'$가 단사인 연접
$\mathcal{O}_X$-가군 $\mathcal{F}'$, 동형
$\alpha' : \mathcal{F}'|_U \to \mathcal{F}$ 및 사상
$\alpha'_0 : \mathcal{F}'/f\mathcal{F}' \to \mathcal{F}_0$
가 존재하여 $\alpha \circ (\alpha' \bmod f) = \alpha'_0|_{U_0}$를 만족한다.
\end{lemma}

\begin{proof}
$\mathcal{F}$가 유한 $A$-가군 $M$에 대응하는 연접
$\mathcal{O}_X$-가군의 $U$로의 제한이 되도록 $M$을 택한다. 국소
코호몰로지의 보조정리
\ref{local-cohomology-lemma-finiteness-pushforwards-and-H1-local}를 보라.
$\mathcal{F}$는 $f$-꼬임이 없으므로 $M$을 그 $f$-멱 꼬임에 의한
몫으로 바꾸어도 된다. 한편 $M_0 = \Gamma(X_0, \mathcal{F}_0)$로 놓으면
$\mathcal{F}_0$는 유한 $A/fA$-가군 $M_0$에 대응하는 연접
$\mathcal{O}_{X_0}$-가군이다. 스킴의 코호몰로지의 보조정리
\ref{coherent-lemma-homs-over-open}에 의해 동형 $\alpha_0$가
$A/fA$-가군 준동형 $\mathfrak m^n M/fM \to M_0$에 대응하는 $n$이
존재한다(그 핵과 여핵은 $\mathfrak m$의 어떤 멱에 의해 소멸하지만 여기서는
이를 쓰지 않는다). 따라서 $M' = \mathfrak m^n M$으로 놓고
$\mathcal{F}'$을 $M'$에 대응하는 연접 $\mathcal{O}_X$-가군으로 잡으면
보조정리는 명백하다.
\end{proof}

\noindent
$(\mathcal{F}, \mathcal{F}_0, \alpha)$를 연접 삼중항이라 하자.
보조정리 \ref{algebraization-lemma-prepare-chi-triple}와 같이
$\mathcal{F}', \alpha', \alpha'_0$를 택하고 다음과 같이 놓는다.
\begin{equation}
\label{algebraization-equation-chi-triple}
\chi(\mathcal{F}, \mathcal{F}_0, \alpha) =
\text{length}_A(\Coker(\alpha'_0)) - 
\text{length}_A(\Ker(\alpha'_0))
\end{equation}
우변은 의미가 있다. 실제로 $\alpha'_0$는 $U_0$ 위에서 동형이므로 그 핵과
여핵은 $\{\mathfrak m\}$에 지지된 연접 가군이며, 따라서 길이가 유한하다
(대수학, 보조정리 \ref{algebra-lemma-support-point}).

\begin{lemma}
\label{algebraization-lemma-well-defined-chi-triple}
(\ref{algebraization-equation-chi-triple})의 양 $\chi(\mathcal{F}, \mathcal{F}_0, \alpha)$는
보조정리 \ref{algebraization-lemma-prepare-chi-triple}와 같은
$\mathcal{F}', \alpha', \alpha'_0$의 선택에 의존하지 않는다.
\end{lemma}

\begin{proof}
$\mathcal{F}', \alpha', \alpha'_0$와
$\mathcal{F}'', \alpha'', \alpha''_0$를 그러한 두 선택이라 하자.
$n > 0$에 대해 $\mathcal{F}'_n = \mathfrak m^n \mathcal{F}'$로 놓는다.
스킴의 코호몰로지의 보조정리 \ref{coherent-lemma-homs-over-open}에 의해
어떤 $n$에 대해 $\alpha'$와 $\alpha''$가 정하는 식별
$\mathcal{F}''|_U = \mathcal{F}'|_U$와 일치하는 $\mathcal{O}_X$-가군 사상
$\mathcal{F}'_n \to \mathcal{F}''$가 존재한다. 그러면 그림식
$$
\xymatrix{
\mathcal{F}'_n/f\mathcal{F}'_n \ar[r] \ar[d] &
\mathcal{F}'/f\mathcal{F}' \ar[d]^{\alpha_0'} \\
\mathcal{F}''/f\mathcal{F}'' \ar[r]^{\alpha_0''} &
\mathcal{F}_0
}
$$
은 $U_0$로 제한하면 가환한다. 따라서 스킴의 코호몰로지의 보조정리
\ref{coherent-lemma-homs-over-open}에 의해 어떤 $l > 0$에 대해
$\mathfrak m^l(\mathcal{F}'_n/f\mathcal{F}'_n)$으로 제한하면 가환한다.
$\mathcal{F}'_{n + l}/f\mathcal{F}'_{n + l} \to
\mathcal{F}'_n/f\mathcal{F}'_n$가
$\mathfrak m^l(\mathcal{F}'_n/f\mathcal{F}'_n)$을 거쳐 분해되므로 $n$을
$n + l$로 바꾸면 그림식이 가환한다. 달리 말해 $U$와 $X_0$ 위의 사상들과
호환되는 $X$ 위의 사상 $\mathcal{F}''' \to \mathcal{F}''$ 및
$\mathcal{F}''' \to \mathcal{F}'$가 존재하는 세 번째 선택
$\mathcal{F}''', \alpha''', \alpha'''_0$를 찾았다. 이로써 다음 문단에서
다루는 경우로 환원된다.

\medskip\noindent
$U$ 위에서 $\alpha', \alpha''$와 호환되고 $X_0$ 위에서
$\alpha'_0, \alpha''_0$와 호환되는 $X$ 위의 사상
$\mathcal{F}'' \to \mathcal{F}'$가 있다고 가정하자. 이 사상은 $U$ 위에서
동형이고 $f : \mathcal{F}'' \to \mathcal{F}''$가 단사이므로
$\mathcal{F}'' \to \mathcal{F}'$는 단사이다. 분명히
$\mathcal{F}'/\mathcal{F}''$는 $\{\mathfrak m\}$에 지지되므로 길이가
유한하다. 다음과 같은 연접 $\mathcal{O}_{X_0}$-가군의 사상들이 있다.
$$
\mathcal{F}''/f\mathcal{F}'' \to
\mathcal{F}'/f\mathcal{F}' \xrightarrow{\alpha'_0}
\mathcal{F}_0
$$
그 합성은 $\alpha''_0$이고 이 사상들은 $U_0$ 위에서 동형이다. 초등적인
호몰로지 대수학으로부터 다음 $6$항 완전열을 얻는다.
$$
\begin{matrix}
0 \to
\Ker(\mathcal{F}''/f\mathcal{F}'' \to \mathcal{F}'/f\mathcal{F}') \to
\Ker(\alpha''_0) \to
\Ker(\alpha'_0) \to \\
\Coker(\mathcal{F}''/f\mathcal{F}'' \to \mathcal{F}'/f\mathcal{F}') \to
\Coker(\alpha''_0) \to
\Coker(\alpha'_0) \to 0
\end{matrix}
$$
길이의 가법성(대수학, 보조정리 \ref{algebra-lemma-length-additive})에 의해
다음을 보이면 충분하다.
$$
\text{length}_A(
\Coker(\mathcal{F}''/f\mathcal{F}'' \to \mathcal{F}'/f\mathcal{F}')) -
\text{length}_A(
\Ker(\mathcal{F}''/f\mathcal{F}'' \to \mathcal{F}'/f\mathcal{F}')) = 0
$$
이는 다음 그림식에 뱀 보조정리를 적용하고
$$
\xymatrix{
0 \ar[r] &
\mathcal{F}'' \ar[r]_f \ar[d] &
\mathcal{F}'' \ar[r] \ar[d] &
\mathcal{F}''/f\mathcal{F}'' \ar[r] \ar[d] &
0 \\
0 \ar[r] &
\mathcal{F}' \ar[r]^f &
\mathcal{F}' \ar[r] &
\mathcal{F}'/f\mathcal{F}' \ar[r] &
0
}
$$
$\mathcal{F}'/\mathcal{F}''$의 길이가 유한하다는 사실을 사용하면 따른다.
\end{proof}

\begin{lemma}
\label{algebraization-lemma-ses-chi-triple}
다음 짧은 연접 삼중항 완전열이 있으면
$\chi(\mathcal{G}, \mathcal{G}_0, \beta) =
\chi(\mathcal{F}, \mathcal{F}_0, \alpha) +
\chi(\mathcal{H}, \mathcal{H}_0, \gamma)$이다.
$$
0 \to
(\mathcal{F}, \mathcal{F}_0, \alpha) \to
(\mathcal{G}, \mathcal{G}_0, \beta) \to
(\mathcal{H}, \mathcal{H}_0, \gamma)
\to 0
$$
\end{lemma}

\begin{proof}
삼중항 $(\mathcal{G}, \mathcal{G}_0, \beta)$에 대해 보조정리
\ref{algebraization-lemma-prepare-chi-triple}과 같이 $\mathcal{G}', \beta', \beta'_0$를
택한다. 포함 사상을 $j : U \to X$로 나타내고,
$\mathcal{F}' \subset \mathcal{G}'$을 다음 합성의 핵이라 하자.
$$
\mathcal{G}' \xrightarrow{\beta'} j_*\mathcal{G} \to j_*\mathcal{H}
$$
$\mathcal{H}' = \mathcal{G}'/\mathcal{F}'$은 $j_*\mathcal{H}$의 연접
부분층이므로 $f : \mathcal{H}' \to \mathcal{H}'$는 단사이다. 따라서 뱀
보조정리에 의해 다음 짧은 완전열을 얻는다.
$$
0 \to \mathcal{F}'/f\mathcal{F}' \to
\mathcal{G}'/f\mathcal{G}' \to
\mathcal{H}'/f\mathcal{H}' \to 0
$$
구성에 의해 다음 동형들이 있다.
$\alpha' : \mathcal{F}'|_U \to \mathcal{F}$,
$\beta' : \mathcal{G}'|_U \to \mathcal{G}$, 그리고
$\gamma' : \mathcal{H}'|_U \to \mathcal{H}$.
증명을 마치려면 보조정리 \ref{algebraization-lemma-prepare-chi-triple}과 같은 사상
$\alpha'_0 : \mathcal{F}'/f\mathcal{F}' \to \mathcal{F}_0$와
$\gamma'_0 : \mathcal{H}'/f\mathcal{H}' \to \mathcal{H}_0$를 구성하여
다음 가환 그림식에 맞추어야 한다.
$$
\xymatrix{
0 \ar[r] &
\mathcal{F}'/f\mathcal{F}' \ar[r] \ar@{..>}[d]^{\alpha'_0} &
\mathcal{G}'/f\mathcal{G}' \ar[r] \ar[d]^{\beta'_0} &
\mathcal{H}'/f\mathcal{H}' \ar[r] \ar@{..>}[d]^{\gamma'_0} &
0 \\
0 \ar[r] &
\mathcal{F}_0 \ar[r] &
\mathcal{G}_0 \ar[r] &
\mathcal{H}_0 \ar[r] &
0
}
$$
그러나 처음 선택한 $\mathcal{G}'$로는 이것이 가능하지 않을 수 있다.
위 그림식에서 장애물은 정확히 다음 합성임을 알 수 있다.
$$
\delta :
\mathcal{F}'/f\mathcal{F}' \to
\mathcal{G}'/f\mathcal{G}' \xrightarrow{\beta'_0}
\mathcal{G}_0 \to
\mathcal{H}_0
$$
선택한 $\mathcal{F}'$와 $\mathcal{H}'$에 의해 $\delta$의 $U_0$로의 제한은
영이다. 따라서 스킴의 코호몰로지의 보조정리
\ref{coherent-lemma-homs-over-open}에 의해 $\delta$가
$\mathfrak m^k \cdot (\mathcal{F}'/f\mathcal{F}')$ 위에서 영이 되는
$k > 0$이 존재한다. $n > k$에 대해
$\mathcal{G}'_n = \mathfrak m^n \mathcal{G}'$,
$\mathcal{F}'_n = \mathcal{G}'_n \cap \mathcal{F}'$, 그리고
$\mathcal{H}'_n = \mathcal{G}'_n/\mathcal{F}'_n$으로 놓는다.
$\beta'_0$를 $\mathcal{G}'_n/f\mathcal{G}'_n \to
\mathcal{G}'/f\mathcal{G}'$와 합성하면 보조정리
\ref{algebraization-lemma-prepare-chi-triple}과 같은 사상
$\beta'_{n, 0} : \mathcal{G}'_n/f\mathcal{G}'_n \to \mathcal{G}_0$를 얻는다.
Artin--Rees 정리(대수학, 보조정리 \ref{algebra-lemma-Artin-Rees})에 의해
$\mathcal{F}'_n \subset \mathfrak m^k \mathcal{F}'$가 되도록 $n$을 택할 수
있다. 위와 마찬가지로 사상
$f : \mathcal{F}'_n \to \mathcal{F}'_n$,
$f : \mathcal{G}'_n \to \mathcal{G}'_n$, 그리고
$f : \mathcal{H}'_n \to \mathcal{H}'_n$은 단사이고, 뱀 보조정리를 써서
다음 짧은 완전열을 얻는다.
$$
0 \to \mathcal{F}'_n/f\mathcal{F}'_n \to
\mathcal{G}'_n/f\mathcal{G}'_n \to
\mathcal{H}'_n/f\mathcal{H}'_n \to 0
$$
위와 마찬가지로 다음 동형들이 있다.
$\alpha'_n : \mathcal{F}'_n|_U \to \mathcal{F}$,
$\beta'_n : \mathcal{G}'_n|_U \to \mathcal{G}$, 그리고
$\gamma'_n : \mathcal{H}'_n|_U \to \mathcal{H}$.
앞에서와 같이 다음 장애물을 생각한다.
$$
\delta_n :
\mathcal{F}'_n/f\mathcal{F}'_n \to
\mathcal{G}'_n/f\mathcal{G}'_n
\xrightarrow{\beta'_{n, 0}}
\mathcal{G}_0 \to
\mathcal{H}_0
$$
그러나 가환 그림식
$$
\xymatrix{
\mathcal{F}'_n/f\mathcal{F}'_n \ar[r] \ar[d] &
\mathcal{G}'_n/f\mathcal{G}'_n \ar[r]_{\beta'_{n, 0}} \ar[d] &
\mathcal{G}_0 \ar[r] \ar[d] &
\mathcal{H}_0 \ar[d] \\
\mathcal{F}'/f\mathcal{F}' \ar[r] &
\mathcal{G}'/f\mathcal{G}' \ar[r]^{\beta'_0} &
\mathcal{G}_0 \ar[r] &
\mathcal{H}_0
}
$$
과 $n$의 선택 및 $\delta$에 관한 관찰로부터 $\delta_n = 0$임을 알 수 있다.
이로써 원하는 사상
$\alpha'_{n, 0} : \mathcal{F}'_n/f\mathcal{F}'_n \to \mathcal{F}_0$와
$\gamma'_{n, 0} : \mathcal{H}'_n/f\mathcal{H}'_n \to \mathcal{H}_0$를 얻는다.
따라서
$\mathcal{F}'_n, \alpha'_n, \alpha'_{n, 0}$,
$\mathcal{G}'_n, \beta'_n, \beta'_{n, 0}$, 그리고
$\mathcal{H}'_n, \gamma'_n, \gamma'_{n, 0}$
을 사용하여
$\chi(\mathcal{F}, \mathcal{F}_0, \alpha)$,
$\chi(\mathcal{G}, \mathcal{G}_0, \beta)$, 그리고
$\chi(\mathcal{H}, \mathcal{H}_0, \gamma)$를 계산해도 된다. 이제 마지막으로
다음 그림식에 뱀 보조정리를 적용하고
$$
\xymatrix{
0 \ar[r] &
\mathcal{F}'_n/f\mathcal{F}'_n \ar[r] \ar[d] &
\mathcal{G}'_n/f\mathcal{G}'_n \ar[r] \ar[d] &
\mathcal{H}'_n/f\mathcal{H}'_n \ar[r] \ar[d] &
0 \\
0 \ar[r] &
\mathcal{F}_0 \ar[r] &
\mathcal{G}_0 \ar[r] &
\mathcal{H}_0 \ar[r] &
0
}
$$
길이의 가법성(대수학, 보조정리 \ref{algebra-lemma-length-additive})을 쓰면
보조정리가 따른다.
\end{proof}

\begin{proposition}
\label{algebraization-proposition-hilbert-triple}
$(\mathcal{F}, \mathcal{F}_0, \alpha)$를 연접 삼중항이라 하고
$(\mathcal{L}, \mathcal{L}_0, \lambda)$를 가역 연접 삼중항이라 하자. 그러면 함수
$$
\mathbf{Z} \longrightarrow \mathbf{Z},\quad
n \longmapsto
\chi((\mathcal{F}, \mathcal{F}_0, \alpha) \otimes
(\mathcal{L}, \mathcal{L}_0, \lambda)^{\otimes n})
$$
는 차수가 $\leq \dim(\text{Supp}(\mathcal{F}))$인 다항식이다.
\end{proposition}

\noindent
더 정확히 말해 $\mathcal{F} = 0$이면 이 함수는 상수이다.
$\mathcal{F}$가 $U$에서 유한 지지를 가지면 이 함수는 상수이다.
$U$에서 $\mathcal{F}$의 지지의 차원이 $1$, 즉 $X$에서
$\mathcal{F}$의 지지의 폐포의 차원이 $2$이면 이 함수는 일차식인 식이다.

\begin{proof}
$\mathcal{F}$의 지지의 차원에 관한 귀납법으로 증명한다.

\medskip\noindent
기초 단계는 $\mathcal{F} = 0$인 경우이다. 이때 $\mathcal{F}_0$는 영이거나
그 지지가 $\{\mathfrak m\}$이다. 이 경우
$$
(\mathcal{F}, \mathcal{F}_0, \alpha) \otimes
(\mathcal{L}, \mathcal{L}_0, \lambda)^{\otimes n} =
(0, \mathcal{F}_0 \otimes \mathcal{L}_0^{\otimes n}, 0) \cong
(0, \mathcal{F}_0, 0)
$$
이므로 보조정리의 함수는 상수이고 그 값은 $\mathcal{F}_0$의 길이와 같다.

\medskip\noindent
귀납 단계로 $\mathcal{F}$의 지지가 공집합이 아니라고 하자.
$\mathcal{G}_0 \subset \mathcal{F}_0$를 $\{\mathfrak m\}$에 지지된 단면들의
부분가군이라 하면 다음 짧은 완전열을 얻는다.
$$
0 \to (0, \mathcal{G}_0, 0) \to
(\mathcal{F}, \mathcal{F}_0, \alpha) \to
(\mathcal{F}, \mathcal{F}_0/\mathcal{G}_0, \alpha) \to 0
$$
이 완전열은 가역 연접 삼중항
$(\mathcal{L}, \mathcal{L}_0, \lambda)$와 텐서곱해도 완전하다. 위의 논의를
보라. 따라서 $\chi$의 가법성(보조정리 \ref{algebraization-lemma-ses-chi-triple})과 위에서
설명한 기초 단계에 의해, 다음 삼중항에 대해 귀납 단계를 증명하면 충분하다.
$(\mathcal{F}, \mathcal{F}_0/\mathcal{G}_0, \alpha)$.
이와 같이 하여 $\mathfrak m$이 $\mathcal{F}_0$의 연관점이 아니라고 가정해도 된다.

\medskip\noindent
$T = \text{Ass}(\mathcal{F}) \cup \text{Ass}(\mathcal{F}/f\mathcal{F})$로 놓는다.
$U$는 준아핀이므로 어떤 $u \in T$에서도 영이 되지 않는
$s \in \Gamma(U, \mathcal{L})$를 찾을 수 있다. 성질의 보조정리
\ref{properties-lemma-quasi-affine-invertible-nonvanishing-section}를 보라.
$s$에 $\mathfrak m$의 적당한 원소를 곱하여, 어떤
$s_0 \in \Gamma(X_0, \mathcal{L}_0)$에 대해
$\lambda(s \bmod f) = s_0|_{U_0}$라고 가정할 수 있다. 세부사항은 생략한다.
연접 삼중항의 범주에서 사상
$$
(s, s_0) :
(\mathcal{O}_U, \mathcal{O}_{X_0}, 1)
\longrightarrow
(\mathcal{L}, \mathcal{L}_0, \lambda)
$$
을 얻는다. 다음과 같이 놓는다.
$\mathcal{G} = \Coker(s : \mathcal{F} \to \mathcal{F} \otimes \mathcal{L})$
및
$\mathcal{G}_0 = \Coker(s_0 : \mathcal{F}_0 \to
\mathcal{F}_0 \otimes \mathcal{L}_0)$. $s$의 선택에 의해 이 사상은 $U_0$
위에서 단사이고 $\mathfrak m$은 $\mathcal{F}_0$의 연관점이 아니므로
$s_0 : \mathcal{F}_0 \to \mathcal{F}_0 \otimes \mathcal{L}_0$는 단사이다.
따라서 다음 짧은 완전열을 주는 동형
$\beta : \mathcal{G}/f\mathcal{G} \to \mathcal{G}_0|_{U_0}$가 존재한다.
$$
0 \to
(\mathcal{F}, \mathcal{F}_0, \alpha) \to
(\mathcal{F}, \mathcal{F}_0, \alpha) \otimes
(\mathcal{L}, \mathcal{L}_0, \lambda) \to
(\mathcal{G}, \mathcal{G}_0, \beta) \to 0
$$
지지의 차원에 관한 귀납가정으로 명제는 연접 삼중항
$(\mathcal{G}, \mathcal{G}_0, \beta)$에 대해 성립한다. 보조정리
\ref{algebraization-lemma-ses-chi-triple}의 가법성을 쓰면 함수
$$
n \longmapsto 
\chi((\mathcal{F}, \mathcal{F}_0, \alpha) \otimes
(\mathcal{L}, \mathcal{L}_0, \lambda)^{\otimes n + 1})
-
\chi((\mathcal{F}, \mathcal{F}_0, \alpha) \otimes
(\mathcal{L}, \mathcal{L}_0, \lambda)^{\otimes n})
$$
가 다항식임을 알 수 있다. 모든 정수에서 정의된 함수에 대한 대수학의 보조정리
\ref{algebra-lemma-numerical-polynomial}의 변형을 적용하면 결론을 얻는다
(세부사항은 생략한다).
\end{proof}

\begin{lemma}
\label{algebraization-lemma-nonnegative-chi-triple}
$\text{depth}(A) \geq 3$, 또는 동치로 $\text{depth}(A/fA) \geq 2$라고
가정하자. $(\mathcal{L}, \mathcal{L}_0, \lambda)$를 가역 연접 삼중항이라 하자.
그러면
$$
\chi(\mathcal{L}, \mathcal{L}_0, \lambda) =
\text{length}_A \Coker(\Gamma(U, \mathcal{L}) \to \Gamma(U_0, \mathcal{L}_0))
$$
이며, 특히 이는 $\geq 0$이다. 또한
$\chi(\mathcal{L}, \mathcal{L}_0, \lambda) = 0$일 필요충분조건은
$\mathcal{L} \cong \mathcal{O}_U$인 것이다.
\end{lemma}

\begin{proof}
깊이 조건들의 동치는 대수학의 보조정리
\ref{algebra-lemma-depth-drops-by-one}에서 따른다. 깊이 조건에 의해
$\Gamma(U, \mathcal{O}_U) = A$이고
$\Gamma(U_0, \mathcal{O}_{U_0}) = A/fA$임을 안다. 쌍대화 복합체의 보조정리
\ref{dualizing-lemma-depth}와 국소 코호몰로지의 보조정리
\ref{local-cohomology-lemma-finiteness-pushforwards-and-H1-local}를 보라.
국소 코호몰로지의 보조정리
\ref{local-cohomology-lemma-finiteness-for-finite-locally-free}
에 의해 $M = \Gamma(U, \mathcal{L})$은 유한 $A$-가군이다. 이는 다시 예컨대
국소 코호몰로지의 보조정리
\ref{local-cohomology-lemma-finiteness-pushforwards-and-H1-local}
의 (4) 또는 인자의 보조정리 \ref{divisors-lemma-depth-pushforward}에 의해
$\text{depth}(M) \geq 2$임을 뜻한다. 또한 $X_0$는 국소 스킴이므로
$\mathcal{L}_0 \cong \mathcal{O}_{X_0}$이다. 따라서
$M_0 = \Gamma(X_0, \mathcal{L}_0) = \Gamma(U_0, \mathcal{L}_0|_{U_0})$
이고 이 가군은 $A/fA$와 동형이다.

\medskip\noindent
위에 의해 $\mathcal{F}' = \widetilde{M}$은 그 $U$로의 제한이
$\mathcal{L}$과 동형인 연접 $\mathcal{O}_X$-가군이다. 동형
$\lambda : \mathcal{L}/f\mathcal{L} \to \mathcal{L}_0|_{U_0}$는 전역단면
위의 사상 $M/fM \to M_0$를 정하며, 이는 $U_0$ 위에서 동형이다.
$\text{depth}(M) \geq 2$이므로 $H^0_\mathfrak m(M/fM) = 0$이고, 따라서
$M/fM \to M_0$는 단사이다. 그러므로 정의에 의해
$$
\chi(\mathcal{L}, \mathcal{L}_0, \lambda) =
\text{length}_A \Coker(M/fM \to M_0)
$$
이며, 이는 보조정리의 첫째 진술을 준다.

\medskip\noindent
마지막으로 이 길이가 $0$이면 $M \to M_0$는 전사이다. 따라서
$\mathcal{L}_0$의 자명화 단면으로 가는 $s \in M = \Gamma(U, \mathcal{L})$를
찾을 수 있다. 다음 완전열로 정의되는 유한 $A$-가군 $K$, $Q$를 생각하자.
$$
0 \to K \to A \xrightarrow{s} M \to Q \to 0
$$
$s$는 $U_0$의 점들에서 영이 아니므로 $K$와 $Q$의 지지는 $U_0$와 만나지
않는다. 대수학의 보조정리 \ref{algebra-lemma-depth-in-ses}를 쓰면
$\text{depth}(K) \geq 2$임을 알 수 있다($As \subset M$은 $M$의
부분가군으로서 $\text{depth} \geq 1$임에 유의하라). 따라서 대수학의
보조정리 \ref{algebra-lemma-bound-depth}에 의해 $K$의 지지가 공집합이
아니면 그 차원은 $\geq 2$이다. 이는 $K = 0$이 아닌 한
$\text{Supp}(M) \cap V(f) \subset \{\mathfrak m\}$와 모순이다.
$K = 0$일 때 $\text{depth}(Q) \geq 2$임을 얻고 앞에서처럼 $Q = 0$을
결론짓는다. 따라서 $A \cong M$이고 $\mathcal{L}$은 자명하다.
\end{proof}







\section{천공 스펙트럼 위의 가역 가군, I}
\label{algebraization-section-local-lefschetz-for-pic}

\noindent
이 절에서는 Picard 군에 대한 몇 가지 국소 Lefschetz 정리를 증명한다.
일부 발상은
\cite{Kollar-pic}, \cite{Bhatt-local}, 그리고 \cite{Kollar-map-pic}에서
가져왔다.

\begin{lemma}
\label{algebraization-lemma-injective-torsion-in-pic}
$(A, \mathfrak m)$을 Noether 국소환이라 하자. $f \in \mathfrak m$은
영인자가 아니고 $\text{depth}(A/fA) \geq 2$, 또는 동치로
$\text{depth}(A) \geq 3$이라고 가정하자. $U$와 $U_0$를 각각 $A$와
$A/fA$의 천공 스펙트럼이라 하자. 사상
$$
\Pic(U) \to \Pic(U_0)
$$
은 꼬임 부분에서 단사이다.
\end{lemma}

\begin{proof}
$\mathcal{L}$을 가역 $\mathcal{O}_U$-가군이라 하자. $\mathcal{L}$이
$\Pic(U_0)$에서 $0$으로 가는 것은 $\mathcal{L}$을 절
\ref{algebraization-section-coherent-triples}과 같은 가역 연접 삼중항
$(\mathcal{L}, \mathcal{L}_0, \lambda)$으로 확장할 수 있는 것과 동치이다.
명제 \ref{algebraization-proposition-hilbert-triple}에 의해 함수
$$
n \longmapsto \chi((\mathcal{L}, \mathcal{L}_0, \lambda)^{\otimes n})
$$
는 다항식이다. 보조정리 \ref{algebraization-lemma-nonnegative-chi-triple}에 의해 이
다항식의 값이 영일 필요충분조건은 $\mathcal{L}^{\otimes n}$이 자명한 것이다.
따라서 $\mathcal{L}$이 꼬임이면 이 다항식은 영점을 무한히 많이 가지므로
항등적으로 영이고, 따라서 $\mathcal{L}$은 자명하다.
\end{proof}

\begin{proposition}[Koll\'ar]
\label{algebraization-proposition-injective-pic}
\begin{reference}
\cite[정리 1.9]{Kollar-map-pic}
\end{reference}
$(A, \mathfrak m)$을 Noether 국소환이라 하고 $f \in \mathfrak m$이라 하자.
다음을 가정한다.
\begin{enumerate}
\item $A$는 쌍대화 복합체를 갖는다.
\item $f$는 영인자가 아니다.
\item $\text{depth}(A/fA) \geq 2$, 또는 동치로 $\text{depth}(A) \geq 3$이다.
\item $f \in \mathfrak p \subset A$가 $\dim(A/\mathfrak p) = 2$인 소
아이디얼이면 $\text{depth}(A_\mathfrak p) \geq 2$이다.
\end{enumerate}
$U$와 $U_0$를 각각 $A$와 $A/fA$의 천공 스펙트럼이라 하자. 사상
$$
\Pic(U) \to \Pic(U_0)
$$
은 단사이다. 마지막으로 (1), (2), (3)이 성립하고 $A$가 $(S_2)$이며
$\dim(A) \geq 4$이면 (4)가 성립한다.
\end{proposition}

\begin{proof}
$\mathcal{L}$을 가역 $\mathcal{O}_U$-가군이라 하자. $\mathcal{L}$이
$\Pic(U_0)$에서 $0$으로 가는 것은 $\mathcal{L}$을 절
\ref{algebraization-section-coherent-triples}과 같은 가역 연접 삼중항
$(\mathcal{L}, \mathcal{L}_0, \lambda)$으로 확장할 수 있는 것과 동치이다.
명제 \ref{algebraization-proposition-hilbert-triple}에 의해 함수
$$
n \longmapsto \chi((\mathcal{L}, \mathcal{L}_0, \lambda)^{\otimes n})
$$
는 다항식 $P$이다. 보조정리 \ref{algebraization-lemma-nonnegative-chi-triple}에 의해 모든
$n \in \mathbf{Z}$에 대해 $P(n) \geq 0$이고, 등호가 성립할 필요충분조건은
$\mathcal{L}^{\otimes n}$이 자명한 것이다. 특히 $P(0) = 0$이고, $P$는
항등적으로 영이어서 결론이 성립하거나 $P$가 차수가 $\geq 2$인 짝수차
다항식이다.

\medskip\noindent
$M = \Gamma(U, \mathcal{L})$ 및
$M_0 = \Gamma(X_0, \mathcal{L}_0) = \Gamma(U_0, \mathcal{L}_0)$로 놓는다.
그러면 $M$은 깊이가 $\geq 2$인 유한 $A$-가군이고 $M_0 \cong A/fA$이다.
보조정리 \ref{algebraization-lemma-nonnegative-chi-triple}의 증명을 보라. 국소 코호몰로지의
보조정리
\ref{local-cohomology-lemma-local-finiteness-for-finite-locally-free}
와 $\text{depth}(A) \geq 3$이므로 $i = 0, 1, 2$에 대해
$H^i_\mathfrak m(A) = 0$이라는 사실에 의해 $H^2_\mathfrak m(M)$은 유한
$A$-가군임에 유의하라. 짧은 완전열
$$
0 \to M/fM \to M_0 \to Q \to 0
$$
을 생각하자. 보조정리 \ref{algebraization-lemma-nonnegative-chi-triple}에 의해 $Q$의 길이는
유한하고 $\chi(\mathcal{L}, \mathcal{L}_0, \lambda)$와 같다. 위 짧은
완전열에 딸린 국소 코호몰로지의 긴 완전열로부터
$Q = H^1_\mathfrak m(M/fM)$ 및 $i > 1$에 대해
$H^i_\mathfrak m(M/fM) = H^i_\mathfrak m(M_0) \cong H^i_\mathfrak m(A/fA)$를
얻는다. 완전열 $0 \to M \to M \to M/fM \to 0$에 딸린 국소
코호몰로지의 긴 완전열을 생각하자. 이는 다음과 같이 시작한다.
$$
0 \to Q \to H^2_\mathfrak m(M) \to H^2_\mathfrak m(M) \to
H^2_\mathfrak m(A/fA)
$$
길이의 가법성을 쓰면 $\chi(\mathcal{L}, \mathcal{L}_0, \lambda)$는 사상
$H^2_\mathfrak m(M) \to H^2_\mathfrak m(A/fA)$의 상의 길이와 같다.

\medskip\noindent
나머지 증명을 밝히기 위해 먼저 특별한 경우에 보조정리를 증명하자. 잠시
$H^2_\mathfrak m(A/fA)$가 유한 길이 가군이라고 가정한다. 그러면
$P(1) \leq \text{length}_A H^2_\mathfrak m(A/fA)$이다.
$\mathcal{L}^{\otimes n}$에 똑같은 논증을 적용하면 모든 $n$에 대해
$P(n) \leq \text{length}_A H^2_\mathfrak m(A/fA)$이다. 따라서 $P$는 양의
차수를 가질 수 없고 결론이 성립한다. 나머지 증명에서는 이 논증을 수정하여
충분한 조건인 $P(n)$의 일차 상계를 얻을 것이다.

\medskip\noindent
사상 $H^2_\mathfrak m(M) \to H^2_\mathfrak m(M_0) \cong
H^2_\mathfrak m(A/fA)$를 살펴보자. $A$의 정규화된 쌍대화 복합체
$\omega_A^\bullet$을 택한다. 국소 쌍대성(쌍대화 복합체의 보조정리
\ref{dualizing-lemma-special-case-local-duality})에 의해 이 사상은 다음
사상의 Matlis 쌍대이다.
$$
\text{Ext}^{-2}_A(M, \omega_A^\bullet)
\longleftarrow
\text{Ext}^{-2}_A(M_0, \omega_A^\bullet)
$$
따라서 두 상의 (유한한) 길이는 같다. 유한 $A$-가군
$\text{Ext}^{-2}_A(M_0, \omega_A^\bullet)$의 지지는 공집합이 아니라면
$\mathfrak m$과 $f$를 포함하고 $\dim(A/\mathfrak p_i) = 1$인 유한 개의
소아이디얼 $\mathfrak p_1, \ldots, \mathfrak p_r$로 이루어진다. 실제로
국소 코호몰로지의 보조정리
\ref{local-cohomology-lemma-sitting-in-degrees}에 의해 이 지지는
$\text{depth}_{A_\mathfrak p}(M_{0, \mathfrak p}) + \dim(A/\mathfrak p) \leq 2$
인 소아이디얼 $\mathfrak p \subset A$들의 집합에 포함된다. 따라서 $f$를
포함하고 $\dim(A/\mathfrak p) = 2$ 및
$\text{depth}_{A_\mathfrak p}(M_{0, \mathfrak p}) = 0$을 만족하는
소아이디얼 $\mathfrak p$가 없음을 보이면 충분하다. 그러나
$M_{0, \mathfrak p} \cong (A/fA)_\mathfrak p$이므로 그러한 소아이디얼은
$\text{depth}(A_\mathfrak p) = 1$을 주어 가정 (4)와 모순이다.
$\mathfrak p_1, \ldots, \mathfrak p_r$에서 영이 되지 않는 단면
$t \in \Gamma(U, \mathcal{L}^{\otimes -1})$를 택한다. 성질의 보조정리
\ref{properties-lemma-quasi-affine-invertible-nonvanishing-section}를 보라.
전역단면에 $t$를 곱하면 사상 $t : M \to A$가 정해지며, 이는
$i = 1, \ldots, r$에 대해 동형 $M_{\mathfrak p_i} \to A_{\mathfrak p_i}$를
정한다. $t$에 호환되는 사상 $t_0 : M_0 \to A/fA$를 마찬가지로 정하는
$\Gamma(U_0, \mathcal{L}_0^{\otimes -1})$의 대응 단면을
$t_0 = t|_{U_0}$라 쓰자. 그러면 다음 가환 그림식이 있다.
$$
\xymatrix{
\text{Ext}^{-2}_A(M, \omega_A^\bullet) &
\text{Ext}^{-2}_A(M_0, \omega_A^\bullet) \ar[l] \\
\text{Ext}^{-2}_A(A, \omega_A^\bullet) \ar[u]^t &
\text{Ext}^{-2}_A(A/fA, \omega_A^\bullet) \ar[l] \ar[u]_{t_0}
}
$$
따라서 위쪽 수평 사상의 상의 길이는
$\text{Ext}^{-2}_A(A/fA, \omega_A^\bullet)$의 길이와 $t_0$의 여핵의 길이의
합 이하이다.

\medskip\noindent
한편 $n > 1$에 대해 $\mathcal{L}$을 $\mathcal{L}^n$으로 바꾸면 $t$ 대신
$$
t^n :
M_n = \Gamma(U, \mathcal{L}^{\otimes n})
\longrightarrow
\Gamma(U, \mathcal{O}_U) = A
$$
을 쓸 수 있다. 이는 $t_0$를 그 $n$제곱으로 바꾼다. 따라서 사상
$\text{Ext}^{-2}_A(M_n, \omega_A^\bullet) \leftarrow
\text{Ext}^{-2}_A(M_{n, 0}, \omega_A^\bullet)$
의 상의 길이는 $\text{Ext}^{-2}_A(A/fA, \omega_A^\bullet)$의 길이와 다음
사상의 여핵의 길이의 합 이하이다.
$$
t_0^n : 
\text{Ext}^{-2}_A(A/fA, \omega_A^\bullet)
\longrightarrow
\text{Ext}^{-2}_A(M_{n, 0}, \omega_A^\bullet)
$$
동형 $M_0 \cong A/fA$를 통해 $t_0$는 어떤 $g \in A/fA$에 대한
$g : A/fA \to A/fA$가 되고, 대응하는 동형 $M_{n, 0} \cong A/fA$를 통해
$t_0^n$은 $g^n : A/fA \to A/fA$가 된다. 따라서 위 여핵의 길이는
$\text{Ext}^{-2}_A(A/fA, \omega_A^\bullet)$를 $g^n$으로 나눈 몫의 길이이다.
$\text{Ext}^{-2}_A(A/fA, \omega_A^\bullet)$는 차원이 $1$인 지지 $T$를 갖는
유한 $A$-가군이고, $t$의 선택에 의해 $V(g) \cap T$는 닫힌점 하나로
이루어지므로 대수학의 보조정리
\ref{algebra-lemma-support-dimension-d}에 의해 이 길이는 $n$에 대해
일차적으로 증가한다.

\medskip\noindent
증명을 마치기 위해 마지막 진술을 증명한다. $f \in \mathfrak m \subset A$가
(1), (2), (3)을 만족하고 $A$가 $(S_2)$이며 $\dim(A) \geq 4$라고 가정하자.
조건 (1)에 의해 $A$는 사슬적이다. 쌍대화 복합체의 보조정리
\ref{dualizing-lemma-universally-catenary}를 보라. 그러면 국소 코호몰로지의
보조정리 \ref{local-cohomology-lemma-catenary-S2-equidimensional}에 의해
$\Spec(A)$는 등차원이다. 따라서 $A$의 모든 소아이디얼 $\mathfrak p$에 대해
$\dim(A_\mathfrak p) + \dim(A/\mathfrak p) \geq 4$이다. 그러므로
$\text{depth}(A_\mathfrak p) \geq \min(2, \dim(A_\mathfrak p))
\geq \min(2, 4 - \dim(A/\mathfrak p))$이고 따라서 (4)가 성립한다.
\end{proof}

\begin{remark}
\label{algebraization-remark-compare-SGA2}
SGA2에는 다음 결과가 있다. $(A, \mathfrak m)$을 Noether 국소환이라 하고
$f \in \mathfrak m$이라 하자. $A$가 정칙환의 몫이고 $f$가 영인자가
아니며 다음을 가정한다.
\begin{enumerate}
\item[(a)] $\mathfrak p \subset A$가 $\dim(A/\mathfrak p) = 1$인
소아이디얼이면 $\text{depth}(A_\mathfrak p) \geq 2$이다.
\item[(b)] $\text{depth}(A/fA) \geq 3$, 또는 동치로
$\text{depth}(A) \geq 4$이다.
\end{enumerate}
$U$와 $U_0$를 각각 $A$와 $A/fA$의 천공 스펙트럼이라 하자. 그러면 사상
$$
\Pic(U) \to \Pic(U_0)
$$
은 단사이다. 이는 \cite[Expos\'e XI, 보조정리 3.16]{SGA2}이다\footnote{조건
(a)는 조건 (b)에서 따른다. 대수학의 보조정리
\ref{algebra-lemma-depth-localization}를 보라.}. SGA2의 이 결과는 명제
\ref{algebraization-proposition-injective-pic}에서 따른다. 실제로
\begin{enumerate}
\item 정칙환의 몫은 쌍대화 복합체를 갖는다(쌍대화 복합체의 보조정리
\ref{dualizing-lemma-regular-gorenstein} 및 명제
\ref{dualizing-proposition-dualizing-essentially-finite-type}를 보라).
\item $\text{depth}(A) \geq 4$이면 $\dim(A/\mathfrak p) = 2$인 모든
소아이디얼 $\mathfrak p$에 대해 $\text{depth}(A_\mathfrak p) \geq 2$이다.
대수학의 보조정리 \ref{algebra-lemma-depth-localization}를 보라.
\end{enumerate}
\end{remark}






\section{천공 스펙트럼 위의 가역 가군, II}
\label{algebraization-section-local-lefschetz-for-pic-surjective}

\noindent
이제 Picard 군에 대한 국소 Lefschetz 정리의 전사성으로 넘어간다. 먼저
$U_0$ 위의 가역 가군을 열린 근방으로 확장하는 다음 간단한 판정법이 있다.

\begin{lemma}
\label{algebraization-lemma-surjective-Pic-first}
$(A, \mathfrak m)$을 Noether 국소환이라 하고 $f \in \mathfrak m$이라 하자.
다음을 가정한다.
\begin{enumerate}
\item $A$는 $f$-진 완비이다.
\item $f$는 영인자가 아니다.
\item $H^1_\mathfrak m(A/fA)$와 $H^2_\mathfrak m(A/fA)$는 유한
$A$-가군이다.
\item $H^3_\mathfrak m(A/fA) = 0$이다\footnote{$\text{depth}(A/fA) \geq 4$,
또는 동치로 $\text{depth}(A) \geq 5$이면 (3)과 (4)가 성립함에 유의하라.}.
\end{enumerate}
$U$와 $U_0$를 각각 $A$와 $A/fA$의 천공 스펙트럼이라 하자. 그러면
$$
\colim_{U_0 \subset U' \subset U\text{ 열린 부분스킴}} \Pic(U')
\longrightarrow
\Pic(U_0)
$$
은 전사이다.
\end{lemma}

\begin{proof}
$U_0 \subset U_n \subset U$를 $U_0$의 $n$차 무한소 근방이라 하자.
$U_0 \subset U$는 영인자가 아닌 $f$로 잘라낸 주 닫힌 부분스킴이므로
$U_{n + 1}$ 안에서 $U_n$의 아이디얼층은 $\mathcal{O}_{U_0}$와 동형이다.
따라서 Abel 군의 완전열
$$
\Pic(U_{n + 1}) \to \Pic(U_n) \to
H^2(U_0, \mathcal{O}_{U_0}) = H^3_\mathfrak m(A/fA) = 0
$$
을 얻는다. 사상에 관하여 더 보기의 보조정리
\ref{more-morphisms-lemma-picard-group-first-order-thickening}를 보라. 따라서
모든 가역 $\mathcal{O}_{U_0}$-가군은 가역 연접 형식 가군, 즉
$\textit{Coh}(U, f\mathcal{O}_U)$의 가역 대상의 제한이다. 보조정리
\ref{algebraization-lemma-equivalence}를 적용하면 결론을 얻는다.
\end{proof}

\begin{remark}
\label{algebraization-remark-surjective-Pic-second}
$(A, \mathfrak m)$을 Noether 국소환이라 하고 $f \in \mathfrak m$이라 하자.
다음을 가정하면 보조정리 \ref{algebraization-lemma-surjective-Pic-first}의 결론이 성립한다.
\begin{enumerate}
\item $A$는 쌍대화 복합체를 갖는다.
\item $A$는 $f$-진 완비이다.
\item $f$는 영인자가 아니다.
\item 다음 중 하나가 성립한다.
\begin{enumerate}
\item $A_f$는 $(S_2)$이고 $f \not \in \mathfrak p$인 극소
$\mathfrak p \subset A$에 대해 $\dim(A/\mathfrak p) \geq 4$이다. 또는
\item $\mathfrak p \not \in V(f)$이고
$V(\mathfrak p) \cap V(f) \not = \{\mathfrak m\}$이면
$\text{depth}(A_\mathfrak p) + \dim(A/\mathfrak p) > 3$.
\end{enumerate}
\item $H^3_{\mathfrak m}(A/fA) = 0$이다.
\end{enumerate}
증명은 보조정리 \ref{algebraization-lemma-surjective-Pic-first}의 증명과 정확히 같되,
보조정리 \ref{algebraization-lemma-equivalence-better}를 보조정리
\ref{algebraization-lemma-equivalence} 대신 쓴다. 여기서 두 가지를
언급해야 한다. (a) $\text{depth}(A/fA) \geq 4$를 가정하지 않고
$H^3_{\mathfrak m}(A/fA) = 0$임을 아는 예를 찾기는 어려워 보인다. (b)
보조정리 \ref{algebraization-lemma-equivalence-better}의 증명은 보조정리
\ref{algebraization-lemma-equivalence}의 증명보다 훨씬 어렵다.
\end{remark}

\begin{lemma}
\label{algebraization-lemma-surjective-Pic-first-better}
$(A, \mathfrak m)$을 Noether 국소환이라 하고 $f \in \mathfrak m$이라 하자.
다음을 가정한다.
\begin{enumerate}
\item 보조정리 \ref{algebraization-lemma-surjective-Pic-first}의 조건들이 성립한다.
\item 모든 극대 아이디얼 $\mathfrak p \subset A_f$에 대해
$(A_f)_\mathfrak p$의 천공 스펙트럼의 Picard 군은 자명하다.
\end{enumerate}
$U$와 $U_0$를 각각 $A$와 $A/fA$의 천공 스펙트럼이라 하자. 그러면
$$
\Pic(U) \longrightarrow \Pic(U_0)
$$
은 전사이다.
\end{lemma}

\begin{proof}
$\mathcal{L}_0 \in \Pic(U_0)$라 하자. 보조정리
\ref{algebraization-lemma-surjective-Pic-first}에 의해 그 $U_0$로의 제한이
$\mathcal{L}_0$인 열린집합 $U_0 \subset U' \subset U$와
$\mathcal{L}' \in \Pic(U')$가 존재한다. $U' \supset U_0$이므로
$U \setminus U'$은 (2)와 같은 소아이디얼
$\mathfrak p_1, \ldots, \mathfrak p_n$에 대응하는 점들로 이루어진다.
자명한 가역 가군은 언제나 확장되므로, 가정에 의해 천공 스펙트럼
$U' \times_U \Spec(A_{\mathfrak p_i})$ 위에서 $\mathcal{L}'$과 일치하는
$\Spec(A_{\mathfrak p_i})$ 위의 가역 가군 $\mathcal{L}'_i$를 찾을 수 있다.
극한의 보조정리 \ref{limits-lemma-glueing-near-closed-point-modules}를 $n$번
적용하면 $\mathcal{L}'$이 $U$ 위의 가역 가군으로 확장됨을 알 수 있다.
\end{proof}

\begin{lemma}
\label{algebraization-lemma-local-pic-to-completion}
$(A, \mathfrak m)$을 깊이가 $\geq 2$인 Noether 국소환이라 하고
$A^\wedge$를 그 완비라 하자. $U$와 $U^\wedge$를 각각 $A$와
$A^\wedge$의 천공 스펙트럼이라 하자. 그러면
$\Pic(U) \to \Pic(U^\wedge)$는 단사이다.
\end{lemma}

\begin{proof}
$\mathcal{L}$을 $U^\wedge$ 위의 당김이 $\mathcal{L}^\wedge$인 가역
$\mathcal{O}_U$-가군이라 하자. 깊이에 대한 가정과 쌍대화 복합체의 보조정리
\ref{dualizing-lemma-depth} 및 국소 코호몰로지의 보조정리
\ref{local-cohomology-lemma-finiteness-pushforwards-and-H1-local}에 의해
$H^0(U, \mathcal{O}_U) = A$이다. 따라서 $\mathcal{L}$이 자명할
필요충분조건은 $M = H^0(U, \mathcal{L})$이 $A$-가군으로서 $A$와 동형인
것이다(세부사항은 생략한다). $A \to A^\wedge$가 평탄하므로 평탄 밑변환에
의해 $M \otimes_A A^\wedge = \Gamma(U^\wedge, \mathcal{L}^\wedge)$이다.
스킴의 코호몰로지의 보조정리
\ref{coherent-lemma-flat-base-change-cohomology}를 보라. 마지막으로
$M \cong A$일 필요충분조건이
$M \otimes_A A^\wedge \cong A^\wedge$임은 쉽게 알 수 있다.
\end{proof}

\begin{lemma}
\label{algebraization-lemma-trivial-local-pic-regular}
$(A, \mathfrak m)$을 정칙 국소환이라 하자. 그러면 $A$의 천공 스펙트럼의
Picard 군은 자명하다.
\end{lemma}

\begin{proof}
인자의 보조정리 \ref{divisors-lemma-extend-invertible-module}와 대수학에
관하여 더 보기의 보조정리 \ref{more-algebra-lemma-regular-local-UFD}를 결합하라.
\end{proof}

\noindent
이제 앞의 결과들을 반복 적용하여 차원이 $\geq 4$인 완전교차의 천공
스펙트럼의 Picard 군이 자명함을 증명할 수 있다. Noether 국소환의 완비가
정칙 국소환을 정칙열이 생성하는 아이디얼로 나눈 몫이면 그 환을 완전교차라
부름을 상기하라. 분할멱 대수의 절 \ref{dpa-section-lci}의 논의를 보라.

\begin{proposition}[Grothendieck]
\label{algebraization-proposition-trivial-local-pic-complete-intersection}
$(A, \mathfrak m)$을 Noether 국소환이라 하자. $A$가 차원이 $\geq 4$인
완전교차이면 $A$의 천공 스펙트럼의 Picard 군은 자명하다.
\end{proposition}

\begin{proof}
보조정리 \ref{algebraization-lemma-local-pic-to-completion}에 의해 $A$가 완비 국소환이라고
가정해도 된다. 가정에 의해 $A = B/(f_1, \ldots, f_r)$로 쓸 수 있는데,
여기서 $B$는 완비 정칙 국소환이고 $f_1, \ldots, f_r$은 정칙열이다.
$r$에 관한 귀납법으로 증명을 마친다. 기초 단계 $r = 0$은 보조정리
\ref{algebraization-lemma-trivial-local-pic-regular}에서 따른다.

\medskip\noindent
$A = B/(f_1, \ldots, f_r)$이고 명제가 $r - 1$에 대해 성립한다고 가정하자.
$A' = B/(f_1, \ldots, f_{r - 1})$로 놓고 보조정리
\ref{algebraization-lemma-surjective-Pic-first-better}를 $f_r \in A'$에 적용한다. 이는
허용된다.
\begin{enumerate}
\item 우리 국소환들이 완비이므로 보조정리
\ref{algebraization-lemma-surjective-Pic-first}의 조건 (1)이 성립한다.
\item $f_1, \ldots, f_r$이 정칙열이므로 보조정리
\ref{algebraization-lemma-surjective-Pic-first}의 조건 (2)가 성립한다.
\item $A = A'/f_r A'$가 차원이 $\dim(A) \geq 4$인 Cohen--Macaulay
환이므로 보조정리 \ref{algebraization-lemma-surjective-Pic-first}의 조건 (3), (4)가 성립한다.
\item $A'_{f_r}$의 극대 소아이디얼 $\mathfrak p$에 대해
$\dim((A'_{f_r})_\mathfrak p) \geq 4$이고
$(A'_{f_r})_\mathfrak p = B_\mathfrak q/(f_1, \ldots, f_{r - 1})$
가 되는 어떤 소아이디얼 $\mathfrak q \subset B$에 대해 $B_\mathfrak q$가
정칙이므로, 귀납가정에 의해 보조정리
\ref{algebraization-lemma-surjective-Pic-first-better}의 조건 (2)가 성립한다.
\end{enumerate}
이로써 증명이 끝난다.
\end{proof}

\begin{example}
\label{algebraization-example-grothendieck-sharp}
명제 \ref{algebraization-proposition-trivial-local-pic-complete-intersection}의 차원 상계는
최적이다. 예를 들어 $A = k[[x, y, z, w]]/(xy - zw)$의 천공 스펙트럼의
Picard 군은 자명하지 않다. 실제로 $I_x, I_y, I_z, I_w$가 각각
$A_x, A_y, A_z, A_w$의 가역 아이디얼임은 쉽게 알 수 있으므로, 아이디얼
$I = (x, z)$는 $A$의 천공 스펙트럼 $U$ 위에서 유효 Cartier 인자 $D$를
잘라낸다. 한편 $A/I$의 깊이는 $\geq 1$(실제로는 $2$)이고, 따라서 $I$의
깊이는 $\geq 2$(실제로는 $3$)이므로
$I = \Gamma(U, \mathcal{O}_U(-D))$이다. 그러므로
$\mathcal{O}_U(-D)$가 자명하다면 $I \cong \Gamma(U, \mathcal{O}_U) = A$일
것인데, $I$는 $1$개의 원소로 생성되지 않으므로 이는 참이 아니다.
\end{example}

\begin{example}
\label{algebraization-example-grothendieck-sharp-bis}
명제 \ref{algebraization-proposition-trivial-local-pic-complete-intersection}은 다음과 같은
몫으로 확장할 수 없다.
$$
A = B/(f_1, \ldots, f_r)
$$
여기서 $B$는 정칙이고 $\dim(B) - r \geq 4$이다. 달리 말해 $A$의 천공
스펙트럼의 Picard 군이 소멸하려면 $f_1, \ldots, f_r$이 정칙열이라는 조건이
(일반적으로) 필요하다. 실제로 $k$를 체라 하고 다음과 같이 놓자.
$$
A = k[[a, b, x, y, z, u, v, w]]/(a^3, b^3, xa^2 + yab + zb^2, w^2)
$$
$A = A_0[w]/(w^2)$이며
$A_0 = k[[a, b, x, y, z, u, v]]/(a^3, b^3, xa^2 + yab + zb^2)$임에 유의하라.
아래에서 $A_0$의 깊이가 $2$임을 보인다. $U$를 $A$의 천공 스펙트럼,
$U_0$를 $A_0$의 천공 스펙트럼이라 하자. 첫째 화살표가 $w$를 곱하는
사상인 짧은 완전열 $0 \to A_0 \to A \to A_0 \to 0$이 있다. 사상에
관하여 더 보기의 보조정리
\ref{more-morphisms-lemma-picard-group-first-order-thickening}에 의해 완전열
$$
H^0(U, \mathcal{O}_U^*) \to
H^0(U_0, \mathcal{O}_{U_0}^*) \to
H^1(U_0, \mathcal{O}_{U_0}) \to
\Pic(U)
$$
을 얻는다. $A_0$, 따라서 $A$의 깊이가 $2$이므로
$H^0(U_0, \mathcal{O}_{U_0}) = A_0$이고 $H^0(U, \mathcal{O}_U) = A$이며
$H^1(U_0, \mathcal{O}_{U_0})$는 영이 아니다. 쌍대화 복합체의 보조정리
\ref{dualizing-lemma-depth}와 국소 코호몰로지의 보조정리
\ref{local-cohomology-lemma-local-cohomology}를 보라. 따라서 위에 표시한
마지막 화살표는 영이 아니고, $\Pic(U)$는 영이 아니다.

\medskip\noindent
$A_0$의 깊이가 $2$임을 보이려면
$A_1 = k[[a, b, x, y, z]]/(a^3, b^3, xa^2 + yab + zb^2)$의 깊이가
$0$임을 보이면 충분하다. 실제로 $a^2b^2$는 $A_1$의 영이 아닌 원소로
가며 각 변수 $a, b, x, y, z$에 의해 소멸된다. 예를 들어 $A_1$에서
$ya^2b^2 = (yab)(ab) = - (xa^2 + zb^2)(ab) = -xa^3b - yab^3 = 0$이다.
다른 경우도 비슷하다.
\end{example}










\section{Lefschetz 정리에의 응용}
\label{algebraization-section-lefschetz}

\noindent
이 절에서는 사영 스킴 $P$ 위의 연접층과 풍부한 인자 $Q$를 따른 형식적
완비 위의 연접 가군 사이의 관계를 논한다.

\medskip\noindent
$k$를 체라 하고 $P$를 $k$ 위의 고유 스킴이라 하자. $\mathcal{L}$을 풍부한
가역 $\mathcal{O}_P$-가군이라 하자. $s \in \Gamma(P, \mathcal{L})$를
단면이라 하고\footnote{$s$가 정칙 단면이라고 요구하지 않는다. 이에 따라
$Q$는 $P$의 국소 주 닫힌 부분스킴일 뿐이며 반드시 유효 Cartier 인자는
아니다.} $Q = Z(s)$를 영점 스킴이라 하자. 인자의 정의
\ref{divisors-definition-zero-scheme-s}를 보라. 모든 $n \geq 1$에 대해
$Q_n = Z(s^n)$으로 $Q$의 $n$차 무한소 근방을 나타낸다. $\mathcal{F}$가
연접 $\mathcal{O}_P$-가군이면 $\mathcal{F}_n = \mathcal{F}|_{Q_n}$으로
그 제한, 즉 닫힌 몰입 $Q_n \to P$에 의한 $\mathcal{F}$의 당김을 나타낸다.

\begin{proposition}
\label{algebraization-proposition-lefschetz}
위 상황에서 모든 점 $p \in P \setminus Q$에 대해
$$
\text{depth}(\mathcal{F}_p) + \dim(\overline{\{p\}}) > \sigma
$$
가 되는 정수 $\sigma$가 존재한다고 가정하자. 그러면 사상
$$
H^i(P, \mathcal{F}) \longrightarrow \lim H^i(Q_n, \mathcal{F}_n)
$$
은 $0 \leq i < \sigma$에 대해 동형이다.
\end{proposition}

\begin{proof}
사상에 관하여 더 보기의 보조정리
\ref{more-morphisms-lemma-apply-proj-spec}와 같은 글의 절
\ref{more-morphisms-section-proj-spec}의 표기와 결과를 더 언급하지 않고
사용한다. 적어도 그 보조정리의 진술을 이해하지 않으면 이 증명은 의미가
없다. 우리 경우
$A = \bigoplus_{m \geq 0} \Gamma(P, \mathcal{L}^{\otimes m})$
는 각 등급 성분이 유한 차원 $k$-벡터공간인 유한형 $k$-대수이다. 스킴의
코호몰로지의 보조정리 \ref{coherent-lemma-coherent-proper-ample}를 보라.

\medskip\noindent
우리는 $s$를 $f \in A_1 \subset A$, 즉 $A$의 차수 $1$인 동차 원소로
간주할 수 있고 그렇게 한다. $f$가 정의하는 닫힌 부분스킴을
$Y = V(f) \subset X$라 쓰자. 그러면 스킴론적으로
$U \cap Y = (\pi|_U)^{-1}(Q)$이다. 표기
$\mathcal{F}_U = \pi^*\mathcal{F}|_U = (\pi|_U)^*\mathcal{F}$를 상기하라.
이는 연접 $\mathcal{O}_U$-가군이다. 다음을 만족하는 유한 $A$-가군 $M$을
택한다.
$\mathcal{F}_U = \widetilde{M}|_U$
(존재성은 국소 코호몰로지의 보조정리
\ref{local-cohomology-lemma-finiteness-pushforwards-and-H1-local}를 보라.)
$i \leq \sigma + 1$에 대해 $H^i_Z(M)$이 $f$의 어떤 멱에 의해 소멸된다고
주장한다.

\medskip\noindent
이 주장을 증명하기 위해 국소 코호몰로지의 명제
\ref{local-cohomology-proposition-annihilator}를 적용한다. 기하학적 언어로
옮기면 $u \in U$, $u \not \in Y$ 및 $z \in \overline{\{u\}} \cap Z$에 대해
$$
\text{depth}(\mathcal{F}_{U, u}) +
\dim(\mathcal{O}_{\overline{\{u\}}, z}) > \sigma + 1
$$
임을 보이면 충분하다. 이는 조금만 생각하면 된다.

\medskip\noindent
$A_0$는 유한 차원 $k$-대수이므로 $Z = \Spec(A_0)$는 $X$의 유한한
닫힌점들의 집합이다. ($Z$가 한 점이기를 원하는 독자는 유한 $k$-대수
$A_0$를 $k$로 바꾸어도 된다. 증명의 다른 부분에는 아무 영향이 없다.)

\medskip\noindent
$\pi : L \to P$와 그 제한 $\pi|_U : U \to P$는 상대 차원 $1$인 매끄러운
사상이다. $u \in U$, $u \not \in Y$ 및
$z \in \overline{\{u\}} \cap Z$라 하자. 그 상을
$p = \pi(u) \in P \setminus Q$라 하자. 그러면 $u$는 $p$ 위의 $\pi$의
올의 일반점이거나 그 올의 닫힌점이다. $u$가 올의 일반점이면
$\text{depth}(\mathcal{F}_{U, u}) = \text{depth}(\mathcal{F}_p)$
이고
$\dim(\overline{\{u\}}) = \dim(\overline{\{p\}}) + 1$.
$u$가 올의 닫힌점이면
$\text{depth}(\mathcal{F}_{U, u}) = \text{depth}(\mathcal{F}_p) + 1$
이고
$\dim(\overline{\{u\}}) = \dim(\overline{\{p\}})$.
두 경우 모두 $Z$의 모든 점이 닫혀 있으므로
$\dim(\overline{\{u\}}) = \dim(\mathcal{O}_{\overline{\{u\}}, z})$
이다. 따라서 원하는 부등식은 보조정리의 진술에 있는 가정에서 따른다.

\medskip\noindent
$A'$을 $A$의 $f$-진 완비라 하자. 대수학의 보조정리
\ref{algebra-lemma-completion-flat}에 의해 $A \to A'$는 평탄하다.
$U$의 역상을 $U' \subset X' = \Spec(A')$라 하고, $Y'$와 $Z'$도 마찬가지로
정의한다. $\mathcal{F}'$을 $U'$ 위에서 $\mathcal{F}_U$의 당김이라 하고
$M' = M \otimes_A A'$로 놓는다. 국소 코호몰로지의 평탄 밑변환(국소
코호몰로지의 보조정리 \ref{local-cohomology-lemma-torsion-change-rings})에 의해
$$
H^i_{Z'}(M') = H^i_Z(M) \otimes_A A'
$$
이고, $i \leq \sigma + 1$에 대해 이들은 $f$의 어떤 멱에 의해 소멸됨을
얻는다. 그림식
$$
\xymatrix{
& H^i(U, \mathcal{F}_U) \ar[ld] \ar[d] \ar[r] &
\lim H^i(U, \mathcal{F}_U/f^n\mathcal{F}_U) \ar@{=}[d] \\
H^i(U, \mathcal{F}_U) \otimes_A A' \ar@{=}[r] &
H^i(U', \mathcal{F}') \ar[r] &
\lim H^i(U', \mathcal{F}'/f^n\mathcal{F}')
}
$$
을 생각하자. 보조정리 \ref{algebraization-lemma-alternative-higher}와 방금 증명한 꼬임
성질에 의해 아래쪽 수평 화살표는 $i < \sigma$에 대해 동형이다. 수평
등호는 평탄 밑변환(스킴의 코호몰로지의 보조정리
\ref{coherent-lemma-flat-base-change-cohomology})이고, 수직 등호는
$U \cap Y$와 $U' \cap Y'$ 및 그 각각의 $n$차 무한소 근방이 서로
동형으로 사상되기 때문이다(우리가 $f$에 관해 완비화하고 있으므로).

\medskip\noindent
사상에 관하여 더 보기의 식
(\ref{more-morphisms-equation-cohomology-torsor})을 적용하면 호환되는 직합 분해
$$
\lim H^i(U, \mathcal{F}_U/f^n\mathcal{F}_U) =
\lim
\left(
\bigoplus\nolimits_{m \in \mathbf{Z}}
H^i(Q_n, \mathcal{F}_n \otimes \mathcal{L}^{\otimes m})
\right)
$$
및
$$
H^i(U, \mathcal{F}_U) =
\bigoplus\nolimits_{m \in \mathbf{Z}}
H^i(P, \mathcal{F} \otimes \mathcal{L}^{\otimes m})
$$
를 얻는다. 따라서 대수학의 보조정리 \ref{algebra-lemma-daniel-litt}에 의해
결론을 얻는다.
\end{proof}

\begin{lemma}
\label{algebraization-lemma-lefschetz-addendum}
$k$를 체라 하고 $X$를 $k$ 위의 고유 스킴이라 하자. $\mathcal{L}$을 풍부한
가역 $\mathcal{O}_X$-가군이라 하고 $s \in \Gamma(X, \mathcal{L})$라 하자.
$Y = Z(s)$를 $s$의 영점 스킴이라 하고 그 $n$차 무한소 근방을
$Y_n = Z(s^n)$라 하자. $\mathcal{F}$를 연접 $\mathcal{O}_X$-가군이라 하자.
모든 $x \in X \setminus Y$에 대해
$$
\text{depth}(\mathcal{F}_x) + \dim(\overline{\{x\}}) > 1
$$
라고 가정하자. 그러면 $Y$를 포함하는 임의의 열린 부분스킴 $V \subset X$에
대해 $\Gamma(V, \mathcal{F}) \to \lim \Gamma(Y_n, \mathcal{F}|_{Y_n})$는
동형이다.
\end{lemma}

\begin{proof}
명제 \ref{algebraization-proposition-lefschetz}에 의해 이는 $V = X$일 때 참이다. 따라서
사상 $\Gamma(V, \mathcal{F}) \to \lim
\Gamma(Y_n, \mathcal{F}|_{Y_n})$가 단사임을 보이면 충분하다.
$\sigma \in \Gamma(V, \mathcal{F})$가 영으로 가면 그 지지는 $Y$와 서로소이다
(세부사항은 생략한다. 힌트: Krull 교차 정리를 사용하라). 그러면
$\text{Supp}(\sigma)$의 폐포 $T \subset X$는 $Y$와 서로소이다. 따라서
$T$는 $X$에서 닫혀 있으므로 $k$ 위에서 고유이고, 아핀 스킴
$X \setminus Y$에서 닫혀 있으므로 아핀이다(사상의 보조정리
\ref{morphisms-lemma-proper-ample-delete-affine}를 보라). 따라서 이는 $k$
위에서 유한이다(사상의 보조정리 \ref{morphisms-lemma-finite-proper}). 그러므로
$T$는 $X$의 유한한 닫힌점들의 집합이다. 가정에 의해 $x \in T$에 대해
$\text{depth}(\mathcal{F}_x) \geq 2$는 적어도 $1$이다. 따라서
$\Gamma(V, \mathcal{F}) \to \Gamma(V \setminus T, \mathcal{F})$는 단사이고,
원하는 대로 $\sigma = 0$이다.
\end{proof}


\begin{example}
\label{algebraization-example-lefschetz}
$k$를 체라 하고 $X$를 $k$ 위의 고유 다양체라 하자. $Y \subset X$는
$\mathcal{O}_X(Y)$가 풍부한 유효 Cartier 인자라 하고 그 $n$차 무한소
근방을 $Y_n$이라 쓰자. $\mathcal{E}$를 유한 국소 자유
$\mathcal{O}_X$-가군이라 하자. 다음은 명제 \ref{algebraization-proposition-lefschetz}의
몇 가지 특별한 경우이다.
\begin{enumerate}
\item $X$가 곡선이면 아무것도 얻지 못한다.
\item $X$가 Cohen--Macaulay(예를 들어 정규) 곡면이면
$$
H^0(X, \mathcal{E}) \to \lim H^0(Y_n, \mathcal{E}|_{Y_n})
$$
은 동형이다.
\item $X$가 Cohen--Macaulay 삼차원체이면 다음 두 사상
$$
H^0(X, \mathcal{E}) \to \lim H^0(Y_n, \mathcal{E}|_{Y_n})
\quad\text{이고}\quad
H^1(X, \mathcal{E}) \to \lim H^1(Y_n, \mathcal{E}|_{Y_n})
$$
은 동형들이다.
\end{enumerate}
그 양상은 분명할 것이다. $X$가 정규 3차원 다양체이면 $H^0$에 대해서는
결론을 얻지만 $H^1$에 대해서는 그렇지 않다.
\end{example}

\noindent
다음 주요 결과를 증명하기 전에 보조정리가 하나 필요하다.

\begin{lemma}
\label{algebraization-lemma-Gm-equivariant-extend-canonically}
상황 \ref{algebraization-situation-algebraize}에서 $(\mathcal{F}_n)$을
$\textit{Coh}(U, I\mathcal{O}_U)$의 대상이라 하자. 다음을 가정한다.
\begin{enumerate}
\item $A$는 등급환이고 $\mathfrak a = A_+$이며 $I$는 동차 아이디얼이다.
\item $(\mathcal{F}_n) = (\widetilde{M_n}|_U)$이며, 여기서 $(M_n)$은 등급
$A$-가군의 역계이다.
\item $(\mathcal{F}_n)$은 $X$로 표준적으로 확장된다.
\end{enumerate}
그러면 다음을 만족하는 유한 등급 $A$-가군 $N$이 존재한다.
\begin{enumerate}
\item[(a)] 역계 $(N/I^nN)$과 $(M_n)$은 $A_+$-멱 꼬임 가군을 법으로 한
등급 $A$-가군의 범주에서 프로-동형이다.
\item[(b)] $(\mathcal{F}_n)$은 $N$에 대응하는 연접 가군의 완비이다.
\end{enumerate}
\end{lemma}

\begin{proof}
$(\mathcal{G}_n)$을 보조정리 \ref{algebraization-lemma-canonically-algebraizable}과 같은
표준 확장이라 하자. $A$와 $M_n$의 등급은 군 스킴 $\mathbf{G}_m$의 $X$
위의 작용
$$
a : \mathbf{G}_m \times X \longrightarrow X
$$
을 정하며, 이에 의해 $(\widetilde{M_n})$은 $\mathbf{G}_m$-등변 준연접
$\mathcal{O}_X$-가군의 역계가 된다. 군체의 예
\ref{groupoids-example-Gm-on-affine}를 보라. $\mathfrak a$와 $I$는 동차
아이디얼이므로 닫힌 부분스킴 $Z$, $Y$와 열린 부분스킴 $U$는
$\mathbf{G}_m$-불변인 닫힌 또는 열린 부분스킴이다. $(\widetilde{M_n})$의
제한 $(\mathcal{F}_n)$은 $\mathbf{G}_m$-등변 연접 $\mathcal{O}_U$-가군의
역계이다. 달리 말해 $(\mathcal{F}_n)$은 $\mathbf{G}_m$-등변 연접 형식
가군이다. 즉 적절한 쌍대순환 조건을 만족하는 동형
$$
\alpha : (a^*\mathcal{F}_n) \longrightarrow (p^*\mathcal{F}_n)
$$
이 $\mathbf{G}_m \times U$ 위에 존재한다. $a$와 $p$는 아핀 스킴의 평탄
사상이므로 보조정리 \ref{algebraization-lemma-canonically-extend-base-change}에 의해
$\mathbf{G}_m \times X$ 위에 유일한 동형
$$
\beta : (a^*\mathcal{G}_n) \longrightarrow (p^*\mathcal{G}_n)
$$
이 존재하며, 그 $\mathbf{G}_m \times U$로의 제한은 $\alpha$이다. 유일성에
의해 $\beta$는 대응하는 쌍대순환 조건을 만족한다. 이와 같이 각
$\mathcal{G}_n$은 전이사상들과 호환되는 방식으로 $\mathbf{G}_m$-등변 연접
$\mathcal{O}_X$-가군이 된다.

\medskip\noindent
군체의 보조정리 \ref{groupoids-lemma-Gm-equivariant-module}에 의해
$\mathbf{G}_m$-등변 구조를 갖춘 $\mathcal{G}_n$은 등급 $A$-가군 $N_n$에
대응한다. 전이사상 $N_{n + 1} \to N_n$은 등급 가군 사상이다.
$(\mathcal{G}_n)$은 $\textit{Coh}(X, I\mathcal{O}_X)$의 대상이므로 $N_n$은
유한 $A$-가군이고 $N_n = N_{n + 1}/I^n N_{n + 1}$임에 유의하라.
$N$을 대수학의 보조정리 \ref{algebra-lemma-finiteness-graded}에서 얻은 유한
등급 $A$-가군이라 하자. 그러면 $N_n = N/I^nN$이고, 따라서
$(\mathcal{G}_n)$은 $N$에 대응하는 연접 가군의 완비이다. 특히 (b)가 참이다.

\medskip\noindent
(a)를 보이려면 위에서 설명한 상황을 조금 더 풀어 써야 한다. 먼저
$M_n \to H^0(U, \mathcal{F}_n)$의 핵과 여핵은 $A_+$-멱 꼬임이다(국소
코호몰로지의 보조정리
\ref{local-cohomology-lemma-finiteness-pushforwards-and-H1-local}).
$H^0(U, \mathcal{F}_n)$에는 이 사상들과 계의 전이사상들이 등급
$A$-가군 사상이 되게 하는 자연스러운 등급이 있다. 예를 들어
$(U \to X)_*\mathcal{F}_n$이 $X$ 위의 $\mathbf{G}_m$-등변 가군임과 군체의
보조정리 \ref{groupoids-lemma-Gm-equivariant-module}를 쓸 수 있다. 다음으로
정의 \ref{algebraization-definition-canonically-algebraizable}와 보조정리
\ref{algebraization-lemma-canonically-algebraizable}에 의해 $(N_n)$과
$(H^0(U, \mathcal{F}_n))$이 프로-동형임을 상기하라. 이 프로-동형을 정의하는
사상들이 등급 가군 사상임을 확인하는 일은 생략한다. 따라서 $(N_n)$과
$(M_n)$은 $A_+$-멱 꼬임 가군을 법으로 한 등급 $A$-가군의 범주에서
프로-동형이다.
\end{proof}

\noindent
$k$를 체라 하고 $P$를 $k$ 위의 고유 스킴이라 하자. $\mathcal{L}$을 풍부한
가역 $\mathcal{O}_P$-가군이라 하자. $s \in \Gamma(P, \mathcal{L})$를
단면이라 하고 $Q = Z(s)$를 영점 스킴이라 하자. 인자의 정의
\ref{divisors-definition-zero-scheme-s}를 보라. $\mathcal{I} \subset
\mathcal{O}_P$를 $Q$의 아이디얼층이라 하자. 스킴의 코호몰로지의 절
\ref{coherent-section-coherent-formal}에서 도입한 연접 형식 가군의 범주를
$\textit{Coh}(P, \mathcal{I})$로 나타낸다.

\begin{proposition}
\label{algebraization-proposition-lefschetz-existence}
위 상황에서 $(\mathcal{F}_n)$을 $\textit{Coh}(P, \mathcal{I})$의 대상이라
하자. 모든 $q \in Q$와 $\mathfrak p \not \in V(\mathcal{I}_q^\wedge)$인
모든 소아이디얼 $\mathfrak p \in \mathcal{O}_{P, q}^\wedge$에 대해
$$
\text{depth}((\mathcal{F}_q^\wedge)_\mathfrak p) +
\dim(\mathcal{O}_{P, q}^\wedge/\mathfrak p) +
\dim(\overline{\{q\}}) > 2
$$
라고 가정하자. 그러면 $(\mathcal{F}_n)$은 연접 $\mathcal{O}_P$-가군의
완비이다.
\end{proposition}

\begin{proof}
스킴의 코호몰로지의 보조정리 \ref{coherent-lemma-existence-easy}에 의해
보조정리를 증명할 때 $(\mathcal{F}_n)$을 그것과 $\mathcal{I}$-꼬임만큼
다른 대상으로 바꾸어도 된다(더 정확한 내용은 아래를 보라).
$T' = \{q \in Q \mid \dim(\overline{\{q\}}) = 0\}$ 및
$T = \{q \in Q \mid \dim(\overline{\{q\}}) \leq 1\}$로 놓는다. 명제의
가정은 정확히 $Q \subset P$, $(\mathcal{F}_n)$ 및 $T' \subset T \subset Q$가
$d = 1$에 대해 보조정리
\ref{algebraization-lemma-improvement-formal-coherent-module-better}의 조건을 만족한다는
것이다. 자명한 부등식 조작 외에, $\mathcal{I}_y^\wedge$가 $1$개의 원소로
생성되므로
$V(\mathfrak p) \cap V(\mathcal{I}^\wedge_y) = \{\mathfrak m^\wedge_y\}
\Leftrightarrow \dim(\mathcal{O}_{P, q}^\wedge/\mathfrak p) = 1$
임을 사용한다. 이 두 관찰을 결합하면 $(\mathcal{F}_n)$을 보조정리
\ref{algebraization-lemma-improvement-formal-coherent-module-better}에서 얻은
$\textit{Coh}(P, \mathcal{I})$의 대상 $(\mathcal{H}_n)$으로 바꾸어도 된다.
따라서 $(\mathcal{F}_n)$이 모든 $q \in Q$에 대해
$\text{depth}(\mathcal{F}''_{n, q}) + \dim(\overline{\{q\}}) \geq 2$를
만족하는 연접 $\mathcal{O}_P$-가군의 역계 $(\mathcal{F}_n'')$과
프로-동형이라고 가정할 수 있고 그렇게 한다.

\medskip\noindent
사상에 관하여 더 보기의 보조정리
\ref{more-morphisms-lemma-apply-proj-spec}와 같은 글의 절
\ref{more-morphisms-section-proj-spec}의 표기와 결과를 더 언급하지 않고
사용한다. 적어도 그 보조정리의 진술을 이해하지 않으면 이 증명은 의미가
없다. 우리 경우
$A = \bigoplus_{m \geq 0} \Gamma(P, \mathcal{L}^{\otimes m})$
는 각 등급 성분이 유한 차원 $k$-벡터공간인 유한형 $k$-대수이다. 스킴의
코호몰로지의 보조정리 \ref{coherent-lemma-coherent-proper-ample}를 보라.

\medskip\noindent
스킴의 코호몰로지의 보조정리
\ref{coherent-lemma-inverse-systems-pullback}에 의해 $\pi|_U : U \to P$에
의한 당김은 $\textit{Coh}(U, f\mathcal{O}_U)$의 대상
$(\pi|_U^*\mathcal{F}_n)$이며, 이는 연접 $\mathcal{O}_U$-가군의 역계
$(\pi|_U^*\mathcal{F}_n'')$과 프로-동형이다. 모든 $y \in U \cap Y$에 대해
$$
\text{depth}(\pi|_U^*\mathcal{F}''_{n, y}) + \delta_Z^Y(y) \geq 3
$$
라고 주장한다. $Z$의 모든 점이 닫혀 있으므로 모든 $y \in U \cap Y$에 대해
$\delta_Z^Y(y) \geq \dim(\overline{\{y\}})$이다. 보조정리
\ref{algebraization-lemma-discussion}을 보라. $q \in Q$를 $y$의 상이라 하자. 사상
$\pi : U \to P$는 상대 차원 $1$로 매끄러우므로 $y$는 $\pi$의 한 올의
닫힌점이거나 일반점이다. 따라서
$$
\text{depth}(\pi^*\mathcal{F}''_{n, y}) + \delta_Z^Y(y)
\geq
\text{depth}(\pi^*\mathcal{F}''_{n, y}) + \dim(\overline{\{y\}}) =
\text{depth}(\mathcal{F}''_{n, q}) + \dim(\overline{\{q\}}) + 1
$$
이다. 깊이 또는 차원 중 하나가 $1$만큼 증가하기 때문이다. 이로써 주장을
증명했다.

\medskip\noindent
보조정리 \ref{algebraization-lemma-cd-1-canonical}에 의해 $(\pi|_U^*\mathcal{F}_n)$은
$X$로 표준적으로 확장된다. 다음에 유의하라.
$$
M_n = \Gamma(U, \pi|_U^*\mathcal{F}_n) =
\bigoplus\nolimits_{m \in \mathbf{Z}}
\Gamma(P, \mathcal{F}_n \otimes_{\mathcal{O}_P} \mathcal{L}^{\otimes m})
$$
이는 표준적으로 등급 $A$-가군이다. 사상에 관하여 더 보기의 식
(\ref{more-morphisms-equation-cohomology-torsor})을 보라. 성질의 보조정리
\ref{properties-lemma-quasi-coherent-quasi-affine}에 의해
$\pi|_U^*\mathcal{F}_n = \widetilde{M_n}|_U$이다. 따라서 보조정리
\ref{algebraization-lemma-Gm-equivariant-extend-canonically}를 적용하여, $(M_n)$과
$(N/I^nN)$이 $A_+$-꼬임 가군을 법으로 한 등급 $A$-가군의 범주에서
프로-동형이 되는 유한 등급 $A$-가군 $N$을 찾을 수 있다. $\mathcal{F}$를
$N$에 대응하는 연접 $\mathcal{O}_P$-가군이라 하자. 스킴의 코호몰로지의
명제 \ref{coherent-proposition-coherent-modules-on-proj-general}를 보라.
같은 명제에 의해 $(\mathcal{F}/\mathcal{I}^n\mathcal{F})$은
$(\mathcal{F}_n)$과 프로-동형이다. 둘 다 $\textit{Coh}(P, \mathcal{I})$의
대상이므로 보조정리 \ref{algebraization-lemma-recognize-formal-coherent-modules}에 의해
결론이 성립한다.
\end{proof}

\begin{example}
\label{algebraization-example-lefschetz-existence}
$k$를 체라 하고 $X$를 $k$ 위의 고유 다양체라 하자. $Y \subset X$는
$\mathcal{O}_X(Y)$가 풍부한 유효 Cartier 인자라 하고, 대응하는 아이디얼층을
$\mathcal{I} \subset \mathcal{O}_X$라 쓰자. $(\mathcal{E}_n)$을 각
$\mathcal{E}_n$이 유한 국소 자유인 $\textit{Coh}(X, \mathcal{I})$의 대상이라
하자. 다음은 명제 \ref{algebraization-proposition-lefschetz-existence}의 몇 가지 특별한
경우이다.
\begin{enumerate}
\item $X$가 곡선 또는 곡면이면 아무것도 얻지 못한다.
\item $X$가 Cohen--Macaulay 3차원 다양체이면 $(\mathcal{E}_n)$은 연접
$\mathcal{O}_X$-가군 $\mathcal{E}$의 완비이다.
\item 더 일반적으로 $\dim(X) \geq 3$이고 $X$가 $(S_3)$이면
$(\mathcal{E}_n)$은 연접 $\mathcal{O}_X$-가군 $\mathcal{E}$의 완비이다.
\end{enumerate}
물론 $\mathcal{E}$가 존재하면 $\mathcal{E}$는 $Y$의 어떤 열린 근방에서 유한 국소 자유이다.
\end{example}

\begin{proposition}
\label{algebraization-proposition-lefschetz-equivalence}
$k$를 체라 하고 $X$를 $k$ 위의 고유 스킴이라 하자. $\mathcal{L}$을 풍부한
가역 $\mathcal{O}_X$-가군이라 하고 $s \in \Gamma(X, \mathcal{L})$라 하자.
$Y = Z(s)$를 $s$의 영점 스킴이라 하고, 대응하는 아이디얼층을
$\mathcal{I} \subset \mathcal{O}_X$라 쓰자. $\mathcal{V}$를 $Y$를 포함하는
$X$의 열린 부분스킴들의 역포함으로 순서매긴 집합이라 하자. 모든
$x \in X \setminus Y$에 대해
$$
\text{depth}(\mathcal{O}_{X, x}) + \dim(\overline{\{x\}}) > 2
$$
라고 가정하자. 그러면 완비화 함자
$$
\colim_\mathcal{V}
\textit{Coh}(\mathcal{O}_V)
\longrightarrow
\textit{Coh}(X, \mathcal{I})
$$
는 유한 국소 자유 대상들로 이루어진 충만한 부분범주 사이의 동치이다.
\end{proposition}

\begin{proof}
완전충실성을 증명하려면 모든 $m$에 대해
$$
\colim_\mathcal{V} \Gamma(V, \mathcal{L}^{\otimes m})
\longrightarrow
\lim \Gamma(Y_n, \mathcal{L}^{\otimes m}|_{Y_n})
$$
가 동형임을 보이면 충분하다. 보조정리
\ref{algebraization-lemma-completion-fully-faithful-general}를 보라. 이는 보조정리
\ref{algebraization-lemma-lefschetz-addendum}에서 따른다.

\medskip\noindent
본질적 전사성을 보이자. $(\mathcal{F}_n)$을
$\textit{Coh}(X, \mathcal{I})$의 유한 국소 자유 대상이라 하자. 그러면
$y \in Y$에 대해 $\mathcal{F}_y^\wedge = \lim \mathcal{F}_{n, y}$는 유한 자유
$\mathcal{O}_{X, y}^\wedge$-가군이다. $\mathfrak p \subset
\mathcal{O}_{X, y}^\wedge$를 $\mathfrak p \not \in
V(\mathcal{I}_y^\wedge)$인 소아이디얼이라 하자. 그러면 $\mathfrak p$는
$x \not \in Y$인 특수화 $x \leadsto y$에 대응하는 소아이디얼
$\mathfrak p_0 \subset \mathcal{O}_{X, y}$ 위에 놓인다. 국소 코호몰로지의
보조정리
\ref{local-cohomology-lemma-change-completion}
와 약간의 차원론(다양체의 절 \ref{varieties-section-algebraic-schemes}를
보라)에 의해
$$
\text{depth}((\mathcal{O}_{X, y}^\wedge)_\mathfrak p) +
\dim(\mathcal{O}_{X, y}^\wedge/\mathfrak p) =
\text{depth}(\mathcal{O}_{X, x}) +
\dim(\overline{\{x\}}) - \dim(\overline{\{y\}})
$$
이다. 따라서 우리의 가정은 명제 \ref{algebraization-proposition-lefschetz-existence}의
가정을 함의하며, $(\mathcal{F}_n)$은 어떤 연접 $\mathcal{O}_X$-가군
$\mathcal{F}$의 완비이다. 그러면 모든 $y \in Y$에 대해 $\mathcal{F}_y$는
유한 자유이고, 따라서 $\mathcal{F}$는 $Y$의 어떤 열린 근방 $V$에서 유한
국소 자유이다. 이로써 증명이 끝난다.
\end{proof}
